\chapter{Une vision riemannienne : géométrie du cône SPD et des matrices orthogonales}
\label{ch:context-riemann}


Nous avons jusqu'ici défini notre objet d'étude qui est la covariance au travers d'un approche statistique en explorant sa définition formelle et en regardant les problématique d'inférence statistique sans nous soucier de la structure particulière de son espace d'existence que nous avions fait remarquer être l'intérieur d'un cône. Sans se soucier de la géométrie de cet espace, nous avons systématiquement traité l'ensemble $\SPD{\Nvar}$ comme un sous-ensemble de l'espace vectoriel des matrices symétriques, dont on se contentait de vérifier \emph{a posteriori} qu'on n'en était pas sorti. Trois symptômes montrent que cette manière de faire est insuffisante.

Le premier est algorithmique. L'ensemble $\SPD{\Nvar}$ n'est pas un espace vectoriel: une itération de descente de gradient euclidienne appliquée à une fonction de coût définie sur les matrices de covariance peut parfaitement produire un itéré à valeurs propres négatives, sur lequel la fonction de coût n'est même plus définie. La contrainte est alors subie --- on projette, on pénalise, on réduit le pas --- plutôt qu'intégrée au problème.

Le deuxième est métrologique et nous l'avons annoncé dès la \cref{sec:context-covariance-estimation}: mesurer l'erreur d'un estimateur par $\norm{\hat{\MCov}-\MCov}_\frob$ revient à traiter les matrices de covariance comme des tableaux de nombres. Or le problème d'estimation, lui, possède des symétries que cette norme ignore. Si l'on applique aux observations un filtrage linéaire inversible $\x\mapsto\mat{A}\x$ --- blanchiment, changement de base d'antennes, simple changement d'unité --- alors toutes les covariances subissent $\MCov\mapsto\mat{A}\MCov\mat{A}^\T$, et il n'y a aucune raison pour que l'écart entre la vraie matrice et son estimée s'en trouve modifié par cette opération de filtrage. La norme de Frobenius, elle, est modifiée.

Le troisième est plus visuel. La moyenne arithmétique de deux matrices de covariance a un déterminant supérieur à celui que l'on attendrait: moyenner deux ellipsoïdes de concentration selon la règle euclidienne \emph{crée} de la puissance. Ce phénomène, connu sous le nom d'effet de gonflement (\textit{swelling effect}) \sidecite[0cm]{arsigny2006logeuclidean}, est une conséquence directe du fait que le déterminant n'est pas une fonction affine. \\

Ces trois symptômes ont une cause commune: nous avons plaqué une structure d'espace vectoriel sur un ensemble qui n'en est pas un. Regardons alors dans ce chapitre la structure géométrique pertinente : une variété riemannienne. Plutôt que d'être une référence dans le domaine de la géométrie riemannienne, on veut ici présenter uniquement les outils essentiels en tant qu'utilisateur sans avoir à subir le formalisme mathématique qui amène à l'existence de ces outils. En suivant l'approche dite des \emph{sous-variétés plongées}, qui évite tout l'appareillage des cartes et des atlas et permet de tout lire dans un espace ambiant euclidien\footnote{C'est le point de vue adopté dans les ouvrages orientés optimisation \sidecite[0.5cm]{absil2008optimization,boumal2023introduction}; pour une présentation en géométrie différentielle classique, on renverra à \sidecite[2cm]{gallot2004riemannian,lee2018riemannian}, et pour tout ce qui concerne spécifiquement le cône des matrices définies positives aux ouvrages de référence \sidecite[4cm]{bhatia2007positive,pennec2020riemannian}.}. Nous nous plaçons dans tout ce chapitre dans le \textbf{cas réel}, où les dimensions et les figures sont plus lisibles; la \cref{sec:context-riemann-complexe} établira le dictionnaire vers le cas complexe, et nous verrons que celui-ci est nettement plus mécanique que ne l'étaient ceux des \cref{sec:context-definition-complex,sec:context-distributions-complex}.




\section{Notion de géométrie riemannienne}
\label{sec:context-riemann-notion}

\subsection{Variétés et espaces tangents}

L'idée directrice est de remplacer «~ensemble de matrices vérifiant des contraintes~» par «~surface lisse~», et de ne travailler qu'avec les déplacements qui restent sur cette surface.

\begin{definition}[{Sous-variété plongée}]
\label{def:context-riemann-manifold}
  Soit $\ambient$ un espace vectoriel euclidien de dimension finie, appelé \textbf{espace ambiant}. Un sous-ensemble $\Manifold\subset\ambient$ est une \textbf{sous-variété plongée} de dimension $d$ si, au voisinage de chacun de ses points, il peut être identifié de manière lisse et inversible à un ouvert de $\R^d$. On dira simplement \textbf{variété} dans la suite.
\end{definition}

Deux régimes très différents se cachent derrière cette définition, et la distinction gouverne tout le coût algorithmique de ce qui suivra:
\begin{itemize}
  \item $\Manifold$ est un \textbf{ouvert} de $\ambient$. C'est le cas du cône $\SPD{\Nvar}$, qui est un ouvert de l'espace $\Sym{\Nvar}$ des matrices symétriques réelles: toute matrice symétrique suffisamment proche d'une matrice définie positive est encore définie positive, que l'on caractérise cette dernière par le signe des valeurs propres ou, de façon équivalente, par la positivité stricte des mineurs principaux dominants (critère de Sylvester). Dans les deux cas, la contrainte reste une inégalité \emph{stricte}: elle ne fait que découper un ouvert, sans réduire la dimension. Celle-ci est donc celle de l'espace ambiant, $\Nvar(\Nvar+1)/2$;
  \item $\Manifold$ est \textbf{définie par une submersion}. C'est le cas de la sphère unité, ou de l'ensemble des matrices orthogonales $\Orthgroup{\Nvar}=\{\Orth\in\R^{\Nvar\times\Nvar};\,\Orth^\T\Orth=\Id{\Nvar}\}$: on écrit $\Manifold=h^{-1}(0)$ pour une application lisse $h:\ambient\to\R^k$ dont la différentielle est surjective en tout point de $\Manifold$ --- une \textbf{submersion} --- ici $h(\Orth)=\Orth^\T\Orth-\Id{\Nvar}$, à valeurs dans $\Sym{\Nvar}$. Cette condition garantit alors que $\Manifold$ est une variété lisse, de dimension $\dim\ambient-k$, strictement inférieure à celle de l'ambiant.
\end{itemize}

Cette distinction n'est pas cosmétique: dans le premier cas les opérations de projection que nous allons introduire seront triviales, dans le second elles ne le seront pas. \\

\begin{figure}[htb]
  \centering
  % Variete, courbe tracee dessus, et espace tangent en un point.
% Utilise dans la \cref{sec:context-riemann-notion}.
% Styles et briques communs aux figures de geometrie riemannienne du
% \cref{ch:context-riemann}. Ce fichier peut etre \input plusieurs fois.
\definecolor{rmanifold}{RGB}{90,90,90}
\definecolor{rtangent}{RGB}{50,110,120}
\definecolor{rgeodesic}{RGB}{40,70,140}
\definecolor{rvector}{RGB}{175,95,45}
\tikzset{
  riemann/.style={
    >={Stealth[length=4.5pt,width=3.5pt]},
    line join=round,
    font=\footnotesize,
  },
  % La nappe qui represente la variete, centree sur l'origine.
  rsurface/.pic={
    \fill[rmanifold!7]
      (-3.2,0.15) .. controls (-2.95,1.75) and (-1.5,2.35) .. (0,2.35)
      .. controls (1.5,2.35) and (2.95,1.7) .. (3.2,0.15)
      .. controls (1.9,-0.7) and (-1.9,-0.7) .. (-3.2,0.15) -- cycle;
    \draw[rmanifold!70, thick]
      (-3.2,0.15) .. controls (-2.95,1.75) and (-1.5,2.35) .. (0,2.35)
      .. controls (1.5,2.35) and (2.95,1.7) .. (3.2,0.15);
    \draw[rmanifold!70, thick]
      (-3.2,0.15) .. controls (-1.9,-0.7) and (1.9,-0.7) .. (3.2,0.15);
  },
  % Le plan tangent, a placer en un point donne via un shift.
  rplane/.pic={
    \fill[rtangent!10, draw=rtangent!65]
      (-1.35,-0.42) -- (1.55,-0.88) -- (2.05,0.22) -- (-0.85,0.68) -- cycle;
  },
}

\begin{tikzpicture}[riemann]
  \coordinate (P) at (-1.04,1.38);

  \path (0,0) pic {rsurface};
  \node[rmanifold!85] at (2.55,1.55) {$\Manifold$};

  % plan tangent
  \path (P) pic[scale=0.85] {rplane};
  \node[rtangent!90, anchor=east, inner sep=1pt] at ($(P)+(-1.20,0.50)$)
    {$\Tangent{\param}{\Manifold}$};

  % courbe tracee sur la variete
  \draw[rgeodesic, very thick]
    (-2.6,0.85) .. controls (-2.0,1.30) and (-0.2,1.55) .. (0.9,1.62);
  \node[rgeodesic, anchor=south west, inner sep=1pt] at (0.72,1.62) {$\geod$};

  % vitesse : un vecteur tangent
  \draw[rvector, very thick, ->] (P) -- ($(P)+(1.20,0.15)$);
  \node[rvector, anchor=north west, inner sep=1pt] at ($(P)+(0.55,0.10)$)
    {$\dot{\geod}(0)$};

  \fill (P) circle (1.7pt);
  \node[anchor=south east, inner sep=2pt] at (P) {$\param$};
\end{tikzpicture}

  \caption{Une variété $\Manifold$, une courbe $\geod$ tracée dessus et
    l'espace tangent en $\param$. Ce dernier rassemble les vitesses en $\param$
    de toutes les courbes passant par ce point: c'est l'ensemble construit à partir des directions des
    déplacements infinitésimaux qui ne quittent pas $\Manifold$.}
  \label{fig:context-riemann-tangent}
\end{figure}


Une fois la variété définie, l'outil central de toute la géométrie riemannienne est, en chacun de ses points, un \textbf{plan tangent} (notion illustrée en \cref{fig:context-riemann-tangent}) : une approximation lisse de $\Manifold$ au premier ordre. Intuitivement, si $\geod$ est une courbe lisse tracée sur $\Manifold$ et passant par $\param$ à la vitesse $\tvec$, elle s'écrit au premier ordre $\geod(t)=\param+t\tvec+o(t)$: le vecteur $\tvec$ est donc la meilleure approximation linéaire d'un chemin qui reste sur $\Manifold$, autrement dit la direction d'un \emph{déplacement admissible} depuis $\param$. Il ne s'agit toutefois que d'une admissibilité \emph{au premier ordre}: rien ne garantit que le point $\param+t\tvec$ lui-même reste sur $\Manifold$ pour $t$ fini --- c'est le cas pour $\SPD{\Nvar}$, ouvert dans $\Sym{\Nvar}$, mais pas pour $\Orthgroup{\Nvar}$, où $\Orth+t\mat{A}$ n'est orthogonale, en général, pour aucun $t\neq0$. L'ensemble de ces directions admissibles est l'espace tangent :


\begin{definition}[{Espace tangent}]
\label{def:context-riemann-tangent}
  Soit $\Manifold$ une variété et $\param\in\Manifold$. L'\textbf{espace tangent} à $\Manifold$ en $\param$ est l'ensemble des vitesses en $\param$ des courbes tracées sur $\Manifold$ et passant par $\param$:
  \begin{equation}
    \Tangent{\param}{\Manifold}
    \eqdef
    \left\{
      \dot{\geod}(0);\;
      \geod:\R\to\Manifold \text{ lisse},\;
      \geod(0)=\param
    \right\}.
    \label{eq:context-riemann-tangent}
  \end{equation}
  C'est un espace vectoriel de dimension $d$. Ses éléments sont appelés \textbf{vecteurs tangents}.
\end{definition}

  L'espace tangent est l'objet qui joue, localement, le rôle que jouait l'espace vectoriel tout entier dans le cas euclidien: c'est là que vivront les gradients, les directions de descente et les erreurs d'estimation. Pour le cône, $\SPD{\Nvar}$ étant ouvert dans $\Sym{\Nvar}$, on a simplement $\Tangent{\MCov}{\SPD{\Nvar}}=\Sym{\Nvar}$ en tout point $\MCov$: une perturbation admissible d'une matrice de covariance est une matrice symétrique quelconque, sans autre contrainte. Pour une variété définie par une submersion, en revanche, l'espace tangent s'obtient en différentiant l'équation $h=0$ qui la définit. Ainsi, en dérivant $\Orth^\T\Orth=\Id{\Nvar}$ le long d'une courbe, on trouve que $\Tangent{\Orth}{\Orthgroup{\Nvar}}=\{\Orth\mat{A};\,\mat{A}^\T=-\mat{A}\}$: les perturbations admissibles d'une matrice orthogonale sont les matrices antisymétriques, transportées en $\Orth$.


\subsection{Métrique riemannienne}

À ce stade nous savons quels déplacements sont admissibles, mais pas comment les mesurer. C'est l'objet de la métrique, et c'est le point où se joue l'essentiel.

\begin{definition}[{Métrique riemannienne}]
\label{def:context-riemann-metric}
  Une \textbf{métrique riemannienne} sur $\Manifold$ est la donnée, en chaque point $\param\in\Manifold$, d'un produit scalaire $\metric{\cdot}{\cdot}{\param}$ sur $\Tangent{\param}{\Manifold}$, dépendant de manière lisse du point $\param$ (\cref{fig:context-riemann-metrique}). Le couple $(\Manifold, \metric{\cdot}{\cdot}{\cdot})$ est appelé \textbf{variété riemannienne}.
\end{definition}

La métrique induit immédiatement une norme $\normriem{\tvec}{\param}^2=\metric{\tvec}{\tvec}{\param}$ sur chaque espace tangent, donc une longueur pour toute courbe $\geod:[0,1]\to\Manifold$:
\begin{equation}
  L(\geod)
  \eqdef
  \int_0^1 \normriem{\dot{\geod}(t)}{\geod(t)}\diff t.
  \label{eq:context-riemann-longueur}
\end{equation}

\begin{figure}[htb]
  \centering
  % La metrique : produit scalaire sur l'espace tangent, donc norme et angle,
% et par integration longueur d'une courbe.
% Utilise dans la \cref{sec:context-riemann-notion}.
% Styles et briques communs aux figures de geometrie riemannienne du
% \cref{ch:context-riemann}. Ce fichier peut etre \input plusieurs fois.
\definecolor{rmanifold}{RGB}{90,90,90}
\definecolor{rtangent}{RGB}{50,110,120}
\definecolor{rgeodesic}{RGB}{40,70,140}
\definecolor{rvector}{RGB}{175,95,45}
\tikzset{
  riemann/.style={
    >={Stealth[length=4.5pt,width=3.5pt]},
    line join=round,
    font=\footnotesize,
  },
  % La nappe qui represente la variete, centree sur l'origine.
  rsurface/.pic={
    \fill[rmanifold!7]
      (-3.2,0.15) .. controls (-2.95,1.75) and (-1.5,2.35) .. (0,2.35)
      .. controls (1.5,2.35) and (2.95,1.7) .. (3.2,0.15)
      .. controls (1.9,-0.7) and (-1.9,-0.7) .. (-3.2,0.15) -- cycle;
    \draw[rmanifold!70, thick]
      (-3.2,0.15) .. controls (-2.95,1.75) and (-1.5,2.35) .. (0,2.35)
      .. controls (1.5,2.35) and (2.95,1.7) .. (3.2,0.15);
    \draw[rmanifold!70, thick]
      (-3.2,0.15) .. controls (-1.9,-0.7) and (1.9,-0.7) .. (3.2,0.15);
  },
  % Le plan tangent, a placer en un point donne via un shift.
  rplane/.pic={
    \fill[rtangent!10, draw=rtangent!65]
      (-1.35,-0.42) -- (1.55,-0.88) -- (2.05,0.22) -- (-0.85,0.68) -- cycle;
  },
}

\begin{tikzpicture}[riemann]
  \coordinate (P) at (-1.70,1.15);

  \path (0,0) pic {rsurface};
  \node[rmanifold!85] at (2.85,1.15) {$\Manifold$};

  % plan tangent et deux vecteurs
  \path (P) pic[scale=0.75] {rplane};
  \draw[rvector, very thick, ->] (P) -- ($(P)+(1.10,-0.16)$);
  \node[rvector, anchor=north west, inner sep=1pt] at ($(P)+(0.85,-0.14)$) {$\tvec$};
  \draw[rvector, very thick, ->] (P) -- ($(P)+(0.40,0.46)$);
  \node[rvector, anchor=south east, inner sep=1pt] at ($(P)+(0.42,0.44)$) {$\tvecbis$};
  \draw[rtangent!80] ($(P)+(0.46,-0.07)$) arc[start angle=-8, end angle=48, radius=0.47];
  \node[rtangent!90, font=\scriptsize, anchor=west] at ($(P)+(0.80,0.42)$)
    {$\metric{\tvec}{\tvecbis}{\param}$};
  \node[rtangent!90, anchor=north, inner sep=2pt] at ($(P)+(0.10,-0.50)$)
    {$\Tangent{\param}{\Manifold}$};

  % longueur d'une courbe : vitesses echantillonnees le long du chemin
  \draw[rgeodesic, very thick]
    (0.55,1.85) .. controls (1.45,1.95) and (2.05,1.55) .. (2.45,0.95);
  \node[rgeodesic, anchor=south east, inner sep=1pt] at (0.60,1.88) {$\geod$};
  \foreach \x/\y/\dx/\dy in {0.85/1.91/0.30/0.05, 1.60/1.89/0.28/-0.13, 2.22/1.42/0.18/-0.28}
    \draw[rvector!85, thick, ->] (\x,\y) -- ++(\dx,\dy);
  \node[rgeodesic!85, anchor=north, font=\scriptsize] at (1.35,1.32)
    {$L(\geod)=\int\normriem{\dot{\geod}}{\geod}$};

  \fill (P) circle (1.7pt);
  \node[anchor=south east, inner sep=2pt] at (P) {$\param$};
\end{tikzpicture}

  \caption{La métrique est un produit scalaire attaché à chaque point: elle
    donne un sens à la norme d'un vecteur tangent et à l'angle entre deux
    d'entre eux, et par intégration le long d'une courbe à la longueur
    \cref{eq:context-riemann-longueur}. Rien n'oblige deux points distincts à
    porter le même produit scalaire, et c'est précisément cette liberté que la
    suite exploite.}
  \label{fig:context-riemann-metrique}
\end{figure}

\begin{remark}[{La métrique est un choix de modélisation}]
\label{rem:context-riemann-choixmetrique}
  Il faut insister sur ce point, car il structure toute la suite: la métrique n'est pas une propriété de l'ensemble $\Manifold$, c'est un choix que l'on fait. Un même ensemble porte une infinité de métriques, et chacune d'elles définit sa propre notion de longueur, d'angle, de droite et de distance. Sur le cône, choisir la métrique euclidienne $\metric{\tvec}{\tvecbis}{\MCov}=\trace{\tvec\tvecbis}$ revient exactement à renoncer à toute structure: la géométrie obtenue est celle de $\Sym{\Nvar}$ tout entier, et la spécificité de $\SPD{\Nvar}$ n'est pas exploitée \sidecite[0cm]{pennec2006riemannian}. C'est précisément ce que nous voulons éviter, et le travail de la \cref{sec:context-riemann-spd} consistera à choisir mieux.
\end{remark}

\subsection{La boîte à outils}

De la métrique découle une série d'objets dont nous n'aurons besoin que d'une poignée. Nous les énumérons ici dans l'ordre où ils se construisent.

\paragraph{Projection orthogonale.} Beaucoup de calculs produisent naturellement un élément de l'espace ambiant $\ambient$ --- un gradient euclidien, une différence de deux points --- qui n'appartient pas à $\Tangent{\param}{\Manifold}$. La projection orthogonale $\projman_\param:\ambient\to\Tangent{\param}{\Manifold}$, prise au sens de la métrique, est l'opération qui le rend admissible (\cref{fig:context-riemann-projection}). Lorsque $\Manifold$ est un ouvert de $\ambient$, elle est l'identité.

\begin{figure}[htb]
  \centering
  % Projection orthogonale d'un vecteur de l'espace ambiant sur l'espace tangent.
% Utilise dans la \cref{sec:context-riemann-notion}.
% Styles et briques communs aux figures de geometrie riemannienne du
% \cref{ch:context-riemann}. Ce fichier peut etre \input plusieurs fois.
\definecolor{rmanifold}{RGB}{90,90,90}
\definecolor{rtangent}{RGB}{50,110,120}
\definecolor{rgeodesic}{RGB}{40,70,140}
\definecolor{rvector}{RGB}{175,95,45}
\tikzset{
  riemann/.style={
    >={Stealth[length=4.5pt,width=3.5pt]},
    line join=round,
    font=\footnotesize,
  },
  % La nappe qui represente la variete, centree sur l'origine.
  rsurface/.pic={
    \fill[rmanifold!7]
      (-3.2,0.15) .. controls (-2.95,1.75) and (-1.5,2.35) .. (0,2.35)
      .. controls (1.5,2.35) and (2.95,1.7) .. (3.2,0.15)
      .. controls (1.9,-0.7) and (-1.9,-0.7) .. (-3.2,0.15) -- cycle;
    \draw[rmanifold!70, thick]
      (-3.2,0.15) .. controls (-2.95,1.75) and (-1.5,2.35) .. (0,2.35)
      .. controls (1.5,2.35) and (2.95,1.7) .. (3.2,0.15);
    \draw[rmanifold!70, thick]
      (-3.2,0.15) .. controls (-1.9,-0.7) and (1.9,-0.7) .. (3.2,0.15);
  },
  % Le plan tangent, a placer en un point donne via un shift.
  rplane/.pic={
    \fill[rtangent!10, draw=rtangent!65]
      (-1.35,-0.42) -- (1.55,-0.88) -- (2.05,0.22) -- (-0.85,0.68) -- cycle;
  },
}

\begin{tikzpicture}[riemann]
  \coordinate (P) at (-0.45,1.35);

  \path (0,0) pic {rsurface};
  \node[rmanifold!85] at (2.55,1.80) {$\Manifold$};

  \path (P) pic {rplane};
  \node[rtangent!90, anchor=east] at ($(P)+(-1.30,0.50)$)
    {$\Tangent{\param}{\Manifold}$};

  % vecteur ambiant, hors du plan tangent
  \draw[rmanifold!85, very thick, ->] (P) -- ($(P)+(1.10,1.15)$);
  \node[rmanifold, anchor=south west, inner sep=1pt] at ($(P)+(1.08,1.13)$)
    {$\mat{Z}\in\ambient$};

  % sa projection
  \draw[rvector, very thick, ->] (P) -- ($(P)+(1.32,-0.18)$);
  \node[rvector, anchor=north, inner sep=2pt] at ($(P)+(1.05,-0.22)$)
    {$\projman_{\param}(\mat{Z})$};

  % trait de projection
  \draw[rmanifold!60, dashed] ($(P)+(1.10,1.15)$) -- ($(P)+(1.32,-0.18)$);

  \fill (P) circle (1.7pt);
  \node[anchor=north east, inner sep=2pt] at (P) {$\param$};
\end{tikzpicture}

  \caption{Projection orthogonale. Un objet $\mat{Z}$ calculé dans l'espace
    ambiant n'est en général pas un déplacement admissible; $\projman_\param$ en
    retient la seule partie qui le soit.}
  \label{fig:context-riemann-projection}
\end{figure}

\paragraph{Connexion de Levi-Civita.} Comparer un vecteur du plan tangent $\tvecbis$ en deux points voisins $\param$ et $\param+t\tvec$, c'est par définition en dériver le champ le long de $\tvec$. La difficulté propre à la variété est que $\tvecbis(\param)$ et $\tvecbis(\param+t\tvec)$ n'habitent pas le même espace tangent: la dérivée directionnelle ambiante $\dirderiv\tvecbis(\param)[\tvec]$ n'a donc aucune raison de rester tangente à $\Manifold$, ni d'être compatible avec la métrique choisie. La \textbf{connexion} $\connection_{\tvec}\tvecbis$ est cette dérivée directionnelle, corrigée pour rester tangente et compatible avec la métrique; parmi toutes les corrections possibles, une seule --- la \textbf{connexion de Levi-Civita} --- est de plus sans torsion, et c'est celle que l'on retient toujours. \\

En partant de cette connexion, on peut alors définir la notion de géodésique :

\paragraph{Géodésiques.} Une \textbf{géodésique} est une courbe qui ne tourne pas: sa vitesse ne change jamais de direction, seulement de position, si bien qu'elle reste parallèle à elle-même tout au long du trajet --- c'est l'\textbf{accélération covariante} $\connection_{\dot\geod(t)}\dot\geod(t)$, généralisation de $\ddot\geod$ sur la variété, qui doit rester nulle:
\begin{equation}
  \connection_{\dot{\geod}(t)}\dot{\geod}(t) = \vector{0}.
  \label{eq:context-riemann-geodesique}
\end{equation}
C'est la généralisation naturelle de la ligne droite. Une géodésique est déterminée soit par un point de départ et une vitesse initiale, soit par ses deux extrémités. Localement, c'est aussi le chemin le plus court (\cref{fig:context-riemann-geodesique}), ce qui donne la \textbf{distance géodésique}
\begin{equation}
  \geodist(\param_0,\param_1)
  \eqdef
  \inf\left\{ L(\geod);\; \geod(0)=\param_0,\, \geod(1)=\param_1 \right\}.
  \label{eq:context-riemann-distance}
\end{equation}

\begin{figure}[htb]
  \centering
  % Geodesique entre deux points et distance geodesique, comparees a une autre
% courbe tracee sur la variete.
% Utilise dans la \cref{sec:context-riemann-notion}.
% Styles et briques communs aux figures de geometrie riemannienne du
% \cref{ch:context-riemann}. Ce fichier peut etre \input plusieurs fois.
\definecolor{rmanifold}{RGB}{90,90,90}
\definecolor{rtangent}{RGB}{50,110,120}
\definecolor{rgeodesic}{RGB}{40,70,140}
\definecolor{rvector}{RGB}{175,95,45}
\tikzset{
  riemann/.style={
    >={Stealth[length=4.5pt,width=3.5pt]},
    line join=round,
    font=\footnotesize,
  },
  % La nappe qui represente la variete, centree sur l'origine.
  rsurface/.pic={
    \fill[rmanifold!7]
      (-3.2,0.15) .. controls (-2.95,1.75) and (-1.5,2.35) .. (0,2.35)
      .. controls (1.5,2.35) and (2.95,1.7) .. (3.2,0.15)
      .. controls (1.9,-0.7) and (-1.9,-0.7) .. (-3.2,0.15) -- cycle;
    \draw[rmanifold!70, thick]
      (-3.2,0.15) .. controls (-2.95,1.75) and (-1.5,2.35) .. (0,2.35)
      .. controls (1.5,2.35) and (2.95,1.7) .. (3.2,0.15);
    \draw[rmanifold!70, thick]
      (-3.2,0.15) .. controls (-1.9,-0.7) and (1.9,-0.7) .. (3.2,0.15);
  },
  % Le plan tangent, a placer en un point donne via un shift.
  rplane/.pic={
    \fill[rtangent!10, draw=rtangent!65]
      (-1.35,-0.42) -- (1.55,-0.88) -- (2.05,0.22) -- (-0.85,0.68) -- cycle;
  },
}

\begin{tikzpicture}[riemann]
  \coordinate (P) at (-1.95,1.05);
  \coordinate (Q) at ( 1.85,1.35);

  \path (0,0) pic {rsurface};
  \node[rmanifold!85] at (2.65,0.55) {$\Manifold$};

  % une courbe quelconque joignant les deux points
  \draw[rmanifold!60, thick, dashed]
    (P) .. controls (-1.20,-0.15) and (1.05,-0.05) .. (Q);
  \node[rmanifold, anchor=north, font=\scriptsize] at (-0.15,0.05)
    {$L(\geod')>\geodist(\param_0,\param_1)$};

  % la geodesique
  \draw[rgeodesic, very thick]
    (P) .. controls (-0.85,2.05) and (0.75,2.15) .. (Q);
  \node[rgeodesic, anchor=east, inner sep=3pt] at (-1.35,1.72) {$\geod$};
  \node[rgeodesic!85, anchor=south, font=\scriptsize] at (0.05,2.06)
    {$L(\geod)=\geodist(\param_0,\param_1)$};

  \fill (P) circle (1.7pt);
  \node[anchor=east, inner sep=3pt] at (P) {$\param_0$};
  \fill (Q) circle (1.7pt);
  \node[anchor=west, inner sep=3pt] at (Q) {$\param_1$};
\end{tikzpicture}

  \caption{Géodésique et distance géodésique. Parmi toutes les courbes tracées
    sur $\Manifold$ et joignant $\param_0$ à $\param_1$, la géodésique $\geod$
    est celle de longueur minimale; cette longueur \emph{est} la distance
    $\geodist(\param_0,\param_1)$.}
  \label{fig:context-riemann-geodesique}
\end{figure}

\paragraph{Applications exponentielle et logarithme.} L'\textbf{exponentielle riemannienne} $\exp_\param:\Tangent{\param}{\Manifold}\to\Manifold$ associe à un vecteur tangent $\tvec$ le point $\geod(1)$ atteint en suivant la géodésique partant de $\param$ à la vitesse $\tvec$. Elle réalise donc «~un déplacement de $\tvec$ en restant sur la variété~». Son inverse local, le \textbf{logarithme riemannien} $\log_{\param_0}(\param_1)\in\Tangent{\param_0}{\Manifold}$, est le vecteur tangent qui pointe de $\param_0$ vers $\param_1$ et dont la norme vaut $\geodist(\param_0,\param_1)$. Ces deux applications, illustrées en \cref{fig:context-riemann-explog} permettent de faire des opérations qui nous sortent de la variété et d'en revenir.

\begin{figure}[htb]
  \centering
  % Applications exponentielle et logarithme riemanniennes.
% Utilise dans la \cref{sec:context-riemann-notion}.
% Styles et briques communs aux figures de geometrie riemannienne du
% \cref{ch:context-riemann}. Ce fichier peut etre \input plusieurs fois.
\definecolor{rmanifold}{RGB}{90,90,90}
\definecolor{rtangent}{RGB}{50,110,120}
\definecolor{rgeodesic}{RGB}{40,70,140}
\definecolor{rvector}{RGB}{175,95,45}
\tikzset{
  riemann/.style={
    >={Stealth[length=4.5pt,width=3.5pt]},
    line join=round,
    font=\footnotesize,
  },
  % La nappe qui represente la variete, centree sur l'origine.
  rsurface/.pic={
    \fill[rmanifold!7]
      (-3.2,0.15) .. controls (-2.95,1.75) and (-1.5,2.35) .. (0,2.35)
      .. controls (1.5,2.35) and (2.95,1.7) .. (3.2,0.15)
      .. controls (1.9,-0.7) and (-1.9,-0.7) .. (-3.2,0.15) -- cycle;
    \draw[rmanifold!70, thick]
      (-3.2,0.15) .. controls (-2.95,1.75) and (-1.5,2.35) .. (0,2.35)
      .. controls (1.5,2.35) and (2.95,1.7) .. (3.2,0.15);
    \draw[rmanifold!70, thick]
      (-3.2,0.15) .. controls (-1.9,-0.7) and (1.9,-0.7) .. (3.2,0.15);
  },
  % Le plan tangent, a placer en un point donne via un shift.
  rplane/.pic={
    \fill[rtangent!10, draw=rtangent!65]
      (-1.35,-0.42) -- (1.55,-0.88) -- (2.05,0.22) -- (-0.85,0.68) -- cycle;
  },
}

\begin{tikzpicture}[riemann]
  \coordinate (P) at (-1.55,1.20);
  \coordinate (Q) at ( 1.75,1.45);

  \path (0,0) pic {rsurface};
  \node[rmanifold!85] at (2.75,0.75) {$\Manifold$};

  \path (P) pic[scale=0.8] {rplane};
  \node[rtangent!90, anchor=east, inner sep=1pt] at ($(P)+(-1.15,0.40)$)
    {$\Tangent{\param_0}{\Manifold}$};

  % la geodesique
  \draw[rgeodesic, very thick]
    (P) .. controls (-0.55,2.10) and (0.75,2.20) .. (Q);

  % le vecteur tangent
  \draw[rvector, very thick, ->] (P) -- ($(P)+(1.60,-0.23)$);
  \node[rvector, anchor=north, inner sep=2pt] at ($(P)+(0.80,-0.34)$)
    {$\tvec=\log_{\param_0}(\param_1)$};

  % le retour sur la variete
  \draw[rvector!60, dotted, very thick, ->]
    ($(P)+(1.63,-0.21)$) .. controls (0.75,0.55) and (1.60,0.70) .. ($(Q)+(0.02,-0.12)$);
  \node[rvector!85, anchor=north, inner sep=2pt, font=\scriptsize] at (1.05,0.62)
    {$\exp_{\param_0}(\tvec)$};

  \fill (P) circle (1.7pt);
  \node[anchor=east, inner sep=3pt] at (P) {$\param_0$};
  \fill (Q) circle (1.7pt);
  \node[anchor=west, inner sep=3pt] at (Q) {$\param_1$};
\end{tikzpicture}

  \caption{Exponentielle et logarithme riemanniens. Le logarithme lit un point
    de la variété comme un vecteur de l'espace tangent, l'exponentielle fait le
    trajet inverse en suivant la géodésique. Ce sont les analogues de la
    soustraction et de l'addition du cas euclidien.}
  \label{fig:context-riemann-explog}
\end{figure}

\paragraph{Transport parallèle.} Le \textbf{transport parallèle} $\transport_{\param_0\to\param_1}$ déplace un vecteur tangent de $\Tangent{\param_0}{\Manifold}$ vers $\Tangent{\param_1}{\Manifold}$ le long de la géodésique qui joint les deux points, en préservant la norme et les angles (\cref{fig:context-riemann-transport}). Nous en aurons besoin en \cref{sec:context-riemann-optimization} pour combiner des gradients calculés en des itérés successifs, ce que fait toute méthode de gradient conjugué.

\begin{figure}[htb]
  \centering
  % Transport parallele : deux vecteurs tangents deplaces le long d'une
% geodesique, l'angle qu'ils forment etant conserve.
% Utilise dans la \cref{sec:context-riemann-notion}.
% Styles et briques communs aux figures de geometrie riemannienne du
% \cref{ch:context-riemann}. Ce fichier peut etre \input plusieurs fois.
\definecolor{rmanifold}{RGB}{90,90,90}
\definecolor{rtangent}{RGB}{50,110,120}
\definecolor{rgeodesic}{RGB}{40,70,140}
\definecolor{rvector}{RGB}{175,95,45}
\tikzset{
  riemann/.style={
    >={Stealth[length=4.5pt,width=3.5pt]},
    line join=round,
    font=\footnotesize,
  },
  % La nappe qui represente la variete, centree sur l'origine.
  rsurface/.pic={
    \fill[rmanifold!7]
      (-3.2,0.15) .. controls (-2.95,1.75) and (-1.5,2.35) .. (0,2.35)
      .. controls (1.5,2.35) and (2.95,1.7) .. (3.2,0.15)
      .. controls (1.9,-0.7) and (-1.9,-0.7) .. (-3.2,0.15) -- cycle;
    \draw[rmanifold!70, thick]
      (-3.2,0.15) .. controls (-2.95,1.75) and (-1.5,2.35) .. (0,2.35)
      .. controls (1.5,2.35) and (2.95,1.7) .. (3.2,0.15);
    \draw[rmanifold!70, thick]
      (-3.2,0.15) .. controls (-1.9,-0.7) and (1.9,-0.7) .. (3.2,0.15);
  },
  % Le plan tangent, a placer en un point donne via un shift.
  rplane/.pic={
    \fill[rtangent!10, draw=rtangent!65]
      (-1.35,-0.42) -- (1.55,-0.88) -- (2.05,0.22) -- (-0.85,0.68) -- cycle;
  },
}

\begin{tikzpicture}[riemann]
  \coordinate (P) at (-1.85,1.05);
  \coordinate (Q) at ( 1.65,1.50);

  \path (0,0) pic {rsurface};
  \node[rmanifold!85] at (2.75,0.75) {$\Manifold$};

  % la geodesique support
  \draw[rgeodesic, very thick]
    (P) .. controls (-0.75,1.90) and (0.60,2.05) .. (Q);
  \node[rgeodesic!85, anchor=north, inner sep=3pt] at (-0.15,1.55)
    {$\transport_{\param_0\to\param_1}$};

  % le couple de vecteurs en theta_0
  \draw[rvector, very thick, ->] (P) -- ++(0.95,-0.20);
  \node[rvector, anchor=north, inner sep=1pt] at ($(P)+(0.95,-0.22)$) {$\tvec$};
  \draw[rvector, very thick, ->] (P) -- ++(0.42,0.52);
  \node[rvector, anchor=east, inner sep=1pt] at ($(P)+(0.42,0.52)$) {$\tvecbis$};
  \draw[rtangent!85] ($(P)+(0.52,-0.11)$) arc[start angle=-12, end angle=51, radius=0.53];
  \node[rtangent!90, font=\scriptsize, anchor=west, inner sep=1pt] at ($(P)+(0.60,0.24)$) {$\phi$};

  % le couple transporte en theta_1 : meme ecart angulaire, tourne dans le
  % meme sens que la geodesique (qui vire vers le bas a droite)
  \draw[rvector, very thick, ->] (Q) -- ++(0.82,-0.51);
  \node[rvector, anchor=north, inner sep=1pt] at ($(Q)+(0.82,-0.51)$)
    {$\transport(\tvec)$};
  \draw[rvector, very thick, ->] (Q) -- ++(0.57,0.34);
  \node[rvector, anchor=south west, inner sep=1pt] at ($(Q)+(0.57,0.34)$)
    {$\transport(\tvecbis)$};
  \draw[rtangent!85] ($(Q)+(0.46,-0.29)$) arc[start angle=-32, end angle=31, radius=0.54];
  \node[rtangent!90, font=\scriptsize, anchor=west, inner sep=1pt] at ($(Q)+(0.75,-0.01)$) {$\phi$};

  \fill (P) circle (1.7pt);
  \node[anchor=north east, inner sep=2pt] at (P) {$\param_0$};
  \fill (Q) circle (1.7pt);
  \node[anchor=north west, inner sep=2pt] at (Q) {$\param_1$};
\end{tikzpicture}

  \caption{Transport parallèle le long d'une géodésique. Les normes et l'angle
    $\phi$ que forment les deux vecteurs sont les mêmes en $\param_0$ et en
    $\param_1$: c'est ce qui permet de comparer, et donc de combiner, des
    vecteurs attachés à deux points distincts de la variété.}
  \label{fig:context-riemann-transport}
\end{figure}

Ces six objets --- projection, connexion, géodésique, distance, exponentielle et logarithme, transport --- suffisent à tout ce qui suit.

\begin{remark}[{Ce dont nous n'aurons pas besoin}]
  Nous n'aborderons ni la courbure, ni le hessien riemannien, ni les méthodes du second ordre. La courbure n'apparaîtra qu'en \cref{sec:context-geometry-icrb}, sous forme de termes que l'on négligera, et une fois en \cref{sec:context-riemann-moyennes} sous forme d'un énoncé d'unicité admis. Le lecteur intéressé trouvera ces développements dans \sidecite[0cm]{absil2008optimization,boumal2023introduction}.
\end{remark}

\begin{remark}[{Calcul différentiel matriciel}]
\label{rem:context-riemann-calculdiff}
  Tout ce qui suit repose sur une seule notion de calcul différentiel, la \textbf{dérivée directionnelle} d'une fonction $f:\ambient\to\R$ en $\param$ dans la direction $\tvec$:
  \begin{equation}
    \dirderiv f(\param)[\tvec]
    \eqdef
    \lim_{t\to 0}\frac{f(\param+t\tvec)-f(\param)}{t}.
    \label{eq:context-riemann-derivdir}
  \end{equation}
  On l'obtient en pratique en développant $f(\param+t\tvec)$ au premier ordre en $t$, sans jamais passer par des coordonnées. Nous utiliserons librement les identités usuelles $\dirderiv\log\Det(\MCov)[\tvec]=\trace{\inv{\MCov}\tvec}$ et $\dirderiv\inv{\MCov}[\tvec]=-\inv{\MCov}\tvec\inv{\MCov}$, qui s'obtiennent toutes deux directement de \cref{eq:context-riemann-derivdir}. Pour un traitement systématique du calcul différentiel matriciel, on se reportera à \sidecite[-1cm]{magnus2019matrix} et, pour tout ce qui concerne les fonctions de matrices, à \sidecite[1cm]{higham2008functions}.
\end{remark}


\section{Le cône des matrices symétriques définies positives}
\label{sec:context-riemann-spd}

\subsection{Ce que l'on attend d'une distance entre covariances}

Plutôt que de choisir une métrique et d'en subir les conséquences, procédons dans l'autre sens: écrivons d'abord ce que nous voudrions qu'une distance entre matrices de covariance vérifie, et voyons ce que cela impose.

\begin{itemize}
  \item \textbf{Invariance par congruence.} Comme nous l'avons dit en ouverture de ce chapitre, un filtrage linéaire inversible $\x\mapsto\mat{A}\x$ transforme toutes les covariances selon $\MCov\mapsto\mat{A}\MCov\mat{A}^\T$. Une distance qui prétend mesurer un écart entre modèles de données, et non entre tableaux de nombres, doit être insensible à cette opération:
  \begin{equation}
    \geodist(\mat{A}\MCov_0\mat{A}^\T, \mat{A}\MCov_1\mat{A}^\T)
    =
    \geodist(\MCov_0,\MCov_1),
    \qquad
    \forall \mat{A} \text{ inversible.}
    \label{eq:context-riemann-invcongruence}
  \end{equation}
  \item \textbf{Invariance par inversion.} On manipule indifféremment la matrice de covariance et son inverse, dite matrice de précision, qui est l'objet qui intervient réellement dans le filtre adapté et dans la forme quadratique de la \cref{eq:context-definition-gaussian}. Il serait incohérent que l'écart entre deux modèles dépende de celle des deux représentations que l'on a choisie: $\geodist(\inv{\MCov_0},\inv{\MCov_1})=\geodist(\MCov_0,\MCov_1)$.
  \item \textbf{Complétude.} Une suite de matrices dont la plus petite valeur propre tend vers zéro quitte le cône. On voudrait que le bord soit à distance infinie, de sorte qu'aucun algorithme itératif ne puisse l'atteindre en un nombre fini d'étapes. C'est la traduction géométrique du premier symptôme évoqué en ouverture.
\end{itemize}

\subsection{La métrique affine invariante}

Imposons la première de ces conditions au niveau de la métrique elle-même, en demandant que la congruence $\MCov\mapsto\mat{A}\MCov\mat{A}^\T$ soit une isométrie. Elle transporte les vecteurs tangents selon $\tvec\mapsto\mat{A}\tvec\mat{A}^\T$, et la condition s'écrit donc
$\metric{\mat{A}\tvec\mat{A}^\T}{\mat{A}\tvecbis\mat{A}^\T}{\mat{A}\MCov\mat{A}^\T}=\metric{\tvec}{\tvecbis}{\MCov}$. En prenant $\mat{A}=\MCov^{-1/2}$, tout se ramène au point $\Id{\Nvar}$: une telle métrique est entièrement déterminée par un produit scalaire sur $\Sym{\Nvar}$, évalué non pas sur le vecteur tangent $\tvec$ lui-même mais sur sa version \emph{blanchie} $\MCov^{-1/2}\tvec\MCov^{-1/2}$. Le choix canonique du produit scalaire de Frobenius conduit à la métrique suivante \sidecite[0cm]{pennec2006riemannian,bhatia2007positive}.

\begin{definition}[{Métrique affine invariante}]
\label{def:context-riemann-metriquenaturelle}
  On appelle \textbf{métrique affine invariante}, ou \textbf{métrique affine invariante naturelle} lorsqu'il faudra la distinguer de la famille de la \cref{def:context-riemann-affineinvariante}, la métrique riemannienne définie sur $\SPD{\Nvar}$, pour $\tvec,\tvecbis\in\Tangent{\MCov}{\SPD{\Nvar}}=\Sym{\Nvar}$, par
  \begin{equation}
    \metric{\tvec}{\tvecbis}{\MCov}
    \eqdef
    \trace{\inv{\MCov}\tvec\inv{\MCov}\tvecbis}
    =
    \trace{
      \left(\MCov^{-1/2}\tvec\MCov^{-1/2}\right)
      \left(\MCov^{-1/2}\tvecbis\MCov^{-1/2}\right)
    },
    \label{eq:context-riemann-metriqueaffine}
  \end{equation}
  soit, pour la norme associée,
  \begin{equation}
    \normriem{\tvec}{\MCov}
    =
    \norm{\MCov^{-1/2}\tvec\MCov^{-1/2}}_\frob.
    \label{eq:context-riemann-normeaffine}
  \end{equation}
\end{definition}

\begin{remark}[{Ce que la métrique mesure}]
\label{rem:context-riemann-lecturemetrique}
  L'écriture \cref{eq:context-riemann-normeaffine} est celle qu'il faut retenir: la métrique n'est rien d'autre que la norme de Frobenius appliquée à la perturbation \emph{relative} $\MCov^{-1/2}\tvec\MCov^{-1/2}$, c'est-à-dire à la perturbation lue dans le repère où $\MCov$ devient l'identité. Une même perturbation absolue $\tvec$ coûte donc d'autant plus cher que la puissance de $\MCov$ est faible dans la direction concernée. C'est exactement la lecture que réclame un problème d'estimation de covariance, où une erreur d'un décibel est une erreur d'un décibel quelle que soit la puissance du canal.

  Cette dépendance quadratique en $\inv{\MCov}$ règle du même coup l'exigence de complétude: la norme d'un vecteur tangent explose lorsque $\MCov$ approche du bord du cône, qui devient littéralement inatteignable, alors qu'il est à distance finie pour la métrique euclidienne. Le cas scalaire $\Nvar=1$, où $\SPD{1}=\R_*^+$ et où \cref{eq:context-riemann-metriqueaffine} se réduit à $\metric{\zeta}{\eta}{\sigma}=\zeta\eta/\sigma^2$, en donne l'image la plus simple: $0$ est à distance infinie de tout point, et la géométrie de la demi-droite est logarithmique.
\end{remark}

\subsection{La boîte à outils du cône}

La géométrie associée à \cref{eq:context-riemann-metriqueaffine} est entièrement connue en forme close, ce qui explique une bonne partie de son succès. Nous reprenons ci-dessous, un par un, les objets de la \cref{sec:context-riemann-notion}; les démonstrations, toutes élémentaires une fois le blanchiment $\MCov\mapsto\Id{\Nvar}$ effectué, se trouvent dans \sidecite[0cm]{bhatia2007positive,pennec2006riemannian}.

\begin{property}[{Espace tangent et projection}]
\label{prop:context-riemann-projspd}
  Le cône étant un ouvert de $\Sym{\Nvar}$, l'espace tangent est $\Tangent{\MCov}{\SPD{\Nvar}}=\Sym{\Nvar}$ en tout point, et la projection orthogonale de $\mat{Z}\in\R^{\Nvar\times\Nvar}$ sur celui-ci est la simple symétrisation
  \begin{equation}
    \projman_{\MCov}(\mat{Z})=\symop(\mat{Z})\eqdef\tfrac{1}{2}(\mat{Z}+\mat{Z}^\T).
    \label{eq:context-riemann-projspd}
  \end{equation}
\end{property}

\begin{property}[{Connexion et géodésiques}]
\label{prop:context-riemann-geodspd}
  La connexion de Levi-Civita associée à \cref{eq:context-riemann-metriqueaffine} est :
  \begin{equation}
    \connection_{\tvec}\tvecbis = \dirderiv\tvecbis[\tvec] - \symop\!\left(\tvecbis\inv{\MCov}\tvec\right),
    \label{eq:context-riemann-connexionspd}
  \end{equation}
  et la géodésique joignant $\MCov_0$ à $\MCov_1$, que l'on note $\MCov_0\geodmean_t\MCov_1$, s'écrit
  \begin{equation}
    \geod(t)
    =
    \MCov_0\geodmean_t\MCov_1
    =
    \MCov_0^{1/2}
    \left(
      \MCov_0^{-1/2}\MCov_1\MCov_0^{-1/2}
    \right)^{t}
    \MCov_0^{1/2}.
    \label{eq:context-riemann-geodspd}
  \end{equation}
\end{property}

\begin{property}[{Exponentielle et logarithme}]
\label{prop:context-riemann-explogspd}
  L'exponentielle et le logarithme riemanniens sont
  \begin{equation}
    \exp_{\MCov}(\tvec)
    =
    \MCov^{1/2}\expm\!\left(\MCov^{-1/2}\tvec\MCov^{-1/2}\right)\MCov^{1/2}, 
  \end{equation}
  \begin{equation}
    \log_{\MCov_0}(\MCov_1)
    =
    \MCov_0^{1/2}\logm\!\left(\MCov_0^{-1/2}\MCov_1\MCov_0^{-1/2}\right)\MCov_0^{1/2},
    \label{eq:context-riemann-explogspd}
  \end{equation}
  où $\expm$ et $\logm$ désignent l'exponentielle et le logarithme de matrice. Toutes deux sont définies globalement: l'exponentielle est un difféomorphisme de $\Sym{\Nvar}$ sur $\SPD{\Nvar}$, de sorte qu'il n'y a ici ni rayon d'injectivité ni point de coupure à surveiller.
\end{property}

\begin{property}[{Transport parallèle}]
\label{prop:context-riemann-transportspd}
  Le transport parallèle de $\tvec\in\Tangent{\MCov_0}{\SPD{\Nvar}}$ vers $\Tangent{\MCov_1}{\SPD{\Nvar}}$ le long de la géodésique \cref{eq:context-riemann-geodspd} est la congruence \sidecite[0cm]{thanwerdas2023invariant}
  \begin{equation}
    \transport_{\MCov_0\to\MCov_1}(\tvec)=\mat{E}\tvec\mat{E}^\T,
    \qquad
    \mat{E}=(\MCov_1\inv{\MCov_0})^{1/2}.
    \label{eq:context-riemann-transportspd}
  \end{equation}
\end{property}

\begin{property}[{Distance géodésique}]
\label{prop:context-riemann-distspd}
  La distance associée à \cref{eq:context-riemann-metriqueaffine} est la norme de Frobenius du logarithme du rapport des deux matrices:
  \begin{equation}
    \geodist(\MCov_0,\MCov_1)
    =
    \norm{\logm\!\left(\MCov_0^{-1/2}\MCov_1\MCov_0^{-1/2}\right)}_\frob
    =
    \left(
      \sum_{\ivar=1}^{\Nvar}\log^2\eig_\ivar\!\left(\inv{\MCov_0}\MCov_1\right)
    \right)^{1/2},
    \label{eq:context-riemann-distspd}
  \end{equation}
  où les $\eig_\ivar(\inv{\MCov_0}\MCov_1)$ sont les \textbf{valeurs propres généralisées} du couple $(\MCov_0,\MCov_1)$, c'est-à-dire les solutions de $\Det(\MCov_1-\eig\MCov_0)=0$. Elle vérifie les trois exigences de la \cref{sec:context-riemann-spd}: invariance par congruence, invariance par inversion, et complétude.
\end{property}

\begin{remark}[{Lecture de la distance}]
\label{rem:context-riemann-lecturedistance}
  La seconde écriture de \cref{eq:context-riemann-distspd} dit exactement ce que la distance mesure, et pourquoi elle est la bonne pour nos problèmes. Les $\eig_\ivar$ sont les rapports de puissance entre les deux modèles dans les directions propres relatives, et valent tous $1$ si et seulement si $\MCov_0=\MCov_1$. La distance ne dépend donc que de \emph{rapports}: l'invariance par congruence en devient évidente --- une congruence ne change pas le spectre de $\inv{\MCov_0}\MCov_1$ --- et l'invariance par inversion tout autant, puisqu'elle change $\eig_\ivar$ en $\inv{\eig_\ivar}$ et laisse $\log^2\eig_\ivar$ inchangé. On notera que ces quantités sont exactement celles qui gouvernent les pertes de rapport signal-à-bruit après blanchiment, ce qui explique qu'elles interviennent également dans les statistiques de détection du \cref{ch:robust-detection}.
\end{remark}

\subsection{Une famille à deux paramètres}
\label{subsec:context-riemann-famille}

L'argument d'invariance de la \cref{eq:context-riemann-metriqueaffine} ne détermine pas la métrique de manière unique. Une fois ramené au point $\Id{\Nvar}$, il ne reste plus qu'à choisir un produit scalaire sur $\Sym{\Nvar}$ invariant par similitude orthogonale $\tvec\mapsto\Orth\tvec\Orth^\T$; or il en existe non pas un mais une famille à deux paramètres, obtenue en pondérant séparément la partie de trace nulle et la partie scalaire \sidecite[0cm]{thanwerdas2023invariant}. La \cref{def:context-riemann-metriquenaturelle} correspond au choix le plus simple de cette pondération.

\begin{definition}[{Famille des métriques affines invariantes}]
\label{def:context-riemann-affineinvariante}
  Soient $\alpha\in\R_*^+$ et $\beta>-\alpha/\Nvar$. La métrique
  \begin{equation}
    \metric{\tvec}{\tvecbis}{\MCov}^{\alpha,\beta}
    \eqdef
    \alpha\,\trace{\inv{\MCov}\tvec\inv{\MCov}\tvecbis}
    +
    \beta\,\trace{\inv{\MCov}\tvec}\trace{\inv{\MCov}\tvecbis}
    \label{eq:context-riemann-metriqueaffineab}
  \end{equation}
  est invariante par congruence, la condition $\beta>-\alpha/\Nvar$ assurant qu'elle est définie positive. Le cas $(\alpha,\beta)=(1,0)$ redonne la \cref{def:context-riemann-metriquenaturelle}.
\end{definition}

\begin{remark}[{Les paramètres n'agissent que sur la distance}]
\label{rem:context-riemann-alphabeta}
  Un calcul direct montre que la connexion de Levi-Civita \cref{eq:context-riemann-connexionspd} ne fait intervenir ni $\alpha$ ni $\beta$. Il en va donc de même des géodésiques, de l'exponentielle, du logarithme et du transport parallèle: \emph{toutes les métriques de la famille \cref{eq:context-riemann-metriqueaffineab} partagent les mêmes droites}, et toute la boîte à outils qui précède reste valable telle quelle. Seule la distance, qui intègre la métrique le long de ces droites, en dépend:
  \begin{equation}
    \geodist_{\alpha,\beta}^2(\MCov_0,\MCov_1)
    =
    \alpha\,\norm{\logm\!\left(\inv{\MCov_0}\MCov_1\right)}_\frob^2
    +
    \beta\,\left(\log\Det\!\left(\inv{\MCov_0}\MCov_1\right)\right)^2.
    \label{eq:context-riemann-distspdab}
  \end{equation}
  Nous n'aurions aucune raison de traîner ces deux paramètres si le seul argument d'invariance était en jeu, et toute la suite de ce chapitre s'en tiendra au cas $(1,0)$. Ils ne réapparaîtront qu'en \cref{sec:context-geometry-elliptical}, où le modèle probabiliste --- et non plus l'invariance --- viendra sélectionner un point précis de cette famille; le découplage ci-dessus signifie alors que changer de modèle ne change aucune des constructions géométriques, seulement la manière de mesurer.
\end{remark}

\subsection{Sous-variétés et retour sur la normalisation}

Nous avions laissé en suspens, en \cref{subsec:context-distributions-elliptical}, la question du choix de normalisation de la matrice de forme. La géométrie tranche.

\begin{property}[{La contrainte de déterminant unité}]
\label{prop:context-riemann-detun}
  L'ensemble des matrices de forme normalisées par le déterminant,
  \begin{equation}
    \SPDdet{\Nvar}
    \eqdef
    \left\{ \MShape\in\SPD{\Nvar};\; \Det(\MShape)=1 \right\},
    \label{eq:context-riemann-spddet}
  \end{equation}
  est une sous-variété de $\SPD{\Nvar}$ de dimension $\Nvar(\Nvar+1)/2-1$, dont l'espace tangent est $\Tangent{\MShape}{\SPDdet{\Nvar}}=\{\tvec\in\Sym{\Nvar};\,\trace{\inv{\MShape}\tvec}=0\}$ et la projection orthogonale
  \begin{equation}
    \projman_{\MShape}^{\SPDdet{\Nvar}}(\mat{Z})
    =
    \symop(\mat{Z}) - \frac{1}{\Nvar}\trace{\inv{\MShape}\mat{Z}}\,\MShape.
    \label{eq:context-riemann-projdet}
  \end{equation}
  Elle est de plus \textbf{totalement géodésique}: toute géodésique de $\SPD{\Nvar}$ issue d'un point de $\SPDdet{\Nvar}$ dans une direction tangente à $\SPDdet{\Nvar}$ y reste. Géodésiques, exponentielle, logarithme et distance sont donc ceux des \cref{prop:context-riemann-geodspd,prop:context-riemann-explogspd,prop:context-riemann-distspd}, sans modification.
\end{property}

\begin{remark}[{Pourquoi le déterminant et pas la trace}]
\label{rem:context-riemann-normalisation}
  Le caractère totalement géodésique se lit directement sur \cref{eq:context-riemann-geodspd}: le déterminant de $\MCov_0\geodmean_t\MCov_1$ vaut $\Det(\MCov_0)^{1-t}\Det(\MCov_1)^{t}$, qui vaut $1$ pour tout $t$ dès que les deux extrémités sont de déterminant unité. Rien de tel pour la contrainte $\trace{\MShape}=\Nvar$, qui n'est pas préservée le long des géodésiques: la sous-variété correspondante est courbe dans $\SPD{\Nvar}$, et il faudrait lui construire une géométrie propre. C'est ce qui justifie l'annonce faite en \cref{subsec:context-distributions-elliptical}: nous conserverons la normalisation par la trace au \cref{ch:robust-detection} par cohérence avec la littérature de détection, mais toute la partie géométrique de ce mémoire utilisera la normalisation par le déterminant. Nous verrons en \cref{sec:context-geometry-compound-gaussian} que ce choix a une seconde vertu, statistique celle-là.
\end{remark}

\begin{remark}[{Le cône positif et la partie échelle}]
\label{rem:context-riemann-rplus}
  Le cas $\Nvar=1$, soit $\SPD{1}=\R_*^+$, mérite d'être noté à part car il resservira. La métrique naturelle y est $\metric{\zeta}{\eta}{\sigma}=\zeta\eta/\sigma^2$, la géodésique $\sigma_0\geodmean_t\sigma_1=\sigma_0^{1-t}\sigma_1^{t}$, et la distance $\geodist(\sigma_0,\sigma_1)=\abs{\log(\sigma_1/\sigma_0)}$. Autrement dit, la géométrie naturelle d'un paramètre de puissance est logarithmique, ce qui est cohérent avec l'usage universel des décibels. Par extension, $\Rpos{\Nobs}$ muni de la métrique produit s'identifie à la sous-variété des matrices diagonales définies positives de $\SPD{\Nobs}$, et hérite donc de toute la géométrie ci-dessus; sa distance est $\geodist^2(\vector{\texture}_0,\vector{\texture}_1)=\sum_\iobs\log^2(\texture_{1,\iobs}/\texture_{0,\iobs})$. Ce sera le point de départ de la \cref{sec:context-geometry-compound-gaussian}.
\end{remark}

\subsection{La métrique log-euclidienne}

L'invariance par congruence n'est pas le seul cahier des charges raisonnable. On peut aussi partir d'une exigence de \emph{simplicité de calcul}, et l'on aboutit alors à une géométrie différente, tout aussi utile en pratique.

L'idée est la suivante: l'application $\logm$ réalise une bijection lisse entre $\SPD{\Nvar}$ et l'espace vectoriel $\Sym{\Nvar}$ tout entier. Pourquoi ne pas simplement transporter la géométrie euclidienne de $\Sym{\Nvar}$ sur le cône par cette bijection ?

\begin{definition}[{Métrique log-euclidienne}]
\label{def:context-riemann-logeuclidienne}
  On appelle \textbf{métrique log-euclidienne} \sidecite[0cm]{arsigny2006logeuclidean} la métrique riemannienne obtenue en tirant en arrière le produit scalaire de Frobenius par le logarithme de matrice:
  \begin{equation}
    \metric{\tvec}{\tvecbis}{\MCov}^{\mathrm{LE}}
    \eqdef
    \trace{
      \dirderiv\logm(\MCov)[\tvec]\,
      \dirderiv\logm(\MCov)[\tvecbis]
    }.
    \label{eq:context-riemann-metriqueLE}
  \end{equation}
\end{definition}

Par construction, $(\SPD{\Nvar},\metric{\cdot}{\cdot}{\cdot}^{\mathrm{LE}})$ est isométrique à l'espace euclidien $\Sym{\Nvar}$. Tout s'en déduit immédiatement: les géodésiques sont les images des segments,
\begin{equation}
  \geod^{\mathrm{LE}}(t) = \expm\!\left( (1-t)\logm\MCov_0 + t\logm\MCov_1 \right),\\
  \geodist_{\mathrm{LE}}(\MCov_0,\MCov_1) = \norm{\logm\MCov_0 - \logm\MCov_1}_\frob,
  \label{eq:context-riemann-geodLE}
\end{equation}
la courbure est nulle, et le transport parallèle est trivial dans les coordonnées logarithmiques.

\begin{remark}[{Ce que l'on gagne et ce que l'on perd}]
\label{rem:context-riemann-LEbilan}
  Cette métrique conserve la complétude et l'absence d'effet de gonflement --- puisque $\log\Det$ est linéaire le long de \cref{eq:context-riemann-geodLE} --- ainsi que l'invariance par similitude orthogonale $\MCov\mapsto\Orth\MCov\Orth^\T$ et par inversion. En revanche elle perd l'invariance par congruence: $\logm(\mat{A}\MCov\mat{A}^\T)\ne\mat{A}\logm(\MCov)\mat{A}^\T$ en général. Ce n'est pas anodin, puisque c'était notre exigence de départ.

  Ce qu'elle apporte en échange est de nature algorithmique, et c'est décisif à grande échelle: une fois payée une décomposition spectrale par matrice, toutes les opérations ultérieures --- distances, moyennes, barycentres, produits scalaires --- sont euclidiennes, donc de coût linéaire en le nombre de matrices. C'est la raison pour laquelle cette métrique est celle qui structure les architectures d'apprentissage profond sur variété \textsc{spd} \sidecite[1cm]{huang2017riemannian}, où la couche dite \textit{LogEig} n'est rien d'autre que le plongement $\MCov\mapsto\logm\MCov$ de \cref{def:context-riemann-logeuclidienne}. Nous y reviendrons au \cref{ch:spdnet}.
\end{remark}

\subsection{Les autres métriques, et un bilan}

D'autres géométries ont été proposées sur $\SPD{\Nvar}$, que nous mentionnons sans les développer. La métrique de \textbf{Bures-Wasserstein} \sidecite[0cm]{bhatia2019bures} est celle induite par le transport optimal entre lois gaussiennes centrées; sa distance est $\geodist^2_{\mathrm{BW}}(\MCov_0,\MCov_1)=\trace{\MCov_0}+\trace{\MCov_1}-2\trace{(\MCov_0^{1/2}\MCov_1\MCov_0^{1/2})^{1/2}}$, elle est de courbure positive et n'est pas complète, le bord du cône restant à distance finie. La métrique \textbf{log-Cholesky} exploite de la même manière que la log-euclidienne un difféomorphisme vers un espace vectoriel, en passant cette fois par la factorisation de Cholesky, avec un coût numérique moindre. Enfin, un certain nombre de \textbf{divergences} --- Stein, Jensen-Bregman, Kullback-Leibler --- sont utilisées comme substituts bon marché de la distance affine invariante, sans être issues d'une métrique riemannienne; nous y reviendrons en \cref{sec:context-geometry-divergences}. Une classification complète des métriques invariantes sur le cône est donnée dans \sidecite[1cm]{thanwerdas2023invariant}.

Le \cref{tab:context-riemann-metriques} rassemble les propriétés qui gouvernent le choix pratique, et la \cref{fig:context-riemann-interpolation} met en image l'écart entre elles sur un exemple concret.

\begin{table}
  \centering
  \footnotesize
  \begin{tabular}{@{}lccccc@{}}
    \toprule
    Métrique & Géod. & Dist. & Moyenne & Inv. congr. & Complète \\
    \midrule
    Euclidienne          & oui & oui & arithm.   & non & non \\
    Affine invariante    & oui & oui & itérative & oui & oui \\
    Log-euclidienne      & oui & oui & close     & non & oui \\
    Bures-Wasserstein    & oui & oui & itérative & non & non \\
    \bottomrule
  \end{tabular}
  \vspace{0.5em}
  \caption{Comparaison des principales géométries sur $\SPD{\Nvar}$. Les colonnes «~géod.~» et «~dist.~» indiquent la disponibilité d'une forme close, la colonne «~moyenne~» le statut de la moyenne de Fréchet introduite en \cref{sec:context-riemann-moyennes}. La courbure sectionnelle est nulle pour les métriques euclidienne et log-euclidienne, négative pour l'affine invariante et positive pour Bures-Wasserstein. L'invariance par inversion, vérifiée par les trois premières lignes, n'est pas reportée.}
  \label{tab:context-riemann-metriques}
\end{table}

\begin{figure}[h]
  \begin{center}
    % This file was created with matplot2tikz v0.5.3.
\begin{tikzpicture}

\definecolor{crimson2143940}{RGB}{214,39,40}
\definecolor{darkgray176}{RGB}{176,176,176}
\definecolor{darkorange25512714}{RGB}{255,127,14}
\definecolor{forestgreen4416044}{RGB}{44,160,44}
\definecolor{gray127}{RGB}{127,127,127}

\begin{groupplot}[group style={group size=2 by 2, horizontal sep=1.4cm, vertical sep=1.5cm}]
\nextgroupplot[
height=4.6cm,
label style={font=\footnotesize},
legend style={draw=none, fill=white, fill opacity=0.9, text opacity=1, font=\footnotesize},
tick align=outside,
tick label style={font=\scriptsize},
tick pos=left,
title style={font=\footnotesize, yshift=-2pt},
title={euclidienne},
width=0.45\textwidth,
x grid style={darkgray176},
xmin=-2.29992834785656, xmax=2.29992834785656,
xtick style={color=black},
y grid style={darkgray176},
ylabel={\(\displaystyle x_2\)},
ymin=-2.29992834785656, ymax=2.29992834785656,
ytick style={color=black}
]
\addplot [darkorange25512714]
table {%
0.0770027631002538 -0.84691612780764
0.0179228597895658 -0.851862151073466
-0.0411749094175406 -0.855959019621099
-0.100231634579482 -0.859202649602875
-0.159188446668272 -0.861589807697886
-0.217986576251448 -0.863118114335028
-0.27656741207472 -0.863786046065008
-0.334872559486944 -0.863592937078952
-0.392843898649179 -0.862538979872096
-0.450423642469809 -0.860625225051906
-0.507554394207967 -0.857853580290806
-0.564179204687863 -0.854226808424571
-0.620241629066952 -0.84974852469827
-0.675685783101388 -0.844423193162514
-0.730456398852649 -0.838256122223592
-0.784498879779822 -0.831253459351934
-0.837759355162621 -0.823422184954183
-0.890184733800895 -0.814770105414971
-0.941722756937088 -0.805305845315343
-0.992322050348907 -0.795038838835594
-1.04193217556026 -0.783979320351064
-1.09050368011944 -0.7721383142303
-1.13798814689438 -0.759527623845714
-1.18433824233593 -0.746159819807737
-1.22950776366094 -0.73204822743415
-1.27345168490822 -0.717206913467127
-1.3161262018214 -0.701650672051193
-1.35748877551395 -0.685395009986104
-1.39749817487288 -0.668456131269327
-1.43611451765879 -0.65085092094355
-1.47329931026137 -0.632596928265297
-1.50901548607069 -0.613712349211451
-1.54322744242594 -0.594216008341108
-1.57590107610502 -0.574127340030847
-1.60700381731928 -0.553466369102122
-1.63650466217987 -0.532253690860081
-1.66437420360308 -0.510510450563723
-1.69058466062389 -0.488258322347839
-1.71510990608875 -0.465519487617766
-1.73792549269961 -0.44231661293848
-1.75900867738355 -0.418672827440069
-1.77833844396356 -0.394611699762115
-1.79589552410794 -0.370157214559959
-1.81166241653737 -0.34533374859628
-1.8256234044706 -0.320166046441802
-1.83776457129123 -0.29467919580937
-1.84807381442013 -0.268898602545975
-1.85654085737955 -0.242849965307649
-1.86315726003693 -0.216559249942487
-1.8679164270182 -0.190052663607328
-1.87081361428223 -0.163356628643885
-1.87184593384974 -0.136497756240387
-1.87101235668213 -0.109502819904963
-1.86831371370725 -0.0823987287772281
-1.86375269499112 -0.0552125008046694
-1.85733384705636 -0.0279712358105662
-1.84906356835022 -0.000702088480298374
-1.83895010286635 0.0265677587070353
-1.82700353192709 0.0538111225747037
-1.81323576413417 0.0810008463451237
-1.79766052349794 0.108109826710329
-1.780293335757 0.135111040849134
-1.76115151290177 0.161977573364073
-1.74025413591756 0.188682643111251
-1.71762203576424 0.215199629896371
-1.69327777261141 0.241502101010323
-1.66724561335003 0.267563837577884
-1.63955150740259 0.293358860693263
-1.61022306085617 0.318861457316436
-1.57928950894411 0.344046205904467
-1.5467816869037 0.368888001752256
-1.51273199923892 0.39336208201745
-1.47717438741892 0.417444050404588
-1.44014429604447 0.441109901483854
-1.40167863751597 0.464336044620215
-1.36181575523837 0.487099327489081
-1.32059538539965 0.509377059155038
-1.27805861736092 0.531147032690675
-1.23424785269762 0.552387547312924
-1.18920676293277 0.57307743001488
-1.14298024600416 0.593196056671508
-1.09561438150921 0.612723372598224
-1.04715638477172 0.631639912541839
-0.997654559776701 0.649926820083947
-0.947158251019851 0.667565866437413
-0.895717794319916 0.684539468617223
-0.843384466642844 0.70083070696759
-0.79021043498779 0.716423342027828
-0.736248704385918 0.731301830720207
-0.681553065063832 0.745451341843627
-0.626178038824313 0.75885777085768
-0.570178824697791 0.771507753942367
-0.513611243918759 0.783388681319441
-0.45653168428195 0.794488709822112
-0.398997043933756 0.804796774700568
-0.341064674654921 0.814302600651561
-0.282792324691033 0.822996712061055
-0.224238081187824 0.83087044244972
-0.165460312288635 0.83791594311187
-0.106517608951783 0.844126190939221
-0.0474687265458208 0.84949499542168
0.0116274737190959 0.854017004818182
0.0707120834653392 0.857687711491416
0.129726205868962 0.86050345640114
0.188611014369336 0.862461432751594
0.247307811308738 0.86355968878937
0.305758086443464 0.863797129748964
0.363903575268102 0.863173518944065
0.421686317094864 0.861689478003485
0.479048712830059 0.859346486251508
0.535933582390121 0.856146879233266
0.592284221699958 0.85209384638662
0.648044459216813 0.847191427862855
0.703158711923275 0.841444510499371
0.757572040733641 0.834858822948379
0.81123020525839 0.827440929966452
0.864079717872185 0.819198225870637
0.9160678970315 0.810138927167636
0.967142919788727 0.800272064363414
1.01725387345042 0.789607472961396
1.06635080632817 0.778155783658222
1.11438477753156 0.765928411746835
1.16130790575344 0.752937545737474
1.20707341699915 0.739196135207898
1.25163569121178 0.724717877894966
1.29495030774727 0.709517206040435
1.33697408965387 0.69360927200459
1.37766514671185 0.677009933162038
1.41698291719062 0.659735736094732
1.45488820828147 0.641803900097984
1.4913432351659 0.623232300015891
1.52631165868031 0.604039448423301
1.55975862153965 0.584244477172084
1.59165078308395 0.563867118320076
1.621956352513 0.542927684461739
1.65064512057613 0.52144704848013
1.67768848968541 0.499446622740345
1.70305950242238 0.476948337745222
1.72673286840979 0.453974620274527
1.74868498952159 0.43054837102946
1.76889398340604 0.406692941804736
1.78733970529853 0.382432112211003
1.80400376810235 0.357790065970818
1.81886956071727 0.332791366811779
1.83192226459792 0.30746093398087
1.84314886852522 0.281824017404417
1.8525381815762 0.255906172518404
1.86008084427944 0.229733234794266
1.86576933794471 0.203331293985526
1.86959799215781 0.176726668120964
1.87156299043296 0.149945877270228
1.87166237401714 0.123015617108058
1.86989604384265 0.095962732303455
1.8662657606258 0.0688141897603267
1.86077514311185 0.0415970517362991
1.85342966446777 0.0143384488664617
1.8442366468264 -0.0129344468810334
1.83320525398765 -0.0401944492905829
1.82034648228382 -0.0674143849989453
1.80567314961819 -0.0945671205822355
1.78919988268786 -0.121625589603117
1.77094310240357 -0.148562819591211
1.75092100752096 -0.175351958929844
1.72915355649962 -0.201966303622322
1.70566244760814 -0.228379323911056
1.68047109729472 -0.254564690723002
1.65360461684522 -0.280496301915051
1.62508978735163 -0.306148308293214
1.59495503301612 -0.331495139379664
1.56323039281717 -0.356511528901939
1.52994749056607 -0.381172539978916
1.49513950338361 -0.405453589978434
1.4588411286284 -0.429330475021797
1.4210885493098 -0.45277939411073
1.38191939801992 -0.475776972852723
1.34137271942058 -0.498300286761145
1.29948893132282 -0.520326884106861
1.25630978439748 -0.541834808298612
1.21187832055726 -0.562802619769816
1.16623883005161 -0.583209417349996
1.1194368073172 -0.603034859099527
1.07151890562813 -0.622259182586915
1.02253289059092 -0.640863224588424
0.972527592530673 -0.658828440190393
0.921552857816001 -0.676136921275208
0.869659499171045 -0.692771414372505
0.816899245024265 -0.708715337857803
0.763324687944402 -0.72395279848143
0.708989232215077 -0.738468607211257
0.653947040600229 -0.752248294373457
0.598252980353496 -0.765278124076191
0.541962568525351 -0.777545107901843
0.485131916622485 -0.789037017854158
0.427817674674649 -0.799742398547373
0.370076974764668 -0.809650578625196
0.311967374077936 -0.818751681398246
0.253546797528166 -0.827036634689351
0.19487348001657 -0.834497179876893
0.136005908382054 -0.841125880127179
0.0770027631002543 -0.84691612780764
};
\addplot [darkorange25512714]
table {%
0.270700171232457 -1.01711070684762
0.216488448730679 -1.03099610447947
0.162060925835092 -1.04385378277262
0.107471857084187 -1.05567092491373
0.0527756580488805 -1.06643575131853
-0.00197314891011602 -1.07613753137394
-0.0567199889913676 -1.08476659413461
-0.111410289354066 -1.09231433796312
-0.165989533517424 -1.09877323910427
-0.220403315703889 -1.10413685918495
-0.274597395072006 -1.108399851632
-0.328517749784866 -1.11155796700186
-0.382110630860251 -1.11360805721643
-0.435322615748785 -1.11454807870121
-0.488100661586696 -1.11437709442233
-0.540392158070094 -1.11309527482065
-0.592144979898066 -1.11070389764182
-0.643307538732313 -1.10720534666261
-0.693828834621523 -1.10260310931474
-0.743658506839234 -1.09690177320847
-0.7927468840845 -1.09010702155963
-0.841045033995328 -1.08222562752443
-0.888504811925513 -1.07326544744784
-0.935078908936278 -1.06323541303224
-0.980720898954844 -1.05214552243405
-1.02538528505295 -1.0400068302974
-1.0690275447992 -1.02683143673461
-1.11160417463995 -1.0126324752645
-1.15307273326464 -0.997424099720676
-1.19339188391219 -0.981221470142615
-1.23252143557638 -0.964040737663851
-1.27042238306914 -0.945899028412162
-1.30705694590175 -0.926814426437868
-1.34238860594526 -0.906805955687263
-1.37638214383257 -0.885893561039128
-1.40900367406583 -0.864098088423244
-1.44022067879426 -0.841441264040714
-1.47000204022864 -0.817945672706817
-1.49831807166023 -0.793634735337967
-1.52514054705309 -0.76853268560524
-1.55044272918042 -0.74266454577772
-1.57419939627682 -0.716056101779759
-1.59638686717986 -0.688733877487004
-1.61698302493603 -0.660725108286824
-1.63596733884737 -0.632057713929485
-1.65332088493688 -0.602760270697129
-1.66902636481245 -0.572861982918322
-1.68306812291017 -0.542392653856537
-1.69543216210018 -0.511382656001622
-1.70610615763935 -0.479862900793833
-1.71507946945678 -0.44786480781064
-1.7223431527601 -0.415420273447014
-1.72788996695183 -0.382561639120399
-1.73171438284695 -0.349321659032092
-1.73381258818453 -0.315733467517136
-1.73418249142785 -0.281830546015297
-1.73282372384932 -0.247646689696037
-1.72973763989801 -0.213215973770754
-1.72492731584951 -0.17857271952587
-1.71839754673946 -0.143751460110624
-1.7101548415837 -0.108786906113684
-1.70020741689 -0.0737139109628693
-1.68856518846762 -0.0385674361825026
-1.67523976154305 -0.00338251654299782
-1.66024441919164 0.0318057748625586
-1.64359410909676 0.0669623615800311
-1.62530542864962 0.102052198759206
-1.60539660840463 0.137040308087263
-1.58388749390677 0.171891812655913
-1.56079952590915 0.206571971727465
-1.53615571900038 0.241046215365143
-1.50998063866312 0.275280178893155
-1.48230037678672 0.30923973715215
-1.45314252565818 0.342891038515912
-1.42253615045762 0.376200538635407
-1.39051176028547 0.409135033876511
-1.35710127775029 0.441661694418121
-1.32233800714769 0.47374809697763
-1.28625660126182 0.505362257131163
-1.24889302682275 0.536472661196344
-1.21028452865405 0.567048297645824
-1.17046959254633 0.597058688020242
-1.12948790689381 0.626473917309816
-1.08738032313198 0.655264663774278
-1.04418881501608 0.683402228171422
-0.999956436780722 0.710858562365131
-0.954727280222474 0.737606297284369
-0.908546430748223 0.763618770205268
-0.861459922433051 0.788870051329112
-0.813514692132465 0.813334969629719
-0.764758532694728 0.836989137944472
-0.715240045319898 0.859808977283974
-0.665008591113101 0.881771740336092
-0.614114241880306 0.902855534140978
-0.562607730215662 0.923039341914444
-0.510540398930142 0.942303043997951
-0.457964149871913 0.960627437914321
-0.404931392189447 0.977994257509187
-0.351494990088939 0.994386191159088
-0.297708210138119 1.00978689902808
-0.243624668168979 1.02418102935562
-0.189298275832344 1.03755423375957
-0.134783186857566 1.04989318153893
-0.0801337430709103 1.06118557296223
-0.02540442022644 1.07142015152811
0.0293502262965976 1.08058671518606
0.0840756158757643 1.0886761265061
0.138717197052579 1.09568032178708
0.193220501910598 1.10159231909484
0.24753120037022 1.10640622522187
0.301595154346092 1.11011724156186
0.355358471713123 1.11272166889302
0.408767560027339 1.11421691106554
0.461769179947985 1.11460147758951
0.514310498307663 1.11387498512066
0.566339140777582 1.11203815784247
0.617803244075427 1.10909282674432
0.668651507663807 1.10504192779631
0.718833244887748 1.0998894990226
0.768298433500249 1.09364067647625
0.816997765525549 1.08630168911948
0.86488269641039 1.07787985261448
0.911905493414281 1.06838356203103
0.958019283190533 1.05782228347812
1.00317809851063 1.04620654466789
1.04733692408536 1.03354792442139
1.09045174143705 1.01985904112659
1.1324795727781 1.00515354016006
1.17337852385224 0.989446080284962
1.21310782569558 0.972752319038908
1.25162787527601 0.955088897126158
1.2889002749704 0.93647342182981
1.32488787084014 0.916924449460485
1.35955478966703 0.896461466858982
1.39286647471246 0.875104871971382
1.42478972016431 0.852875953515925
1.45529270423727 0.829796869761953
1.48434502089345 0.805890626442069
1.51191771015179 0.781181053819517
1.53798328695598 0.755692782933664
1.56251576857221 0.72945122104724
1.58549070048928 0.702482526319823
1.6068851807954 0.674813581732825
1.62667788300732 0.646471968291945
1.64484907732903 0.617485937533823
1.66138065031886 0.587884383364298
1.67625612294532 0.55769681325633
1.68946066701379 0.526953318836313
1.70098111994758 0.495684545888081
1.71080599790866 0.463921663804533
1.71892550724501 0.431696334517297
1.72533155425314 0.399040680935432
1.73001775324608 0.365987254924614
1.73297943291877 0.332569004858722
1.73421364100448 0.298819242776193
1.73371914721777 0.264771611173846
1.73149644448078 0.230460049471321
1.72754774843196 0.195918760179518
1.72187699521739 0.161182174806795
1.71448983756721 0.126284919536891
1.70539363916081 0.0912617807127907
1.69459746728658 0.0561476701609467
1.68211208380345 0.0209775903904073
1.66794993441319 -0.0142134002984475
1.65212513625432 -0.0493902227607758
1.63465346382978 -0.0845178119749357
1.61555233328263 -0.119561151996055
1.59484078503518 -0.154485310860689
1.57253946480917 -0.189255475407766
1.5486706030456 -0.223836985981092
1.52325799274496 -0.258195370978865
1.49632696574986 -0.292296381215703
1.46790436749371 -0.326106024062984
1.43801853024057 -0.359590597333439
1.40669924484295 -0.392716722876221
1.37397773104567 -0.425451379848973
1.3398866063653 -0.457761937633722
1.30445985357631 -0.489616188363791
1.26773278683633 -0.520982379029297
1.22974201648423 -0.551829243129241
1.19052541254609 -0.582126031838645
1.15012206698559 -0.611842544659636
1.10857225473627 -0.640949159525972
1.06591739355458 -0.669416862330954
1.02220000273379 -0.697217275849317
0.977463660719833 -0.724322688024276
0.931752961671327 -0.750706079591494
0.885113471007187 -0.77634115101248
0.837591679986007 -0.80120234869054
0.789234959362565 -0.825264890443154
0.740091512167622 -0.848504790205401
0.690210325658098 -0.870898881939783
0.639641122485503 -0.892424842728642
0.588434311131321 -0.913061215026124
0.53664093565874 -0.932787428047529
0.484312624830812 -0.951583818274721
0.431501540645792 -0.969431649057141
0.378260326340909 -0.986313129288919
0.324642053916456 -1.00221143114343
0.270700171232457 -1.01711070684762
};
\addplot [darkorange25512714]
table {%
0.589698790492918 -0.962300966163674
0.542968983004148 -0.990277298074362
0.49569793222482 -1.01726650003812
0.447932758966384 -1.04324166863149
0.399721076592324 -1.0681769112415
0.351110943556121 -1.09204737187605
0.302150815495544 -1.11482925594085
0.252889496931009 -1.13649985395841
0.20337609261618 -1.15703756420533
0.15365995858928 -1.17642191424541
0.103790652973914 -1.19463358133695
0.0538178865784602 -1.21165441169414
0.00379147334325258 -1.22746743858311
-0.0462387193150393 -1.24205689923477
-0.0962228202123341 -1.25540825055745
-0.146111004109821 -1.26750818363379
-0.19585354138084 -1.27834463698735
-0.245400847582456 -1.28790680860578
-0.294703532882309 -1.29618516670843
-0.343712451291475 -1.30317145924793
-0.392378749654254 -1.30885872213592
-0.440653916346062 -1.31324128618509
-0.488489829630868 -1.31631478276028
-0.53583880562999 -1.31807614813325
-0.582653645854425 -1.3185236265367
-0.628887684253322 -1.31765677191442
-0.674494833731726 -1.31547644836596
-0.719429632091186 -1.31198482928524
-0.763647287347472 -1.30718539519411
-0.807103722380195 -1.30108293027282
-0.849755618869837 -1.29368351759115
-0.8915604604784 -1.28499453304455
-0.932476575230612 -1.27502463800181
-0.972463177053457 -1.26378377067111
-1.01148040643262 -1.25128313619348
-1.04948937014533 -1.23753519547322
-1.08645218002995 -1.22255365275662
-1.12233199075373 -1.20635344197129
-1.15709303654105 -1.18895071183964
-1.19070066682548 -1.17036280978155
-1.22312138079027 -1.15060826462202
-1.25432286076268 -1.12970676812129
-1.28427400442889 -1.10767915534562
-1.31294495583755 -1.08454738389848
-1.34030713516076 -1.06033451203272
-1.36633326718311 -1.03506467566557
-1.39099740849019 -1.00876306431949
-1.41427497332955 -0.981455896012674
-1.43614275811835 -0.953170391124348
-1.45657896457311 -0.92393474526093
-1.47556322143882 -0.893778101150025
-1.49307660479536 -0.862730519590318
-1.50910165692132 -0.830822949486319
-1.52362240369625 -0.798087196997824
-1.53662437052397 -0.764555893834839
-1.54809459676118 -0.730262464729583
-1.5580216486369 -0.695241094117991
-1.56639563064996 -0.659526692063919
-1.57320819543296 -0.623154859460042
-1.57845255207316 -0.586161852540099
-1.58212347288182 -0.5485845467379
-1.58421729860524 -0.510460399929086
-1.5847319420724 -0.471827415092308
-1.5836668902755 -0.432724102427027
-1.58102320488134 -0.393189440965714
-1.57680352117302 -0.353262839718699
-1.57101204542301 -0.312984098390414
-1.5636545507003 -0.272393367706184
-1.55473837111563 -0.231531109389109
-1.54427239451073 -0.190438055826937
-1.53226705359868 -0.149155169469136
-1.51873431556439 -0.107723601994633
-1.50368767013545 -0.0661846532909256
-1.48714211613523 -0.0245797302854578
-1.4691141465318 0.0170496943297099
-1.44962173199732 0.0586621234387791
-1.42868430299452 0.100216076867441
-1.40632273040797 0.141670132731217
-1.38255930473952 0.182982968725694
-1.35741771388869 0.224113403317495
-1.33092301954001 0.265020436794929
-1.303101632181 0.305663292137393
-1.27398128477563 0.346001455662795
-1.2435910051195 0.385994717412474
-1.21196108690427 0.425603211233357
-1.17912305952031 0.464787454517412
-1.14510965662747 0.503508387558767
-1.10995478352545 0.541727412489279
-1.0736934833563 0.579406431753721
-1.03636190217263 0.616507886086263
-0.99799725290643 0.652994791950354
-0.95863777827444 0.688830778404725
-0.918322712656945 0.72398012335873
-0.87709224298811 0.758407789180902
-0.834987468696767 0.792079457625231
-0.792050360737609 0.824961564040329
-0.748323719753633 0.857021330827415
-0.703851133411522 0.88822680011373
-0.658676932952517 0.918546865608844
-0.612846149002067 0.947951303612081
-0.566404466682327 0.976410803140162
-0.519398180072231 1.00389699514503
-0.471874146060548 1.03038248079273
-0.423879737637916 1.05584085877518
-0.375462796674418 1.08024675162755
-0.32667158622976 1.10357583102506
-0.277554742443609 1.12580484203402
-0.228161226054025 1.14691162629283
-0.17854027359234 1.16687514409991
-0.128741348303106 1.18567549538656
-0.0788140908380613 1.20329393955375
-0.0288082697732482 1.21971291415318
0.0212262680013873 1.23491605239388
0.0712396469704575 1.24888819945703
0.121182012710132 1.26161542760258
0.171003581584114 1.27308505005272
0.220654690369046 1.28328563363836
0.270085845759896 1.29220701019593
0.31924777370596 1.29984028670327
0.368091468528304 1.30617785414436
0.416568241769693 1.31121339509413
0.464629770728306 1.31494189001585
0.51222814662685 1.31735962226471
0.559315922369076 1.31846418179262
0.605846159836076 1.31825446755061
0.65177247667522 1.31673068858641
0.697049092535094 1.31389436383603
0.741630874700345 1.30974832060965
0.785473383080966 1.30429669177334
0.82853291451113 1.29754491162928
0.870766546313463 1.28949971049876
0.912132179085303 1.28016910801321
0.952588578664295 1.26956240512013
0.992095417231511 1.25769017481162
1.03061331351108 1.24456425158504
1.06810387202633 1.23019771964614
1.10452972137319 1.21460489986638
1.13985455147285 1.19780133550763
1.17404314976645 1.17980377672828
1.20706143631567 1.16063016388633
1.23887649777442 1.14029960965603
1.26945662019754 1.11883237997598
1.29877132065394 1.0962498738476
1.32679137761269 1.0725746020041
1.35348886007163 1.04783016447135
1.37883715539965 1.02204122704281
1.4028109958647 0.99523349669215
1.42538648382122 0.967433695947998
1.44654111553186 0.938669536256312
1.46625380359967 0.908969690357028
1.48450489798846 0.878363763702439
1.50127620561045 0.846882264945841
1.51655100846145 0.814556575529844
1.53031408028579 0.781418918404666
1.54255170175417 0.747502325907595
1.55325167413935 0.712840606835638
1.56240333147618 0.677468312744178
1.56999755119359 0.641420703505229
1.57602676320821 0.604733712159631
1.58048495747037 0.567443909098215
1.5833676899551 0.529588465607636
1.58467208709195 0.491205116817233
1.58439684862948 0.452332124083829
1.58254224893139 0.413008236851976
1.57911013670299 0.37327265402767
1.57410393314838 0.333164984904025
1.56752862856014 0.292725209677871
1.55939077734489 0.251993639596625
1.54969849148969 0.211010876775158
1.53846143247587 0.169817773722731
1.52569080164824 0.128455392620318
1.51139932904933 0.0869649643889382
1.49560126072983 0.0453878475897729
1.47831234454782 0.00376548719704961
1.45954981447095 -0.037860626715201
1.43933237339722 -0.0794490003313633
1.4176801745116 -0.120958177456165
1.3946148011968 -0.162346780839122
1.37015924551858 -0.203573553420279
1.34433788530667 -0.244597399456137
1.3171764598545 -0.2853774254848
1.28870204426165 -0.32587298108945
1.25894302244483 -0.366043699419578
1.22792905884419 -0.405849537429533
1.19569106885312 -0.445250815794303
1.16226118800114 -0.484208258462731
1.1276727399205 -0.522683031808729
1.09196020312848 -0.560636783341491
1.05515917665841 -0.598031679936086
1.01730634457382 -0.634830445546333
0.978439439400922 -0.670996398362389
0.938597204516073 -0.706493487375961
0.897819355525476 -0.741286328316742
0.856146540675828 -0.775340238924219
0.813620300335257 -0.808621273519705
0.77028302558498 -0.841096256844134
0.726177915962953 -0.872732817127871
0.681348936401628 -0.903499418359605
0.635840773402758 -0.933365391722127
0.589698790492918 -0.962300966163674
};
\addplot [darkorange25512714]
table {%
-0.783356969605672 0.702954216759224
-0.745592515470949 0.744252656534965
-0.707084838501919 0.784809209076924
-0.667872323988434 0.824583446731517
-0.627994059817742 0.863535721673654
-0.587489797510827 0.901627205428552
-0.54639991259723 0.938819927576803
-0.50476536436784 0.975076813604117
-0.462627655045787 1.010361721858
-0.420028788416119 1.04463947957454
-0.377011227955517 1.07787591793938
-0.333617854503781 1.11003790614794
-0.289891923519269 1.14109338443095
-0.245877021960925 1.17101139601224
-0.201617024839842 1.19976211796725
-0.157156051483702 1.22731689095105
-0.112538421557668 1.2536482477667
-0.0678086108855781 1.27872994074506
-0.0230112071154799 1.30253696790912
0.0218091347263109 1.32504559789644
0.066607736748309 1.34623339361515
0.111339942729754 1.36607923460968
0.155961162634904 1.38456333811414
0.200426917061351 1.40166727877217
0.244692881578056 1.41737400700374
0.288714930908908 1.43166786600057
0.332449182917758 1.44453460733316
0.375852042351086 1.45596140515393
0.418880244294699 1.46593686898229
0.461490897301138 1.47445105505897
0.503641526144802 1.48149547625811
0.545290114162185 1.48706311054747
0.586395145134991 1.49114840798811
0.626915644674409 1.49374729626669
0.666811221065271 1.49485718475485
0.706042105529397 1.49447696709161
0.744569191867974 1.49260702228621
0.782354075443463 1.48924921434029
0.819359091462187 1.48440689038985
0.855547352519404 1.4780848773687
0.890882785369483 1.47028947719692
0.925330166884497 1.46102846049892
0.958855159165402 1.45031105885755
0.991424343770801 1.43814795561186
1.02300525502918 1.42455127520766
1.05356641240138 1.4095345711117
1.0830773518611 1.39311281230118
1.1115086562621 1.37530236834243
1.13883198466183 1.35612099307332
1.16502010057233 1.33558780690586
1.19004689911013 1.31372327776658
1.21388743301816 1.29054920069363
1.23651793753377 1.26608867611101
1.25791585407787 1.24036608680158
1.27805985274192 1.21340707360178
1.29692985354995 1.18523850984226
1.31450704647478 1.15588847455997
1.33077391018824 1.12538622450838
1.34571422952682 1.0937621649937
1.35931311165528 1.06104781956618
1.37155700091218 1.02727579859683
1.38243369232242 0.992479766770592
1.39193234376346 0.956694409528736
1.40004348677298 0.919955398493603
1.40675903598723 0.882299355910382
1.4120722972007 0.84376381814125
1.41597797403906 0.804387198248313
1.4184721732387 0.764208747702628
1.41955240852762 0.723268517257487
1.41921760310378 0.681607317024951
1.41746809070852 0.639266675795446
1.41430561529384 0.596288799640958
1.40973332928401 0.552716529843106
1.40375579043315 0.508593300188021
1.396378957282 0.463963093670597
1.38761018321826 0.418870398651295
1.37745820914663 0.373360164509169
1.36593315477564 0.327477756835353
1.35304650853013 0.28126891221165
1.33881111609934 0.234779692619317
1.32324116763201 0.188056439523474
1.3063521835914 0.141145727678928
1.28816099928409 0.0940943187034436
1.26868574807822 0.0469491144647443
1.24794584332773 -0.000242889672286029
1.22596195902072 -0.0474346516916088
1.20275600917116 -0.0945791298185329
1.17835112597454 -0.14162932941198
1.15277163674917 -0.188538349809764
1.12604303968622 -0.235259431080191
1.09819197843252 -0.281746000633377
1.06924621553161 -0.327951719645817
1.03923460474947 -0.373830529251934
1.0081870623124 -0.419336696456557
0.976134537085922 -0.464424859722563
0.943108979724237 -0.509050074188247
0.909143310821158 -0.553167856469337
0.874271388094119 -0.596734229000995
0.838527972634077 -0.639705763875616
0.801948694254899 -0.682039626132708
0.764570015976786 -0.72369361645771
0.726429197679144 -0.764626213247189
0.687564258959127 -0.804796613998478
0.64801394123287 -0.844164775982497
0.60781766911721 -0.882691456159215
0.567015511130366 -0.920338250295965
0.525648139750772 -0.957067631249624
0.483756790873861 -0.99284298637448
0.441383222707233 -1.02762865401851
0.398569674145161 -1.06138995907171
0.355358822663946 -1.09409324753096
0.31179374178008 -1.12570592004712
0.267917858113625 -1.15619646442071
0.223774908099618 -1.18553448701401
0.179408894390639 -1.21369074304806
0.134864041994008 -1.24063716575451
0.0901847541873421 -1.2663468943532
0.0454155682563954 -1.29079430082752
0.000601111099328699 -1.31395501547103
-0.0442139452583546 -1.33580595117968
-0.0889849281938548 -1.35632532646556
-0.133667209017719 -1.37549268716912
-0.178216247460685 -1.39328892684837
-0.222587636072398 -1.40969630582458
-0.266737144487715 -1.42469846886556
-0.3106207635165 -1.43828046148892
-0.35419474901294 -1.45042874486905
-0.397415665480658 -1.46113120933287
-0.440240429370169 -1.47037718643106
-0.482626352025487 -1.47815745957256
-0.524531182237116 -1.4844642732119
-0.565913148358969 -1.48929134058009
-0.606730999947255 -1.4926338499514
-0.64694404887982 -1.4944884694398
-0.686512209914947 -1.49485335032024
-0.7253960406492 -1.49372812887154
-0.763556780834474 -1.49111392673889
-0.800956391015045 -1.48701334981584
-0.837557590446131 -1.48143048564665
-0.873323894256149 -1.47437089935173
-0.908219649815627 -1.46584162808023
-0.942210072276524 -1.45585117399524
-0.975261279246518 -1.44440949579868
-1.00734032456372 -1.43152799880418
-1.0384152311381 -1.4172195235681
-1.068455022827 -1.40149833308973
-1.09742975531275 -1.38438009859364
-1.12531054595188 -1.36588188390825
-1.15206960256595 -1.34602212845626
-1.17768025114537 -1.32482062887381
-1.20211696243867 -1.30229851927675
-1.22535537740058 -1.27847825019374
-1.2473723314737 -1.25338356618704
-1.26814587767939 -1.22703948218341
-1.28765530849501 -1.19947225853872
-1.30588117649565 -1.170709374861
-1.3228053137397 -1.14077950261825
-1.33841084987911 -1.10971247655801
-1.35268222897606 -1.07753926496749
-1.36560522500951 -1.04429193880371
-1.37716695605596 -1.01000363972442
-1.38735589713046 -0.974708547051855
-1.39616189167497 -0.938441843701964
-1.40357616168262 -0.901239681113363
-1.40959131644781 -0.863139143210781
-1.41420135993347 -0.824178209438996
-1.41740169674798 -0.784395716904101
-1.41918913672602 -0.743831321659831
-1.41956189810854 -0.702525459177545
-1.41851960931888 -0.660519304039272
-1.41606330933318 -0.617854728893987
-1.41219544664469 -0.574574262718051
-1.40691987682305 -0.530721048421397
-1.40024185867098 -0.486338799841739
-1.39216804898222 -0.441471758169672
-1.38270649590584 -0.396164647848087
-1.37186663092373 -0.350462631989877
-1.35965925944903 -0.304411267358358
-1.34609655005512 -0.258056458955306
-1.33119202234567 -0.211444414261846
-1.31496053347804 -0.16462159717783
-1.29741826335334 -0.117634681705614
-1.27858269848793 -0.070530505424395
-1.25847261458249 -0.0233560228014958
-1.23710805780602 0.0238417416128691
-1.21451032481332 0.0710157400606938
-1.19070194151611 0.118118948474399
-1.16570664062861 0.165104413351456
-1.13954933801036 0.211925298558652
-1.11225610782945 0.258534932019365
-1.0838541565713 0.304886852237314
-1.0543717959186 0.350934854610379
-1.02383841452961 0.39663303748836
-0.992284448742942 0.441935847928736
-0.959741352237906 0.486798127104827
-0.926241564680838 0.531175155321105
-0.891818479388517 0.575022696590745
-0.856506410040983 0.618297042731031
-0.820340556476916 0.660955056932614
-0.783356969605672 0.702954216759224
};
\addplot [darkorange25512714]
table {%
-0.719649272415747 0.45310184040373
-0.687703256543099 0.50304531035938
-0.655071723108398 0.552487333809609
-0.62178719993383 0.601378625870194
-0.587882865755957 0.649670450637938
-0.553392517152425 0.697314669771679
-0.518350534852808 0.744263790477635
-0.482791849467142 0.790471012851252
-0.446751906666338 0.83589027652837
-0.410266631849164 0.880476306599194
-0.373372394331026 0.924184658739306
-0.336105971090247 0.966971763512741
-0.298504510107968 1.00879496980294
-0.260605493338242 1.04961258732832
-0.222446699345198 1.08938392820004
-0.184066165644554 1.12806934748058
-0.145502150786988 1.1656302827027
-0.106793096221187 1.2020292923093
-0.0679775879745738 1.237230092976
-0.0290943181899148 1.27119759577918
0.00981795344384906 1.30389794117332
0.0487204383281028 1.33529853274299
0.0875743576198756 1.3653680696956
0.126340980887437 1.39407657806279
0.164981664717629 1.42139544057908
0.203457891236467 1.44729742520818
0.241731306504589 1.47175671228852
0.279763758749299 1.49474892027079
0.317517336395083 1.51625113002205
0.354954405854691 1.53624190767201
0.392037649043118 1.55470132597875
0.42873010057708 1.57161098419264
0.464995184622914 1.58695402639864
0.500796751356162 1.6007151583186
0.536099112996505 1.61288066255696
0.570867079382116 1.62343841227451
0.60506599304798 1.63237788327673
0.638661763773208 1.63969016450451
0.671620902562913 1.6453679669169
0.703910555030766 1.64940563075698
0.735498534148961 1.65179913119362
0.766353352332951 1.65254608233353
0.796444252828949 1.65164573959958
0.825741240372934 1.64909900047297
0.854215111090582 1.64490840359866
0.881837481608321 1.63907812625475
0.908580817346498 1.63161398018849
0.934418459966444 1.622523405823
0.959324653944088 1.61181546484051
0.983274572243621 1.59950083114946
1.00624434106563 1.58559178024454
1.02821106364501 1.57010217697017
1.04915284307497 1.55304746169976
1.06904880413434 1.53444463494439
1.08787911409645 1.51431224040634
1.10562500249882 1.49267034649431
1.12226877985394 1.46954052631883
1.13779385528259 1.44494583618765
1.15218475305196 1.41891079262278
1.16542712800216 1.39146134792185
1.17750777984587 1.36262486428829
1.1884146663266 1.33243008655612
1.19813691522271 1.30090711353644
1.20666483518506 1.26808736801431
1.21398992539756 1.23400356542582
1.220104884051 1.19868968124663
1.2250036156216 1.16218091712449
1.2286812369472 1.12451366578946
1.23113408209488 1.08572547477683
1.23235970601524 1.04585500899897
1.23235688697972 1.00494201220324
1.23112562779839 0.963027267354619
1.22866715581718 0.920152555982317
1.22498392169442 0.876360616531054
1.22007959695799 0.831695101758447
1.21395907034545 0.786200535220991
1.20662844293083 0.739922266892011
1.19809502204296 0.692906427955821
1.18836731398136 0.645199884823156
1.17745501553698 0.596850192413716
1.16536900432624 0.547905546752388
1.15212132794794 0.498414736926398
1.13772519197401 0.448427096451289
1.12219494678587 0.3979924540942
1.10554607326967 0.347161084203469
1.0877951673846 0.295983656594066
1.06895992361969 0.244511186038819
1.04905911735555 0.192794981415786
1.02811258614865 0.140886594562438
1.00614120995685 0.0888377688876772
0.983166890325748 0.0367003877928772
0.959212528556801 -0.0154735770466222
0.934302002878767 -0.0676321174880076
0.908460144645389 -0.119723240763841
0.881712713582953 -0.171695021309582
0.854086372112435 -0.223495652524118
0.825608658771823 -0.275073498411712
0.796307960765096 -0.326377145053886
0.76621348566525 -0.377355451859934
0.735355232299548 -0.427957602544973
0.70376396084604 -0.478133155784721
0.671471162171146 -0.52783209549651
0.638509026438879 -0.577004880696399
0.604910411022986 -0.625602494882708
0.570708807754002 -0.673576494896733
0.535938309533873 -0.720879059211936
0.500633576351401 -0.76746303560349
0.464829800732417 -0.813281988150635
0.428562672659105 -0.858290243525017
0.391868343993455 -0.902442936518856
0.354783392440295 -0.945696054767558
0.317344785085846 -0.988006482622198
0.279589841548118 -1.02933204412814
0.241556196775913 -1.06963154506694
0.203281763533481 -1.10886481401964
0.164804694608258 -1.14699274241052
0.126163344779334 -1.18397732349136
0.0873962325845776 -1.21978169022742
0.0485420019245182 -1.25437015204727
0.00963938354126559 -1.28770823041996
-0.0292728435891347 -1.31976269322394
-0.0681558909124304 -1.35050158787356
-0.106970998961433 -1.37989427317016
-0.145679475992304 -1.4079114498458
-0.184242736553318 -1.43452518976946
-0.222622339947689 -1.45970896378641
-0.260780028552087 -1.48343766816301
-0.298677765952671 -1.5056876496107
-0.336277774860613 -1.52643672886408
-0.373542574769322 -1.54566422278972
-0.410435019315826 -1.56335096500351
-0.446918333309079 -1.57947932497618
-0.482956149388264 -1.59403322560774
-0.518512544274577 -1.60699815925353
-0.553552074580326 -1.61836120218574
-0.588039812139675 -1.62811102747606
-0.621941378825798 -1.63623791628659
-0.655222980819742 -1.64273376755784
-0.687851442296844 -1.64759210608399
-0.719794238497111 -1.65080808896752
-0.751019528146607 -1.65237851044678
-0.781496185197537 -1.65230180509147
-0.811193829855353 -1.65057804936315
-0.840082858862008 -1.64720896153904
-0.86813447500512 -1.64219789999913
-0.89532071582366 -1.63554985987855
-0.921614481481543 -1.62727146808828
-0.946989561781333 -1.61737097670929
-0.97142066229114 -1.60585825476669
-0.994883429558656 -1.59274477839208
-1.01735447538721 -1.57804361938392
-1.03881140014963 -1.56176943217721
-1.05923281511666 -1.54393843923569
-1.07859836377776 -1.52456841488091
-1.09688874213287 -1.50367866757438
-1.11408571793506 -1.48129002067054
-1.13017214886487 -1.45742479165951
-1.14513199961809 -1.43210676992058
-1.15895035789017 -1.40536119300844
-1.17161344924105 -1.37721472149581
-1.18310865082588 -1.34769541239759
-1.19342450397774 -1.31683269120309
-1.20255072562985 -1.28465732254398
-1.21047821856599 -1.25120137952747
-1.21719908048881 -1.21649821176516
-1.22270661189698 -1.18058241212938
-1.2269953227634 -1.1434897822703
-1.23006093800783 -1.10525729692803
-1.23190040175832 -1.06592306707543
-1.23251188039739 -1.02552630192824
-1.23189476438984 -0.984107269860503
-1.23004966889033 -0.941707258264145
-1.22697843313018 -0.898368532392829
-1.22268411858398 -0.854134293231016
-1.21717100591786 -0.809048634430295
-1.21044459072244 -0.763156498355864
-1.20251157803465 -0.716503631287012
-1.19337987565408 -0.669136537816219
-1.18305858626025 -0.621102434492382
-1.17155799833888 -0.572449202754325
-1.15888957592613 -0.523225341201534
-1.14506594718097 -0.473479917249712
-1.13010089179717 -0.423262518219298
-1.11400932726738 -0.372623201905756
-1.09680729401308 -0.321612446680876
-1.07851193939513 -0.270281101174828
-1.05914150062088 -0.21868033358915
-1.03871528656498 -0.166861580691166
-1.0172536585218 -0.114876496540701
-0.994778009908937 -0.0627769010001914
-0.971310744941706 -0.0106147280795132
-0.94687525630019 0.0415580258329586
-0.92149590181088 0.0936893538014921
-0.895197980166271 0.145727290184633
-0.868007705706579 0.197619962435689
-0.839952182288709 0.249315642810421
-0.811059376268545 0.300762799930368
-0.781358088623479 0.351910150150462
-0.750877926242968 0.402706708679665
-0.719649272415747 0.45310184040373
};
\addplot [line width=0.52pt, gray127, dashed]
table {%
0 0.5
0.0631370995296211 0.499750794563087
0.126211262625347 0.499003426665747
0.18915961558969 0.497758641300553
0.251919410135435 0.496017679296629
0.314428085934502 0.493782276082762
0.376623332979436 0.49105465995749
0.438443153695383 0.487837549867887
0.499825924740617 0.484134152699254
0.560710458434029 0.479948160078428
0.621036063748338 0.475283744693891
0.680742606808226 0.470145556136338
0.739770570833094 0.464538716263865
0.798061115464682 0.458468814096394
0.855556135420419 0.451941900244412
0.912198318414032 0.44496448087759
0.967931202285678 0.437543511239297
1.02269923128465 0.429686388713456
1.07644781144858 0.421400945450675
1.12912336502384 0.412695440560988
1.18067338387306 0.403578551880993
1.23104648181636 0.394059367323596
1.28019244585421 0.384147375818985
1.32806228622086 0.373852457855855
1.37460828521837 0.36318487563232
1.41978404478267 0.352155262826336
1.46354453273415 0.340774613995817
1.50584612766675 0.329054273619019
1.54664666243068 0.317005924786119
1.58590546616556 0.304641577553258
1.62358340484204 0.291973556970653
1.65964292027145 0.27901449079672
1.6940480675446 0.265777296910451
1.72676455086245 0.252275170434589
1.75775975772292 0.238521570582445
1.78700279142975 0.22453020724146
1.81446450189096 0.210315027306892
1.84011751467635 0.195890200779247
1.86393625830487 0.181270106639312
1.88589698973487 0.166469318514881
1.90597781803168 0.151502590153436
1.92415872618893 0.136384840715302
1.94042159108197 0.121131139901892
1.95475020153341 0.10575669293391
1.96713027447274 0.0902768253944512
1.97754946917401 0.0747069679521305
1.98599739955734 0.059062640979454
1.99246564454201 0.0433594390817754
1.99694775644076 0.0276130155522541
1.99943926738696 0.011839066768312
1.99993769378831 -0.00394668345485666
1.99844253880255 -0.0197284995381263
1.99495529283268 -0.035490649824008
1.98947943204131 -0.0512174222583066
1.98202041488558 -0.0668931400522395
1.97258567667601 -0.0825021773093997
1.96118462216481 -0.098028974601989
1.94782861617108 -0.113458054480795
1.93253097225204 -0.128774036903448
1.91530693943186 -0.143961654565583
1.89617368700102 -0.159005768119619
1.87515028740165 -0.173891381265991
1.85225769721568 -0.188603655701788
1.82751873627487 -0.203127925911895
1.8009580649135 -0.217449713787896
1.77260215938642 -0.231554743060174
1.74247928547692 -0.245428953528797
1.71061947032082 -0.259058515079039
1.67705447247475 -0.272429841467535
1.64181775025853 -0.285529603865347
1.60494442840316 -0.298344744144428
1.56647126303778 -0.310862487894248
1.52643660505034 -0.323070357155605
1.48488036185857 -0.334956182858928
1.44184395762943 -0.346508116954666
1.39737029198661 -0.357714644223686
1.35150369724722 -0.368564593755889
1.30428989423037 -0.379047150085623
1.25577594668172 -0.389151863972765
1.20601021435923 -0.39886866281876
1.15504230482718 -0.4081878607072
1.10292302400622 -0.417100168058959
1.04970432552894 -0.425596700892247
0.995439258951366 -0.433668989678357
0.940181916872006 -0.441308987784273
0.883987381011158 -0.448509079493732
0.826911667304268 -0.455262087598731
0.769011670064022 -0.46156128055393
0.710345105266857 -0.467400379186799
0.650970453020423 -0.472773562956834
0.59094689926934 -0.477675475757597
0.530334276797357 -0.482101231255806
0.469193005584738 -0.486046417762128
0.407584032580305 -0.489507102628857
0.345568770948199 -0.492479836170055
0.283209038849903 -0.494961655100279
0.22056699782255 -0.496950085488444
0.157705090814953 -0.498443145223897
0.0946859799431164 -0.499439345992222
0.0315724840272736 -0.499937694758829
-0.031572484027274 -0.499937694758829
-0.0946859799431159 -0.499439345992222
-0.157705090814953 -0.498443145223897
-0.22056699782255 -0.496950085488444
-0.283209038849904 -0.494961655100279
-0.345568770948199 -0.492479836170055
-0.407584032580304 -0.489507102628857
-0.469193005584737 -0.486046417762128
-0.530334276797357 -0.482101231255806
-0.59094689926934 -0.477675475757597
-0.650970453020424 -0.472773562956834
-0.710345105266857 -0.467400379186799
-0.769011670064021 -0.46156128055393
-0.826911667304268 -0.455262087598731
-0.883987381011158 -0.448509079493732
-0.940181916872006 -0.441308987784273
-0.995439258951366 -0.433668989678357
-1.04970432552894 -0.425596700892247
-1.10292302400621 -0.417100168058959
-1.15504230482718 -0.4081878607072
-1.20601021435923 -0.39886866281876
-1.25577594668172 -0.389151863972765
-1.30428989423037 -0.379047150085623
-1.35150369724722 -0.36856459375589
-1.39737029198661 -0.357714644223686
-1.44184395762943 -0.346508116954666
-1.48488036185857 -0.334956182858928
-1.52643660505034 -0.323070357155605
-1.56647126303778 -0.310862487894248
-1.60494442840316 -0.298344744144428
-1.64181775025852 -0.285529603865347
-1.67705447247475 -0.272429841467535
-1.71061947032082 -0.259058515079039
-1.74247928547692 -0.245428953528797
-1.77260215938642 -0.231554743060174
-1.8009580649135 -0.217449713787897
-1.82751873627487 -0.203127925911894
-1.85225769721568 -0.188603655701788
-1.87515028740165 -0.173891381265991
-1.89617368700102 -0.159005768119619
-1.91530693943186 -0.143961654565583
-1.93253097225204 -0.128774036903448
-1.94782861617108 -0.113458054480795
-1.96118462216481 -0.098028974601989
-1.97258567667601 -0.0825021773093994
-1.98202041488558 -0.0668931400522396
-1.98947943204131 -0.0512174222583065
-1.99495529283268 -0.0354906498240081
-1.99844253880255 -0.0197284995381264
-1.99993769378831 -0.00394668345485689
-1.99943926738696 0.0118390667683119
-1.99694775644076 0.0276130155522542
-1.99246564454201 0.0433594390817752
-1.98599739955734 0.059062640979454
-1.97754946917401 0.0747069679521302
-1.96713027447274 0.090276825394451
-1.95475020153341 0.10575669293391
-1.94042159108197 0.121131139901891
-1.92415872618893 0.136384840715302
-1.90597781803168 0.151502590153436
-1.88589698973487 0.166469318514881
-1.86393625830487 0.181270106639312
-1.84011751467635 0.195890200779246
-1.81446450189096 0.210315027306892
-1.78700279142975 0.22453020724146
-1.75775975772292 0.238521570582445
-1.72676455086245 0.252275170434589
-1.6940480675446 0.265777296910451
-1.65964292027145 0.27901449079672
-1.62358340484204 0.291973556970653
-1.58590546616556 0.304641577553258
-1.54664666243068 0.317005924786119
-1.50584612766675 0.329054273619019
-1.46354453273415 0.340774613995817
-1.41978404478267 0.352155262826336
-1.37460828521837 0.36318487563232
-1.32806228622086 0.373852457855854
-1.28019244585421 0.384147375818985
-1.23104648181636 0.394059367323596
-1.18067338387306 0.403578551880993
-1.12912336502384 0.412695440560988
-1.07644781144858 0.421400945450675
-1.02269923128465 0.429686388713456
-0.967931202285677 0.437543511239297
-0.912198318414033 0.44496448087759
-0.855556135420419 0.451941900244412
-0.798061115464683 0.458468814096394
-0.739770570833094 0.464538716263865
-0.680742606808225 0.470145556136338
-0.621036063748338 0.475283744693891
-0.560710458434029 0.479948160078428
-0.499825924740618 0.484134152699254
-0.438443153695383 0.487837549867887
-0.376623332979435 0.49105465995749
-0.314428085934502 0.493782276082762
-0.251919410135435 0.496017679296629
-0.189159615589691 0.497758641300553
-0.126211262625347 0.499003426665747
-0.0631370995296205 0.499750794563087
-4.89858719658941e-16 0.5
};
\addplot [line width=0.52pt, gray127, dashed]
table {%
-0.445503262094184 0.226995249869774
-0.416617575056473 0.283137680568154
-0.387316594559492 0.338997873468965
-0.357629528457827 0.394520145917093
-0.327585969464592 0.449649152103885
-0.297215865652769 0.504329938237046
-0.266549490602317 0.55850799731978
-0.235617413222806 0.612129323484547
-0.204450467281669 0.66514046582729
-0.173079720668428 0.717488581688452
-0.141536444425547 0.769121489327695
-0.109852081576777 0.819987719939789
-0.0780582157840671 0.870036568959839
-0.0461865398642803 0.919218146606698
-0.014268824197108 0.967483427614181
0.0176631149443375 1.01478430010053
0.0495774471086752 1.06107361352735
0.0814423593955418 1.10630522570036
0.113226088167383 1.15043404876484
0.144896950712137 1.19341609415028
0.176423376825253 1.23520851641907
0.207773940279555 1.2757696559758
0.238917390151591 1.31505908059447
0.269822681973231 1.35303762572215
0.300459008677471 1.38966743351911
0.330795831307583 1.42491199059634
0.360802909459014 1.45873616441286
0.390450331423681 1.4911062382966
0.419708544006606 1.52198994505395
0.44854838198519 1.55135649913435
0.476941097181737 1.57917662731806
0.504858387120264 1.60542259789628
0.532272423239025 1.63006824831473
0.559155878630625 1.65308901125312
0.585481955282079 1.67446193911427
0.611224410787649 1.69416572689891
0.636357584507851 1.71218073344284
0.660856423148531 1.72848900099576
0.684696505734541 1.74307427312192
0.707854067953093 1.75592201090485
0.730306025842545 1.76701940744013
0.752029998802995 1.77635540060155
0.773004331905752 1.78392068406809
0.793208117479436 1.78970771660067
0.812621215951201 1.79371072955939
0.831224275922303 1.7959257326539
0.848998753457988 1.79635051792092
0.865926930572497 1.79498466192528
0.881991932890733 1.79182952618192
0.89717774646901 1.78688825579878
0.911469233758097 1.78016577634164
0.924852148692657 1.77166878892421
0.937313150892037 1.76140576352835
0.948839818958248 1.74938693056096
0.959420662857882 1.73562427065613
0.969045135375633 1.72013150273255
0.977703642627984 1.70292407031823
0.985387553626604 1.68401912615602
0.992089208881904 1.66343551510539
0.997801928038179 1.64119375535748
1.00252001653274 1.61731601798212
1.00623877127236 1.59182610482722
1.00895448532144 1.56474942479254
1.01066445159716 1.53611296850159
1.01136696556796 1.5059452813968
1.01106132695262 1.47427643528474
1.00974784041836 1.44113799835993
1.00742781527712 1.40656300373699
1.00410356418043 1.37058591652248
0.999778400814071 1.33324259945935
0.994456636594945 1.29457027717819
0.988143576373364 1.25460749909087
0.98084551314505 1.21339410096369
0.972569721778146 1.17097116520819
0.963324451761466 1.12738097992934
0.953118918981224 1.08266699677181
0.941963296534443 1.03687378760643
0.929868704588185 0.990047000099996
0.916847199294733 0.942233312212662
0.90291176077375 0.893480385668352
0.888076280173416 0.843836818444506
0.872355545823431 0.793352096328566
0.855765228493684 0.742076543589459
0.838321865773289 0.690061272813279
0.82004284558556 0.637358133953147
0.800946388855343 0.584019662644047
0.781051531346012 0.530099027834167
0.760378104684199 0.475649978784932
0.738946716591197 0.420726791492574
0.71677873034074 0.365384214584636
0.693896243463625 0.309677414745351
0.670322065720412 0.253661921724291
0.646079696364157 0.197393572983105
0.621193300715849 0.140928458035523
0.595687686075892 0.0843228625361073
0.569588276995645 0.0276332121734919
0.542921089933682 -0.0290839835759755
0.515712707322013 -0.0857721877781061
0.487990251068145 -0.142374892398116
0.45978135551937 -0.198835674629064
0.431114139916248 -0.255098253135332
0.402017180362737 -0.311106544155087
0.372519481340911 -0.36680471740579
0.34265044679866 -0.422137251737043
0.312439850839197 -0.47704899047527
0.281917808041584 -0.531485196405095
0.251114743441865 -0.585391606332588
0.220061362204727 -0.638714485175996
0.188788619015925 -0.691400679530039
0.157327687225979 -0.743397670650385
0.125709927775898 -0.794653626805474
0.0939668579359139 -0.845117454943511
0.062130119888381 -0.89473885162314
0.0302314491861592 -0.943468353157001
-0.00169735688207434 -0.991257384918214
-0.0336244709880548 -1.03805830976062
-0.0655180674901006 -1.08382447550452
-0.0973463541576082 -1.12851026144061
-0.129077603862237 -1.17207112380564
-0.1606801862042 -1.21446364018466
-0.192122599042139 -1.25564555279542
-0.223373499895137 -1.29557581061181
-0.25440173718559 -1.33421461028444
-0.285176381291777 -1.37152343581751
-0.315666755379174 -1.40746509696234
-0.345842465979799 -1.44200376628946
-0.375673433289073 -1.47510501490214
-0.405129921150029 -1.5067358467559
-0.434182566694959 -1.53686473154964
-0.462802409614956 -1.56546163615578
-0.490960921028181 -1.59249805455788
-0.51863003191808 -1.6179470362661
-0.545782161113175 -1.64178321318202
-0.572390242780593 -1.66398282488611
-0.598427753405864 -1.68452374232264
-0.623868738232139 -1.70338548985838
-0.648687837132466 -1.72054926569319
-0.6728603098893 -1.73599796060206
-0.696362060856105 -1.74971617498991
-0.719169662976409 -1.76169023424232
-0.741260381136405 -1.7719082023566
-0.762612194827812 -1.78035989383994
-0.783203820098389 -1.78703688386244
-0.803014730768248 -1.79193251665526
-0.822025178890792 -1.79504191214514
-0.840216214437899 -1.79636197081906
-0.857569704189724 -1.7958913768138
-0.874068349810285 -1.79363059922771
-0.889695705090821 -1.78958189165306
-0.904436192343729 -1.78374928992961
-0.918275117930745 -1.77613860812159
-0.93119868690988 -1.76675743272216
-0.943194016786517 -1.75561511509095
-0.954249150354966 -1.74272276213252
-0.964353067617662 -1.72809322522467
-0.973495696770141 -1.71174108740794
-0.981667924240834 -1.69368264884893
-0.988861603775671 -1.67393591059193
-0.995069564558446 -1.65252055661507
-1.00028561835884 -1.62945793420894
-1.00450456570099 -1.6047710326971
-1.00772220104644 -1.57848446051987
-1.00993531698629 -1.55062442070406
-1.01114170743844 -1.52121868474325
-1.01134016984666 -1.49029656491456
-1.01053050637931 -1.4578888850595
-1.00871352412653 -1.42402794985816
-1.00589103429575 -1.38874751262713
-1.00206585040621 -1.35208274167352
-0.99724178548441 -1.31407018523838
-0.99142364826317 -1.27474773506466
-0.984617238388205 -1.23415458862585
-0.976829340636902 -1.19233121005314
-0.968067718155106 -1.14931928979983
-0.958341104718653 -1.10516170308347
-0.947659196027347 -1.05990246714679
-0.936032640040078 -1.01358669738047
-0.923473026360705 -0.966260562351072
-0.909992874685284 -0.917971237779223
-0.895605622322169 -0.868766859513837
-0.880325610797404 -0.818696475549247
-0.86416807155878 -0.767809997133103
-0.847149110792799 -0.716158149013752
-0.829285693369675 -0.663792418876705
-0.810595625932384 -0.610765006020596
-0.791097539146611 -0.557128769323779
-0.770810869129303 -0.502937174553459
-0.749755838074319 -0.44824424106984
-0.727953434094514 -0.393104487978462
-0.705425390300325 -0.337572879784369
-0.682194163135738 -0.281704771602288
-0.658282909993217 -0.225555853977458
-0.633715466129905 -0.169182097372071
-0.608516320908132 -0.112639696372711
-0.582710593383877 -0.0559850136743582
-0.556324007267548 0.000725476103169143
-0.529382865282035 0.0574352427103378
-0.501914022943573 0.114087756618493
-0.473944861791589 0.170626545369709
-0.445503262094184 0.226995249869773
};

\nextgroupplot[
height=4.6cm,
label style={font=\footnotesize},
legend style={draw=none, fill=white, fill opacity=0.9, text opacity=1, font=\footnotesize},
tick align=outside,
tick label style={font=\scriptsize},
tick pos=left,
title style={font=\footnotesize, yshift=-2pt},
title={affine invariante},
width=0.45\textwidth,
x grid style={darkgray176},
xmin=-2.29992834785656, xmax=2.29992834785656,
xtick style={color=black},
y grid style={darkgray176},
ymin=-2.29992834785656, ymax=2.29992834785656,
ytick style={color=black}
]
\addplot [forestgreen4416044]
table {%
0.0277035868646075 -0.61262564284864
-0.0237349791137063 -0.614645786139378
-0.0751498855486594 -0.616053237143389
-0.126489880943364 -0.616846592882904
-0.177703788473886 -0.617025062523668
-0.228740557003418 -0.616588468163261
-0.279549311971158 -0.615537245008438
-0.330079406105184 -0.6138724409413
-0.380280469908752 -0.611595715474748
-0.430102461869708 -0.608709338098239
-0.479495718342953 -0.605216186015515
-0.528411003056247 -0.601119741276539
-0.576799556189996 -0.596424087306516
-0.624613142982099 -0.591133904835441
-0.671804101809407 -0.585254467232252
-0.71832539169787 -0.578791635248215
-0.764130639213998 -0.571751851174804
-0.809174184690909 -0.56414213242188
-0.853411127742882 -0.555970064522592
-0.896797372023031 -0.547243793571944
-0.939289669179511 -0.537972018106595
-0.980845661966407 -0.528163980433969
-1.02142392646636 -0.51782945741932
-1.06098401338285 -0.50697875073994
-1.09948648836088 -0.495622676616225
-1.13689297129606 -0.48377255502983
-1.17316617459271 -0.471440198439663
-1.2082699403329 -0.458637900006974
-1.24216927631954 -0.44537842134126
-1.27483039095728 -0.431674979779218
-1.30622072693675 -0.417541235209415
-1.33630899368836 -0.402991276455813
-1.36506519857349 -0.388039607233726
-1.39246067678172 -0.372701131692198
-1.41846811990457 -0.356991139557224
-1.44306160315716 -0.340925290890621
-1.46621661122052 -0.324519600479733
-1.48791006267904 -0.307790421873547
-1.50812033302853 -0.290754431081119
-1.52682727623199 -0.273428609948563
-1.54401224480158 -0.255830229231179
-1.55965810838692 -0.237976831377591
-1.57374927085091 -0.219886213043049
-1.58627168581637 -0.201576407349336
-1.59721287066772 -0.183065665908964
-1.60656191899396 -0.164372440631565
-1.61430951146039 -0.14551536533063
-1.62044792509836 -0.126513237148915
-1.62497104100366 -0.107384997821045
-1.62787435043597 -0.0881497147919752
-1.62915495931332 -0.0688265622101513
-1.62881159109694 -0.0494348018143002
-1.62684458806372 -0.0299937637329096
-1.62325591096507 -0.0105228272155363
-1.61804913707236 0.00895859868485228
-1.61122945661105 0.0284310944592434
-1.60280366758693 0.047875249500368
-1.59278016900975 0.0672716814516162
-1.58116895252086 0.0866010555277925
-1.56798159243339 0.105844103788449
-1.55323123419466 0.124981644344588
-1.53693258128256 0.143994600479587
-1.51910188054875 0.162864019665284
-1.49975690602352 0.181571092454268
-1.47891694119817 0.200097171229549
-1.45660275980286 0.218423788792904
-1.43283660509889 0.236532676773386
-1.40764216770612 0.25440578383763
-1.38104456198767 0.272025293683816
-1.35307030101536 0.289373642801347
-1.32374727014089 0.306433537978535
-1.29310469919915 0.323187973540857
-1.26117313337125 0.339620248302573
-1.22798440273646 0.355713982214838
-1.19357159054324 0.371453132693687
-1.15796900023119 0.386822010611634
-1.1212121212366 0.401805295936932
-1.08333759361582 0.416388053004918
-1.04438317152169 0.430555745406207
-1.0043876855693 0.444294250476897
-0.963391004128853 0.457589873376356
-0.921433993584003 0.470429360738527
-0.878558477595321 0.48279991388318
-0.834807195409588 0.49468920157391
-0.790223759256375 0.506085372310181
-0.744852610874422 0.516977066141165
-0.698738977211155 0.527353425989582
-0.651928825339471 0.537204108474274
-0.604468816636767 0.54651929422071
-0.556406260271861 0.555289697649156
-0.507789066046187 0.563506576230737
-0.458665696636258 0.571161739202186
-0.40908511928502 0.578247555730572
-0.359096756990237 0.584756962519878
-0.308750439238574 0.590683470851854
-0.258096352334482 0.596021173054105
-0.207184989373393 0.600764748388995
-0.15606709990911 0.604909468357468
-0.104793639365544 0.608451201412521
-0.0534157182432418 0.61138641707762
-0.0019845511713292 0.613712189465958
0.0494485941443503 0.615426200197036
0.100832448025996 0.616526740707683
0.15211578993058 0.617012713955183
0.203247499507601 0.616883635510835
0.254176607556956 0.616139634042838
0.304852346836153 0.614781451188037
0.355224202666202 0.612810440812639
0.40524196328574 0.610228567662649
0.454855769903214 0.607038405405375
0.504016166397203 0.603243134063933
0.552674148615359 0.598846536847333
0.60078121322281 0.593852996379298
0.648289406051339 0.588267490329563
0.695151369901143 0.582095586452031
0.74132039174752 0.575343437034711
0.786750449305425 0.568017772766986
0.831396256905486 0.560125896030313
0.875213310635747 0.551675673619054
0.918157932704132 0.542675528898681
0.960187314977417 0.533134433409183
1.00125956165331 0.52306189792204
1.04133373102311 0.51246796295968
1.08036987628326 0.501363188786867
1.11832908535522 0.489758644883998
1.15517351967386 0.477665898912809
1.19086645190578 0.46509700518547
1.22537230255989 0.452064492648596
1.25865667545384 0.438581352394107
1.29068639200087 0.424661024709437
1.32142952428296 0.410317385679957
1.35085542687725 0.395564733356991
1.37893476740409 0.380417773505211
1.40563955576619 0.364891604943606
1.43094317204974 0.349001704494658
1.45482039305974 0.332763911556702
1.47724741746301 0.316194412314873
1.49820188951388 0.299309723606362
1.51766292133892 0.282126676456076
1.53561111375836 0.264662399299105
1.55202857562367 0.246934300906728
1.56689894165182 0.228960053032966
1.58020738873855 0.210757572798996
1.59194065073439 0.192345004832978
1.60208703166861 0.17374070318309
1.61063641740803 0.154963213021817
1.61758028573901 0.136031252159718
1.6229117148626 0.116963692387105
1.62662539029429 0.0977795406622352
1.62871761016165 0.0784979201647582
1.62918628889442 0.0591380512333148
1.62803095930344 0.0397192322062883
1.62525277304637 0.0202608201847991
1.6208544994797 0.000782211737123797
1.61484052289817 -0.0186971764362218
1.60721683816443 -0.0381579268574765
1.5979910447332 -0.0575806406273945
1.58717233907604 -0.0769459567624907
1.57477150551401 -0.0962345714944926
1.5608009054677 -0.115427257512752
1.54527446513506 -0.134504883130452
1.52820766160949 -0.153448431355482
1.50961750745189 -0.172239018846993
1.48952253373223 -0.190857914738723
1.46794277155733 -0.20928655931033
1.44489973210343 -0.227506582488125
1.42041638517339 -0.245499822156763
1.39451713629989 -0.263248342263629
1.36722780241745 -0.280734450697893
1.33857558612757 -0.297940716926386
1.30858904858255 -0.314849989368737
1.27729808101517 -0.331445412494436
1.2447338749424 -0.347710443624801
1.21092889107305 -0.363628869423075
1.17591682695022 -0.379184822056234
1.13973258336085 -0.394362795012387
1.10241222954588 -0.409147658558007
1.06399296724561 -0.423524674819576
1.02451309361625 -0.437479512474617
0.984011963054422 -0.45099826103747
0.942529947967866 -0.464067444725568
0.900108398531364 -0.476674035892386
0.85678960146796 -0.488805468013686
0.81261673789665 -0.500449648214098
0.7676338402885 -0.511594969321564
0.721885748574115 -0.522230321437622
0.675418065446232 -0.532345103011987
0.628277110901953 -0.541929231410414
0.580509876069976 -0.550973152965282
0.532163976368806 -0.559467852498901
0.483287604042672 -0.567404862310038
0.43392948012246 -0.574776270614699
0.384138805859513 -0.581574729432778
0.333965213680768 -0.587793461912673
0.283458717714067 -0.593426269086605
0.232669663932987 -0.598467536049884
0.18164867997088 -0.602912237557964
0.130446624654145 -0.60675594303572
0.0791145373050461 -0.609994820993941
0.0277035868646079 -0.61262564284864
};
\addplot [forestgreen4416044]
table {%
0.129442481370526 -0.724924354440328
0.087176143775836 -0.732098600944183
0.0448229071051537 -0.738543075641191
0.0024249899858759 -0.744251354538822
-0.0399753444161577 -0.749217747500529
-0.0823358305255072 -0.7534373039178
-0.124614242488378 -0.756905817645033
-0.166768436264269 -0.759619831192312
-0.20875639163607 -0.761576639171911
-0.250536254096727 -0.76277429099508
-0.292066376570724 -0.763211592816434
-0.333305360928798 -0.762888108724009
-0.374212099254489 -0.761804161173782
-0.414745814821408 -0.759960830668245
-0.454866102740366 -0.757359954679333
-0.494532970235842 -0.754004125816796
-0.533706876511655 -0.749896689243826
-0.572348772166091 -0.745041739342524
-0.610420138117188 -0.739444115632537
-0.647883023999406 -0.733109397946914
-0.684700085993366 -0.726043900870009
-0.720834624050989 -0.718254667442967
-0.756250618478896 -0.709749462143064
-0.790912765843628 -0.700536763143912
-0.824786514162872 -0.690625753864229
-0.857838097347627 -0.680026313813605
-0.890034568860981 -0.668749008744395
-0.921343834559933 -0.656805080119548
-0.951734684687529 -0.644206433906869
-0.981176824983434 -0.630965628710895
-1.0096409068819 -0.6170958632542
-1.03709855676707 -0.602610963220626
-1.0635224042564 -0.587525367473541
-1.08888610948404 -0.571854113662861
-1.11316438935704 -0.555612823235197
-1.13633304275801 -0.538817685862057
-1.15836897466939 -0.52148544330163
-1.17925021919501 -0.503633372710237
-1.19895596145621 -0.485279269420082
-1.21746655834055 -0.466441429200484
-1.2347635580825 -0.447138630020253
-1.25082971865655 -0.427390113329402
-1.26564902496443 -0.407215564878853
-1.27920670479932 -0.38663509509725
-1.29148924357114 -0.365669219044449
-1.3024843977781 -0.344338835961657
-1.31218120721138 -0.322665208438617
-1.32057000588044 -0.300669941218594
-1.32764243164835 -0.2783749596623
-1.33339143456729 -0.255802487892212
-1.33781128390612 -0.23297502663909
-1.34089757386289 -0.209915330812754
-1.34264722795666 -0.186646386819497
-1.34305850209418 -0.163191389648739
-1.34213098630845 -0.139573719751749
-1.33986560516736 -0.11581691973551
-1.33626461685211 -0.0919446708949326
-1.33133161090616 -0.0679807696068262
-1.32507150465713 -0.0439491036091534
-1.31749053831506 -0.0198736281892134
-1.30859626875207 0.00422165769551044
-1.29839756196949 0.0283127353400324
-1.28690458426004 0.0523755902342317
-1.27412879207384 0.0763862360010396
-1.2600829205984 0.10032073830658
-1.24478097106386 0.124155238718436
-1.2282381967863 0.147865978488253
-1.21047108796288 0.171429322234974
-1.19149735523408 0.194821781505107
-1.17133591202931 0.218020038186518
-1.15000685571363 0.241000967752442
-1.12753144755422 0.263741662312504
-1.10393209152672 0.286219453447812
-1.07923231198246 0.308411934807321
-1.05345673019885 0.330296984442979
-1.02663103983636 0.35185278686136
-0.998781981326549 0.373057854769825
-0.969937315216589 0.393891050495523
-0.94012579449696 0.414331607055882
-0.909377135939849 0.434359148859588
-0.877721990476814 0.453953712017419
-0.84519191264528 0.473095764242683
-0.811819329134281 0.491766224321427
-0.777637506460836 0.509946481133008
-0.74268051780916 0.527618412202055
-0.706983209065775 0.544764401763358
-0.670581164084368 0.561367358321637
-0.633510669215036 0.577410731688718
-0.595808677133259 0.592878529481124
-0.557512770004676 0.607755333061627
-0.518661122022359 0.622026312908883
-0.479292461353949 0.635677243399818
-0.439446031536574 0.648694516990045
-0.399161552358045 0.661065157778154
-0.358479180263298 0.672776834440374
-0.317439468325584 0.683817872522707
-0.27608332582228 0.694177266078275
-0.234451977455632 0.703844688638292
-0.192586922259076 0.712810503505703
-0.150529892230099 0.721065773361266
-0.108322810730875 0.728602269172453
-0.0660077506981477 0.735412478396341
-0.023626892704011 0.741489612468268
0.0187775170906025 0.746827613568829
0.0611632090478137 0.751421160662437
0.103487932188089 0.75526567480145
0.14570949630694 0.758357323690566
0.187785814031033 0.760693025506935
0.229674942771799 0.762270451972191
0.271335126534719 0.763088030673329
0.312724837542608 0.763144946630119
0.353802817631403 0.762441143107496
0.394528119377202 0.760977321672119
0.434860146913545 0.758754941493028
0.474758696398258 0.755776217887119
0.514183996089514 0.752044120110861
0.553096745991173 0.747562368400483
0.591458157027872 0.742335430263556
0.62922998971082 0.73636851602569
0.666374592255748 0.729667573636762
0.702854938115033 0.722239282741877
0.73863466288656 0.714091048022946
0.773678100562551 0.705230991817542
0.80795031908222 0.695667946022375
0.841417155152804 0.685411443289468
0.874045248304285 0.674471707523801
0.905802074143823 0.662859643691898
0.936655976776786 0.650586826951521
0.966576200362029 0.637665491113303
0.995532919769989 0.624108516445816
1.02349727031302 0.60992941683624
1.05044137651833 0.595142326319429
1.0763383799149 0.579761984988794
1.10116246580654 0.563803724303061
1.12488888900457 0.547283451803534
1.14749399849434 0.53021763525712
1.16895526101111 0.512623286240892
1.18925128350167 0.494517943184587
1.20836183444941 0.47591965388791
1.22626786404155 0.456846957530103
1.24295152315835 0.437318866189681
1.2583961811656 0.417354845892782
1.2725864424923 0.39697479720901
1.28550816197735 0.376199035414109
1.29714845896976 0.355048270239262
1.30749573016833 0.333543585227174
1.31653966118811 0.311706416715542
1.32427123684197 0.289558532468845
1.33068275012716 0.267122009979764
1.33576780990779 0.244419214461859
1.3395213472857 0.221472776555434
1.34193962065318 0.198305569768825
1.34302021942276 0.174940687677591
1.34276206643008 0.15140142090433
1.34116541900767 0.127711233902087
1.33823186872838 0.10389374156447
1.33396433981895 0.0799726856858167
1.32836708624499 0.0559719112948541
1.3214456874706 0.0319153428854544
1.31320704289659 0.00782696056818344
1.30365936498306 -0.0162692238336006
1.29281217106298 -0.0403491907192499
1.28067627385515 -0.0643889366540895
1.26726377068584 -0.0883644982965646
1.2525880314299 -0.112251976285417
1.23666368518339 -0.136027559063091
1.21950660568094 -0.159667546611611
1.20113389547256 -0.183148374077274
1.18156386887532 -0.20644663526061
1.16081603371734 -0.229539105948189
1.13891107189193 -0.252402767063021
1.11587081874142 -0.275014827610478
1.09171824129122 -0.297352747396845
1.06647741535579 -0.319394259497883
1.04017350153935 -0.341117392454978
1.01283272015525 -0.362500492176772
0.984482325088956 -0.383522243524429
0.955150576630839 -0.404161691559032
0.924866713305652 -0.424398262429922
0.893660922726964 -0.444211783883155
0.861564311505473 -0.46358250536965
0.828608874241274 -0.482491117732961
0.794827461630932 -0.500918772457068
0.760253747721204 -0.518847100454981
0.724922196342015 -0.536258230379455
0.688868026752156 -0.553134806437528
0.652127178531968 -0.569460005691161
0.614736275757983 -0.585217554826703
0.576732590495252 -0.600391746376486
0.538154005643736 -0.61496745437637
0.499038977175808 -0.628930149443633
0.459426495802505 -0.642265913260175
0.419356048106737 -0.654961452446605
0.378867577182208 -0.667004111813361
0.338001442817263 -0.678381886975689
0.29679838126338 -0.689083436319865
0.255299464628386 -0.699098092308771
0.213546059934885 -0.708415872115521
0.171579787884716 -0.717027487574564
0.129442481370526 -0.724924354440328
};
\addplot [forestgreen4416044]
table {%
0.422611165574894 -0.689638743542854
0.389121969517002 -0.709688149151102
0.355244888217406 -0.729030122178225
0.321013691087492 -0.747645382124868
0.286462500529009 -0.765515372895076
0.251625757920111 -0.78262228129342
0.216538189283449 -0.798949054781574
0.181234770670512 -0.814479418476655
0.14575069329674 -0.829197891374383
0.110121328462154 -0.843089801780881
0.0743821922924765 -0.85614130193774
0.038568910335884 -0.868339381825765
0.00271718205068533 -0.879671882133642
-0.0331372547806738 -0.89012750637861
-0.0689586596758099 -0.899695832167036
-0.104711325079292 -0.908367321583686
-0.14035961195671 -0.916133330699324
-0.175867985320438 -0.922986118187172
-0.211201049651682 -0.92891885303963
-0.246323584183506 -0.933925621377573
-0.281200578009656 -0.938001432345438
-0.315797264984194 -0.941142223086211
-0.350079158377147 -0.943344862791378
-0.384012085251629 -0.944607155821778
-0.417562220528171 -0.944927843896267
-0.450696120702293 -0.944306607345998
-0.483380757181716 -0.942744065433074
-0.515583549209983 -0.940241775733257
-0.547272396343667 -0.936802232583337
-0.578415710450784 -0.932428864594728
-0.608982447198541 -0.927126031235753
-0.638942136998992 -0.920899018486027
-0.66826491538179 -0.913754033567283
-0.696921552763743 -0.905698198755873
-0.724883483585493 -0.896739544283135
-0.752122834786291 -0.886887000330674
-0.778612453588468 -0.876150388128574
-0.804325934563915 -0.864540410165373
-0.829237645955588 -0.852068639519595
-0.8533227552278 -0.838747508323451
-0.876557253819834 -0.824590295370221
-0.898917981078198 -0.809611112877662
-0.920382647343668 -0.793824892420646
-0.940929856170103 -0.777247370047038
-0.960539125652893 -0.759895070591661
-0.979190908845761 -0.741785291203984
-0.996866613245587 -0.722936084105946
-1.01354861932582 -0.703366238597108
-1.02922029810001 -0.683095262325073
-1.04386602769793 -0.662143361839834
-1.0574712089378 -0.640531422451449
-1.07002227987908 -0.618280987411109
-1.08150672934131 -0.595414236436358
-1.0919131093755 -0.571953963601866
-1.10123104667572 -0.547923554617799
-1.10945125291942 -0.523346963518438
-1.11656553402625 -0.498248688784274
-1.12256679832609 -0.472653748921391
-1.12744906362816 -0.446587657522473
-1.13120746318423 -0.420076397834303
-1.1338382505399 -0.393146396857097
-1.13533880326911 -0.365824499001493
-1.13570762558826 -0.338137939329459
-1.13494434984726 -0.310114316405796
-1.13304973689595 -0.281781564787281
-1.13002567532573 -0.253167927176899
-1.12587517958694 -0.224301926270898
-1.120602386984 -0.19521233632675
-1.11421255355124 -0.165928154480342
-1.10671204881359 -0.136478571841001
-1.09810834943729 -0.106892944393159
-1.088410031777 -0.0772007637336755
-1.07762676332667 -0.0474316276739667
-1.06576929308281 -0.0176152107362664
-1.05284944082962 0.0122187654265853
-1.0388800853568 0.0420405616583305
-1.0238751516217 0.0718204509439311
-1.00784959686858 0.101528748042104
-0.990819395718942 0.131135839076218
-0.972801524247654 0.160612211054043
-0.953813943060846 0.189928481286951
-0.933875579392381 0.219055426679214
-0.913006308236776 0.247964012858222
-0.891226932537377 0.276625423116577
-0.86855916244953 0.305011087137216
-0.845025593699432 0.333092709472922
-0.820649685060211 0.36084229775184
-0.795455734967718 0.388232190580887
-0.769468857299314 0.415235085119224
-0.742714956339812 0.441824064294326
-0.715220700959525 0.467972623633503
-0.687013498030149 0.493654697684136
-0.658121465105007 0.518844685996286
-0.628573402390851 0.543517478641782
-0.598398764039196 0.567648481244339
-0.567627628785778 0.591213639495768
-0.536290669967418 0.61418946313384
-0.504419124946172 0.636553049357881
-0.472044763971247 0.658282105658788
-0.43919985850973 0.679354972040684
-0.405917149077685 0.699750642612073
-0.372229812603695 0.719448786524973
-0.338171429357376 0.738429768241143
-0.303775949475839 0.756674667105216
-0.269077659121451 0.774165296205221
-0.234111146304643 0.790884220501691
-0.198911266405831 0.806814774207289
-0.163513107430812 0.821941077399631
-0.12795195503428 0.836248051850731
-0.0922632573463159 0.849721436057315
-0.0564825896369187 0.86234779945699
-0.0206456188538082 0.874114555816126
0.0152119319311694 0.885009975776085
0.0510543191315744 0.895023198545301
0.0868458142763573 0.904144242725554
0.122550739624777 0.912364016261644
0.158133503730748 0.919674325504547
0.193558636921168 0.926067883379031
0.228790826652858 0.93153831664756
0.263794952712878 0.936080172263282
0.298536122227112 0.939688922805747
0.332979704442252 0.942360970993931
0.367091365246475 0.944093653272089
0.400837101394446 0.944885242464845
0.43418327440248 0.944734949498875
0.467096644080124 0.943642924189479
0.499544401664693 0.941610255091231
0.531494202525764 0.938638968412896
0.562914198407006 0.934732025997653
0.593773069173215 0.929893322370669
0.6240400540309 0.924127680856948
0.653684982191322 0.917440848773342
0.682678302945381 0.909839491699493
0.710991115120414 0.901331186833446
0.738595195889508 0.891924415438525
0.765463028904627 0.881628554389034
0.79156783172551 0.870453866823175
0.816883582516986 0.85841149191254
0.84138504598811 0.845513433758331
0.865047798547244 0.831772549425419
0.887848252648026 0.81720253612614
0.909763680301954 0.801817917566613
0.930772235734132 0.785634029469199
0.950852977159619 0.768667004285514
0.969985887658654 0.750933755115254
0.988151895129954 0.73245195884685
1.00533289130221 0.713240038536755
1.02151174978478 0.693317145044951
1.03667234313968 0.672703137944946
1.05079955895775 0.651418565727324
1.06387931492303 0.629484645316559
1.07589857285031 0.60692324092153
1.08684535168191 0.583756842240795
1.09670873943062 0.560008542044373
1.10547890405702 0.53570201315437
1.11314710327029 0.510861484847393
1.11970569324269 0.485511718702282
1.12514813622908 0.45967798391723
1.12946900708396 0.433386032120891
1.13266399866927 0.406662071702606
1.13472992614793 0.379532741687301
1.13566473015849 0.352025085181137
1.13546747886798 0.324166522414345
1.13413836890079 0.29598482340815
1.13167872514263 0.267508080293006
1.1280909994199 0.238764679305753
1.12337876805561 0.209783272493597
1.11754672830447 0.180592749153126
1.11060069367054 0.151222207032831
1.1025475881122 0.121700923327837
1.09339543914021 0.0920583254957609
1.0831533698157 0.0623239619227776
1.07183158965612 0.0325274724691507
1.05944138445814 0.00269855892356946
1.04599510504778 -0.0271330446042332
1.03150615496876 -0.0569376013230953
1.01598897712164 -0.0866854014026983
0.999459039366739 -0.11634679158898
0.981932819105514 -0.145892204763136
0.963427786855472 -0.175292189414771
0.943962388835199 -0.204517438999804
0.92355602857677 -0.233538821153867
0.902229047583879 -0.262327406732084
0.880002705054993 -0.290854498646265
0.856899156691716 -0.319091660470792
0.832941432613502 -0.347010744788668
0.80815341440074 -0.374583921249468
0.782559811289062 -0.401783704311247
0.756186135538653 -0.428582980638716
0.729058677003066 -0.45495503613041
0.701204476922927 -0.480873582547886
0.672651300970641 -0.506312783720406
0.64342761157296 -0.531247281299
0.613562539539023 -0.555652220034214
0.583085855022129 -0.579503272552365
0.552027937844199 -0.602776663605601
0.520419747212516 -0.625449193771577
0.488292790858903 -0.647498262579151
0.455679093632139 -0.668901891037018
0.422611165574894 -0.689638743542854
};
\addplot [forestgreen4416044]
table {%
-0.587146672544441 0.444442865354003
-0.560980575804343 0.478402377835106
-0.534255281426272 0.511885008421855
-0.506997429764878 0.544857380899911
-0.47923419203949 0.577286627691318
-0.450993243249258 0.609140422617612
-0.422302734586111 0.640387013123268
-0.393191265373036 0.670995251927328
-0.363687854555649 0.700934628071695
-0.333821911775485 0.730175297335117
-0.30362320805382 0.758688111982558
-0.273121846115274 0.786444649820288
-0.242348230380763 0.813417242527753
-0.211333036659702 0.839579003237951
-0.180107181571702 0.864903853338849
-0.148701791728202 0.88936654846911
-0.117148172704793 0.912942703682223
-0.0854777778351341 0.935608817753944
-0.0537221768575883 0.957342296608833
-0.0219130244458194 0.978121475842515
0.00991797134527378 0.997925642317231
0.0417390806868376 1.01673505480914
0.0735185836050454 1.03453096368682
0.105224801600259 1.05129562960128
0.136826129224846 1.06701234116894
0.168291065588174 1.08166543162991
0.199588245757388 1.09524029446494
0.230686472022655 1.10772339795555
0.261554744995721 1.11910229867269
0.292162294510773 1.12936565388067
0.322478610296813 1.1385032328438
0.352473472390957 1.14650592702467
0.38211698126235 1.15336575916366
0.411379587616674 1.15907589123093
0.440232121851522 1.16363063124263
0.468645823133292 1.16702543893488
0.496592368066617 1.16925693028953
0.524043898927735 1.17032288090748
0.550973051433677 1.17022222822596
0.577352982019579 1.16895507257776
0.603157394596932 1.16652267709119
0.628360566766094 1.16292746643096
0.652937375456936 1.15817302438125
0.676863321972068 1.1522640902733
0.700114556407669 1.14520655426112
0.722667901427595 1.13700745145009
0.744500875367048 1.12767495488423
0.765591714642789 1.11721836739904
0.785919395447551 1.10564811234836
0.805463654707029 1.09297572321404
0.824205010278551 1.07921383210917
0.842124780371304 1.06437615718608
0.859205102168746 1.04847748896182
0.875428949634651 1.03153367557464
0.890780150485037 1.01356160698622
0.905243402309042 0.99457919814536
0.918804287822706 0.97460537113003
0.931449289240432 0.953660036285379
0.943165801749808 0.931764072376691
0.953942146076358 0.908939305776978
0.963767580125696 0.885208488709969
0.972632309691483 0.860595276570208
0.9805274982185 0.835124204342835
0.987445275611121 0.808820662146582
0.993378746078397 0.781710869924351
0.998321995007925 0.753821851306602
1.00227009486167 0.725181406673616
1.00521911008784 0.695818085443464
1.00716610104391 0.665761157613332
1.00810912692696 0.635040584582549
1.00804724770829 0.603686989286412
1.00698052507044 0.57173162567058
1.00491002234574 0.539206347536458
1.00183780345633 0.50614357678864
0.997766930856814 0.472576271116044
0.992701462481531 0.438537891138976
0.986646447699521 0.40406236705485
0.979607922281199 0.369184064815836
0.971592902381775 0.333937751872127
0.962609377547392 0.298358562514989
0.952666302750975 0.262481962854139
0.941773589465721 0.226343715464349
0.929942095785123 0.18997984373654
0.917183615599388 0.153426595968878
0.903510866839033 0.116720409233686
0.888937478797374 0.0798978730561752
0.873477978544545 0.042995692941211
0.857147776446608 0.00605065378446773
0.839963150804156 -0.0309004167955557
0.821941231625756 -0.0678206851681061
0.80309998355238 -0.10467334840674
0.783458187949859 -0.141421670975282
0.763035424187214 -0.178029021346607
0.74185205011951 -0.21445890851774
0.719929181794705 -0.250675018384885
0.697288672404718 -0.286641249942111
0.673953090501687 -0.32232175126762
0.649945697501146 -0.357680955261724
0.625290424494542 -0.392683615100902
0.600011848394203 -0.427294839372601
0.574135167434536 -0.461480126855754
0.547686176053878 -0.495205400912347
0.520691239182041 -0.528437043455756
0.493177265959175 -0.561141928461983
0.465171682912153 -0.593287454990399
0.436702406615213 -0.624841579681072
0.407797815862106 -0.655772848696281
0.378486723377503 -0.686050429074389
0.348798347095839 -0.71564413946481
0.318762281036245 -0.744524480213433
0.288408465802582 -0.772662662768519
0.257767158738001 -0.800030638377758
0.226868903763759 -0.826601126047874
0.195744500932379 -0.852347639738911
0.164424975725486 -0.877244514766097
0.132941548126933 -0.901266933382955
0.101325601502041 -0.924390949520176
0.0696086513139768 -0.946593512655583
0.0378223137084586 -0.967852490791388
0.00599827399808911 -0.988146692515852
-0.0258317449222498 -1.00745588812735
-0.0576360141974711 -1.02576082979975
-0.0893828306402917 -1.04304327076913
-0.121040548333663 -1.05928598352246
-0.152577610176108 -1.07447277697043
-0.183962579338525 -1.08858851258705
-0.215164170601103 -1.10161911950008
-0.246151281539109 -1.11355160851717
-0.276893023526466 -1.12437408507375
-0.307358752526198 -1.13407576108984
-0.33751809963708 -1.14264696572378
-0.367341001366018 -1.1500791550124
-0.396797729595989 -1.15636492038779
-0.425858921219681 -1.16149799606231
-0.454495607409269 -1.1654732652745
-0.482679242493178 -1.16828676538955
-0.510381732411031 -1.16993569184937
-0.53757546271841 -1.17041840096821
-0.564233326113549 -1.1697344115711
-0.590328749458467 -1.16788440547356
-0.61583572026766 -1.1648702268019
-0.640728812637909 -1.16069488015497
-0.664983212593373 -1.15536252760912
-0.6885747428207 -1.14887848456932
-0.7114798867695 -1.14124921447069
-0.73367581209415 -1.13248232233559
-0.755140393413578 -1.12258654719275
-0.77585223436632 -1.11157175336605
-0.795790688938882 -1.09944892064152
-0.81493588204613 -1.08623013332245
-0.833268729343212 -1.07192856818354
-0.850770956249246 -1.05655848133593
-0.86742511616382 -1.04013519401646
-0.883214607858141 -1.02267507731509
-0.898123692023505 -1.00419553585588
-0.912137506960579 -0.984714990447624
-0.925242083393865 -0.964252859721654
-0.937424358396579 -0.942829540774878
-0.948672188412066 -0.920466388837528
-0.95897436135875 -0.897185695985802
-0.968320607806591 -0.873010668920633
-0.97670161121387 -0.847965405834753
-0.984109017214123 -0.822074872391079
-0.990535441943957 -0.795364876836407
-0.995974479403441 -0.767862044275187
-1.00042070784175 -0.739593790129035
-1.00386969516168 -0.710588292808453
-1.00631800333766 -0.680874465623974
-1.00776319184286 -0.650481927964739
-1.00820382008195 -0.619440975773255
-1.00763944882711 -0.587782551345732
-1.00607064065589 -0.555538212488137
-1.0034989593904 -0.522740101058682
-0.999926968538429 -0.489420910928121
-0.995358228738159 -0.455613855389792
-0.989797294208779 -0.421352634051878
-0.983249708210764 -0.386671399244912
-0.97572199752023 -0.351604721977981
-0.967221665922906 -0.316187557477606
-0.95775718673419 -0.280455210343601
-0.947337994352766 -0.244443299356682
-0.935974474856194 -0.208187721972898
-0.923677955647838 -0.171724618540255
-0.910460694165465 -0.135090336273244
-0.896335865662765 -0.09832139302114
-0.881317550075971 -0.0614544408662115
-0.865420717988674 -0.0245262295881388
-0.848661215708817 0.0124264299689763
-0.831055749472752 0.0493667025904527
-0.812621868792101 0.086257765409176
-0.79337794896002 0.123062844611437
-0.773343172734308 0.159745252093859
-0.752537511215615 0.196268422034912
-0.730981703939819 0.232595947344503
-0.708697238204399 0.268691615955363
-0.685706327649436 0.304519446920029
-0.66203189011457 0.340043726277421
-0.637697524793997 0.375229042653307
-0.612727488712285 0.410040322559126
-0.587146672544442 0.444442865354003
};
\addplot [forestgreen4416044]
table {%
-0.53327627941799 0.3028102983899
-0.507592116898936 0.347423422518515
-0.48140197551882 0.391690227423921
-0.454731962180144 0.435566586992287
-0.427608662132214 0.479008764314188
-0.400059112470387 0.521973455282505
-0.372110775184902 0.564417831758899
-0.343791509786177 0.606299584265833
-0.315129545533839 0.647576964161579
-0.286153453297185 0.688208825256169
-0.256892117075118 0.728154664826817
-0.227374705203951 0.767374663991912
-0.19763064128177 0.805829727403342
-0.167689574838357 0.843481522217591
-0.137581351779891 0.880292516306751
-0.107335984637899 0.91622601567136
-0.0769836226521137 0.951246201017784
-0.0465545217170562 0.985318163463661
-0.0160790142222995 1.01840793933584
0.0144125211835147 1.05048254402611
0.0448896898747747 1.08151000487096
0.0753221115469223 1.11145939302265
0.105679450500201 1.14030085427975
0.13593144587894 1.16800563884643
0.166047941836238 1.19454612999094
0.195998917593961 1.21989587157449
0.22575451736811 1.24402959442338
0.255285080129715 1.26692324151783
0.284561169171596 1.28855399197253
0.313553601451509 1.30890028378501
0.34223347668244 1.3279418353291
0.37057220614104 1.34565966557212
0.398541541165483 1.36203611299555
0.426113601314348 1.37705485320047
0.453260902158454 1.39070091518
0.479956382677927 1.4029606962428
0.506173432237223 1.41382197557249
0.53188591711118 1.42327392640961
0.557068206535691 1.43130712684403
0.581695198257001 1.43791356920683
0.605742343554184 1.44308666805261
0.629185671709838 1.44682126672393
0.65200181390462 1.4491136424916
0.674168026511793 1.44996151026562
0.695662213768568 1.44936402487295
0.716462949801638 1.44732178190003
0.736549499984958 1.44383681709905
0.755901841608471 1.43891260435873
0.774500683837178 1.43255405224141
0.792327486940664 1.42476749909013
0.809364480773905 1.4155607067104
0.825594682490933 1.40494285263311
0.841001913473708 1.3929245209661
0.855570815459315 1.37951769184376
0.869286865849415 1.36473572948491
0.882136392186688 1.34859336887112
0.894106585783834 1.33110670105849
0.905185514491551 1.31229315713781
0.915362134592763 1.29217149085879
0.924626301811235 1.27076175993599
0.932968781423611 1.2480853060548
0.940381257464782 1.22416473359761
0.946856341017428 1.19902388711126
0.952387577577445 1.17268782753828
0.956969453487938 1.1451828072356
0.960597401435356 1.11653624380564
0.963267805002285 1.0867766927658
0.964978002272374 1.05593381908378
0.965726288483791 1.02403836760681
0.965511917728569 0.991122132414579
0.964335103696136 0.957217925126179
0.962197019460315 0.92235954219277
0.959099796309969 0.886581731208549
0.955046521624493 0.849920156273587
0.950041235796241 0.812411362443084
0.944088928202981 0.774092739298458
0.937195532234371 0.735002483676599
0.929367919377426 0.695179561594433
0.920613892366875 0.654663669406746
0.910942177407223 0.613495194235996
0.900362415474276 0.571715173713552
0.888885152704817 0.529365255072487
0.876521829883975 0.486487653632714
0.863284771040813 0.443125110719835
0.849187171163461 0.399320851059661
0.834243083046065 0.355118539690858
0.81846740328066 0.310562238438698
0.801875857407921 0.265696361993262
0.7844849842416 0.220565633635919
0.766312119382279 0.175215040658183
0.747375377936869 0.129689789517403
0.727693636461072 0.0840352607739838
0.707286514142822 0.0382969638550555
0.686174353245453 -0.00747950831031258
0.664378198830087 -0.0532485247391359
0.641919777777462 -0.0989644618804682
0.6188214771301 -0.144581749093953
0.59510632177642 -0.190054914075628
0.570797951499028 -0.235338628185704
0.545920597410066 -0.280387751633142
0.520499057797114 -0.325157378471976
0.494558673403718 -0.36960288136454
0.468125302169182 -0.413679956066965
0.441225293452819 -0.457344665592614
0.413885461768326 -0.500553484009428
0.386133060054506 -0.543263339827518
0.357995752508936 -0.585431658933765
0.329501587011698 -0.627016407030623
0.300678967166641 -0.667976131536818
0.27155662398805 -0.708270002908184
0.242163587260949 -0.747857855337442
0.212529156603577 -0.786700226792346
0.182682872260892 -0.824758398352301
0.152654485658214 -0.861994432804219
0.122473929744355 -0.898371212459163
0.0921712891538056 -0.933852476152062
0.0617767702177193 -0.968402855387634
0.0313206708535802 -1.00198790959647
0.000833350363583285 -1.03457416046613
-0.0296548008281796 -1.06612912531308
-0.0601133914695544 -1.09662134946211
-0.0905120597749879 -1.12602043760106
-0.120820503690812 -1.15429708407951
-0.151008511100996 -1.18142310212129
-0.181045989943232 -1.2073714519217
-0.210902998205369 -1.23211626760134
-0.240549773772249 -1.2556328829898
-0.26995676409324 -1.27789785621343
-0.299094655640857 -1.2988889930627
-0.327934403131126 -1.31858536911591
-0.356447258476548 -1.33696735059704
-0.384604799442831 -1.3540166139472
-0.412378957980777 -1.36971616408989
-0.439742048205134 -1.3840503513721
-0.466666793992484 -1.39700488716423
-0.493126356170673 -1.40856685810325
-0.519094359272692 -1.41872473896509
-0.544544917828303 -1.42746840415319
-0.56945266216725 -1.43478913779195
-0.59379276370829 -1.44067964241487
-0.617540959708867 -1.44513404623882
-0.64067357745074 -1.44814790901721
-0.663167557837463 -1.44971822646607
-0.685000478380197 -1.44984343325883
-0.706150575548929 -1.44852340458662
-0.726596766466836 -1.44575945628274
-0.746318669926155 -1.44155434351096
-0.765296626704613 -1.43591225801914
-0.783511719162167 -1.42883882396081
-0.800945790098519 -1.42034109228887
-0.81758146085261 -1.41042753372704
-0.833402148626046 -1.39910803032612
-0.848392083013191 -1.38639386561324
-0.862536321721453 -1.37229771334431
-0.875820765466082 -1.35683362487047
-0.88823217202465 -1.34001701513141
-0.899758169437178 -1.32186464728944
-0.910387268338783 -1.3023946160196
-0.920108873412523 -1.28162632947249
-0.928913293951037 -1.25958048992779
-0.936791753516456 -1.23627907315782
-0.94373639868895 -1.21174530652155
-0.949740306895179 -1.18600364581113
-0.954797493308872 -1.15907975087377
-0.958902916816632 -1.13100046003347
-0.962052485043021 -1.10179376333801
-0.964243058429936 -1.07148877465782
-0.965472453366185 -1.04011570266469
-0.96573944436416 -1.00770582071907
-0.965043765281426 -0.974291435696124
-0.963386109586023 -0.93990585578152
-0.960768129665199 -0.904583357269128
-0.957192435178273 -0.868359150393646
-0.952662590455274 -0.831269344232277
-0.947183110943933 -0.793350910710422
-0.940759458708594 -0.75464164774725
-0.933398036985505 -0.715180141577921
-0.925106183799932 -0.675005728289976
-0.91589216465146 -0.634158454612282
-0.905765164274759 -0.592679037995574
-0.894735277484043 -0.550608826024412
-0.882813499110339 -0.507989755201022
-0.8700117130416 -0.464864309142061
-0.856342680376589 -0.42127547623004
-0.841820026704336 -0.377266706761561
-0.826458228521855 -0.332881869635116
-0.810272598803657 -0.288165208621621
-0.793279271737444 -0.243161298261254
-0.775495186641201 -0.197914999430593
-0.756938071077715 -0.152471414624313
-0.737626423183352 -0.10687584299604
-0.717579493228718 -0.0611737352031705
-0.696817264429563 -0.0154106481006644
-0.675360433027085 0.030367800671014
-0.653230387657462 0.0761159781585538
-0.630449188031194 0.121788281583723
-0.607039542943507 0.167339183801198
-0.583024787637729 0.212723278681028
-0.558428860544213 0.257895326370436
-0.53327627941799 0.302810298389899
};
\addplot [line width=0.52pt, gray127, dashed]
table {%
0 0.5
0.0631370995296211 0.499750794563087
0.126211262625347 0.499003426665747
0.18915961558969 0.497758641300553
0.251919410135435 0.496017679296629
0.314428085934502 0.493782276082762
0.376623332979436 0.49105465995749
0.438443153695383 0.487837549867887
0.499825924740617 0.484134152699254
0.560710458434029 0.479948160078428
0.621036063748338 0.475283744693891
0.680742606808226 0.470145556136338
0.739770570833094 0.464538716263865
0.798061115464682 0.458468814096394
0.855556135420419 0.451941900244412
0.912198318414032 0.44496448087759
0.967931202285678 0.437543511239297
1.02269923128465 0.429686388713456
1.07644781144858 0.421400945450675
1.12912336502384 0.412695440560988
1.18067338387306 0.403578551880993
1.23104648181636 0.394059367323596
1.28019244585421 0.384147375818985
1.32806228622086 0.373852457855855
1.37460828521837 0.36318487563232
1.41978404478267 0.352155262826336
1.46354453273415 0.340774613995817
1.50584612766675 0.329054273619019
1.54664666243068 0.317005924786119
1.58590546616556 0.304641577553258
1.62358340484204 0.291973556970653
1.65964292027145 0.27901449079672
1.6940480675446 0.265777296910451
1.72676455086245 0.252275170434589
1.75775975772292 0.238521570582445
1.78700279142975 0.22453020724146
1.81446450189096 0.210315027306892
1.84011751467635 0.195890200779247
1.86393625830487 0.181270106639312
1.88589698973487 0.166469318514881
1.90597781803168 0.151502590153436
1.92415872618893 0.136384840715302
1.94042159108197 0.121131139901892
1.95475020153341 0.10575669293391
1.96713027447274 0.0902768253944512
1.97754946917401 0.0747069679521305
1.98599739955734 0.059062640979454
1.99246564454201 0.0433594390817754
1.99694775644076 0.0276130155522541
1.99943926738696 0.011839066768312
1.99993769378831 -0.00394668345485666
1.99844253880255 -0.0197284995381263
1.99495529283268 -0.035490649824008
1.98947943204131 -0.0512174222583066
1.98202041488558 -0.0668931400522395
1.97258567667601 -0.0825021773093997
1.96118462216481 -0.098028974601989
1.94782861617108 -0.113458054480795
1.93253097225204 -0.128774036903448
1.91530693943186 -0.143961654565583
1.89617368700102 -0.159005768119619
1.87515028740165 -0.173891381265991
1.85225769721568 -0.188603655701788
1.82751873627487 -0.203127925911895
1.8009580649135 -0.217449713787896
1.77260215938642 -0.231554743060174
1.74247928547692 -0.245428953528797
1.71061947032082 -0.259058515079039
1.67705447247475 -0.272429841467535
1.64181775025853 -0.285529603865347
1.60494442840316 -0.298344744144428
1.56647126303778 -0.310862487894248
1.52643660505034 -0.323070357155605
1.48488036185857 -0.334956182858928
1.44184395762943 -0.346508116954666
1.39737029198661 -0.357714644223686
1.35150369724722 -0.368564593755889
1.30428989423037 -0.379047150085623
1.25577594668172 -0.389151863972765
1.20601021435923 -0.39886866281876
1.15504230482718 -0.4081878607072
1.10292302400622 -0.417100168058959
1.04970432552894 -0.425596700892247
0.995439258951366 -0.433668989678357
0.940181916872006 -0.441308987784273
0.883987381011158 -0.448509079493732
0.826911667304268 -0.455262087598731
0.769011670064022 -0.46156128055393
0.710345105266857 -0.467400379186799
0.650970453020423 -0.472773562956834
0.59094689926934 -0.477675475757597
0.530334276797357 -0.482101231255806
0.469193005584738 -0.486046417762128
0.407584032580305 -0.489507102628857
0.345568770948199 -0.492479836170055
0.283209038849903 -0.494961655100279
0.22056699782255 -0.496950085488444
0.157705090814953 -0.498443145223897
0.0946859799431164 -0.499439345992222
0.0315724840272736 -0.499937694758829
-0.031572484027274 -0.499937694758829
-0.0946859799431159 -0.499439345992222
-0.157705090814953 -0.498443145223897
-0.22056699782255 -0.496950085488444
-0.283209038849904 -0.494961655100279
-0.345568770948199 -0.492479836170055
-0.407584032580304 -0.489507102628857
-0.469193005584737 -0.486046417762128
-0.530334276797357 -0.482101231255806
-0.59094689926934 -0.477675475757597
-0.650970453020424 -0.472773562956834
-0.710345105266857 -0.467400379186799
-0.769011670064021 -0.46156128055393
-0.826911667304268 -0.455262087598731
-0.883987381011158 -0.448509079493732
-0.940181916872006 -0.441308987784273
-0.995439258951366 -0.433668989678357
-1.04970432552894 -0.425596700892247
-1.10292302400621 -0.417100168058959
-1.15504230482718 -0.4081878607072
-1.20601021435923 -0.39886866281876
-1.25577594668172 -0.389151863972765
-1.30428989423037 -0.379047150085623
-1.35150369724722 -0.36856459375589
-1.39737029198661 -0.357714644223686
-1.44184395762943 -0.346508116954666
-1.48488036185857 -0.334956182858928
-1.52643660505034 -0.323070357155605
-1.56647126303778 -0.310862487894248
-1.60494442840316 -0.298344744144428
-1.64181775025852 -0.285529603865347
-1.67705447247475 -0.272429841467535
-1.71061947032082 -0.259058515079039
-1.74247928547692 -0.245428953528797
-1.77260215938642 -0.231554743060174
-1.8009580649135 -0.217449713787897
-1.82751873627487 -0.203127925911894
-1.85225769721568 -0.188603655701788
-1.87515028740165 -0.173891381265991
-1.89617368700102 -0.159005768119619
-1.91530693943186 -0.143961654565583
-1.93253097225204 -0.128774036903448
-1.94782861617108 -0.113458054480795
-1.96118462216481 -0.098028974601989
-1.97258567667601 -0.0825021773093994
-1.98202041488558 -0.0668931400522396
-1.98947943204131 -0.0512174222583065
-1.99495529283268 -0.0354906498240081
-1.99844253880255 -0.0197284995381264
-1.99993769378831 -0.00394668345485689
-1.99943926738696 0.0118390667683119
-1.99694775644076 0.0276130155522542
-1.99246564454201 0.0433594390817752
-1.98599739955734 0.059062640979454
-1.97754946917401 0.0747069679521302
-1.96713027447274 0.090276825394451
-1.95475020153341 0.10575669293391
-1.94042159108197 0.121131139901891
-1.92415872618893 0.136384840715302
-1.90597781803168 0.151502590153436
-1.88589698973487 0.166469318514881
-1.86393625830487 0.181270106639312
-1.84011751467635 0.195890200779246
-1.81446450189096 0.210315027306892
-1.78700279142975 0.22453020724146
-1.75775975772292 0.238521570582445
-1.72676455086245 0.252275170434589
-1.6940480675446 0.265777296910451
-1.65964292027145 0.27901449079672
-1.62358340484204 0.291973556970653
-1.58590546616556 0.304641577553258
-1.54664666243068 0.317005924786119
-1.50584612766675 0.329054273619019
-1.46354453273415 0.340774613995817
-1.41978404478267 0.352155262826336
-1.37460828521837 0.36318487563232
-1.32806228622086 0.373852457855854
-1.28019244585421 0.384147375818985
-1.23104648181636 0.394059367323596
-1.18067338387306 0.403578551880993
-1.12912336502384 0.412695440560988
-1.07644781144858 0.421400945450675
-1.02269923128465 0.429686388713456
-0.967931202285677 0.437543511239297
-0.912198318414033 0.44496448087759
-0.855556135420419 0.451941900244412
-0.798061115464683 0.458468814096394
-0.739770570833094 0.464538716263865
-0.680742606808225 0.470145556136338
-0.621036063748338 0.475283744693891
-0.560710458434029 0.479948160078428
-0.499825924740618 0.484134152699254
-0.438443153695383 0.487837549867887
-0.376623332979435 0.49105465995749
-0.314428085934502 0.493782276082762
-0.251919410135435 0.496017679296629
-0.189159615589691 0.497758641300553
-0.126211262625347 0.499003426665747
-0.0631370995296205 0.499750794563087
-4.89858719658941e-16 0.5
};
\addplot [line width=0.52pt, gray127, dashed]
table {%
-0.445503262094184 0.226995249869773
-0.416617575056473 0.283137680568154
-0.387316594559491 0.338997873468965
-0.357629528457826 0.394520145917093
-0.327585969464592 0.449649152103885
-0.297215865652769 0.504329938237046
-0.266549490602316 0.55850799731978
-0.235617413222806 0.612129323484548
-0.204450467281669 0.66514046582729
-0.173079720668428 0.717488581688452
-0.141536444425547 0.769121489327695
-0.109852081576777 0.819987719939789
-0.0780582157840669 0.870036568959839
-0.04618653986428 0.919218146606698
-0.0142688241971077 0.967483427614182
0.0176631149443378 1.01478430010053
0.0495774471086754 1.06107361352735
0.0814423593955421 1.10630522570036
0.113226088167383 1.15043404876484
0.144896950712137 1.19341609415028
0.176423376825253 1.23520851641907
0.207773940279555 1.2757696559758
0.238917390151591 1.31505908059447
0.269822681973231 1.35303762572215
0.300459008677471 1.38966743351911
0.330795831307583 1.42491199059635
0.360802909459015 1.45873616441286
0.390450331423681 1.4911062382966
0.419708544006607 1.52198994505395
0.448548381985191 1.55135649913435
0.476941097181738 1.57917662731806
0.504858387120265 1.60542259789628
0.532272423239025 1.63006824831474
0.559155878630626 1.65308901125312
0.585481955282079 1.67446193911427
0.61122441078765 1.69416572689891
0.636357584507851 1.71218073344284
0.660856423148531 1.72848900099576
0.684696505734541 1.74307427312192
0.707854067953093 1.75592201090485
0.730306025842545 1.76701940744013
0.752029998802996 1.77635540060155
0.773004331905752 1.78392068406809
0.793208117479436 1.78970771660067
0.812621215951201 1.79371072955939
0.831224275922303 1.7959257326539
0.848998753457988 1.79635051792092
0.865926930572497 1.79498466192528
0.881991932890734 1.79182952618192
0.89717774646901 1.78688825579878
0.911469233758097 1.78016577634164
0.924852148692657 1.77166878892421
0.937313150892037 1.76140576352835
0.948839818958248 1.74938693056096
0.959420662857882 1.73562427065613
0.969045135375633 1.72013150273255
0.977703642627984 1.70292407031823
0.985387553626605 1.68401912615602
0.992089208881904 1.66343551510539
0.997801928038179 1.64119375535748
1.00252001653274 1.61731601798212
1.00623877127236 1.59182610482722
1.00895448532144 1.56474942479254
1.01066445159716 1.5361129685016
1.01136696556796 1.5059452813968
1.01106132695262 1.47427643528474
1.00974784041836 1.44113799835993
1.00742781527712 1.40656300373699
1.00410356418043 1.37058591652248
0.999778400814071 1.33324259945935
0.994456636594945 1.29457027717819
0.988143576373364 1.25460749909087
0.98084551314505 1.21339410096369
0.972569721778146 1.17097116520819
0.963324451761466 1.12738097992934
0.953118918981224 1.08266699677181
0.941963296534443 1.03687378760643
0.929868704588185 0.990047000099997
0.916847199294733 0.942233312212662
0.90291176077375 0.893480385668352
0.888076280173416 0.843836818444507
0.872355545823431 0.793352096328567
0.855765228493684 0.74207654358946
0.838321865773289 0.69006127281328
0.82004284558556 0.637358133953148
0.800946388855343 0.584019662644048
0.781051531346012 0.530099027834167
0.760378104684199 0.475649978784933
0.738946716591197 0.420726791492575
0.71677873034074 0.365384214584636
0.693896243463625 0.309677414745351
0.670322065720412 0.253661921724291
0.646079696364157 0.197393572983106
0.621193300715849 0.140928458035523
0.595687686075892 0.0843228625361076
0.569588276995645 0.0276332121734922
0.542921089933682 -0.0290839835759753
0.515712707322013 -0.0857721877781059
0.487990251068145 -0.142374892398115
0.45978135551937 -0.198835674629064
0.431114139916248 -0.255098253135332
0.402017180362737 -0.311106544155086
0.372519481340911 -0.36680471740579
0.34265044679866 -0.422137251737043
0.312439850839196 -0.47704899047527
0.281917808041584 -0.531485196405095
0.251114743441865 -0.585391606332588
0.220061362204726 -0.638714485175996
0.188788619015925 -0.691400679530039
0.157327687225979 -0.743397670650386
0.125709927775897 -0.794653626805474
0.0939668579359136 -0.845117454943511
0.0621301198883808 -0.89473885162314
0.0302314491861589 -0.943468353157002
-0.00169735688207456 -0.991257384918214
-0.033624470988055 -1.03805830976062
-0.0655180674901009 -1.08382447550452
-0.0973463541576085 -1.12851026144061
-0.129077603862237 -1.17207112380564
-0.160680186204201 -1.21446364018466
-0.192122599042139 -1.25564555279542
-0.223373499895138 -1.29557581061181
-0.254401737185591 -1.33421461028444
-0.285176381291777 -1.37152343581751
-0.315666755379174 -1.40746509696234
-0.345842465979799 -1.44200376628946
-0.375673433289074 -1.47510501490214
-0.40512992115003 -1.5067358467559
-0.434182566694959 -1.53686473154964
-0.462802409614956 -1.56546163615578
-0.490960921028182 -1.59249805455788
-0.51863003191808 -1.6179470362661
-0.545782161113175 -1.64178321318202
-0.572390242780594 -1.66398282488611
-0.598427753405864 -1.68452374232264
-0.62386873823214 -1.70338548985838
-0.648687837132466 -1.72054926569319
-0.6728603098893 -1.73599796060206
-0.696362060856105 -1.74971617498991
-0.719169662976409 -1.76169023424232
-0.741260381136405 -1.7719082023566
-0.762612194827812 -1.78035989383994
-0.783203820098389 -1.78703688386244
-0.803014730768248 -1.79193251665526
-0.822025178890792 -1.79504191214514
-0.840216214437899 -1.79636197081906
-0.857569704189725 -1.7958913768138
-0.874068349810286 -1.79363059922772
-0.889695705090822 -1.78958189165307
-0.904436192343729 -1.78374928992961
-0.918275117930746 -1.77613860812159
-0.93119868690988 -1.76675743272216
-0.943194016786517 -1.75561511509095
-0.954249150354966 -1.74272276213252
-0.964353067617662 -1.72809322522467
-0.973495696770141 -1.71174108740794
-0.981667924240834 -1.69368264884893
-0.988861603775671 -1.67393591059193
-0.995069564558446 -1.65252055661507
-1.00028561835884 -1.62945793420894
-1.00450456570099 -1.6047710326971
-1.00772220104644 -1.57848446051987
-1.00993531698629 -1.55062442070406
-1.01114170743844 -1.52121868474325
-1.01134016984666 -1.49029656491456
-1.01053050637931 -1.4578888850595
-1.00871352412653 -1.42402794985816
-1.00589103429575 -1.38874751262713
-1.00206585040621 -1.35208274167352
-0.99724178548441 -1.31407018523838
-0.99142364826317 -1.27474773506466
-0.984617238388205 -1.23415458862585
-0.976829340636902 -1.19233121005314
-0.968067718155106 -1.14931928979984
-0.958341104718653 -1.10516170308347
-0.947659196027347 -1.05990246714679
-0.936032640040078 -1.01358669738047
-0.923473026360705 -0.966260562351073
-0.909992874685284 -0.917971237779224
-0.895605622322169 -0.868766859513838
-0.880325610797404 -0.818696475549247
-0.86416807155878 -0.767809997133104
-0.847149110792799 -0.716158149013753
-0.829285693369675 -0.663792418876705
-0.810595625932384 -0.610765006020596
-0.791097539146611 -0.55712876932378
-0.770810869129303 -0.50293717455346
-0.749755838074319 -0.448244241069841
-0.727953434094514 -0.393104487978463
-0.705425390300325 -0.33757287978437
-0.682194163135738 -0.281704771602289
-0.658282909993217 -0.225555853977459
-0.633715466129905 -0.169182097372072
-0.608516320908132 -0.112639696372711
-0.582710593383876 -0.0559850136743585
-0.556324007267548 0.000725476103168893
-0.529382865282034 0.0574352427103375
-0.501914022943573 0.114087756618492
-0.473944861791589 0.170626545369709
-0.445503262094184 0.226995249869773
};

\nextgroupplot[
height=4.6cm,
label style={font=\footnotesize},
legend style={draw=none, fill=white, fill opacity=0.9, text opacity=1, font=\footnotesize},
tick align=outside,
tick label style={font=\scriptsize},
tick pos=left,
title style={font=\footnotesize, yshift=-2pt},
title={log-euclidienne},
width=0.45\textwidth,
x grid style={darkgray176},
xlabel={\(\displaystyle x_1\)},
xmin=-2.29992834785656, xmax=2.29992834785656,
xtick style={color=black},
y grid style={darkgray176},
ylabel={\(\displaystyle x_2\)},
ymin=-2.29992834785656, ymax=2.29992834785656,
ytick style={color=black}
]
\addplot [crimson2143940]
table {%
0.0539286511369935 -0.593135136454357
0.00111494539502881 -0.597638961374307
-0.0516998717487531 -0.601547046780314
-0.104463153336025 -0.604855497007855
-0.157122303780223 -0.607561014121791
-0.209624831294972 -0.609660901203825
-0.261918400219038 -0.611153065040845
-0.313950883185655 -0.612036018211488
-0.365670413084218 -0.612308880568832
-0.417025434762546 -0.611971380117743
-0.467964756418178 -0.611023853286013
-0.518437600627478 -0.609467244588992
-0.568393654961671 -0.607303105688083
-0.617783122139371 -0.604533593844008
-0.66655676966559 -0.601161469766402
-0.714665978907761 -0.597190094861883
-0.76206279355985 -0.592623427883322
-0.808699967446241 -0.587466020983679
-0.854531011617756 -0.581723015178313
-0.899510240692847 -0.575400135220307
-0.943592818397779 -0.568503683893911
-0.986734802260404 -0.561040535731789
-1.02889318741297 -0.553018130162335
-1.07002594946031 -0.544444464093889
-1.11009208637068 -0.535328083943246
-1.14905165934745 -0.525678077116398
-1.18686583264102 -0.515504062950017
-1.22349691226105 -0.504816183122685
-1.25890838355077 -0.493625091545448
-1.29306494758545 -0.481941943741772
-1.32593255635926 -0.469778385727468
-1.35747844672497 -0.457146542401695
-1.38767117305301 -0.444059005460591
-1.41648063857716 -0.430528820845603
-1.44387812539565 -0.416569475739005
-1.46983632309779 -0.402194885119582
-1.49432935598757 -0.387419377891877
-1.51733280887715 -0.372257682602828
-1.53882375142441 -0.356724912760029
-1.55878076099045 -0.340836551766262
-1.57718394399411 -0.324608437485299
-1.59401495574235 -0.30805674645438
-1.60925701871663 -0.291197977759084
-1.62289493929708 -0.274048936586685
-1.63491512290789 -0.256626717474375
-1.64530558756862 -0.238948687269059
-1.65405597583812 -0.221032467815699
-1.66115756513907 -0.202895918391484
-1.66660327645282 -0.184557117903306
-1.67038768137589 -0.166034346866318
-1.67250700753115 -0.147346069181516
-1.67295914232821 -0.128510913730523
-1.67174363506926 -0.109547655805913
-1.66886169739837 -0.0904751983955907
-1.66431620209369 -0.071312553339885
-1.65811168020379 -0.0520788223801293
-1.65025431653103 -0.0327931781176331
-1.6407519434664 -0.0134748449020145
-1.62961403318201 0.00585692033204868
-1.61685168818907 0.0251828472605505
-1.60247763027061 0.0444836713792352
-1.58650618780017 0.0637401532068738
-1.56895328145895 0.0829330974635969
-1.54983640836572 0.102043372205168
-1.52917462463533 0.121051927894122
-1.50698852638314 0.139939816388755
-1.48330022919439 0.158688209831049
-1.45813334607888 0.177278419414688
-1.43151296393302 0.195691914014465
-1.40346561853267 0.213910338658513
-1.37401926808168 0.231915532824938
-1.34320326534258 0.249689548544624
-1.31104832837706 0.26721466829216
-1.27758650992559 0.284473422647051
-1.24285116545653 0.301448607707617
-1.20687691991664 0.318123302240222
-1.16969963321623 0.334480884546718
-1.13135636448322 0.350505049033324
-1.09188533512174 0.366179822464391
-1.05132589071228 0.381489579884875
-1.00971846179108 0.396419060195626
-0.967104523548157 0.410953381365991
-0.923526554483912 0.425078055268538
-0.879027994065622 0.438779002121145
-0.833653199426053 0.452042564522023
-0.787447401147295 0.464855521063722
-0.740456658173942 0.477205099512507
-0.692727811900527 0.489078989540008
-0.644308439479006 0.500465354994416
-0.595246806392815 0.511352845699018
-0.545591818344793 0.521730608766304
-0.495392972506902 0.531588299416356
-0.444700308180383 0.540916091288752
-0.393564356915482 0.549704686237696
-0.342036092140505 0.557945323600612
-0.290166878350402 0.565629788930961
-0.238008419905511 0.572750422186584
-0.185612709491533 0.579300125365395
-0.133031976292094 0.585272369580823
-0.0803186339255527 0.590661201569954
-0.0275252281979734 0.595461249627865
0.0252956152756827 0.599667728962262
0.0780912435299118 0.603276446463064
0.130809028734289 0.606283804882186
0.183396420654035 0.608686806419352
0.235800999033231 0.610483055710369
0.287970525848473 0.611670762214881
0.339852997380866 0.612248742001214
0.391396696054464 0.612216418926548
0.442550241989477 0.611573825211226
0.493262644218846 0.610321601406638
0.543483351517152 0.608460995756706
0.593162302791179 0.605993862953605
0.642249976981897 0.602922662288968
0.690697442428129 0.599250455202409
0.738456405642702 0.594980902229814
0.785479259452436 0.590118259354438
0.831719130454018 0.584667373764451
0.877129925738425 0.578633679021152
0.921666378837339 0.572023189642681
0.965284094845748 0.564842495108613
1.00793959467575 0.557098753291422
1.04959035839746 0.54879968332136
1.0901948676238 0.539953557891859
1.1297126468969 0.530569195013129
1.16810430403499 0.520655949222175
1.20533156939928 0.510223702257992
1.24135733404207 0.49928285321123
1.27614568669767 0.487844308158158
1.3096619495796 0.475919469289245
1.34187271294812 0.463520223543213
1.37274586841379 0.450658930757876
1.40225064094382 0.437348411349587
1.4303576195393 0.423601933533563
1.4570387865527 0.409433200097844
1.4822675456165 0.394856334744058
1.50601874815499 0.3798858680086
1.52826871845294 0.364536722778286
1.54899527725608 0.348824199414888
1.56817776387984 0.332763960503405
1.58579705680438 0.316372015239259
1.60183559273534 0.299664703469972
1.61627738411129 0.282658679407264
1.62910803504052 0.265370895025759
1.64031475565115 0.247818583164898
1.64988637484035 0.230019240350867
1.65781335140995 0.211990609355675
1.66408778357732 0.19375066151078
1.66870341685202 0.175317578792868
1.67165565027041 0.15670973569967
1.67294154098203 0.137945680933858
1.67255980718305 0.119044118913295
1.67051082939402 0.100023891126071
1.66679665008056 0.0809039573488865
1.66142097161739 0.0617033767475461
1.6543891525977 0.0424412888783596
1.64570820249162 0.0231368946094208
1.63538677465901 0.00380943698076269
1.62343515772359 -0.0155218179775246
1.60986526531706 -0.0348376004500884
1.59469062420326 -0.054118656044888
1.57792636079448 -0.0733457649863883
1.55958918607307 -0.0924997612742515
1.53969737893357 -0.111561551788424
1.51827076796195 -0.130512135321572
1.49533071167001 -0.149332621519899
1.47090007720473 -0.168004249713461
1.44500321755388 -0.186508407617211
1.41766594727034 -0.204826649884128
1.38891551673965 -0.222940716491941
1.35878058501623 -0.240832550945118
1.32729119125537 -0.25848431827397
1.2944787247696 -0.27587842281294
1.26037589373911 -0.292997525740342
1.22501669260752 -0.309824562362077
1.18843636819551 -0.326342759122093
1.15067138456598 -0.342535650322627
1.1117593866759 -0.358387094537573
1.07173916285104 -0.373881290702607
1.03065060612083 -0.389002793866027
0.988534674452212 -0.403736530584624
0.945433349921791 -0.418067813949213
0.901389596867211 -0.431982358224863
0.856447319059366 -0.445466293091234
0.810651315938165 -0.458506177468804
0.764047237955473 -0.47108901291724
0.716681541069754 -0.483202256592521
0.668601440437754 -0.494833833749915
0.619854863349415 -0.505972149780353
0.570490401452902 -0.516606101768184
0.520557262317392 -0.526725089558802
0.470105220381907 -0.536319026325114
0.419184567339059 -0.545378348622305
0.367846062003207 -0.55389402592089
0.316140879712964 -0.561857569608543
0.264120561318499 -0.569261041451727
0.211836961804511 -0.5760970615087
0.159342198600039 -0.582358815486003
0.106688599626683 -0.588040061531093
0.0539286511369939 -0.593135136454357
};
\addplot [crimson2143940]
table {%
0.176582364103728 -0.663478757152825
0.132061364102443 -0.674973744362106
0.0874087224613957 -0.685795903063912
0.0426689499048645 -0.695934445495094
-0.00211335598888132 -0.705379265336065
-0.0468935552434168 -0.714120947785006
-0.091627009982261 -0.722150778942741
-0.136269128924883 -0.729460754498942
-0.180775411836258 -0.736043587710999
-0.225101493885674 -0.741892716667606
-0.269203189870556 -0.747002310829814
-0.313036538261238 -0.751367276843039
-0.356557845022771 -0.754983263614234
-0.399723727170089 -0.757846666649143
-0.442491156013111 -0.759954631645352
-0.484817500048676 -0.761305057337506
-0.52666056745655 -0.761896597591909
-0.567978648157151 -0.76172866274837
-0.608730555389068 -0.760801420207993
-0.648875666764912 -0.759115794266308
-0.688373964764603 -0.756673465191912
-0.727186076625698 -0.753476867551541
-0.765273313591022 -0.749529187783241
-0.802597709474462 -0.744834361020059
-0.839122058506481 -0.739397067167413
-0.874809952421645 -0.733222726238064
-0.909625816751174 -0.726317492949329
-0.943534946284338 -0.71868825058792
-0.976503539663376 -0.710342604148544
-1.00849873307742 -0.701288872753068
-1.03948863302185 -0.691536081357845
-1.06944234809044 -0.681093951757438
-1.0983300197686 -0.669972892893725
-1.12612285219696 -0.658183990480038
-1.15279314087572 -0.645738995950684
-1.1783143002811 -0.632650314746859
-1.2026608903664 -0.618930993950635
-1.22580864192113 -0.604594709279344
-1.24773448076313 -0.589655751453329
-1.26841655073941 -0.574129011950636
-1.28783423551283 -0.558029968162873
-1.30596817911292 -0.541374667967004
-1.32280030523033 -0.524179713728476
-1.33831383523568 -0.506462245751626
-1.35249330490489 -0.488239925193846
-1.36532457983421 -0.46953091646056
-1.37679486952975 -0.450353869098551
-1.38689274015729 -0.430727899205685
-1.39560812593978 -0.410672570375571
-1.40293233919113 -0.390207874196142
-1.40885807897628 -0.369354210321608
-1.41337943838896 -0.34813236613764
-1.41649191043977 -0.326563496040043
-1.41819239254888 -0.304669100347604
-1.41847918963874 -0.282471003870101
-1.41735201582377 -0.25999133415286
-1.41481199469536 -0.237252499419533
-1.41086165820179 -0.2142771662351
-1.4055049441244 -0.191088236911339
-1.39874719215227 -0.167708826677304
-1.39059513855953 -0.144162240637563
-1.3810569094905 -0.120471950541165
-1.37014201285935 -0.0966615713844856
-1.3578613288724 -0.0727548378712874
-1.34422709918253 -0.0487755807534525
-1.32925291468642 -0.024747703075968
-1.31295370197682 -0.00069515634985301
-1.29534570846343 0.0233580833232299
-1.27644648617711 0.0473880391508723
-1.25627487427366 0.0713707575505106
-1.2348509802545 0.0952823320268761
-1.21219615992311 0.11909892700251
-1.18833299609709 0.142796801577586
-1.16328527609711 0.166352333195355
-1.13707796803519 0.189742041189621
-1.10973719592596 0.212942610190785
-1.08129021364565 0.23593091336711
-1.05176537776485 0.258684035478046
-1.02119211928207 0.281179295716648
-0.989600914286255 0.30339427031829
-0.957023253577586 0.325306814913165
-0.923491611276741 0.346895086600275
-0.889039412454001 0.368137565720908
-0.853700999810406 0.389013077309904
-0.817511599444208 0.409500812203315
-0.780507285736726 0.42958034778144
-0.74272494539262 0.449231668326533
-0.704202240670423 0.468435184974904
-0.66497757183998 0.487171755243526
-0.625090038904225 0.50542270211168
-0.584579402623451 0.523169832638611
-0.543486044880911 0.540395456098652
-0.501850928429288 0.557082401615722
-0.459715556058128 0.573214035279626
-0.417121929222954 0.588774276727105
-0.374112506177309 0.603747615171083
-0.330730159649437 0.618119124862165
-0.287018134105822 0.631874479966945
-0.24302000264416 0.644999968848308
-0.198779623558747 0.657482507733487
-0.154341096621581 0.669309653756257
-0.109748719122746 0.680469617360248
-0.0650469417139121 0.690951274051038
-0.02028032409896 0.700744175485285
0.0245065093841042 0.709838559885869
0.0692689142456546 0.718225361772635
0.113962270347049 0.725896220999071
0.158542026378949 0.732843491085877
0.202963744271038 0.739060246843142
0.247183143488852 0.744540291273529
0.291156145173584 0.749278161749571
0.334838916080844 0.75326913545894
0.378187912274595 0.756509234112248
0.42115992253269 0.758995227908693
0.463712111420766 0.760724638755594
0.505802061991528 0.761695742738609
0.547387818066888 0.761907571840167
0.5884279260608 0.761359914904414
0.628881476301088 0.760053317847694
0.668708143809104 0.757989083114367
0.707868228496538 0.755169268378509
0.746322694739335 0.751596684492773
0.784033210289249 0.74727489268647
0.820962184484265 0.742208201015663
0.857072805719787 0.736401660068798
0.892329078143245 0.729861057932169
0.926695857535548 0.722592914420228
0.9601388863436 0.714604474576493
0.992624827828978 0.705903701451531
1.02412129929872 0.696499268165214
1.05459690438508 0.686400549261166
1.08402126434214 0.675617611362013
1.11236504832799 0.664161203134758
1.13960000264234 0.652042744576271
1.16569897889044 0.639274315629591
1.19063596104515 0.625868644142377
1.21438609138036 0.611839093179507
1.23692569524964 0.597199647702492
1.25823230468566 0.581964900628957
1.27828468079677 0.566150038286104
1.29706283493837 0.549770825272654
1.31454804863803 0.532843588744346
1.33072289225448 0.515385202138682
1.34557124235182 0.497413068355107
1.35907829777178 0.478945102407429
1.37123059438775 0.45999971356574
1.3820160185262 0.440595787005655
1.39142381904177 0.420752664983159
1.39944461803432 0.400490127553824
1.40607042019697 0.379828372855624
1.41129462078605 0.358787996974986
1.41511201220478 0.337389973416166
1.41751878919438 0.315655632194406
1.4185125526272 0.293606638573705
1.41809231189825 0.271264971470418
1.41625848591262 0.248652901544194
1.41301290266793 0.225792968998093
1.40835879743215 0.202707961110028
1.40230080951859 0.179420889517906
1.39484497766136 0.155954967281131
1.38599873399579 0.132333585741314
1.37577089664995 0.108580291205285
1.36417166095454 0.0847187614736174
1.35121258927996 0.0607727822380779
1.33690659951069 0.0367662233715387
1.32126795216845 0.0127230151339631
1.304312236197 -0.0113328756817937
1.28605635342279 -0.0353774696406068
1.26651850170678 -0.0593867985683059
1.24571815680446 -0.0833369294436681
1.22367605295199 -0.107203988255379
1.20041416219786 -0.130964183800161
1.17595567250066 -0.154593831398368
1.15032496461484 -0.178069376503384
1.12354758778742 -0.20136741818131
1.09565023428997 -0.224464732437526
1.06666071281117 -0.247338295366865
1.03660792073646 -0.269965306104357
1.00552181534258 -0.292323209553611
0.973433383935422 -0.314389718870239
0.940374612961267 -0.336142837677859
0.906378456121974 -0.357560881994564
0.871478801526013 -0.378622501847987
0.835710437908074 -0.399306702557418
0.799109019950888 -0.419592865661755
0.761711032743869 -0.439460769472441
0.723553755413985 -0.458890609230877
0.684675223965108 -0.477863016850238
0.645114193362903 -0.496359080222004
0.604910098903037 -0.51436036206796
0.564103016901208 -0.531848918318877
0.522733624744216 -0.548807316001547
0.480843160341849 -0.565218650616356
0.438473381020044 -0.581066562988053
0.395666521896272 -0.596335255572929
0.352465253778655 -0.611009508206158
0.308912640630781 -0.625074693273584
0.265052096644599 -0.638516790292849
0.220927342964222 -0.651322399889308
0.176582364103728 -0.663478757152825
};
\addplot [crimson2143940]
table {%
0.381434826773266 -0.622445065599631
0.344373693480989 -0.644729423568221
0.30696928100173 -0.666371101226047
0.269258874867308 -0.687348525678163
0.23128006563067 -0.707640786171666
0.193070711394813 -0.727227654939988
0.154668900075002 -0.746089607366371
0.116112911431912 -0.764207841446434
0.0774411789135223 -0.781564296530417
0.0386922513438281 -0.798141671326429
-9.52455034710642e-05 -0.813923441146756
-0.0388826474079811 -0.828893874380023
-0.0776312902439486 -0.843038048172812
-0.116302548521503 -0.856341863305084
-0.154857873889383 -0.868792058244589
-0.193258833560775 -0.880376222366247
-0.231467148623948 -0.891082808323335
-0.269444732199505 -0.900901143558127
-0.307153727406216 -0.909821440940535
-0.344556545097575 -0.917834808524133
-0.381615901331486 -0.924933258409841
-0.418294854535699 -0.93110971470844
-0.454556842331979 -0.936358020593967
-0.490365717982279 -0.940672944440977
-0.525685786420591 -0.94405018503954
-0.560481839834567 -0.94648637588278
-0.594719192761432 -0.947979088522683
-0.628363716663206 -0.948526834990828
-0.661381873946774 -0.948129069281621
-0.693740751394889 -0.946786187896573
-0.725408092974781 -0.944499529449053
-0.756352331991672 -0.94127137332993
-0.786542622555147 -0.93710493743543
-0.815948870327003 -0.932004374959461
-0.844541762519944 -0.925974770253626
-0.872292797117205 -0.919022133759024
-0.899174311283979 -0.911153396014917
-0.925159508942335 -0.902376400750215
-0.950222487482129 -0.892699897064677
-0.974338263581292 -0.882133530707618
-0.997482798109748 -0.870687834462816
-1.01963302009214 -0.858374217649205
-1.04076684970548 -0.845204954747817
-1.06086322028879 -0.831193173166313
-1.07990209934284 -0.816352840153296
-1.09786450849893 -0.800698748875457
-1.11473254243697 -0.78424650367142
-1.13048938673389 -0.767012504497004
-1.14511933462463 -0.749013930577383
-1.15860780265895 -0.730268723282469
-1.17094134523859 -0.710795568242555
-1.18210766802006 -0.690613876722083
-1.19209564016999 -0.669743766270064
-1.20089530546051 -0.648206040666472
-1.2084978921939 -0.626022169184584
-1.21489582194635 -0.60321426518994
-1.22008271712236 -0.57980506409726
-1.22405340731201 -0.55581790070729
-1.22680393444495 -0.53127668594616
-1.22833155673594 -0.506205883030458
-1.22863475141785 -0.480630483081761
-1.22771321625962 -0.454575980214944
-1.22556786986754 -0.428068346125088
-1.22220085076955 -0.401134004198334
-1.21761551528352 -0.373799803172471
-1.21181643417158 -0.346092990373538
-1.20480938808388 -0.31804118455509
-1.19660136179638 -0.289672348367229
-1.18720053724816 -0.261014760482822
-1.17661628538559 -0.232096987408707
-1.1648591568211 -0.202947855009978
-1.15194087131614 -0.173596419775737
-1.13787430609865 -0.144071939854948
-1.12267348302674 -0.114403845891285
-1.10635355461144 -0.0846217116860227
-1.08893078891231 -0.0547552247182283
-1.0704225533211 -0.024834156551637
-1.05084729724956 0.00511166684228936
-1.03022453373866 0.035052394815541
-1.00857482000757 0.0649581817993328
-0.985919736961762 0.0947992170548192
-0.962281867680708 0.124545754389085
-0.937684774906579 0.154168141806798
-0.91215297755638 0.183636851067963
-0.885711926280966 0.212922507122312
-0.858387978095281 0.241995917390998
-0.830208370105107 0.270828100866385
-0.801201192356529 0.299390317000953
-0.771395359835165 0.327654094356498
-0.740820583643079 0.355591258985081
-0.709507341382109 0.383173962513432
-0.677486846773131 0.410374709902816
-0.64479101854154 0.437166386856683
-0.611452448599977 0.463522286848787
-0.577504369559993 0.48941613774484
-0.542980621605067 0.514822127991136
-0.507915618757972 0.539714932344078
-0.472344314576127 0.564069737114927
-0.436302167309132 0.587862264904627
-0.399825104553215 0.611068798804046
-0.362949487437823 0.633666206035504
-0.325712074380053 0.655631961012035
-0.288149984443068 0.676944167791372
-0.250300660335004 0.697581581902309
-0.212201831085267 0.717523631521644
-0.173891474435412 0.736750437980623
-0.135407778982109 0.75524283558043
-0.0967891061099149 0.772982390696972
-0.0580739517518148 0.789951420155916
-0.0193009080156363 0.806133008859659
0.0194913752854021 0.82151102664867
0.0582642291596638 0.836070144380381
0.0969790039831889 0.849795849209611
0.135597108026478 0.862674459055284
0.174080045923567 0.874693136239026
0.212389457045049 0.885839900282044
0.250487153736774 0.896103639847521
0.288335159386148 0.905474123816648
0.325895746278033 0.913942011487218
0.36313147320256 0.921498861884644
0.400005222777342 0.928137142176107
0.436480238446885 0.933850235179445
0.472520161122323 0.938632445959305
0.508089065424951 0.942479007503981
0.543151495497418 0.945386085477271
0.577672500346904 0.947350782040629
0.611617668685023 0.948371138741794
0.644953163229748 0.948446138467015
0.677645754435149 0.947575706454935
0.709662853615325 0.945760710371113
0.740972545429515 0.943002959443118
0.771543619696009 0.93930520265705
0.801345602503126 0.93467112601729
0.830348786586283 0.929105348872213
0.858524260940834 0.922613419309521
0.885843939641191 0.915201808625784
0.912280589837488 0.906877904875718
0.937807858901874 0.897650005507607
0.962400300697387 0.887527309092226
0.986033400943218 0.876519906153505
1.00868360165108 0.864638769110081
1.03032832460831 0.851895741337744
1.05094599388435 0.83830352536371
1.07051605733808 0.823875670204459
1.08901900710464 0.808626557859786
1.10643639904132 0.79257138897651
1.12275087111303 0.775726167696134
1.13794616069922 0.758107685701577
1.15200712080478 0.739733505478849
1.16491973515886 0.720621942810393
1.17667113218667 0.700792048517507
1.18724959784007 0.680263589470077
1.19664458727446 0.659057028882528
1.20484673536005 0.637193505915645
1.21184786601724 0.614694814604603
1.21764100036673 0.591583382134194
1.22222036368623 0.567882246482923
1.22558139116678 0.543615033458253
1.2277207324631 0.51880593314588
1.22863625503326 0.493479675796537
1.22832704626444 0.467661507174335
1.22679341438268 0.441377163391241
1.22403688814558 0.414652845252763
1.22006021531845 0.38751519214041
1.21486735993525 0.359991255456986
1.20846349834715 0.332108471661158
1.20085501406266 0.303894634918198
1.19204949138437 0.275377869394157
1.18205570784878 0.246586601221082
1.17088362547668 0.21754953016123
1.15854438084274 0.188295600998521
1.14505027397434 0.158853974685749
1.1304147560907 0.129253999276311
1.11465241619432 0.0995251806694221
1.0977789665284 0.0696971531980036
1.07981122691453 0.0397996500885271
1.06076710798631 0.00986247382229433
1.04066559333566 -0.0200845335723311
1.01952672058953 -0.0500115202671003
0.997371561436013 -0.079888654390833
0.974222200619563 -0.109686153766479
0.950101713926491 -0.139374315598634
0.925034145182494 -0.168923546081936
0.899044482285258 -0.198304389900814
0.872158632295974 -0.227487559591182
0.844403395614614 -0.256443964734831
0.815806439264714 -0.285144740957377
0.786396269314279 -0.313561278700908
0.756202202460319 -0.341665251742609
0.725254336805325 -0.369428645430955
0.693583521854823 -0.39682378461133
0.661221327765924 -0.423823361213218
0.628200013877495 -0.450400461471486
0.594552496553358 -0.476528592754607
0.560312316370536 -0.502181709973091
0.525513604685277 -0.527334241541802
0.490191049610182 -0.551961114870261
0.45437986143633 -0.576037781355559
0.418115737534903 -0.599540240852933
0.381434826773266 -0.622445065599631
};
\addplot [crimson2143940]
table {%
-0.510997185558914 0.458549090999412
-0.480033592812941 0.492541885770032
-0.448591492142048 0.526043704077286
-0.416702225715976 0.5590211505801
-0.384397581449015 0.591441352642619
-0.35170976131312 0.623271993102362
-0.318671349238287 0.654481342484678
-0.28531527863218 0.685038290631373
-0.251674799551394 0.714912377711992
-0.21778344555708 0.744073824586841
-0.183675000287956 0.772493562491484
-0.149383463784044 0.800143262013113
-0.114943018594693 0.82699536132993
-0.0803879957046657 0.853023093685366
-0.0457528403122699 0.878200514069766
-0.0110720774936341 0.902502525082942
0.0236197222126391 0.925904901951804
0.0582879772661374 0.948384316678143
0.09289812959627 0.969918361292484
0.127415679050505 0.990485570190841
0.16180621778487 1.0100654415321
0.196035464562442 1.02863845767468
0.230069298925627 1.04618610463218
0.263873795208176 1.06269089052849
0.297415256353026 1.0781363630341
0.330660247502259 1.09250712576609
0.363575629325698 1.10578885363566
0.396128591054908 1.11796830712761
0.428286683189689 1.12903334549783
0.460017849844436 1.13897293887543
0.491290460702149 1.14777717925757
0.522073342544221 1.15543729038597
0.552335810324574 1.16194563649526
0.582047697757189 1.16729572992451
0.611179387386509 1.17148223758423
0.639701840110764 1.17450098627253
0.667586624128785 1.17634896683508
0.694805943281431 1.17702433716466
0.721332664759416 1.17652642403743
0.747140346149878 1.17485572378404
0.772203261794748 1.17201390179482
0.796496428434653 1.16800379085975
0.819995630112767 1.1628293883446
0.842677442313805 1.15649585220635
0.864519255314088 1.14900949585154
0.88549929671941 1.14037778184302
0.90559665316823 1.13060931446101
0.924791291178572 1.11971383112624
0.943064077117827 1.10770219269347
0.960396796275581 1.09458637262509
0.976772171020432 1.0803794450558
0.992173878022708 1.06509557175996
1.00658656452592 1.04874998803486
1.01999586365071 1.03135898751384
1.03238840871611 1.01293990592442
1.04375184656367 0.993511103807719
1.05407484987142 0.973091948216205
1.06334712844517 0.951702793408244
1.07155943947598 0.929364960558512
1.07870359675362 0.90610071650459
1.08477247882674 0.881933251550891
1.08976003610172 0.856886656352069
1.09366129687298 0.830985897898922
1.09647237227896 0.804256794630757
1.09819046017855 0.776725990699007
1.09881384794437 0.748420929407761
1.09834191416993 0.719369825857683
1.0967751292891 0.689601638820585
1.0941150551071 0.659146041872694
1.09036434324374 0.628033393815386
1.08552673249017 0.59629470841287
1.07960704508199 0.563961623476998
1.07261118189236 0.531066369330003
1.06454611654984 0.497641736676612
1.05541988848697 0.463721043917564
1.04524159492635 0.429338103937094
1.03402138181235 0.394527190397523
1.02177043369742 0.359323003574525
1.00850096259308 0.32376063576713
0.994226195796699 0.28787553631697
0.97896036270626 0.251703476271592
0.962718680636185 0.215280512727107
0.945517339648381 0.178642952885689
0.927373486413633 0.14182731786377
0.908305207119432 0.104870306286993
0.888331509441268 0.0678087577082271
0.867472303595368 0.0306796158850957
0.845748382491758 -0.00648010804635935
0.823181401007438 -0.0436333724651868
0.79979385440033 -0.080743142189418
0.775609055885518 -0.117772425393541
0.750651113396123 -0.154684310482757
0.724944905551996 -0.191442002887264
0.698516056860164 -0.228008861739878
0.671390912171769 -0.264348436400453
0.643596510420945 -0.300424502790666
0.615160557671821 -0.336201099502969
0.586111399500512 -0.371642563647703
0.556477992739633 -0.40671356640264
0.526289876613501 -0.441379148229522
0.495577143292791 -0.475604753722488
0.464370407898008 -0.509356266053656
0.432700777981666 -0.542600040981519
0.400599822519606 -0.575302940388254
0.368099540442352 -0.607432365312523
0.335232328737885 -0.63895628844482
0.302030950157615 -0.669843286052995
0.268528500557765 -0.700062569306114
0.234758375908693 -0.729584014965431
0.200754239005073 -0.758378195411891
0.166549985910078 -0.786416407980219
0.132179712167057 -0.813670703570362
0.0976776788123502 -0.840113914507757
0.0630782782231427 -0.865719681624667
0.0284159998343979 -0.890462480535561
-0.00627460424096371 -0.91431764708039
-0.0409589536543594 -0.937261401910347
-0.0756024742919914 -0.959270874191643
-0.110170632739073 -0.980324124403656
-0.144628970703487 -1.00040016620872
-0.178943139364564 -1.01947898737174
-0.213078933612729 -1.03754156970888
-0.247002326145902 -1.05456990804524
-0.280679501388652 -1.07054702816285
-0.314076889200298 -1.08545700372089
-0.347161198338354 -1.09928497213148
-0.379899449643965 -1.11201714937499
-0.412259008916249 -1.12364084374023
-0.444207619442781 -1.13414446847589
-0.475713434153779 -1.14351755334039
-0.506745047367962 -1.15175075503889
-0.537271526098414 -1.1588358665369
-0.567262440887255 -1.16476582524119
-0.596687896138395 -1.16953472003995
-0.625518559918103 -1.17313779719514
-0.653725693193723 -1.1755714650811
-0.681281178481369 -1.17683329776474
-0.708157547874039 -1.17692203742381
-0.734328010422224 -1.17583759560068
-0.759766478839709 -1.17358105329055
-0.784447595507949 -1.17015465986387
-0.808346757753094 -1.16556183082412
-0.831440142370466 -1.15980714440317
-0.853704729372052 -1.15289633699759
-0.875118324933326 -1.14483629745051
-0.895659583516537 -1.13563506018463
-0.915308029148405 -1.12530179719335
-0.934044075831018 -1.11384680889797
-0.951849047065576 -1.10128151387992
-0.968705194469537 -1.08761843749855
-0.984595715468587 -1.07287119940553
-0.999504770045812 -1.05705449996851
-1.01341749653137 -1.04018410561747
-1.02632002641694 -1.02227683312839
-1.03819949818012 -1.00335053285993
-1.04904407010512 -0.983424070959792
-1.05884293208679 -0.962517310558557
-1.06758631640641 -0.940651091969665
-1.07526550746833 -0.917847211915344
-1.08187285048791 -0.894128401799165
-1.08740175912194 -0.869518305046878
-1.09184672203403 -0.844041453538138
-1.09520330838849 -0.817723243152592
-1.09746817226706 -0.79058990845471
-1.09863905600415 -0.762668496542606
-1.09871479243739 -0.733986840086899
-1.09769530607107 -0.704573529586502
-1.09558161315136 -0.674457884868993
-1.09237582065333 -0.643669925863971
-1.08808112418067 -0.612240342678539
-1.08270180478021 -0.580200465004733
-1.07624322467452 -0.547582230889413
-1.06871182191671 -0.514418154897715
-1.06011510397285 -0.480741295701827
-1.05046164023833 -0.44658522312739
-1.03976105349574 -0.411983984690358
-1.02802401032267 -0.376972071657703
-1.015262210459 -0.341584384665762
-1.00148837514438 -0.305856198930526
-0.986716234437389 -0.269823129084544
-0.970960513529149 -0.233521093675469
-0.954236918064913 -0.196986279361685
-0.936562118488336 -0.160255104840642
-0.917953733424023 -0.123364184545922
-0.898430312114894 -0.0863502921491658
-0.878011315931898 -0.0492503239032784
-0.856717098974497 -0.0121012618634463
-0.834568887781252 0.0250598629773862
-0.811588760170753 0.0621960076018098
-0.787799623233965 0.0992701538932378
-0.763225190499945 0.136245345536373
-0.737889958297679 0.17308472485606
-0.711819181337607 0.209751569557826
-0.685038847537181 0.246209329333451
-0.657575652115545 0.282421662295122
-0.629456970983151 0.318352471201809
-0.600710833452872 0.353965939441783
-0.571365894299755 0.389226566735393
-0.541451405197333 0.424099204522518
-0.510997185558915 0.458549090999412
};
\addplot [crimson2143940]
table {%
-0.504004181026227 0.317328497015738
-0.475511920647303 0.36202483972527
-0.446545659644609 0.406360308201374
-0.417134272217063 0.45029070788488
-0.387307076275279 0.493772247998001
-0.357093804216844 0.53676158519593
-0.326524573288414 0.579215866772431
-0.295629855564187 0.621092773376349
-0.264440447570678 0.662350561196454
-0.232987439588072 0.70294810357259
-0.201302184658754 0.742844931991613
-0.169416267333919 0.782001276427294
-0.1373614721894 0.820378104983939
-0.105169752142112 0.857937162804239
-0.0728731965986937 0.894641010202537
-0.040503999468085 0.930453059985523
-0.00809442706993957 0.965337613923145
0.024323214029144 0.999259898333382
0.0567166092195104 1.0321860987454
0.0890534680603536 1.06408339360657
0.121301556467531 1.09491998699967
0.153428728845197 1.1246651403378
0.185402960129235 1.15328920300521
0.217192377710537 1.18076364191373
0.248765293206314 1.20706106994515
0.280090234047761 1.23215527325131
0.311135974852604 1.2560212373846
0.341871568551237 1.27863517223296
0.372266377235434 1.29997453573432
0.402290102698883 1.32001805634707
0.431912816639096 1.33874575425396
0.461104990490594 1.35613896127842
0.48983752485962 1.37218033949345
0.518081778531049 1.38685389850437
0.545809597018576 1.40014501138845
0.572993340629717 1.4120404292753
0.599605912017659 1.42252829455369
0.625620783192481 1.43159815269142
0.651012021964833 1.43924096265664
0.675754317795707 1.44544910593018
0.699823007026532 1.4502163940998
0.723194097464446 1.45353807502898
0.745844292298234 1.45541083759393
0.767751013321096 1.45583281498419
0.788892423437097 1.45480358656354
0.809247448428858 1.45232417828923
0.828795797964793 1.44839706168937
0.847517985824953 1.44302615139918
0.865395349325317 1.43621680125884
0.882410067921156 1.42797579897666
0.898545180970943 1.41831135936294
0.913784604643091 1.40723311614127
0.92811314794866 1.39475211234541
0.941516527884071 1.38088078931138
0.953981383668714 1.36563297427564
0.96549529006326 1.34902386659184
0.976046769755415 1.33107002257972
0.985625304800752 1.31178933902144
0.994221347107227 1.29120103532169
1.00182632795292 1.26932563434934
1.00843266652755 1.24618494197981
1.01403377748913 1.22180202535852
1.01862407752846 1.19620118990703
1.02219899093463 1.16940795509487
1.02475495415621 1.1414490290012
1.02628941935347 1.11235228169158
1.02680085693813 1.08214671743652
1.02628875709808 1.05086244579937
1.02475363030559 1.01853065162245
1.02219700680843 0.985183563941347
1.01862143510449 0.950854423858277
1.01403047940142 0.915577451406653
1.0084287160637 0.879387811439806
1.00182172905085 0.842321578577896
0.994216104351205 0.804415701247942
0.98561942341688 0.765707964852832
0.976040255606385 0.726236954106005
0.965488149642517 0.686042014569377
0.953973624093984 0.645163213432822
0.941508156890268 0.603641299574333
0.928104173880171 0.561517662940646
0.913775036445464 0.51883429328885
0.898535028181966 0.475633738330075
0.882399340661348 0.431959061317015
0.865384058287845 0.387853798117539
0.847506142264969 0.34336191381719
0.828783413688215 0.298527758893832
0.809234535780602 0.253396025008132
0.788878995288772 0.208011700453934
0.767737083058171 0.162420025312957
0.745829873806699 0.116666446358494
0.723179205116967 0.0707965717530803
0.699807655668125 0.0248561255852836
0.67573852272894 -0.0211090977090589
0.650995798934571 -0.0670532789957324
0.625604148370194 -0.112930620115651
0.5995888819853 -0.158695389537462
0.572975932363193 -0.20430196794373
0.545791827870824 -0.249704893705264
0.518063666214742 -0.294858908198258
0.48981908742951 -0.33971900091907
0.461086246325512 -0.384240454351675
0.431893784423634 -0.428378888543054
0.402270801404766 -0.472090305342099
0.372246826102612 -0.515331132257923
0.341851787068705 -0.558058265893869
0.311115982738974 -0.600229114913916
0.280070051231606 -0.641801642498642
0.248744939806305 -0.682734408248446
0.217171874015386 -0.722986609492237
0.185382326577469 -0.762518121960426
0.153407986004791 -0.801289539781665
0.121280725015399 -0.839262214763483
0.0890325687617405 -0.876398294917642
0.0566956629072912 -0.912660762191821
0.0243022415830578 -0.948013469370019
-0.00811540474410177 -0.982421176104881
-0.0405249614593216 -1.01584958404605
-0.0728941220116357 -1.04826537102949
-0.105190620117853 -1.07963622429379
-0.137382261926287 -1.10993087269015
-0.169436958108292 -1.13911911785422
-0.201322755845599 -1.16717186430845
-0.233007870681581 -1.19406114846509
-0.264460718204693 -1.21976016650092
-0.295649945532499 -1.24424330107585
-0.326544462564907 -1.2674861468689
-0.357113472975453 -1.28946553490591
-0.387326504909759 -1.31015955565487
-0.417153441360529 -1.32954758086586
-0.446564550188841 -1.34761028413371
-0.475530513761795 -1.36432966016297
-0.504022458176959 -1.37968904271601
-0.532011982044507 -1.39367312122627
-0.559471184798342 -1.40626795606018
-0.586372694507983 -1.41746099241246
-0.612689695163511 -1.42724107282106
-0.638395953406336 -1.43559844828912
-0.663465844679193 -1.44252478800303
-0.68787437876925 -1.44801318763674
-0.711597224718895 -1.45205817623412
-0.734610735079372 -1.45465572166257
-0.756891969483063 -1.45580323463232
-0.77841871751095 -1.45549957127749
-0.799169520832442 -1.45374503429631
-0.819123694595507 -1.45054137264941
-0.838261348045782 -1.44589177981637
-0.856563404354109 -1.43980089061246
-0.874011619632732 -1.4322747765685
-0.890588601121204 -1.42332093987865
-0.906277824523867 -1.41294830592203
-0.92106365048163 -1.40116721436576
-0.93493134016162 -1.38798940885812
-0.947867069949176 -1.37342802532221
-0.959857945227522 -1.35749757886182
-0.970892013231409 -1.34021394929245
-0.980958274961884 -1.3215943653119
-0.990046696150338 -1.30165738732645
-0.998148217260878 -1.28042288894931
-1.00525476252107 -1.25791203719027
-1.01135924797205 -1.23414727135593
-1.01645558852997 -1.2091522806817
-1.02053870405171 -1.18295198071786
-1.02360452439893 -1.15557248849321
-1.02564999349524 -1.12704109648104
-1.02667307237255 -1.0973862453934
-1.0266727412036 -1.06663749583078
-1.02564900031848 -1.03482549881546
-1.0236028702044 -1.00198196523789
-1.02053639048833 -0.968139634246611
-1.01645261790395 -0.933332240613153
-1.01135562324459 -0.897594481104467
-1.00525048730537 -0.860961979896445
-0.998143295818565 -0.823471253062959
-0.990041133387213 -0.785159672175853
-0.980952076423031 -0.746065427052143
-0.970885185095665 -0.706227487685576
-0.959850494301322 -0.665685565400493
-0.947859003659765 -0.624480073266721
-0.934922666549652 -0.582652085814942
-0.921054378193148 -0.540243298092706
-0.906267962801692 -0.497295984101909
-0.890578159795713 -0.453852954659138
-0.874000609112066 -0.409957514720926
-0.856551835613795 -0.365653420216428
-0.838249232617792 -0.320984834430554
-0.819111044556763 -0.275996283981054
-0.799156348792777 -0.230732614433409
-0.77840503660054 -0.185238945597805
-0.756877793339338 -0.139560626552717
-0.734596077833419 -0.0937431904399643
-0.711582100981376 -0.047832309076288
-0.687858803615835 -0.00187374742668284
-0.663449833635534 0.0440866820151069
-0.638379522432591 0.0900031648934808
-0.612672860638429 0.135829930659733
-0.586355473212588 0.18152129819712
-0.559453593899192 0.227031721356803
-0.531994039076596 0.27231583435923
-0.504004181026227 0.317328497015737
};
\addplot [line width=0.52pt, gray127, dashed]
table {%
0 0.5
0.0631370995296211 0.499750794563087
0.126211262625347 0.499003426665747
0.18915961558969 0.497758641300553
0.251919410135435 0.496017679296629
0.314428085934502 0.493782276082762
0.376623332979436 0.49105465995749
0.438443153695383 0.487837549867887
0.499825924740617 0.484134152699254
0.560710458434029 0.479948160078428
0.621036063748338 0.475283744693891
0.680742606808226 0.470145556136338
0.739770570833094 0.464538716263865
0.798061115464682 0.458468814096394
0.855556135420419 0.451941900244412
0.912198318414032 0.44496448087759
0.967931202285678 0.437543511239297
1.02269923128465 0.429686388713456
1.07644781144858 0.421400945450675
1.12912336502384 0.412695440560988
1.18067338387306 0.403578551880993
1.23104648181636 0.394059367323596
1.28019244585421 0.384147375818985
1.32806228622086 0.373852457855855
1.37460828521837 0.36318487563232
1.41978404478267 0.352155262826336
1.46354453273415 0.340774613995817
1.50584612766675 0.329054273619019
1.54664666243068 0.317005924786119
1.58590546616556 0.304641577553258
1.62358340484204 0.291973556970653
1.65964292027145 0.27901449079672
1.6940480675446 0.265777296910451
1.72676455086245 0.252275170434589
1.75775975772292 0.238521570582445
1.78700279142975 0.22453020724146
1.81446450189096 0.210315027306892
1.84011751467635 0.195890200779247
1.86393625830487 0.181270106639312
1.88589698973487 0.166469318514881
1.90597781803168 0.151502590153436
1.92415872618893 0.136384840715302
1.94042159108197 0.121131139901892
1.95475020153341 0.10575669293391
1.96713027447274 0.0902768253944512
1.97754946917401 0.0747069679521305
1.98599739955734 0.059062640979454
1.99246564454201 0.0433594390817754
1.99694775644076 0.0276130155522541
1.99943926738696 0.011839066768312
1.99993769378831 -0.00394668345485666
1.99844253880255 -0.0197284995381263
1.99495529283268 -0.035490649824008
1.98947943204131 -0.0512174222583066
1.98202041488558 -0.0668931400522395
1.97258567667601 -0.0825021773093997
1.96118462216481 -0.098028974601989
1.94782861617108 -0.113458054480795
1.93253097225204 -0.128774036903448
1.91530693943186 -0.143961654565583
1.89617368700102 -0.159005768119619
1.87515028740165 -0.173891381265991
1.85225769721568 -0.188603655701788
1.82751873627487 -0.203127925911895
1.8009580649135 -0.217449713787896
1.77260215938642 -0.231554743060174
1.74247928547692 -0.245428953528797
1.71061947032082 -0.259058515079039
1.67705447247475 -0.272429841467535
1.64181775025853 -0.285529603865347
1.60494442840316 -0.298344744144428
1.56647126303778 -0.310862487894248
1.52643660505034 -0.323070357155605
1.48488036185857 -0.334956182858928
1.44184395762943 -0.346508116954666
1.39737029198661 -0.357714644223686
1.35150369724722 -0.368564593755889
1.30428989423037 -0.379047150085623
1.25577594668172 -0.389151863972765
1.20601021435923 -0.39886866281876
1.15504230482718 -0.4081878607072
1.10292302400622 -0.417100168058959
1.04970432552894 -0.425596700892247
0.995439258951366 -0.433668989678357
0.940181916872006 -0.441308987784273
0.883987381011158 -0.448509079493732
0.826911667304268 -0.455262087598731
0.769011670064022 -0.46156128055393
0.710345105266857 -0.467400379186799
0.650970453020423 -0.472773562956834
0.59094689926934 -0.477675475757597
0.530334276797357 -0.482101231255806
0.469193005584738 -0.486046417762128
0.407584032580305 -0.489507102628857
0.345568770948199 -0.492479836170055
0.283209038849903 -0.494961655100279
0.22056699782255 -0.496950085488444
0.157705090814953 -0.498443145223897
0.0946859799431164 -0.499439345992222
0.0315724840272736 -0.499937694758829
-0.031572484027274 -0.499937694758829
-0.0946859799431159 -0.499439345992222
-0.157705090814953 -0.498443145223897
-0.22056699782255 -0.496950085488444
-0.283209038849904 -0.494961655100279
-0.345568770948199 -0.492479836170055
-0.407584032580304 -0.489507102628857
-0.469193005584737 -0.486046417762128
-0.530334276797357 -0.482101231255806
-0.59094689926934 -0.477675475757597
-0.650970453020424 -0.472773562956834
-0.710345105266857 -0.467400379186799
-0.769011670064021 -0.46156128055393
-0.826911667304268 -0.455262087598731
-0.883987381011158 -0.448509079493732
-0.940181916872006 -0.441308987784273
-0.995439258951366 -0.433668989678357
-1.04970432552894 -0.425596700892247
-1.10292302400621 -0.417100168058959
-1.15504230482718 -0.4081878607072
-1.20601021435923 -0.39886866281876
-1.25577594668172 -0.389151863972765
-1.30428989423037 -0.379047150085623
-1.35150369724722 -0.36856459375589
-1.39737029198661 -0.357714644223686
-1.44184395762943 -0.346508116954666
-1.48488036185857 -0.334956182858928
-1.52643660505034 -0.323070357155605
-1.56647126303778 -0.310862487894248
-1.60494442840316 -0.298344744144428
-1.64181775025852 -0.285529603865347
-1.67705447247475 -0.272429841467535
-1.71061947032082 -0.259058515079039
-1.74247928547692 -0.245428953528797
-1.77260215938642 -0.231554743060174
-1.8009580649135 -0.217449713787897
-1.82751873627487 -0.203127925911894
-1.85225769721568 -0.188603655701788
-1.87515028740165 -0.173891381265991
-1.89617368700102 -0.159005768119619
-1.91530693943186 -0.143961654565583
-1.93253097225204 -0.128774036903448
-1.94782861617108 -0.113458054480795
-1.96118462216481 -0.098028974601989
-1.97258567667601 -0.0825021773093994
-1.98202041488558 -0.0668931400522396
-1.98947943204131 -0.0512174222583065
-1.99495529283268 -0.0354906498240081
-1.99844253880255 -0.0197284995381264
-1.99993769378831 -0.00394668345485689
-1.99943926738696 0.0118390667683119
-1.99694775644076 0.0276130155522542
-1.99246564454201 0.0433594390817752
-1.98599739955734 0.059062640979454
-1.97754946917401 0.0747069679521302
-1.96713027447274 0.090276825394451
-1.95475020153341 0.10575669293391
-1.94042159108197 0.121131139901891
-1.92415872618893 0.136384840715302
-1.90597781803168 0.151502590153436
-1.88589698973487 0.166469318514881
-1.86393625830487 0.181270106639312
-1.84011751467635 0.195890200779246
-1.81446450189096 0.210315027306892
-1.78700279142975 0.22453020724146
-1.75775975772292 0.238521570582445
-1.72676455086245 0.252275170434589
-1.6940480675446 0.265777296910451
-1.65964292027145 0.27901449079672
-1.62358340484204 0.291973556970653
-1.58590546616556 0.304641577553258
-1.54664666243068 0.317005924786119
-1.50584612766675 0.329054273619019
-1.46354453273415 0.340774613995817
-1.41978404478267 0.352155262826336
-1.37460828521837 0.36318487563232
-1.32806228622086 0.373852457855854
-1.28019244585421 0.384147375818985
-1.23104648181636 0.394059367323596
-1.18067338387306 0.403578551880993
-1.12912336502384 0.412695440560988
-1.07644781144858 0.421400945450675
-1.02269923128465 0.429686388713456
-0.967931202285677 0.437543511239297
-0.912198318414033 0.44496448087759
-0.855556135420419 0.451941900244412
-0.798061115464683 0.458468814096394
-0.739770570833094 0.464538716263865
-0.680742606808225 0.470145556136338
-0.621036063748338 0.475283744693891
-0.560710458434029 0.479948160078428
-0.499825924740618 0.484134152699254
-0.438443153695383 0.487837549867887
-0.376623332979435 0.49105465995749
-0.314428085934502 0.493782276082762
-0.251919410135435 0.496017679296629
-0.189159615589691 0.497758641300553
-0.126211262625347 0.499003426665747
-0.0631370995296205 0.499750794563087
-4.89858719658941e-16 0.5
};
\addplot [line width=0.52pt, gray127, dashed]
table {%
-0.445503262094184 0.226995249869773
-0.416617575056473 0.283137680568154
-0.387316594559491 0.338997873468965
-0.357629528457826 0.394520145917093
-0.327585969464592 0.449649152103885
-0.297215865652769 0.504329938237046
-0.266549490602316 0.55850799731978
-0.235617413222806 0.612129323484547
-0.204450467281669 0.66514046582729
-0.173079720668428 0.717488581688452
-0.141536444425547 0.769121489327695
-0.109852081576777 0.819987719939789
-0.0780582157840671 0.870036568959839
-0.0461865398642802 0.919218146606697
-0.0142688241971079 0.967483427614181
0.0176631149443375 1.01478430010053
0.0495774471086752 1.06107361352735
0.0814423593955418 1.10630522570036
0.113226088167383 1.15043404876484
0.144896950712137 1.19341609415028
0.176423376825253 1.23520851641907
0.207773940279555 1.2757696559758
0.238917390151591 1.31505908059447
0.269822681973231 1.35303762572215
0.300459008677471 1.38966743351911
0.330795831307583 1.42491199059634
0.360802909459014 1.45873616441286
0.390450331423681 1.4911062382966
0.419708544006606 1.52198994505395
0.448548381985191 1.55135649913435
0.476941097181737 1.57917662731806
0.504858387120264 1.60542259789628
0.532272423239025 1.63006824831473
0.559155878630625 1.65308901125311
0.585481955282079 1.67446193911427
0.611224410787649 1.6941657268989
0.636357584507851 1.71218073344284
0.660856423148531 1.72848900099576
0.68469650573454 1.74307427312192
0.707854067953092 1.75592201090485
0.730306025842545 1.76701940744013
0.752029998802995 1.77635540060155
0.773004331905752 1.78392068406809
0.793208117479436 1.78970771660067
0.812621215951201 1.79371072955939
0.831224275922303 1.7959257326539
0.848998753457988 1.79635051792092
0.865926930572497 1.79498466192528
0.881991932890733 1.79182952618192
0.89717774646901 1.78688825579878
0.911469233758096 1.78016577634164
0.924852148692657 1.77166878892421
0.937313150892037 1.76140576352835
0.948839818958248 1.74938693056096
0.959420662857882 1.73562427065613
0.969045135375633 1.72013150273255
0.977703642627984 1.70292407031823
0.985387553626604 1.68401912615602
0.992089208881904 1.66343551510538
0.997801928038178 1.64119375535748
1.00252001653274 1.61731601798212
1.00623877127236 1.59182610482722
1.00895448532144 1.56474942479254
1.01066445159716 1.53611296850159
1.01136696556796 1.5059452813968
1.01106132695262 1.47427643528474
1.00974784041836 1.44113799835993
1.00742781527712 1.40656300373699
1.00410356418043 1.37058591652248
0.999778400814071 1.33324259945935
0.994456636594945 1.29457027717819
0.988143576373363 1.25460749909087
0.98084551314505 1.21339410096369
0.972569721778146 1.17097116520819
0.963324451761466 1.12738097992934
0.953118918981224 1.08266699677181
0.941963296534442 1.03687378760643
0.929868704588185 0.990047000099996
0.916847199294733 0.942233312212661
0.90291176077375 0.893480385668351
0.888076280173416 0.843836818444506
0.872355545823431 0.793352096328566
0.855765228493684 0.742076543589459
0.838321865773289 0.690061272813279
0.82004284558556 0.637358133953147
0.800946388855343 0.584019662644047
0.781051531346012 0.530099027834167
0.760378104684198 0.475649978784932
0.738946716591197 0.420726791492574
0.71677873034074 0.365384214584635
0.693896243463625 0.309677414745351
0.670322065720412 0.253661921724291
0.646079696364157 0.197393572983105
0.621193300715849 0.140928458035523
0.595687686075892 0.0843228625361073
0.569588276995645 0.0276332121734919
0.542921089933682 -0.0290839835759756
0.515712707322013 -0.0857721877781061
0.487990251068145 -0.142374892398116
0.45978135551937 -0.198835674629064
0.431114139916248 -0.255098253135332
0.402017180362737 -0.311106544155086
0.372519481340911 -0.36680471740579
0.34265044679866 -0.422137251737042
0.312439850839197 -0.47704899047527
0.281917808041584 -0.531485196405095
0.251114743441865 -0.585391606332588
0.220061362204727 -0.638714485175995
0.188788619015925 -0.691400679530039
0.157327687225979 -0.743397670650385
0.125709927775897 -0.794653626805474
0.0939668579359138 -0.845117454943511
0.062130119888381 -0.89473885162314
0.0302314491861592 -0.943468353157001
-0.00169735688207434 -0.991257384918213
-0.0336244709880549 -1.03805830976062
-0.0655180674901006 -1.08382447550452
-0.0973463541576082 -1.12851026144061
-0.129077603862237 -1.17207112380564
-0.1606801862042 -1.21446364018466
-0.192122599042139 -1.25564555279542
-0.223373499895137 -1.29557581061181
-0.254401737185591 -1.33421461028444
-0.285176381291776 -1.37152343581751
-0.315666755379174 -1.40746509696234
-0.345842465979799 -1.44200376628946
-0.375673433289073 -1.47510501490214
-0.405129921150029 -1.5067358467559
-0.434182566694959 -1.53686473154964
-0.462802409614956 -1.56546163615578
-0.490960921028182 -1.59249805455788
-0.518630031918079 -1.6179470362661
-0.545782161113175 -1.64178321318202
-0.572390242780593 -1.66398282488611
-0.598427753405864 -1.68452374232264
-0.623868738232139 -1.70338548985838
-0.648687837132466 -1.72054926569319
-0.6728603098893 -1.73599796060206
-0.696362060856105 -1.74971617498991
-0.719169662976409 -1.76169023424232
-0.741260381136405 -1.7719082023566
-0.762612194827812 -1.78035989383994
-0.783203820098389 -1.78703688386244
-0.803014730768248 -1.79193251665526
-0.822025178890792 -1.79504191214514
-0.840216214437899 -1.79636197081905
-0.857569704189724 -1.7958913768138
-0.874068349810285 -1.79363059922771
-0.889695705090821 -1.78958189165306
-0.904436192343729 -1.78374928992961
-0.918275117930745 -1.77613860812159
-0.93119868690988 -1.76675743272215
-0.943194016786517 -1.75561511509095
-0.954249150354966 -1.74272276213252
-0.964353067617662 -1.72809322522467
-0.973495696770141 -1.71174108740794
-0.981667924240834 -1.69368264884893
-0.988861603775671 -1.67393591059193
-0.995069564558445 -1.65252055661507
-1.00028561835884 -1.62945793420894
-1.00450456570099 -1.6047710326971
-1.00772220104644 -1.57848446051987
-1.00993531698629 -1.55062442070406
-1.01114170743844 -1.52121868474325
-1.01134016984666 -1.49029656491456
-1.01053050637931 -1.4578888850595
-1.00871352412653 -1.42402794985816
-1.00589103429575 -1.38874751262713
-1.00206585040621 -1.35208274167352
-0.99724178548441 -1.31407018523838
-0.99142364826317 -1.27474773506466
-0.984617238388205 -1.23415458862585
-0.976829340636901 -1.19233121005314
-0.968067718155106 -1.14931928979983
-0.958341104718653 -1.10516170308347
-0.947659196027347 -1.05990246714679
-0.936032640040078 -1.01358669738047
-0.923473026360704 -0.966260562351072
-0.909992874685284 -0.917971237779223
-0.895605622322169 -0.868766859513837
-0.880325610797403 -0.818696475549246
-0.86416807155878 -0.767809997133103
-0.847149110792799 -0.716158149013752
-0.829285693369675 -0.663792418876705
-0.810595625932384 -0.610765006020596
-0.791097539146611 -0.557128769323779
-0.770810869129303 -0.502937174553459
-0.749755838074319 -0.44824424106984
-0.727953434094514 -0.393104487978462
-0.705425390300325 -0.337572879784369
-0.682194163135738 -0.281704771602288
-0.658282909993217 -0.225555853977458
-0.633715466129905 -0.169182097372071
-0.608516320908132 -0.112639696372711
-0.582710593383877 -0.0559850136743581
-0.556324007267548 0.00072547610316917
-0.529382865282035 0.0574352427103378
-0.501914022943573 0.114087756618493
-0.473944861791589 0.170626545369709
-0.445503262094184 0.226995249869773
};

\nextgroupplot[
height=4.6cm,
label style={font=\footnotesize},
legend cell align={left},
legend columns=3,
legend style={
  fill opacity=0.8,
  draw opacity=1,
  text opacity=1,
  at={(-1.3,-0.45)},
  anchor=north west,
  draw=none
},
legend style={draw=none, fill=white, fill opacity=0.9, text opacity=1, font=\footnotesize},
tick align=outside,
tick label style={font=\scriptsize},
tick pos=left,
title style={font=\footnotesize, yshift=-2pt},
title={déterminant le long du chemin},
width=0.45\textwidth,
x grid style={darkgray176},
xlabel={\(\displaystyle t\)},
xmin=-0.05, xmax=1.05,
xtick style={color=black},
y grid style={darkgray176},
ylabel={\(\displaystyle \det\)},
ymin=0.860448561810231, ymax=3.93058020198513,
ytick style={color=black}
]
\addplot [line width=0.56pt, darkorange25512714]
table {%
0 1
0.01 1.1105247390463
0.02 1.21881665508156
0.03 1.32487574810578
0.04 1.42870201811897
0.05 1.53029546512112
0.06 1.62965608911223
0.07 1.72678389009231
0.08 1.82167886806136
0.09 1.91434102301936
0.1 2.00477035496633
0.11 2.09296686390226
0.12 2.17893054982716
0.13 2.26266141274102
0.14 2.34415945264385
0.15 2.42342466953563
0.16 2.50045706341639
0.17 2.5752566342861
0.18 2.64782338214478
0.19 2.71815730699243
0.2 2.78625840882903
0.21 2.8521266876546
0.22 2.91576214346914
0.23 2.97716477627264
0.24 3.0363345860651
0.25 3.09327157284652
0.26 3.14797573661691
0.27 3.20044707737626
0.28 3.25068559512458
0.29 3.29869128986186
0.3 3.3444641615881
0.31 3.38800421030331
0.32 3.42931143600748
0.33 3.46838583870062
0.34 3.50522741838272
0.35 3.53983617505378
0.36 3.57221210871381
0.37 3.6023552193628
0.38 3.63026550700075
0.39 3.65594297162767
0.4 3.67938761324355
0.41 3.70059943184839
0.42 3.7195784274422
0.43 3.73632460002498
0.44 3.75083794959671
0.45 3.76311847615741
0.46 3.77316617970707
0.47 3.7809810602457
0.48 3.78656311777329
0.49 3.78991235228984
0.5 3.79102876379536
0.51 3.78991235228985
0.52 3.78656311777329
0.53 3.7809810602457
0.54 3.77316617970707
0.55 3.76311847615741
0.56 3.75083794959671
0.57 3.73632460002497
0.58 3.7195784274422
0.59 3.70059943184839
0.6 3.67938761324355
0.61 3.65594297162767
0.62 3.63026550700075
0.63 3.6023552193628
0.64 3.57221210871381
0.65 3.53983617505378
0.66 3.50522741838272
0.67 3.46838583870062
0.68 3.42931143600748
0.69 3.38800421030331
0.7 3.34446416158811
0.71 3.29869128986186
0.72 3.25068559512458
0.73 3.20044707737626
0.74 3.14797573661691
0.75 3.09327157284652
0.76 3.0363345860651
0.77 2.97716477627264
0.78 2.91576214346914
0.79 2.8521266876546
0.8 2.78625840882903
0.81 2.71815730699243
0.82 2.64782338214478
0.83 2.5752566342861
0.84 2.50045706341639
0.85 2.42342466953564
0.86 2.34415945264385
0.87 2.26266141274102
0.88 2.17893054982716
0.89 2.09296686390226
0.9 2.00477035496633
0.91 1.91434102301936
0.92 1.82167886806136
0.93 1.72678389009231
0.94 1.62965608911223
0.95 1.53029546512112
0.96 1.42870201811897
0.97 1.32487574810578
0.98 1.21881665508156
0.99 1.1105247390463
1 1
};
\addlegendentry{euclidienne}
\addplot [line width=0.56pt, forestgreen4416044]
table {%
0 1
0.01 1
0.02 1
0.03 1
0.04 1
0.05 1
0.06 1
0.07 1
0.08 1
0.09 1
0.1 1
0.11 1
0.12 1
0.13 1
0.14 1
0.15 1
0.16 1
0.17 1
0.18 1
0.19 1
0.2 1
0.21 1
0.22 1
0.23 1
0.24 1
0.25 1
0.26 1
0.27 1
0.28 1
0.29 1
0.3 1
0.31 1
0.32 1
0.33 1
0.34 1
0.35 1
0.36 1
0.37 1
0.38 1
0.39 1
0.4 1
0.41 1
0.42 1
0.43 1
0.44 1
0.45 1
0.46 1
0.47 1
0.48 1
0.49 1
0.5 1
0.51 1
0.52 1
0.53 1
0.54 1
0.55 1
0.56 1
0.57 1
0.58 1
0.59 1
0.6 1
0.61 1
0.62 1
0.63 1
0.64 1
0.65 1
0.66 1
0.67 1
0.68 1
0.69 1
0.7 1
0.71 1
0.72 1
0.73 1
0.74 1
0.75 1
0.76 1
0.77 1
0.78 1
0.79 1
0.8 1
0.81 1
0.82 1
0.83 1
0.84 1
0.85 1
0.86 1
0.87 1
0.88 1
0.89 1
0.9 1
0.91 1
0.92 1
0.93 1
0.94 1
0.95 1
0.96 1
0.97 1
0.98 1
0.99 1
1 1
};
\addlegendentry{affine invariante}
\addplot [line width=0.56pt, crimson2143940]
table {%
0 1
0.01 1
0.02 1
0.03 1
0.04 1
0.05 1
0.06 1
0.07 1
0.08 1
0.09 1
0.1 1
0.11 1
0.12 1
0.13 1
0.14 1
0.15 1
0.16 1
0.17 1
0.18 1
0.19 1
0.2 1
0.21 1
0.22 1
0.23 1
0.24 1
0.25 1
0.26 1
0.27 1
0.28 1
0.29 1
0.3 1
0.31 1
0.32 1
0.33 1
0.34 1
0.35 1
0.36 1
0.37 1
0.38 1
0.39 1
0.4 1
0.41 1
0.42 1
0.43 1
0.44 1
0.45 1
0.46 1
0.47 1
0.48 1
0.49 1
0.5 1
0.51 1
0.52 1
0.53 1
0.54 1
0.55 1
0.56 1
0.57 1
0.58 1
0.59 1
0.6 1
0.61 1
0.62 1
0.63 1
0.64 1
0.65 1
0.66 1
0.67 1
0.68 1
0.69 1
0.7 1
0.71 1
0.72 1
0.73 1
0.74 1
0.75 1
0.76 1
0.77 1
0.78 1
0.79 1
0.8 1
0.81 1
0.82 1
0.83 1
0.84 1
0.85 1
0.86 1
0.87 1
0.88 1
0.89 1
0.9 1
0.91 1
0.92 1
0.93 1
0.94 1
0.95 1
0.96 1
0.97 1
0.98 1
0.99 1
1 1
};
\addlegendentry{log-euclidienne}
\end{groupplot}

\end{tikzpicture}

  \end{center}
  \provInterpolationQR[.5cm]
  \caption{Trois manières de joindre les deux mêmes matrices de covariance
    $2\times 2$, de même déterminant et de directions propres différentes: le
    segment euclidien $(1-t)\MCov_0+t\MCov_1$, la géodésique affine invariante
    \cref{eq:context-riemann-geodspd} et la géodésique log-euclidienne
    \cref{eq:context-riemann-geodLE}. Les extrémités, communes aux trois
    panneaux, sont en pointillés. Le dernier panneau suit le déterminant le
    long des trois chemins: seul le chemin euclidien le fait gonfler, jusqu'à
    plus de quatre fois celui des extrémités, alors que les deux autres le
    conservent. \provInterpolation}
  \label{fig:context-riemann-interpolation}
\end{figure}


\section{Moyennes sur le cône}
\label{sec:context-riemann-moyennes}

Nous avons maintenant de quoi mesurer un écart entre deux matrices de covariance. La question immédiatement suivante, et celle dont nous aurons le plus besoin en pratique, est celle de la moyenne d'une collection de matrices --- c'est l'opération élémentaire de tout algorithme de classification, de partitionnement ou d'agrégation.

Or la moyenne arithmétique ne fait pas de sens sur une variété: elle utilise l'addition et la division, qui n'y sont pas définies. La sortie est de caractériser la moyenne non par sa formule mais par la propriété qu'elle vérifie, à savoir minimiser la dispersion.

\begin{definition}[{Moyenne de Fréchet}]
\label{def:context-riemann-frechet}
  Soit $\{\MCov_\iobs\}_{\iobs=1}^{\Nobs}$ une collection de points d'une variété riemannienne munie d'une distance $\geodist$. On appelle \textbf{moyenne de Fréchet}, ou \textbf{barycentre riemannien}, tout minimiseur de la variance de Fréchet \sidecite[-0.5cm]{frechet1948elements,karcher1977riemannian}:
  \begin{equation}
    \Fmean
    \in
    \argmin_{\MCov\in\SPD{\Nvar}}
    \;
    \frac{1}{\Nobs}\sum_{\iobs=1}^{\Nobs}\geodist^2(\MCov,\MCov_\iobs).
    \label{eq:context-riemann-frechet}
  \end{equation}
\end{definition}

Cette définition ne demande rien d'autre qu'une distance, et elle redonne bien les objets attendus dans chacune des géométries de la \cref{tab:context-riemann-metriques}:
\begin{itemize}
  \item pour la métrique \textbf{euclidienne}, on retrouve la moyenne arithmétique $\frac{1}{\Nobs}\sum_\iobs\MCov_\iobs$;
  \item pour la métrique \textbf{log-euclidienne}, la propriété d'isométrie avec $\Sym{\Nvar}$ donne immédiatement une forme close, $\Fmean_{\mathrm{LE}}=\expm\!\left(\frac{1}{\Nobs}\sum_\iobs\logm\MCov_\iobs\right)$;
  \item pour la métrique \textbf{affine invariante}, il n'existe pas de forme close pour $\Nobs>2$. La condition d'optimalité s'écrit
  \begin{equation}
    \sum_{\iobs=1}^{\Nobs}\log_{\Fmean}(\MCov_\iobs) = \vector{0},
    \label{eq:context-riemann-frechetopt}
  \end{equation}
  soit «~la moyenne est le point depuis lequel les autres se compensent~», et il faut recourir à un algorithme itératif \sidecite[0.5cm]{jeuris2012survey}. Ce sera notre premier exemple concret d'optimisation riemannienne, en \cref{sec:context-riemann-optimization}.
\end{itemize}

Le cas de deux points est instructif car il est le seul à admettre une forme close pour la métrique affine invariante: on obtient le milieu de la géodésique, $\Fmean=\MCov_0\geodmean_{1/2}\MCov_1=\MCov_0^{1/2}(\MCov_0^{-1/2}\MCov_1\MCov_0^{-1/2})^{1/2}\MCov_0^{1/2}$. Pour $\Nvar=1$ cette expression se réduit à $\sqrt{\sigma_0\sigma_1}$: c'est la moyenne géométrique. La moyenne riemannienne du cône est donc la généralisation matricielle exacte de la moyenne géométrique scalaire, et l'absence d'effet de gonflement en découle directement puisque $\Det(\Fmean)=\sqrt{\Det(\MCov_0)\Det(\MCov_1)}$. La \cref{fig:context-riemann-moyennes} compare les trois moyennes --- arithmétique, log-euclidienne, de Fréchet --- sur un échantillon, et rend cet effet visible.

\begin{remark}[{Existence et unicité}]
  Rien ne garantit \emph{a priori} qu'un minimiseur de \cref{eq:context-riemann-frechet} existe ni qu'il soit unique --- sur une sphère, la moyenne de deux points antipodaux ne l'est pas. Pour le cône muni de la métrique affine invariante, tout se passe bien: la variété est complète, simplement connexe et de courbure sectionnelle négative ou nulle, ce qui en fait une variété de Hadamard, et la fonction objectif de \cref{eq:context-riemann-frechet} y est strictement convexe le long des géodésiques. Le minimiseur existe donc et est unique. Nous admettons ce résultat, dont la démonstration demande l'appareil de courbure que nous avons choisi de ne pas introduire, et renvoyons à \sidecite[0cm]{moakher2005differential,bhatia2007positive}.
\end{remark}

\begin{remark}[{Au-delà de la moyenne}]
  La même mécanique --- remplacer une formule par le problème de minimisation qu'elle résout --- s'étend à d'autres statistiques descriptives. En remplaçant le carré de la distance par la distance elle-même dans \cref{eq:context-riemann-frechet}, on obtient la \textbf{médiane de Fréchet}, plus robuste aux points aberrants \sidecite[0cm]{arnaudon2013riemannian}. Plus loin encore, il est possible de définir de véritables lois de probabilité sur la variété, dont des analogues riemanniens de la loi gaussienne \sidecite[1.5cm]{said2017riemannian}, et de développer sur cette base une statistique inférentielle intrinsèque. Nous n'irons pas dans cette direction ici, mais elle constitue un prolongement naturel des outils de ce chapitre.
\end{remark}

\begin{figure}[h]
  \begin{center}
    % This file was created with matplot2tikz v0.5.3.
\begin{tikzpicture}

\definecolor{crimson2143940}{RGB}{214,39,40}
\definecolor{darkgray176}{RGB}{176,176,176}
\definecolor{darkorange25512714}{RGB}{255,127,14}
\definecolor{forestgreen4416044}{RGB}{44,160,44}
\definecolor{mediumpurple148103189}{RGB}{148,103,189}
\definecolor{sienna1408675}{RGB}{140,86,75}
\definecolor{steelblue31119180}{RGB}{31,119,180}

\begin{groupplot}[group style={group size=2 by 1, vertical sep=1.5cm, horizontal sep=1.4cm}]
\nextgroupplot[
height=4.6cm,
label style={font=\footnotesize},
legend cell align={left},
legend columns=2,
legend style={
  fill opacity=0.8,
  draw opacity=1,
  text opacity=1,
  at={(0,-0.45)},
  anchor=north west,
  draw=none
},
legend style={draw=none, fill=white, fill opacity=0.9, text opacity=1, font=\footnotesize},
tick align=outside,
tick label style={font=\scriptsize},
tick pos=left,
title style={font=\footnotesize, yshift=-2pt},
title={ellipses de concentration},
width=0.45\textwidth,
x grid style={darkgray176},
xlabel={\(\displaystyle x_1\)},
xmin=-1.77650169622035, xmax=1.77650169622035,
xtick style={color=black},
y grid style={darkgray176},
ylabel={\(\displaystyle x_2\)},
ymin=-1.77650169622035, ymax=1.77650169622035,
ytick style={color=black}
]
\addplot [line width=0.36pt, steelblue31119180, opacity=0.7]
table {%
-1.29428725258679 -0.681523228387517
-1.31939520383936 -0.632275709933187
-1.343187953259 -0.582397925300684
-1.36564178371564 -0.531939593750454
-1.38673431274277 -0.480951013244735
-1.40644451484875 -0.429483010309459
-1.42475274247545 -0.377586889369259
-1.44164074558343 -0.325314381606107
-1.45709168984391 -0.272717593392546
-1.47109017341963 -0.219848954350924
-1.48362224231773 -0.166761165090406
-1.49467540429938 -0.113507144673862
-1.50423864133235 -0.0601399778669939
-1.51230242057398 -0.00671286222228407
-1.51885870387375 0.0467209449494778
-1.52390095578589 0.100108179667243
-1.52742415008403 0.153395624374398
-1.52942477477149 0.206530160987182
-1.52990083558205 0.259458823843948
-1.52885185796795 0.312128852502486
-1.52627888757288 0.364487744332779
-1.52218448918971 0.416483306852769
-1.51657274420377 0.468063709754954
-1.50944924652451 0.519177536571968
-1.50082109700933 0.569773835929635
-1.49069689638533 0.619802172336403
-1.47908673667584 0.669212676458545
-1.46600219114058 0.717956094830987
-1.45145630273908 0.765983838954243
-1.43546357112924 0.813248033728479
-1.41803993821374 0.859701565176458
-1.39920277224878 0.905298127407779
-1.37897085053108 0.949992268777591
-1.35736434068018 0.993739437193783
-1.33440478053501 1.03649602452748
-1.31011505668443 1.07821941008257
-1.28451938165342 1.11886800308096
-1.25764326976749 1.15840128412112
-1.22951351171947 1.19677984556875
-1.20015814786394 1.23396543083918
-1.169606440266 1.26992097253229
-1.1378888435322 1.30461062938218
-1.10503697445276 1.33799982198448
-1.07108358048522 1.37005526726588
-1.0360625071111 1.40074501166138
-1.00000866409797 1.43003846296634
-0.962957990700626 1.45790642083143
-0.924947419836074 1.4843211058702
-0.886014841267994 1.50925618735013
-0.846199063837395 1.53268680943974
-0.80553977677712 1.55458961598536
-0.764077510148755 1.57494277379309
-0.721853594441381 1.59372599439259
-0.678910119372446 1.61092055426108
-0.63528989193182 1.62650931348735
-0.591036393710853 1.64047673285717
-0.54619373755898 1.65280888934316
-0.500806623611078 1.66349348998355
-0.454920294729391 1.67251988413603
-0.408580491404469 1.67987907409463
-0.361833406160049 1.68556372405875
-0.314725637507341 1.68956816744567
-0.267304143494623 1.69188841253914
-0.219616194898431 1.69252214646839
-0.171709328103028 1.69146873751365
-0.123631297715097 1.68872923573589
-0.0754300289609131 1.68430637193005
-0.0271535699134343 1.67820455490296
0.0211499564030598 1.67042986707854
0.0694323999830496 1.66099005843468
0.117645631836745 1.64989453877794
0.165741591966108 1.63715436836361
0.213672337272111 1.62278224687063
0.261390089345447 1.60679250074227
0.308847282093104 1.58920106890521
0.355996609153266 1.57002548688126
0.402791071051335 1.54928486930763
0.449184022050019 1.52699989088296
0.495129216646826 1.50319276575839
0.540580855672581 1.477887225394
0.585493631945037 1.45110849490277
0.629822775432058 1.42288326790561
0.673524097879362 1.39323967992263
0.716554036858331 1.36220728032701
0.758869699189992 1.32981700288955
0.800428903701866 1.29610113494321
0.84119022327508 1.26109328519841
0.881113026139821 1.2248283502411
0.920157516377966 1.18734247974713
0.95828477359252 1.14867304044743
0.995456791704312 1.10885857888008
1.03163651683728 1.06793878296622
1.0667878842546 1.02595444244835
1.10087585430876 0.982947408230138
1.13386644736988 0.938960550658606
1.16572677769733 0.89403771678999
1.19642508622095 0.848223686682057
1.22593077219918 0.801564128756365
1.25421442372258 0.75410555427499
1.28124784703218 0.705895270977088
1.30700409462373 0.656981335921512
1.33145749210949 0.607412507582487
1.35458366381104 0.557238197246108
1.37635955705749 0.506508419756091
1.39676346516487 0.455273743657892
1.41577504907386 0.40358524079088
1.43337535762414 0.351494435378816
1.44954684644539 0.299053252669388
1.46427339544586 0.246313967173992
1.47754032488125 0.193329150559367
1.48933440998776 0.140151619243009
1.49964389416486 0.0868343817446212
1.50845850069453 0.0334305858460663
1.51576944298528 -0.0200065343874967
1.52156943333083 -0.0734237116724868
1.52585269017467 -0.12676769860489
1.52861494387323 -0.179985320738422
1.52985344095196 -0.233023529589985
1.52956694685002 -0.285829455519537
1.52775574715097 -0.338350460431708
1.52442164729806 -0.390534190246592
1.51956797079453 -0.442328627087437
1.51319955589068 -0.493682141133192
1.50532275076097 -0.544543542084237
1.49594540717608 -0.594862130189989
1.48507687267601 -0.644587746787513
1.47272798125231 -0.693670824300769
1.45891104254854 -0.742062435650649
1.44363982958965 -0.789714343026558
1.42692956505286 -0.836579045970905
1.40879690609325 -0.882609828728597
1.38925992773963 -0.927760806814321
1.36833810487689 -0.971986972751192
1.34605229283306 -1.01524424093519
1.32242470659027 -1.05748949158065
1.29747889864031 -1.09868061370298
1.27123973550707 -1.13877654709585
1.24373337295899 -1.17773732326086
1.21498722993647 -1.21552410524902
1.18502996122003 -1.25209922637427
1.15389142886662 -1.28742622776045
1.12160267244249 -1.32146989468431
1.08819587808224 -1.35419629167827
1.05370434640502 -1.3855727963581
1.0181624593197 -1.41556813194157
0.981605645752287 -1.44415239842583
0.944070346329571 -1.47129710239241
0.905593977054322 -1.49697518541007
0.8662148920082 -1.52116105100722
0.825972345119586 -1.54383059018705
0.784906451034395 -1.56496120545988
0.743058145128941 -1.58453183336888
0.700469142704653 -1.60252296548653
0.657181897405356 -1.61891666786107
0.613239558898563 -1.63369659889345
0.568685929862939 -1.64684802562699
0.523565422324842 -1.6583578384335
0.477923013387455 -1.66821456408127
0.431804200396619 -1.67640837717182
0.385254955588101 -1.68293110993405
0.33832168026146 -1.6877762603661
0.291051158526232 -1.69093899871665
0.243490510666507 -1.69241617229932
0.195687146170402 -1.69220630863536
0.147688716471259 -1.69030961692145
0.0995430674476534 -1.68672798782112
0.051298191729588 -1.68146499158015
0.00300218085839254 -1.67452587446766
-0.0452968226519742 -1.6659175535465
-0.0935506733044207 -1.65564860977816
-0.141711270611211 -1.64372927946913
-0.189730607041569 -1.63017144406707
-0.237560815876631 -1.61498861831716
-0.285154218924038 -1.5981959367903
-0.332463374044586 -1.57981013879664
-0.379441122443589 -1.55984955169946
-0.426040635679786 -1.53833407264609
-0.472215462344955 -1.51528514873394
-0.51791957436768 -1.49072575563165
-0.56310741289514 -1.46468037467635
-0.607733933707163 -1.43717496847022
-0.651754652117279 -1.40823695500033
-0.695125687316031 -1.37789518030785
-0.737803806112316 -1.34617988973366
-0.779746466029165 -1.31312269776913
-0.820911857711023 -1.27875655654212
-0.861258946600208 -1.2431157229696
-0.900747513841062 -1.20623572460958
-0.939338196370976 -1.16815332424646
-0.976992526158345 -1.12890648324515
-1.01367296854835 -1.08853432371027
-1.04934295967829 -1.04707708948846
-1.08396694292528 -1.00457610605237
-1.11751040434983 -0.961073739306594
-1.14993990710009 -0.916613353356382
-1.18122312474247 -0.871239267281355
-1.21132887348527 -0.824996710957295
-1.24022714326342 -0.777931779970009
-1.26788912765313 -0.730091389666259
-1.29428725258679 -0.681523228387517
};
\addlegendentry{échantillon}
\addplot [line width=0.36pt, steelblue31119180, opacity=0.7, forget plot]
table {%
-0.106070073624268 -0.678450496630078
-0.136616402509617 -0.673328420673853
-0.16702654919386 -0.667535156304588
-0.19727020018143 -0.661076478374195
-0.227317207943295 -0.653958825033298
-0.257137620968668 -0.646189291313539
-0.286701713621319 -0.637775622055098
-0.315980015770743 -0.62872620418647
-0.344943342168623 -0.61905005836419
-0.373562821541321 -0.608756829980847
-0.401809925369393 -0.597856779550347
-0.429656496325434 -0.586360772480009
-0.457074776341917 -0.574280268239694
-0.48403743428104 -0.561627308938749
-0.510517593178993 -0.548414507322179
-0.536488857037507 -0.534655034197982
-0.561925337135952 -0.520362605308204
-0.586801677837783 -0.50555146765679
-0.611093081865586 -0.490236385307858
-0.634775335019546 -0.474432624668561
-0.657824830314685 -0.458155939271198
-0.68021859151282 -0.441422554069754
-0.701934296025778 -0.424249149266509
-0.722950297167035 -0.406652843684852
-0.743245645729603 -0.388651177704865
-0.76280011086866 -0.37026209577869
-0.781594200268092 -0.351503928543113
-0.799609179570857 -0.332395374547178
-0.816827091053803 -0.312955481613076
-0.833230771528316 -0.293203627848854
-0.848803869448955 -0.273159502331903
-0.863530861213037 -0.252843085482449
-0.877397066634892 -0.232274629146642
-0.890388663579399 -0.211474636409074
-0.902492701740187 -0.190463841154857
-0.913697115548783 -0.169263187401635
-0.923990736201832 -0.147893808422135
-0.933363302794406 -0.126377005678056
-0.941805472548293 -0.104734227586313
-0.949308830125087 -0.0829870481387895
-0.955865896014773 -0.0611571453969076
-0.961470133991473 -0.0392662798824696
-0.966115957628891 -0.0173362728862953
-0.969798735868991 0.00461101528371521
-0.97251479763833 0.0265537070934121
-0.97426143550747 0.0484699295903977
-0.975036908389797 0.0703378362074602
-0.974840443277078 0.0921356285397085
-0.973672236010008 0.113841578073695
-0.971533451082998 0.13543404784687
-0.968426220483377 0.156891514015779
-0.964353641566182 0.178192587311495
-0.959319773966645 0.199316034360908
-0.953329635553465 0.220240798852614
-0.946389197426884 0.240946022526304
-0.938505377966564 0.261411065964729
-0.929686035935199 0.281615529167524
-0.919939962644724 0.301539271886368
-0.909276873192951 0.321162433701224
-0.897707396779342 0.340465453817637
-0.885243066109594 0.359429090565361
-0.87189630589959 0.378034440578871
-0.857680420490167 0.39626295764065
-0.842609580585064 0.414096471168463
-0.826698809125253 0.43151720432819
-0.809963966313748 0.448507791754156
-0.792421733805807 0.46505129685931
-0.774089598080295 0.481131228717981
-0.754985833008785 0.49673155850439
-0.735129481639763 0.511836735470538
-0.714540337216102 0.526431702447521
-0.693238923444729 0.54050191085485
-0.671246474038147 0.554033335202792
-0.648584911548209 0.567012487073279
-0.62527682551324 0.579426428565461
-0.601345449940296 0.591262785192484
-0.576814640145 0.60250975821665
-0.551708848972044 0.613156136410652
-0.526053102420059 0.623191307233174
-0.499872974695162 0.632605267407698
-0.473194562718026 0.641388632893991
-0.446044460109901 0.649532648242317
-0.418449730683519 0.657029195321066
-0.390437881465292 0.663870801409075
-0.362036835275713 0.670050646644609
-0.33327490289528 0.675562570823539
-0.304180754843694 0.680401079539972
-0.274783392800463 0.684561349663193
-0.245112120695394 0.688039234145467
-0.215196515497803 0.690831266155907
-0.185066397733545 0.692934662536285
-0.154751801759267 0.694347326575344
-0.124282945823506 0.695067850098849
-0.093690201944486 0.695095514873281
-0.0630040656346256 0.694430293321792
-0.0322551255019504 0.693072848551691
-0.00147403275870167 0.691024533693446
0.0293085293324577 0.688287390551854
0.0600618760441868 0.684864147570725
0.0907553517716707 0.68075821711311
0.121358360590792 0.675973692059784
0.151840396756817 0.670515341729371
0.182171075113183 0.664388607124188
0.212320161380083 0.657599595506531
0.242257602292655 0.650155074310826
0.271953555558726 0.642062464397701
0.301378419606265 0.633329832656713
0.330502863090868 0.623965883965096
0.35929785413388 0.613979952510549
0.387734689262007 0.603381992486717
0.415785022019558 0.592182568170633
0.443420891224812 0.580392843392015
0.47061474884233 0.568024570404922
0.497339487443435 0.555090078172849
0.523568467227491 0.541602260078947
0.549275542577029 0.52757456107362
0.574435088120273 0.513020964272307
0.599022024275063 0.497955977016806
0.623011842248734 0.482394616414041
0.646380628469006 0.466352394366684
0.669105088421558 0.449845302110552
0.691162569870508 0.432889794274197
0.712531085438647 0.415502772476573
0.733189334524943 0.397701568479135
0.753116724537436 0.379503926909161
0.772293391420391 0.360927987571524
0.790700219455209 0.341992267366539
0.808318860315401 0.322715641831915
0.825131751356592 0.303117326327205
0.841122133123353 0.283216856879522
0.856274066055385 0.2630340707096
0.870572446376422 0.24258908645762
0.884003021150003 0.221902284128517
0.896552402487112 0.20099428477674
0.908208080891512 0.179885929950744
0.918958437729486 0.158598260917677
0.928792756811542 0.137152497688987
0.937701235074546 0.115570017867856
0.945674992353627 0.0938723353395368
0.952706080234119 0.0720810788258486
0.958787489974713 0.0502179703251932
0.963913159493923 0.028304803459594
0.968077979412901 0.00636342175034393
0.971277798148577 -0.015584303156097
0.973509426052046 -0.0375164932902299
0.974770638588082 -0.0594112861679549
0.9750601785526 -0.0812468565835722
0.974377757325867 -0.103001438365618
0.972724055160203 -0.124653346073862
0.970100720501891 -0.146180996615822
0.966510368347969 -0.167562930761261
0.961956577639546 -0.188777834533218
0.956443887694234 -0.209804560454237
0.94997779368126 -0.230622148626638
0.942564741143757 -0.251209847625798
0.934212119573712 -0.271547135185621
0.924928255045955 -0.291613738655578
0.91472240191855 -0.311389655208928
0.903604733607848 -0.330855171781972
0.891586332447404 -0.34999088472446
0.878679178640867 -0.368777719141576
0.86489613831985 -0.387196947908201
0.850250950718697 -0.405230210336526
0.834758214478904 -0.422859530478383
0.818433373096888 -0.440067335044056
0.801292699529563 -0.456836470919726
0.783353279973113 -0.473150222266063
0.764632996831091 -0.488992327180937
0.745150510888854 -0.504346993909643
0.72492524271209 -0.519198916586457
0.703977353287967 -0.533533290491859
0.682327723928232 -0.547335826810203
0.65999793545426 -0.560592766873113
0.637010246684821 -0.573290895874427
0.613387572248009 -0.585417556043002
0.58915345973944 -0.596960659260257
0.564332066249499 -0.607908699109867
0.538948134283022 -0.618250762347618
0.513026967095435 -0.627976539779958
0.486594403469916 -0.63707633654043
0.459676791960731 -0.645541081753723
0.43230096462843 -0.653362337577721
0.404494210293055 -0.660532307614525
0.37628424733206 -0.667043844682072
0.347699196050027 -0.672890457938601
0.318767550647728 -0.67806631935287
0.289518150818499 -0.682566269513657
0.259980153000194 -0.686385822772779
0.230183001311419 -0.68952117171648
0.200156398200998 -0.691969190960748
0.169930274839918 -0.693727440266761
0.139534761285294 -0.694794166973372
0.109000156446066 -0.695168307744203
0.0783568978803953 -0.694849489627595
0.0476355314548396 -0.693838030428381
0.0168666808955679 -0.692134938391088
-0.0139189827380304 -0.689741911194896
-0.0446907716269311 -0.686661334261358
-0.0754180117827546 -0.682896278376556
-0.106070073624268 -0.678450496630078
};
\addplot [line width=0.36pt, steelblue31119180, opacity=0.7, forget plot]
table {%
0.0574480571897439 -0.859006166756709
0.0109631199473195 -0.861684903339426
-0.0355327455714906 -0.86350469367087
-0.0819931912767611 -0.864463723744463
-0.128371904385808 -0.864561037578171
-0.174622653588775 -0.863796538167448
-0.220699335133012 -0.862170987581933
-0.266556018780296 -0.859686006205801
-0.312146993591092 -0.856344071122531
-0.357426813490212 -0.852148513645692
-0.402350342568451 -0.847103515998222
-0.446872800075043 -0.841214107143492
-0.490949805056085 -0.834486157772329
-0.534537420594442 -0.826926374450984
-0.577592197607016 -0.818542292935878
-0.620071218155741 -0.809342270661801
-0.661932138229115 -0.799335478411034
-0.703133229951632 -0.788531891171716
-0.743633423179036 -0.776942278194561
-0.783392346437934 -0.764578192257831
-0.822370367168958 -0.751451958151277
-0.860528631233364 -0.737576660390523
-0.897829101643678 -0.72296613017413
-0.934234596479798 -0.707634931596363
-0.969708825952739 -0.691598347129383
-1.00421642857908 -0.674872362389343
-1.03772300643006 -0.657473650201584
-1.07019515942019 -0.639419553980796
-1.10160051860119 -0.620728070442726
-1.13190777842801 -0.601417831664659
-1.16108672796496 -0.581508086512564
-1.18910828100056 -0.561018681453397
-1.21594450504135 -0.539970040771718
-1.24156864915558 -0.518383146210317
-1.26595517063912 -0.496279516055157
-1.28907976047702 -0.473681183685479
-1.31091936757522 -0.450610675610449
-1.33145222173839 -0.427090989014247
-1.35065785537098 -0.403145568831968
-1.3685171238797 -0.378798284379208
-1.3850122247573 -0.354073405558606
-1.40012671532849 -0.328995578667077
-1.41384552914038 -0.303589801827852
-1.426154990981 -0.277881400071799
-1.43704283051109 -0.251896000092886
-1.44649819449541 -0.225659504702934
-1.45451165762152 -0.19919806701113
-1.46107523189511 -0.172538064354039
-1.46618237460258 -0.145706072002101
-1.46982799483302 -0.118728836668828
-1.4720084585529 -0.0916332498490885
-1.47272159222857 -0.064446321013092
-1.47196668499287 -0.0371951506827532
-1.46974448935375 -0.00990690341730685
-1.46605722044415 0.0173912192649178
-1.46090855381391 0.0446720060015611
-1.45430362176589 0.0719082627111106
-1.44624900824001 0.0990728397005538
-1.4367527422502 0.126138658728782
-1.42582428988091 0.153078739998769
-1.41347454485113 0.179866229051623
-1.39971581765529 0.20647442353569
-1.38456182329189 0.232876799824042
-1.36802766759205 0.259047039453806
-1.35012983216177 0.284959055360977
-1.33088615795264 0.310587017884574
-1.31031582747761 0.335905380514205
-1.28843934568945 0.360888905355386
-1.26527851954099 0.385512688287223
-1.2408564362474 0.409752183787378
-1.21519744027245 0.433583229399587
-1.18832710906135 0.456982069819312
-1.1602722275446 0.479925380573555
-1.13106076143825 0.502390291271191
-1.10072182936697 0.524354408400678
-1.06928567383807 0.545795837652395
-1.03678363109503 0.566693205743361
-1.00324809988091 0.587025681722595
-0.968712509142535 0.606772997735858
-0.933211284707823 0.62591546922909
-0.89677981496936 0.644434014570402
-0.859454415608479 0.662310174071063
-0.821272293395004 0.679526128386518
-0.782271509098728 0.6960647162791
-0.742490939549614 0.711909451724721
-0.701970238884522 0.727044540346504
-0.660749799019105 0.741454895158957
-0.618870709384269 0.755126151607012
-0.576374715967333 0.768044681884925
-0.533304179698724 0.780197608520766
-0.489702034225677 0.791572817212966
-0.445611743115045 0.802158968906119
-0.401077256527867 0.811945511093991
-0.356142967408892 0.820922688338496
-0.31085366723473 0.829081551994124
-0.265254501364732 0.836413969128146
-0.219390924039111 0.842912630627702
-0.17330865306917 0.848571058485677
-0.127053624264793 0.853383612258124
-0.0806719456446307 0.857345494686783
-0.0342098514746247 0.860452756481086
0.0122863438193151 0.862702300254901
0.0587702918185366 0.864091883614072
0.105195656312758 0.864620121391687
0.151516159489008 0.864286487028841
0.197685628062359 0.863091313099522
0.243658039302458 0.861035790979095
0.28938756690999 0.858121969656714
0.334828626697327 0.85435275369284
0.379935922027835 0.849731900323918
0.424664488968551 0.84426401571708
0.468969741111208 0.837954550378614
0.512807514016934 0.830809793720787
0.556134109240328 0.822836867792416
0.59890633788902 0.814043720179459
0.641081563675298 0.804439116082688
0.682617745416885 0.794032629580343
0.723473478944507 0.782834634084487
0.763608038374467 0.770856292000559
0.802981416705094 0.758109543600451
0.841554365696588 0.744607095120179
0.879288434994524 0.730362406094034
0.916146010457997 0.715389675937828
0.952090351654222 0.699703829794601
0.987085628482187 0.683320503656918
1.02109695688889 0.666256028780574
1.0540904336425 0.648527415405236
1.08603317012786 0.630152335798274
1.11689332513057 0.611149106638654
1.14664013657698 0.591536670758478
1.17524395219852 0.571334578260357
1.20267625908971 0.550562967029439
1.22890971213045 0.529242542659535
1.25391816124424 0.507394557813324
1.27767667746513 0.485040791037245
1.30016157778745 0.462203525052161
1.32135044877357 0.438905524541464
1.34122216889608 0.415170013458739
1.35975692959221 0.391020651877625
1.37693625500939 0.366481512406938
1.39274302042245 0.341577056194571
1.40716146930386 0.316332108544091
1.42017722903021 0.290771834168333
1.43177732520915 0.264921712104671
1.44195019461252 0.238807510316958
1.45068569670286 0.212455260009458
1.45797512374173 0.185891229678378
1.46381120946971 0.159141898926857
1.46818813634964 0.132233932069522
1.47110154136561 0.105194151552923
1.47254852037214 0.0780495112183299
1.47252763098909 0.0508270694335606
1.47103889403945 0.0235539621206113
1.46808379352858 -0.00374262429402404
1.46366527516496 -0.031035479979369
1.45778774342376 -0.0582973988233225
1.45045705715646 -0.0855012055522972
1.44168052375055 -0.112619782820126
1.4314668918454 -0.139626098239225
1.41982634261141 -0.166493231327058
1.40677047960122 -0.193194400341069
1.39231231718301 -0.219702988975294
1.3764662675675 -0.245992572892084
1.35924812644157 -0.272036946062455
1.34067505722275 -0.297810146888826
1.32076557395035 -0.323286484084106
1.29953952283029 -0.34844056228133
1.27701806245194 -0.373247307348306
1.25322364269679 -0.397681991382066
1.22817998235993 -0.421720257358165
1.20191204550662 -0.445338143410305
1.17444601658759 -0.468512106716035
1.14580927433777 -0.491219046964751
1.11603036448464 -0.513436329384592
1.08513897129315 -0.535141807305266
1.05316588797584 -0.556313844234342
1.02014298599752 -0.576931335424964
0.986103183305035 -0.596973728913533
0.951080411514004 -0.616421046006343
0.915109582085012 -0.635253901194784
0.878226551523111 -0.653453521479232
0.840468085635287 -0.67100176508239
0.801871822881493 -0.687881139533401
0.76247623685582 -0.704074819104725
0.722320597935183 -0.71956666158439
0.681444934133746 -0.734341224366903
0.639889991202137 -0.748383779846769
0.597697192011192 -0.761680330099293
0.55490859526074 -0.774217620834016
0.511566853554573 -0.785983154606875
0.467715170883404 -0.796965203277933
0.423397259558192 -0.807152819702242
0.378657296636756 -0.816535848642193
0.333539879887124 -0.825104936890479
0.288089983331513 -0.832851542593578
0.242352912415235 -0.839767943766453
0.196374258845255 -0.84584724599
0.150199855143379 -0.851083389283553
0.103875728959401 -0.855471154145602
0.0574480571897443 -0.859006166756709
};
\addplot [line width=0.36pt, steelblue31119180, opacity=0.7, forget plot]
table {%
-0.00359705609146724 -0.722279457905249
-0.048008802975559 -0.721698280122709
-0.092372693640766 -0.720397697799301
-0.136644505196063 -0.718379007383767
-0.180780106536883 -0.715644221150616
-0.224735502335964 -0.71219606519424
-0.268466876898847 -0.708037976711486
-0.311930637840313 -0.703174100575382
-0.35508345953822 -0.697609285203439
-0.397882326321428 -0.691349077724641
-0.440284575348753 -0.684399718449947
-0.482247939136214 -0.676768134651815
-0.523730587690183 -0.668461933658944
-0.564691170204432 -0.659489395273124
-0.605088856279514 -0.649859463515749
-0.644883376623394 -0.639581737712225
-0.684035063192757 -0.628666462923147
-0.722504888734972 -0.617124519731808
-0.760254505691311 -0.60496741339819
-0.797246284422626 -0.592207262390277
-0.833443350719393 -0.578856786304094
-0.868809622558727 -0.564929293184548
-0.903309846071725 -0.55043866625967
-0.936909630685289 -0.535399350101514
-0.969575483403402 -0.519826336227499
-1.00127484219367 -0.50373514815653
-1.03197610844586 -0.487141825934821
-1.06164867847009 -0.470062910146829
-1.09026297400326 -0.452515425427239
-1.11779047169328 -0.434516863490434
-1.1442037315318 -0.41608516569438
-1.16947642420698 -0.397238705156286
-1.19358335734914 -0.377996268437883
-1.21650050064307 -0.35837703681857
-1.23820500978192 -0.338400567175097
-1.25867524923897 -0.318086772486846
-1.27789081383435 -0.297455901986139
-1.29583254907537 -0.276528520973365
-1.31248257025016 -0.255325490317039
-1.3278242802555 -0.233867945659234
-1.34184238614123 -0.212177276347118
-1.35452291435453 -0.190275104111581
-1.36585322466911 -0.168183261514229
-1.37582202278524 -0.145923770184208
-1.38441937158815 -0.123518818866569
-1.3916367010536 -0.100990741304038
-1.3974668167906 -0.0783619939742647
-1.401903907213 -0.055655133704708
-1.40494354933258 -0.0328927951875061
-1.40658271316795 -0.0100976684167202
-1.40681976476497 0.0127075239295442
-1.40565446782544 0.0355000491395967
-1.4030879839427 0.0582571871286243
-1.3991228714437 0.0809562530865648
-1.3937630828388 0.10357462009082
-1.38701396088185 0.126089741661264
-1.3788822332444 0.148479174235065
-1.3693760058094 0.170720599538914
-1.3585047545911 0.19279184683637
-1.34627931628913 0.21467091502813
-1.33271187748626 0.236335994583202
-1.31781596250057 0.25776548927912
-1.30160641990405 0.27893803772953
-1.28409940772129 0.299832534677684
-1.26531237732276 0.320428152034617
-1.24526405602894 0.340704359641043
-1.2239744284425 0.360640945732258
-1.20146471652722 0.380218037085673
-1.17775735845347 0.399416118830871
-1.15287598623136 0.418216053902456
-1.12684540215382 0.436599102116294
-1.09969155407318 0.454546938850136
-1.07144150953572 0.472041673310004
-1.04212342880023 0.489065866364118
-1.01176653676716 0.505602547926606
-0.980401093846701 0.521635233873652
-0.948058365794472 0.537147942475231
-0.914770592545175 0.552125210326044
-0.880570956075115 0.566552107759777
-0.84549354732569 0.580414253731316
-0.80957333222079 0.593697830152089
-0.772846116812013 0.606389595664232
-0.735348511586409 0.618476898839867
-0.697117894972351 0.629947690792317
-0.658192376079908 0.640790537186703
-0.618610756712849 0.65099462963793
-0.578412492690167 0.66054979648473
-0.537637654515658 0.669446512928986
-0.49632688743477 0.677675910530264
-0.454521370918542 0.685229786046066
-0.412262777615011 0.692100609608998
-0.369593231809005 0.698281532232708
-0.326555267431747 0.703766392639104
-0.283191785662103 0.708549723400051
-0.239546012161753 0.712626756387419
-0.195661453986917 0.715993427526061
-0.151581856219571 0.71864638084497
-0.107351158361395 0.72058297182258
-0.0630134505339283 0.721801270022889
-0.0186129295285746 0.722300061019757
0.0258061452497202 0.722078847607469
0.070199495901207 0.721137850296366
0.114522870168505 0.719478007093033
0.158732085548216 0.717100972565271
0.202783073333002 0.714009116192794
0.246631922540234 0.710205520005273
0.290234923683413 0.705693975510109
0.333548612342735 0.700478979912968
0.376529812491375 0.694565731634878
0.419135679534284 0.687960125130319
0.461323743016617 0.680668745011512
0.503051948959202 0.672698859484731
0.544278701778865 0.664058413105195
0.584962905751811 0.654756018857763
0.625064005978734 0.644800949571328
0.664542028810823 0.634203128675453
0.703357621696363 0.622973120308486
0.74147209240822 0.611122118786991
0.778847447613086 0.598661937447019
0.815446430744069 0.585604996868308
0.85123255913884 0.571964312493184
0.886170160406352 0.557753481652486
0.92022440798585 0.542986670011448
0.953361355862747 0.527678597449057
0.985547972406748 0.511844523384958
1.0167521732985 0.495500231568532
1.04694285351194 0.478662014345312
1.07608991832044 0.461346656416417
1.10416431329593 0.4435714181072
1.13113805327095 0.425354018161776
1.15698425023487 0.406712616080598
1.18167714013647 0.387665794018665
1.20519210856607 0.368232538262434
1.22750571529176 0.34843222030387
1.24859571762507 0.328284577530523
1.26844109259303 0.307809693550876
1.28702205789428 0.287027978174558
1.30432009161851 0.265960147067412
1.32031795070952 0.24462720110166
1.33499968815344 0.223050405421782
1.3483506688751 0.201251268246958
1.36035758432656 0.179251519431202
1.37100846575339 0.157073088802578
1.38029269612533 0.134738084303063
1.38820102071965 0.112268769950879
1.39472555634641 0.0896875436472256
1.39985979920659 0.0670169148495708
1.40359863137526 0.0442794821337339
1.4059383259032 0.0214979106671307
1.40687655053202 -0.00130509038435398
1.40641237001901 -0.0241067904933615
1.40454624706939 -0.0468844604293391
1.40128004187512 -0.0696153949155332
1.39661701026055 -0.0922769352621047
1.39056180043701 -0.1148464919528
1.38312044836934 -0.137301567162675
1.37430037175911 -0.159619777184416
1.36411036265051 -0.181778874740903
1.35256057866622 -0.203756771161786
1.33966253288209 -0.225531558401942
1.32542908235068 -0.247081530879902
1.30987441528501 -0.268385207114437
1.29301403691548 -0.289421351137773
1.27486475403394 -0.310168993664059
1.25544465824027 -0.330607452992013
1.23477310790828 -0.35071635562089
1.21287070888892 -0.370475656559226
1.18975929396984 -0.389865659306127
1.16546190111207 -0.40886703548517
1.1400027504852 -0.427460844111342
1.11340722032427 -0.445628550471835
1.08570182163212 -0.463352044601843
1.05691417175271 -0.480613659336971
1.02707296684149 -0.497396187924255
0.996207953260501 -0.513682901174215
0.964349897926514 -0.529457564136889
0.931530557641928 -0.544704452285171
0.897782647438882 -0.559408367189372
0.863139807968202 -0.573554651667341
0.827636571965677 -0.58712920439506
0.791308329829063 -0.600118493963156
0.754191294340178 -0.612509572365301
0.716322464567203 -0.624290087905068
0.677739588983215 -0.635448297508366
0.638481127837685 -0.6459730784292
0.59858621481846 -0.655853939337057
0.55809461804246 -0.6650810307749
0.517046700413946 -0.673645154977315
0.475483379389908 -0.681537775039048
0.433446086192647 -0.688751023424775
0.390976724510234 -0.695277709811634
0.348117628726004 -0.701111328256696
0.304911521718718 -0.706246063682225
0.261401472275473 -0.710676797672283
0.217630852159796 -0.714399113574868
0.173643292877725 -0.717409300904538
0.129482642184981 -0.719704359041099
0.0851929203785619 -0.721282000220687
0.0408182764163637 -0.722140651816265
-0.0035970560914669 -0.722279457905249
};
\addplot [line width=0.36pt, steelblue31119180, opacity=0.7, forget plot]
table {%
0.0727609943744723 -0.613653329082534
0.02925871616423 -0.618501251231454
-0.0142727277705933 -0.622732637882197
-0.0577899443399765 -0.626343271096525
-0.101249554636046 -0.629329551716729
-0.144608237174114 -0.631688502953341
-0.187822771076478 -0.633417773352468
-0.230850079155927 -0.634515639139766
-0.273647270856027 -0.634981005938745
-0.316171685005353 -0.634813409861658
-0.358380932343072 -0.634013017971919
-0.400232937773488 -0.632580628117572
-0.441685982307405 -0.630517668135972
-0.482698744648525 -0.627826194430496
-0.523230342383414 -0.624508889920666
-0.563240372733984 -0.62056906136776
-0.602688952831856 -0.616010636078563
-0.641536759474477 -0.610838157990537
-0.679745068323334 -0.605056783142331
-0.717275792505227 -0.598672274534123
-0.754091520578086 -0.591690996382944
-0.790155553823513 -0.584119907778679
-0.82543194282886 -0.575966555747105
-0.859885523322388 -0.567239067726841
-0.893481951225777 -0.557946143467748
-0.92618773688905 -0.548097046358828
-0.957970278473782 -0.537701594194274
-0.988797894451326 -0.526770149386879
-1.01863985518365 -0.515313608638562
-1.0474664135553 -0.503343392078292
-1.07524883462598 -0.490871431878262
-1.10195942427417 -0.477910160359633
-1.12757155680321 -0.464472497599735
-1.1520597014824 -0.450571838553042
-1.17539944799656 -0.4362220396988
-1.19756753077877 -0.421437405228581
-1.21854185220202 -0.406232672787556
-1.23830150460657 -0.390622998783689
-1.25682679114118 -0.374623943279502
-1.27409924539734 -0.358251454481462
-1.29010164981703 -0.341521852842467
-1.3048180528555 -0.32445181479326
-1.31823378488215 -0.307058356119001
-1.33033547280354 -0.289358814997565
-1.34111105339397 -0.271370834716465
-1.35054978532035 -0.25311234608564
-1.35864225984944 -0.234601549563638
-1.36538041022664 -0.215856897114998
-1.3707575197171 -0.196897073816932
-1.37476822830115 -0.177740979233637
-1.37740853801724 -0.158407708576793
-1.37867581694724 -0.138916533671044
-1.37856880183994 -0.119286883743423
-1.37708759937034 -0.0995383260558766
-1.37423368603326 -0.079690546400193
-1.37000990667158 -0.0597633294747753
-1.36442047164044 -0.0397765391628251
-1.35747095261022 -0.0197500987315911
-1.34916827701263 0.000296029027577083
-1.33952072113528 0.0203418616985724
-1.32853790187166 0.0403674171594389
-1.31623076713488 0.0603527335009844
-1.30261158494449 0.0802778889252457
-1.28769393119748 0.100123021603971
-1.27149267613554 0.11986834947732
-1.25402396952204 0.139494189973056
-1.23530522454365 0.158980979626555
-1.21535510045244 0.178309293582099
-1.19419348396598 0.197459864955992
-1.1718414694438 0.216413604042208
-1.14832133786007 0.235151617341428
-1.12365653459348 0.253655226394485
-1.09787164605631 0.271905986401467
-1.07099237518622 0.289885704607887
-1.04304551582497 0.307576458439621
-1.01405892600976 0.324960613368515
-0.984061500203697 0.342020840490874
-0.953083140493206 0.358740133801281
-0.921154726780951 0.375101827144559
-0.888308086004087 0.391089610828955
-0.85457596040848 0.406687547883996
-0.819991974910516 0.421880089946806
-0.784590603579061 0.436652092761057
-0.748407135270954 0.450988831273085
-0.711477638454311 0.464876014310151
-0.673838925254689 0.478299798826191
-0.635528514759958 0.491246803700864
-0.596584595620461 0.503704123078142
-0.557045987981734 0.515659339231155
-0.516952104787745 0.527100534940441
-0.476342912493209 0.5380163053733
-0.43525889122416 0.548395769452371
-0.393740994426479 0.558228580702131
-0.351830608042612 0.567504937562486
-0.309569509257158 0.576215593159179
-0.266999824852466 0.584351864521274
-0.224163989215742 0.591905641236533
-0.181104702039527 0.598869393536049
-0.137864885757718 0.605236179800085
-0.0944876427595464 0.61099965347763
-0.0510162124241885 0.616154069412781
-0.00749392801880341 0.620694289571636
0.0360358264970062 0.624615788164003
0.0795296597172689 0.627914656154801
0.122944216043143 0.630587605160674
0.166236218900716 0.632631970727922
0.209362513880042 0.634045714988489
0.252280111752388 0.634827428691349
0.29494623132284 0.634976332607284
0.337318342075521 0.634492278305631
0.37935420656894 0.633375748302246
0.421011922539186 0.631627855578517
0.462249964669025 0.629250342471924
0.50302722598123 0.626245578939238
0.543303058814916 0.622616560194093
0.583037315344009 0.618366903721299
0.622190387597474 0.613500845670847
0.660723246941402 0.608023236635227
0.698597482983598 0.601939536814251
0.735775341861902 0.595255810572206
0.772219763878051 0.587978720392767
0.8078944204396 0.580115520237683
0.842763750273047 0.571674048315876
0.876792994872089 0.562662719270139
0.909948233145652 0.55309051578924
0.942196415231179 0.542966979653782
0.973505395439443 0.532302202224748
1.00384396429808 0.521106814384213
1.03318187966185 0.50939197593825
1.06148989685869 0.497169364492592
1.08873979784143 0.484451163812142
1.11490441931615 0.471250051675927
1.13995767981914 0.457579187239621
1.16387460571548 0.4434521979182
1.18663135609331 0.42888316580185
1.20820524652893 0.413886613618627
1.22857477169919 0.398477490257886
1.2477196268184 0.382671155868907
1.26562072787863 0.366483366549558
1.28226023067302 0.34993025864028
1.29762154858333 0.333028332639027
1.31168936911379 0.315794436753215
1.32444966915496 0.298245750105061
1.33588972896224 0.280399765607056
1.34599814483522 0.262274272524657
1.35476484048513 0.243887338743557
1.36218107707909 0.225257292759219
1.36823946195118 0.206402705406643
1.3729339559716 0.187342371348542
1.37625987956662 0.168095290340424
1.37821391738326 0.148680648291217
1.37879412159414 0.129117798138338
1.37799991383907 0.109426240556267
1.37583208580163 0.089625604517844
1.37229279841995 0.069735627727692
1.36738557973265 0.0497761369472344
1.36111532136205 0.0297670282309536
1.35348827363805 0.00972824709356981
1.3445120393677 -0.0103202313720839
1.33419556625652 -0.0303584224066646
1.32254913798928 -0.0503663515055666
1.30958436397895 -0.0703240743299378
1.29531416779422 -0.0902116965876329
1.27975277427699 -0.110009393864277
1.26291569536275 -0.129697431384671
1.24481971461791 -0.149256183684847
1.22548287050964 -0.168666154175154
1.20492443842467 -0.187907994574885
1.1831649114552 -0.206962524199068
1.16022597997093 -0.225810749078179
1.13613050999763 -0.244433880891755
1.1109025204238 -0.262813355696996
1.08456715905808 -0.280930852433704
1.05715067756143 -0.29876831118712
1.02868040527884 -0.316307951190434
0.999184721996884 -0.333532288549048
0.968693029654119 -0.35042415366889
0.937235723032594 -0.366966708371454
0.904844159459674 -0.383143462678448
0.871550627550369 -0.398938291249373
0.837388315021341 -0.41433544945561
0.80239127560867 -0.429319589075005
0.766594395122345 -0.443875773591318
0.730033356671334 -0.457989493083259
0.692744605093881 -0.471646678688298
0.654765310628502 -0.484833716626813
0.61613333186188 -0.497537461772598
0.576887177990604 -0.509745250756217
0.537065970434365 -0.52144491458812
0.496709403838874 -0.532624790788958
0.455857706507372 -0.543273735015
0.414551600300189 -0.553381132167049
0.372832260042302 -0.562936906971815
0.330741272479379 -0.571931534025154
0.288320594823202 -0.580356047287209
0.245612512927809 -0.588202049019946
0.202659599138035 -0.595461718158207
0.159504669852469 -0.602127818105913
0.116190742843141 -0.608193703949665
0.0727609943744727 -0.613653329082534
};
\addplot [line width=0.36pt, steelblue31119180, opacity=0.7, forget plot]
table {%
0.120572155997036 -0.801515271339218
0.0771213662069393 -0.807643061279324
0.0335937001617985 -0.812965775051606
-0.00996745281424766 -0.81747810685922
-0.0535186700165599 -0.821175558711686
-0.0970165386446942 -0.824054444908588
-0.140417699077225 -0.826111895713555
-0.183678888093561 -0.827345860214881
-0.226756981999661 -0.827755108369914
-0.269609039614671 -0.827339232231194
-0.31219234507563 -0.826098646353099
-0.354464450417566 -0.824034587378614
-0.396383217886557 -0.821149112806612
-0.437906861943554 -0.817445098940898
-0.478993990917119 -0.812926238023048
-0.519603648263543 -0.8075970345519
-0.559695353393219 -0.801462800793371
-0.599229142022572 -0.794529651485076
-0.63816560601133 -0.786804497741026
-0.676465932645412 -0.778295040162477
-0.714091943326287 -0.769009761161805
-0.751006131628232 -0.758957916507049
-0.787171700685548 -0.748149526095566
-0.822552599872485 -0.736595363965975
-0.857113560739281 -0.724306947558362
-0.890820132168524 -0.711296526233448
-0.923638714716775 -0.697577069062155
-0.955536594107222 -0.683162251897756
-0.986481973839986 -0.668066443743491
-1.01644400688756 -0.652304692429226
-1.04539282644379 -0.635892709611451
-1.07329957569579 -0.618846855111562
-1.10013643658899 -0.601184120608031
-1.12587665755681 -0.582922112698735
-1.1504945801872 -0.564079035350314
-1.17396566479951 -0.54467367175206
-1.19626651490615 -0.52472536559243
-1.21737490053474 -0.504254001776833
-1.23726978038743 -0.483279986605924
-1.25593132281532 -0.461824227434161
-1.27334092558706 -0.439908111828898
-1.28948123443201 -0.417553486250796
-1.30433616033929 -0.394782634276791
-1.31789089559568 -0.371618254387344
-1.33013192854632 -0.348083437340098
-1.34104705706334 -0.324201643152514
-1.35062540070924 -0.299996677716411
-1.35885741158279 -0.275492669067737
-1.36573488383652 -0.250714043335217
-1.37125096185651 -0.225685500391859
-1.37540014709623 -0.200431989233581
-1.3781783035576 -0.174978683109509
-1.37958266191384 -0.149350954428738
-1.37961182227 -0.12357434946856
-1.37826575555841 -0.0976745629093793
-1.37554580356763 -0.0716774122216982
-1.37145467760496 -0.0456088119306993
-1.36599645579374 -0.0194947477840832
-1.35917657900816 0.00663874915108917
-1.35100184544973 0.0327656284367304
-1.34148040387063 0.0588598462313704
-1.33062174545091 0.0848953912512236
-1.31843669433742 0.110846310698802
-1.30493739685412 0.136686736133229
-1.2901373093943 0.162390909256462
-1.27405118500703 0.187933207589728
-1.25669505869091 0.213288170014562
-1.23808623141013 0.238430522153009
-1.21824325284841 0.26333520156167
-1.19718590291831 0.287977382714491
-1.17493517204421 0.31233250174939
-1.15151324023853 0.336376280954047
-1.12694345499226 0.360084752966455
-1.10125030800169 0.383434284666107
-1.0744594107545 0.406401600732009
-1.04659746899969 0.42896380684402
-1.01769225612675 0.451098412504414
-0.987772585480479 0.472783353456891
-0.956868281639269 0.493997013680715
-0.925010150685282 0.514718246938026
-0.892229949496294 0.534926397852876
-0.85856035408974 0.554601322500953
-0.824034927050557 0.573723408489484
-0.788688084075265 0.592273594507295
-0.752555059665653 0.610233389325541
-0.71567187200626 0.627584890230164
-0.678075287060668 0.644310800867707
-0.639802781922388 0.660394448486698
-0.600892507456879 0.675819800557409
-0.561383250271936 0.69057148075343
-0.521314394054354 0.704634784279127
-0.480725880311415 0.717995692527702
-0.439658168556321 0.730640887055245
-0.398152195977278 0.742557762856847
-0.356249336630409 0.753734440931545
-0.3139913601972 0.764159780123568
-0.271420390347571 0.773823388228083
-0.228578862750088 0.78271563235037
-0.185509482771162 0.790827648508102
-0.142255182905413 0.798151350467156
-0.0988590799796279 0.804679437802147
-0.0553644321729666 0.810405403173648
-0.0118145958962689 0.81532353881485
0.031747017426558 0.819428942221189
0.0752769846319916 0.822717521037265
0.118731914102046 0.825185997136197
0.162068489017982 0.826831909887326
0.205243510539457 0.827653618609025
0.248213940866061 0.827650304204172
0.29093694613834 0.826821969976637
0.333369939135512 0.825169441627992
0.375470621727341 0.822694366434435
0.417197027037826 0.819399211604744
0.458507561278705 0.815287261820917
0.499361045211038 0.810362615963922
0.539716755193574 0.80463018302785
0.579534463776961 0.798095677226518
0.618774479803341 0.790765612297419
0.65739768797136 0.782647295008687
0.695365587827153 0.773748817875549
0.732640332142431 0.764079051093532
0.769184764641423 0.753647633696459
0.804962457039056 0.74246496394805
0.839937745353469 0.730542188976709
0.87407576545664 0.717891193663817
0.907342487827707 0.704524588796634
0.93970475147434 0.690455698497581
0.971130296988329 0.675698546942477
1.00158779870248 0.660267844380922
1.03104689591669 0.644178972472813
1.05947822316221 0.627447968955563
1.08685343947372 0.610091511657337
1.11314525664025 0.592126901872234
1.13832746640666 0.573572047113974
1.16237496659861 0.554445443265305
1.18526378614492 0.534766156140903
1.20697110897251 0.51455380248215
1.22747529674987 0.493828530402747
1.24675591045674 0.472610999304626
1.26479373075804 0.450922359284219
1.28157077716224 0.428784230049575
1.2970703259446 0.406218679369366
1.31127692681782 0.383248201075261
1.3241764183332 0.359895692639574
1.33575594199705 0.336184432350577
1.34600395508837 0.312138056108192
1.35491024216483 0.287780533863207
1.36246592524578 0.263136145723518
1.36866347266201 0.238229457751177
1.37349670656347 0.213085297474419
1.3769608090775 0.187728729139032
1.37905232711138 0.162185028723781
1.37976917579444 0.136479658744755
1.37911064055633 0.110638242873779
1.37707737783929 0.0846865403961859
1.37367141444383 0.0586504205333911
1.36889614550832 0.0325558366558998
1.36275633112468 0.0064288004124175
1.35525809159344 -0.0197046441991279
1.34640890132281 -0.0458184467928051
1.3362175813781 -0.0718865765622763
1.32469429068869 -0.097883048228866
1.31185051592132 -0.123781947944254
1.29769906003002 -0.149557459121965
1.28225402949377 -0.175183888171899
1.2655308202549 -0.200635690112265
1.24754610237208 -0.225887494033373
1.22831780340322 -0.250914128387904
1.20786509053489 -0.275690646082466
1.1862083514761 -0.300192349345388
1.1633691741353 -0.324394814346004
1.13937032510117 -0.348273915540854
1.11423572694836 -0.371805849722553
1.08799043439091 -0.394967159747344
1.06066060930722 -0.41773475791769
1.0322734946613 -0.440085948996594
1.00285738734637 -0.461998452830701
0.972441609977944 -0.483450426559645
0.941056481664368 -0.504420486389483
0.9087332877841 -0.524887728908526
0.875504248799753 -0.544831751924317
0.84140248814004 -0.564232674800978
0.806461999181621 -0.583071158276662
0.770717611363761 -0.601328423741347
0.734204955469596 -0.618986271955764
0.696960428108584 -0.636027101192798
0.659021155435581 -0.652433924783267
0.620424956142676 -0.668190388048605
0.581210303760702 -0.683280784603564
0.54141628830798 -0.697690072012677
0.501082577324538 -0.711403886784887
0.460249376330651 -0.724408558691385
0.418957388749093 -0.736691124392394
0.37724777533109 -0.748239340359306
0.335162113126384 -0.759041695079299
0.292742354038324 -0.769087420530262
0.250030783005301 -0.778366502914594
0.207069975850194 -0.786869692641177
0.163902756839874 -0.794588513545566
0.120572155997037 -0.801515271339218
};
\addplot [line width=0.36pt, steelblue31119180, opacity=0.7, forget plot]
table {%
-0.0281157227512392 -0.564603349556718
-0.0798387339317738 -0.561745580950493
-0.13148216012602 -0.55832785213253
-0.182994522043625 -0.554353569969241
-0.234324471041955 -0.549826696111517
-0.285420840311543 -0.54475174304567
-0.336232695880276 -0.539133769595283
-0.386709387385483 -0.53297837587847
-0.436800598563323 -0.526291697725553
-0.486456397405128 -0.519080400562734
-0.535627285930714 -0.511351672767854
-0.584264249529049 -0.50311321850486
-0.632318805817077 -0.49437325004413
-0.679743052968017 -0.485140479576297
-0.726489717460943 -0.475424110527757
-0.772512201204058 -0.465233828386484
-0.817764627984686 -0.454579791047332
-0.862201889199678 -0.44347261868642
-0.905779688820648 -0.431923383174714
-0.948454587549221 -0.419943597041339
-0.990184046118263 -0.407545201997646
-1.03092646769596 -0.394740557033451
-1.07064123935042 -0.381542426097324
-1.10928877253356 -0.367963965373211
-1.14683054254382 -0.354018710166062
-1.1832291269284 -0.339720561409543
-1.2184482427868 -0.325083771809287
-1.2524527829384 -0.310122931635482
-1.28520885091806 -0.294852954178977
-1.31668379476485 -0.279289060885387
-1.34684623957023 -0.263446766182025
-1.37566611875327 -0.247341862012784
-1.40311470403163 -0.230990402096384
-1.42916463405855 -0.214408685923675
-1.45378994169728 -0.197613242509952
-1.4769660799056 -0.180620813918471
-1.49866994620492 -0.163448338571597
-1.51887990570931 -0.146112934366218
-1.53757581269162 -0.128631881610249
-1.55473903066518 -0.111022605797249
-1.57035245096104 -0.0933026602363096
-1.58440050978231 -0.0754897085545328
-1.59686920371843 -0.0576015070895448
-1.60774610370413 -0.0396558871895925
-1.61702036740896 -0.0216707374388671
-1.62468275004516 -0.00366398582577493
-1.63072561358307 0.0143464181280714
-1.63514293436491 0.0323425212603268
-1.63793030910926 0.0503063846640159
-1.63908495930034 0.0682201015694262
-1.63860573395775 0.0860658151939651
-1.63649311078372 0.103825736542189
-1.63274919568697 0.121482162138262
-1.6273777206835 0.139017491673163
-1.62038404017641 0.156414245549061
-1.61177512561853 0.173655082303353
-1.6015595585631 0.190722815895014
-1.58974752210954 0.207600432836017
-1.57635079075266 0.224271109150748
-1.56138271864562 0.240718227146507
-1.54485822628824 0.25692539197839
-1.52679378565385 0.272876447992025
-1.50720740376974 0.288555494827879
-1.48611860476733 0.303946903271084
-1.46354841042009 0.319035330830977
-1.43951931918861 0.333805737034828
-1.4140552837936 0.348243398420512
-1.38718168733931 0.362333923213174
-1.35892531801114 0.376063265671264
-1.32931434237254 0.38941774008764
-1.298378277288 0.402384034431775
-1.26614796049998 0.414949223619481
-1.23265551988919 0.427100782396911
-1.1979343414488 0.438826597826008
-1.1620190360046 0.450114981358944
-1.12494540471413 0.460954680489518
-1.08675040337931 0.471334889969897
-1.04747210560814 0.481245262581527
-1.00714966486203 0.49067591944946
-0.965823275426855 0.499617459889838
-0.923534132346342 0.50806097078069
-0.880324390358006 0.515998035446734
-0.836237121872373 0.5234207420493
-0.791316274037462 0.530321691473018
-0.74560662493132 0.536694004701424
-0.699153738926253 0.542531329674109
-0.652003921269272 0.547827847618592
-0.604204171924014 0.552578278850599
-0.555802138720161 0.556777888036967
-0.506846069857038 0.560422488915928
-0.457384765808771 0.563508448470065
-0.407467530678901 0.56603269054778
-0.357144123052988 0.567992698929675
-0.306464706398169 0.569386519836769
-0.255479799059124 0.57021276387807
-0.204240223900291 0.570470607435548
-0.15279705764453 0.570159793485138
-0.101201579958731 0.569280631852945
-0.0495052223371319 0.567833998906403
0.00224048316671547 0.565821336680695
0.0539839553082163 0.563244651441298
0.105673615069042 0.560106511684096
0.157257937072213 0.556410045575045
0.20868550094372 0.552158937831952
0.259905042569467 0.547357426051468
0.310865505196456 0.542010296484957
0.36151609032727 0.536122879267457
0.411806308357119 0.52970104310449
0.461686028902978 0.522751189422003
0.511105530774638 0.515280245985288
0.560015551537878 0.507295659993239
0.608367336620321 0.49880539065482
0.656112687911059 0.489817901255149
0.703204011805576 0.480342150719118
0.74959436664809 0.470387584680936
0.795237509524011 0.459964126068518
0.840087942355889 0.449082165212095
0.884100957256887 0.437752549486905
0.92723268109657 0.425986572500297
0.969440119234602 0.413795962834013
1.01068119837873 0.401192872352885
1.05091480852438 0.388189864091592
1.09010084393396 0.374799899731555
1.12820024311517 0.361036326680447
1.16517502775833 0.346912864767212
1.20098834059401 0.332443592565836
1.2356044821331 0.317642933361523
1.26898894625291 0.302525640773248
1.30110845459356 0.287106784047026
1.33193098973062 0.271401733034568
1.36142582709074 0.255426142872271
1.3895635655786 0.23919593837584
1.41631615688455 0.222727298166085
1.44165693344376 0.206036638541723
1.46556063501907 0.189140597115249
1.48800343388088 0.172056016228207
1.50896295855922 0.154799926162373
1.52841831614406 0.137389528163607
1.54635011311187 0.119842177295267
1.56274047465743 0.102175365138312
1.57757306251193 0.0844067023553147
1.5908330912292 0.066553901135761
1.60250734292424 0.0486347575401645
1.61258418044909 0.0306671337605597
1.62105355899295 0.0126689403150842
1.62790703609509 -0.00534188180561657
1.6331377800605 -0.0233473790223606
1.63674057676982 -0.041329603063944
1.63871183487696 -0.0592706288583719
1.63904958938896 -0.0771525724009562
1.63775350362478 -0.0949576085814844
1.63482486955088 -0.112667988952674
1.63026660649342 -0.13026605942221
1.62408325822813 -0.147734277850734
1.61628098845104 -0.165055231538222
1.60686757463435 -0.182211654581349
1.59585240027366 -0.199186445084514
1.58324644553434 -0.21596268220738
1.56906227630622 -0.232523643031942
1.55331403167765 -0.248852819232285
1.53601740984137 -0.264933933530454
1.51718965244619 -0.280750955921988
1.49684952741014 -0.296288119654984
1.4750173102122 -0.311529936946738
1.45171476368128 -0.326461214422298
1.42696511630253 -0.341067068259557
1.40079303906271 -0.355332939025767
1.3732246208576 -0.369244606190699
1.344287342486 -0.382788202301976
1.31401004925631 -0.395950226808454
1.2824229222329 -0.408717559517862
1.24955744815091 -0.421077473675294
1.21544638802967 -0.433017648649521
1.18012374451573 -0.444526182214456
1.14362472798834 -0.455591602413559
1.10598572146094 -0.466202878995328
1.06724424431379 -0.476349434408497
1.02743891489385 -0.486021154345965
0.98660941201916 -0.495208397826963
0.944796435426133 -0.503902006807385
0.902041665199168 -0.512093315308737
0.858387720223044 -0.519774158056562
0.813878115699483 -0.526936878619768
0.768557219770255 -0.533574337042725
0.722470209290045 -0.539679916962529
0.675663024793183 -0.545247532204334
0.628182324699123 -0.550271632848184
0.580075438802312 -0.554747210761294
0.531390321092823 -0.558669804590268
0.482175501954778 -0.56203550420827
0.432480039790198 -0.564840954612724
0.382353472116531 -0.567083359269655
0.331845766186566 -0.568760482901334
0.28100726918 -0.569870653714452
0.229888658016269 -0.570412765066599
0.178540888838686 -0.570386276569388
0.127015146220266 -0.569791214627131
0.0753627921418106 -0.568628172410512
0.0236353147931844 -0.566898309265306
-0.0281157227512388 -0.564603349556718
};
\addplot [line width=0.36pt, steelblue31119180, opacity=0.7, forget plot]
table {%
-0.0330981970129618 -0.541800007498646
-0.0820616827139916 -0.53853781440782
-0.130943367545043 -0.534738795111659
-0.179694525179615 -0.530406736555218
-0.228266559403554 -0.525545957028677
-0.276611052556817 -0.520161301862778
-0.324679813797251 -0.514258138598896
-0.372424927138268 -0.507842351638551
-0.419798799212545 -0.500920336377714
-0.466754206714129 -0.493498992831734
-0.513244343471652 -0.485585718757255
-0.559222867105743 -0.47718840227797
-0.604643945224119 -0.468315414021566
-0.649462301108309 -0.458975598775704
-0.693633258846473 -0.449178266671341
-0.737112787867324 -0.438933183902187
-0.779857546830757 -0.428250562989553
-0.821824926831437 -0.417141052602287
-0.862973093872289 -0.405615726941948
-0.903261030565531 -0.393686074703804
-0.942648577019697 -0.381363987624647
-0.981096470871897 -0.368661748628854
-1.01856638642539 -0.3555920195845
-1.05502097285346 -0.342167828681733
-1.09042389143156 -0.328402557445989
-1.12473985176053 -0.314309927398998
-1.1579346469448 -0.299903986380872
-1.18997518769065 -0.285199094546915
-1.22082953529023 -0.270209910053102
-1.25046693345887 -0.254951374444517
-1.27885783899351 -0.239438697761292
-1.30597395122208 -0.223687343376909
-1.33178824021421 -0.207713012583973
-1.35627497372523 -0.191531628942824
-1.37940974284664 -0.17515932240858
-1.40116948633746 -0.158612413252456
-1.42153251361215 -0.141907395793354
-1.44047852636227 -0.125060921955972
-1.45798863879029 -0.108089784671804
-1.47404539643534 -0.091010901139575
-1.48863279357219 -0.073841295961817
-1.50173628916614 -0.0565980841743694
-1.51334282136781 -0.0392984541857412
-1.52344082053349 -0.0219596506433317
-1.53202022075799 -0.00459895724359007
-1.53907246990861 0.0127663204967486
-1.54459053815 0.0301188724911785
-1.54856892495175 0.0474414013384946
-1.55100366457137 0.0647166395652172
-1.55189233000745 0.0819273668381862
-1.55123403541897 0.0990564271301632
-1.54902943700828 0.116086745821332
-1.54528073236703 0.133001346719652
-1.53999165828553 0.149783368983094
-1.53316748702785 0.166416083926898
-1.5248150210763 0.182882911699084
-1.51494258635059 0.199167437807615
-1.50356002390837 0.215253429482714
-1.49067868013542 0.231124851858049
-1.47631139543534 0.246765883954628
-1.46047249142999 0.262160934451503
-1.44317775668333 0.277294657227532
-1.42444443096308 0.292151966658728
-1.40429118805572 0.30671805265594
-1.38273811715208 0.320978395427866
-1.35980670282194 0.334918779954703
-1.33551980359784 0.348525310157985
-1.30990162918908 0.361784422752493
-1.28297771634907 0.37468290076644
-1.25477490341964 0.38720788671643
-1.22532130357808 0.399346895424078
-1.19464627681324 0.411087826461509
-1.16278040065893 0.422418976213329
-1.1297554397135 0.433329049543038
-1.09560431397623 0.443807171052273
-1.06036106603197 0.453842895921639
-1.02406082711673 0.463426220322333
-0.986739782098093 0.472547591388177
-0.948435133405389 0.481197916738125
-0.909185063945469 0.489368573539743
-0.869028699041168 0.497051417104639
-0.828006067430328 0.504238789007263
-0.786158061364286 0.510923524718995
-0.743526395845583 0.5170989607499
-0.700153567045548 0.522758941291041
-0.656082809943185 0.527897824350724
-0.61135805522761 0.532510487378555
-0.566023885506982 0.536592332371714
-0.52012549086759 0.540139290458341
-0.473708623827391 0.543147825953478
-0.426819553728912 0.545614939883515
-0.379505020616953 0.547538172975631
-0.331812188647099 0.548915608109255
-0.283788599071465 0.549745872227091
-0.235482122848541 0.550028137703808
-0.186940912924379 0.549762123171041
-0.138213356232687 0.548948093797862
-0.0893480254616868 0.547586861026452
-0.0403936306357964 0.545679781763243
0.00860102943957895 0.543228757026318
0.0575871158217561 0.540236230050438
0.106515798114498 0.536705183851575
0.15533830314321 0.532639138253371
0.204005963573106 0.528042146378505
0.252470266421862 0.522918790608451
0.300682901418412 0.517274178015663
0.348595809159681 0.51111393527273
0.39616122901724 0.504444203043589
0.443331746746137 0.497271629862377
0.490060341748453 0.489603365506027
0.536300433944453 0.481447053867217
0.58200593020462 0.472810825334769
0.627131270296297 0.463703288689102
0.67163147229911 0.454133522520812
0.715462177443928 0.444111066180935
0.758579694330647 0.433645910271915
0.800941042480729 0.422748486688756
0.842503995181079 0.411429658220279
0.883227121576548 0.399700707720858
0.923069827969113 0.387573326863429
0.961992398282557 0.375059604484972
0.999956033652317 0.362172014536099
1.03692289110104 0.348923403646746
1.07285612126127 0.335326978320372
1.10771990510772 0.321396291769434
1.14147948966246 0.307145230405245
1.1741012226374 0.292587999995701
1.20555258597967 0.277739111504659
1.2358022282863 0.262613366627095
1.26481999605597 0.247225843034453
1.29257696374671 0.231591879344893
1.31904546260948 0.215727059833422
1.34419910826897 0.199647198897152
1.36801282702416 0.183368325291163
1.39046288084231 0.166906666150693
1.41152689102153 0.150278630815577
1.43118386049839 0.133500794473056
1.44941419477821 0.116589881635275
1.46619972146732 0.0995627494679239
1.48152370838767 0.082436370986647
1.49537088025583 0.0652278181379736
1.50772743390976 0.0479542447816278
1.51858105206802 0.030632869591192
1.52792091560802 0.0132809588901541
1.53573771435064 -0.00408419055953511
1.5420236563409 -0.0214452687992526
1.54677247561514 -0.0387849699286471
1.549979438447 -0.0560860093565343
1.55164134806621 -0.0733311410305564
1.55175654784512 -0.0905031746284185
1.55032492295006 -0.10758499269358
1.54734790045588 -0.124559567698303
1.54282844792335 -0.141409979017066
1.53677107044103 -0.158119429793411
1.52918180613452 -0.17467126368341
1.52006822014755 -0.19104898145908
1.50943939710082 -0.207236257455165
1.4973059320363 -0.223216955842916
1.48367991985583 -0.238975146714639
1.4685749432647 -0.254495121962967
1.45200605923206 -0.269761410939052
1.43398978398186 -0.284758795874035
1.41454407652907 -0.299472327048456
1.39368832077779 -0.313887337694454
1.3714433061989 -0.327989458615926
1.34783120710672 -0.341764632512049
1.32287556055514 -0.355199127989906
1.29660124287535 -0.368279553252234
1.26903444487864 -0.380992869446666
1.24020264574875 -0.393326403663132
1.21013458565008 -0.405267861566501
1.17886023707885 -0.416805339651836
1.14641077498585 -0.427927337110068
1.11281854570061 -0.438622767292255
1.07811703468782 -0.448880968760988
1.04234083316834 -0.458691715917953
1.00552560363789 -0.468045229197024
0.96770804431791 -0.476932184812745
0.928925852573979 -0.48534372405449
0.889217687338246 -0.493271462117011
0.848623130573376 -0.500707496458597
0.807182647816391 -0.507644414678501
0.764937547841733 -0.514075301905779
0.721929941483791 -0.519993747692186
0.678202699659898 -0.52539385240225
0.633799410635665 -0.530270233094157
0.588764336575264 -0.534618028885583
0.54314236941992 -0.538432905799128
0.496978986138667 -0.541711061082522
0.450320203395906 -0.544449226999284
0.403212531680992 -0.546644674086081
0.355702928945573 -0.54829521387351
0.307838753794876 -0.549399201067617
0.259667718279626 -0.549955535189956
0.211237840335634 -0.549963661674577
0.162597395918466 -0.549423572420821
0.113794870880936 -0.548335805801404
0.0648789126413395 -0.546701446125747
0.0158982816906581 -0.544522122559119
-0.0330981970129614 -0.541800007498646
};
\addplot [line width=0.36pt, steelblue31119180, opacity=0.7, forget plot]
table {%
-0.00531888209642684 -0.73632783662905
-0.0514352878701382 -0.735627700696611
-0.0975004218303039 -0.734194275073942
-0.143468365249584 -0.732028988630878
-0.189293296282283 -0.729133999772034
-0.234929535640563 -0.725512194285264
-0.28033159212855 -0.721167182465045
-0.32545420798895 -0.71610329551365
-0.370252404016962 -0.710325581223722
-0.414681524396528 -0.703839798946517
-0.458697281214227 -0.696652413850858
-0.502255798606422 -0.688770590478518
-0.545313656495682 -0.680202185602445
-0.587827933872855 -0.670955740394964
-0.629756251581665 -0.661040471913744
-0.67105681456318 -0.650466263914042
-0.711688453518027 -0.639243656996357
-0.751610665944854 -0.627383838099329
-0.7907836565141 -0.614898629348357
-0.829168376736844 -0.60180047627105
-0.866726563889192 -0.588102435391247
-0.903420779153386 -0.573818161214
-0.93921444493764 -0.558961892614456
-0.97407188133747 -0.543548438644246
-1.00795834170221 -0.527593163769492
-1.04084004727121 -0.51111197255518
-1.07268422084526 -0.494121293811139
-1.10345911945961 -0.47663806421545
-1.13313406602604 -0.458679711431591
-1.16167947991244 -0.440264136736171
-1.18906690642947 -0.42140969717454
-1.21526904519476 -0.402135187262094
-1.24025977734654 -0.382459820249489
-1.26401419157953 -0.362403208970456
-1.28650860897701 -0.341985346291302
-1.30772060661452 -0.321226585181587
-1.32762903991151 -0.300147618425839
-1.34621406370869 -0.278769457996537
-1.36345715205018 -0.257113414108924
-1.37934111665051 -0.235201073978513
-1.39385012402833 -0.213054280302489
-1.40696971128955 -0.19069510948643
-1.41868680054427 -0.168145849638063
-1.4289897119431 -0.14542897835
-1.43786817531989 -0.122567140293578
-1.44531334042927 -0.0995831246461642
-1.45131778576873 -0.076499842374404
-1.45587552597659 -0.0533403033960729
-1.45898201779827 -0.0301275936432904
-1.46063416461519 -0.00688485204995976
-1.46083031953144 0.0163647525136238
-1.45957028701556 0.0395980443360066
-1.45685532309536 0.062791863966631
-1.45268813410591 0.0859230913016777
-1.44707287399183 0.108968668630689
-1.44001514016651 0.131905623620997
-1.43152196793253 0.154711092217042
-1.42160182346867 0.177362341431762
-1.41026459539067 0.199836792007328
-1.39752158489405 0.22211204092264
-1.38338549448879 0.244165883725145
-1.36787041533723 0.265976336664717
-1.3509918132077 0.287521658607542
-1.33276651305786 0.308780372708146
-1.31321268226326 0.329731287817985
-1.29234981250769 0.350353519609246
-1.27019870035343 0.370626511392795
-1.24678142651083 0.39053005460953
-1.22212133382771 0.410044308974719
-1.19624300402076 0.42914982225522
-1.16917223317192 0.447827549659896
-1.14093600601431 0.466058872823874
-1.11156246903322 0.483825618367734
-1.08108090240913 0.501110076013134
-1.04952169083054 0.517895016236796
-1.0169162932059 0.534163707445272
-0.983297211304649 0.549899932653361
-0.948697957358764 0.565088005649548
-0.91315302065703 0.579712786632365
-0.876697833165372 0.593759697302076
-0.839368734207494 0.607214735392636
-0.801202934241059 0.620064488629462
-0.762238477765489 0.632296148099075
-0.722514205398385 0.643897521017306
-0.682069715158352 0.654857042883327
-0.640945322992834 0.665163789007399
-0.599182022590301 0.674807485400837
-0.556821444516848 0.683778519017349
-0.513905814717941 0.692067947335527
-0.470477912426677 0.699667507272948
-0.426581027520517 0.706569623422994
-0.382258917368988 0.71276741560618
-0.337555763215392 0.718254705728472
-0.292516126135978 0.723026023939741
-0.247184902620491 0.727076614086228
-0.201607279818382 0.730402438451584
-0.155828690495264 0.733000181781753
-0.109894767744553 0.734867254589689
-0.0638512994994006 0.736001795736612
-0.0177441828902908 0.736402674287233
0.0283806215062177 0.736069490637096
0.0744771354819969 0.735002576910908
0.120499409029425 0.733202996631475
0.166401566145353 0.730672543659556
0.212137850561305 0.727413740405704
0.257662671354316 0.723429835315874
0.302930648392963 0.718724799633301
0.347896657573261 0.713303323439877
0.392515875799355 0.707170810980974
0.436743825664153 0.700333375278381
0.480536419785377 0.692797832036704
0.523850004752819 0.684571692849326
0.566641404643016 0.675663157710688
0.608867964057944 0.666081106842357
0.65048759064485 0.655835091841025
0.691458797054824 0.644935326157272
0.731740742298288 0.633392674914576
0.771293272456185 0.621218644078722
0.81007696070628 0.608425368988402
0.848053146624671 0.595025602258452
0.885183974723341 0.581032701067761
0.921432432185334 0.566460613844547
0.956762385759932 0.551323866362265
0.991138617781068 0.535637547259994
1.02452686127306 0.519417293001754
1.05689383410869 0.502679272289744
1.08820727218551 0.48544016994702
1.11843596158746 0.46771717028571
1.14754976969955 0.449527939977308
1.17551967524469 0.430890610442159
1.20231779721275 0.411823759775663
1.22791742265296 0.392346394229222
1.25229303330195 0.372477929264401
1.27542033102091 0.352238170199167
1.29727626201651 0.331647292465526
1.31783903982142 0.310725821498206
1.33708816701153 0.289494612274459
1.35500445563823 0.267974828525369
1.37157004635539 0.246187921639381
1.38676842622191 0.224155609279093
1.40058444516221 0.201899853732618
1.41300433106816 0.179442840021094
1.42401570352736 0.156806953784177
1.4336075861643 0.134014758965544
1.44177041758174 0.111088975320668
1.44849606089182 0.0880524557692691
1.45377781182701 0.064928163615024
1.45761040542312 0.041739149655253
1.45999002126749 0.0185085292033765
1.46091428730732 -0.00474054095292639
1.4603822822141 -0.0279848856349124
1.45839453630208 -0.051201334374292
1.45495303099961 -0.0743667445100577
1.45006119687404 -0.0974580242575945
1.44372391021199 -0.120452155727067
1.43594748815864 -0.14332621786816
1.4267396824206 -0.166057409318274
1.41610967153889 -0.188623071131419
1.40406805173953 -0.211000709365153
1.39062682637101 -0.23316801750302
1.37579939393909 -0.255102898690184
1.35960053475089 -0.276783487760044
1.34204639618153 -0.298188173029914
1.3231544765781 -0.319295617844015
1.30294360781691 -0.340084781842318
1.28143393653148 -0.360534941934036
1.25864690402994 -0.380625712954846
1.23460522492185 -0.400337067987266
1.20933286447579 -0.419649358323925
1.18285501473027 -0.438543333053818
1.15519806938174 -0.457000158252036
1.12638959747481 -0.475001435753827
1.09645831592079 -0.492529221494292
1.06543406087206 -0.509566043395418
1.03334775798077 -0.526094918782621
1.00023139157143 -0.542099371313454
0.966117972758292 -0.557563447401571
0.931041506539108 -0.572471732119623
0.895036957898239 -0.586809364565183
0.858140216952786 -0.600562052674416
0.820388063176537 -0.61371608746873
0.781818128737401 -0.62625835672017
0.742468860984829 -0.638176358021983
0.702379484124673 -0.649458211251282
0.661589960119637 -0.660092670411414
0.620140948854325 -0.670069134842214
0.578073767604586 -0.679377659786972
0.53543034985155 -0.688008966305586
0.492253203481429 -0.695954450524006
0.448585368412727 -0.703206192210769
0.404470373693112 -0.709756962672051
0.359952194108719 -0.715600231957395
0.315075206349122 -0.720730175368898
0.26988414477169 -0.725141679267405
0.224424056809405 -0.728830346169888
0.178740258066595 -0.73179249913296
0.132878287147363 -0.734025185418127
0.0868838602617116 -0.735526179435145
0.040802825654631 -0.736293984960535
-0.00531888209642648 -0.73632783662905
};
\addplot [line width=0.36pt, steelblue31119180, opacity=0.7, forget plot]
table {%
0.0107093726494511 -0.725768114411018
-0.0340393370493741 -0.726066613088103
-0.0787541155967528 -0.725641352775019
-0.123390390328987 -0.724492757380492
-0.167903666836691 -0.722621971849392
-0.212249573317781 -0.72003086102142
-0.256383904808258 -0.716722007772202
-0.300262667246674 -0.712698710438615
-0.34384212132837 -0.707964979530939
-0.387078826105768 -0.702525533735091
-0.429929682291247 -0.696385795208934
-0.472351975219457 -0.689551884177356
-0.514303417426221 -0.682030612831496
-0.555742190801604 -0.6738294785382
-0.596626988275107 -0.664956656366488
-0.636917054991461 -0.655420990938465
-0.676572228935949 -0.645231987612805
-0.715552980968776 -0.634399803009612
-0.753820454228575 -0.622935234886073
-0.791336502865775 -0.61084971137302
-0.828063730067206 -0.598155279583124
-0.863965525334064 -0.584864593602065
-0.899006100976048 -0.570990901874673
-0.933150527785306 -0.55654803399858
-0.966364769854634 -0.541550386938584
-0.998615718505206 -0.526012910675438
-1.02987122529002 -0.509951093303384
-1.06010013404019 -0.493380945591287
-1.08927231192206 -0.476318985022747
-1.11735867947429 -0.458782219331117
-1.14433123959489 -0.440788129545821
-1.17016310544934 -0.422354652566886
-1.19482852727198 -0.403500163285051
-1.21830291803392 -0.384243456265272
-1.24056287795194 -0.364603727011895
-1.26158621781388 -0.344600552834157
-1.28135198109737 -0.324253873331101
-1.2998404648597 -0.303583970515348
-1.31703323937818 -0.282611448595546
-1.33291316652127 -0.261357213437649
-1.34746441683224 -0.239842451725491
-1.36067248530832 -0.218088609841444
-1.37252420585961 -0.196117372488191
-1.38300776443332 -0.173950641072949
-1.39211271079028 -0.151610511875667
-1.39982996892193 -0.129119254022976
-1.40615184609755 -0.106499287289833
-1.41107204053247 -0.0837731597510105
-1.41458564766989 -0.0609635253046786
-1.4166891650698 -0.0380931210905095
-1.4173804959003 -0.0151847448248013
-1.41665895102778 0.00773876792478072
-1.41452524970388 0.0306545665021954
-1.41098151884847 0.0535398079410557
-1.40603129092953 0.0763716797349953
-1.39967950044192 0.0991274225776703
-1.39193247898852 0.121784353049731
-1.3827979489688 0.144319886230148
-1.37228501588094 0.166711558209352
-1.36040415924527 0.188937048481748
-1.34716722215807 0.21097420219528
-1.3325873994861 0.232801052235871
-1.31667922471368 0.254395841124716
-1.29945855545537 0.275737042706616
-1.28094255764882 0.296803383607712
-1.26114968844331 0.317573864441251
-1.24009967780131 0.33802778074023
-1.21781350883122 0.358144743596055
-1.19431339687095 0.377904699982657
-1.16962276734316 0.397287952745776
-1.14376623240433 0.416275180237518
-1.11676956641082 0.434847455576591
-1.08865968022638 0.452986265515034
-1.05946459439691 0.470673528892623
-1.02921341121887 0.487891614660568
-0.99793628572956 0.504623359456522
-0.965664395647858 0.520852084713397
-0.932429910295641 0.536561613284919
-0.89826595853068 0.551736285571363
-0.863206595723079 0.566360975129381
-0.827286769808145 0.580421103750368
-0.790542286449521 0.59390265599234
-0.753009773347325 0.60679219315083
-0.714726643726861 0.619076866654882
-0.675731059044304 0.630744430874784
-0.636061890946527 0.641783255328779
-0.595758682523002 0.652182336276583
-0.554861608888385 0.661931307688152
-0.513411437135091 0.671020451576766
-0.471449485695766 0.679440707686128
-0.429017583156181 0.687183682521828
-0.386158026559581 0.694241657718158
-0.342913539244075 0.700607597731949
-0.299327228255087 0.70627515685575
-0.255442541375308 0.711238685543371
-0.211303223815012 0.71549323604147
-0.166953274605868 0.719034567321584
-0.122436902741755 0.721859149307679
-0.0777984831102734 0.723964166395002
-0.0330825122588928 0.725347520256741
0.0116664359601738 0.726007831935683
0.0564037548221577 0.725944443218785
0.101084849194686 0.725157417293299
0.145665179991178 0.723647538683784
0.19010030856838 0.721416312470073
0.234345941023766 0.718465962786981
0.278357972348664 0.714799430607235
0.322092530393093 0.71042037080985
0.36550601959847 0.705333148536866
0.408555164454613 0.699542834842081
0.451197052637713 0.693055201636111
0.493389177786266 0.685876715932826
0.535089481872345 0.678014533402891
0.576256397125954 0.669476491240835
0.616848887470684 0.660271100352768
0.656826489429368 0.650407536872524
0.696149352458955 0.639895633014688
0.734778278674396 0.628745867273634
0.772674761921952 0.616969353978336
0.809801026162962 0.604577832213358
0.846120063129821 0.59158365411708
0.881595669216629 0.577999772568823
0.916192481567724 0.563839728277131
0.949876013328157 0.549117636282102
0.98261268802093 0.533848171885207
1.01436987301677 0.518046556020633
1.04511591206302 0.50172854008272
1.0748201568393 0.484910390224636
1.10345299750844 0.467608871143913
1.13098589223214 0.449841229371038
1.15739139562218 0.431625176077736
1.18264318609852 0.412978869422085
1.20671609212724 0.393920896448073
1.22958611731209 0.374470254557627
1.25123046431459 0.354646332573589
1.27162755757893 0.334468891412517
1.29075706483897 0.313958044386571
1.30859991738584 0.29313423715413
1.32513832907608 0.272018227339112
1.34035581406125 0.250631063839325
1.35423720322138 0.228994065844466
1.36676865928585 0.20712880158469
1.37793769062673 0.185057066830929
1.38773316371069 0.162800863168393
1.39614531419713 0.140382376064915
1.40316575667149 0.117823952755989
1.40878749300403 0.0951480799685605
1.41300491932573 0.0723773615057727
1.41581383161431 0.0495344957149977
1.41721142988491 0.0266422528616355
1.41719632098117 0.0037234524312138
1.41576851996399 -0.0191990596175677
1.41292945009647 -0.0421024336261873
1.40868194142519 -0.0649638390133399
1.40303022795917 -0.0877604870329523
1.39597994344931 -0.110469653490506
1.38753811577253 -0.133068701395005
1.37771315992626 -0.155535103524021
1.36651486964014 -0.177846464879322
1.35395440761348 -0.199980545010687
1.34004429438797 -0.221915280185679
1.32479839586701 -0.243628805383245
1.30823190949378 -0.26509947608925
1.2903613491022 -0.286305889872198
1.27120452845551 -0.307226907717641
1.25078054348914 -0.327841675100007
1.22910975327547 -0.348129642770846
1.20621375972949 -0.368070587242769
1.1821153860755 -0.387644630948657
1.15683865409643 -0.406832262056054
1.13040876018843 -0.425614353916988
1.10285205024453 -0.443972184133824
1.07419599339251 -0.461887453222162
1.04446915461304 -0.479342302852162
1.0137011662655 -0.496319333650107
0.981922698549798 -0.512801622542491
0.949165428933673 -0.528772739625292
0.915462010575865 -0.544216764541669
0.880846039776766 -0.559118302351715
0.845352022488884 -0.573462498878465
0.809015339920563 -0.58723505551487
0.771872213267225 -0.600422243476954
0.733959667605293 -0.613010917488966
0.695315494984793 -0.624988528886877
0.655978216757412 -0.636343138127159
0.615987045177581 -0.647063426688387
0.575381844314845 -0.657138708353782
0.534203090316491 -0.666558939863465
0.492491831060046 -0.675314730925801
0.450289645235861 -0.683397353577839
0.407638600900561 -0.690798750885543
0.364581213542701 -0.697511544975113
0.321160403702396 -0.703529044387413
0.277419454187195 -0.708845250748162
0.233401966926839 -0.713454864747246
0.189151819509905 -0.717353291421181
0.144713121445676 -0.720536644733477
0.100130170194819 -0.723001751448323
0.0554474070127152 -0.724746154293734
0.0107093726494514 -0.725768114411018
};
\addplot [line width=0.36pt, steelblue31119180, opacity=0.7, forget plot]
table {%
-0.0584324126310926 -0.59556308700673
-0.105000090593441 -0.590694505282316
-0.15146310218198 -0.585237106428818
-0.197775132056298 -0.579196330500099
-0.243890015377836 -0.572578199072977
-0.289761783828011 -0.565389309244787
-0.335344711430428 -0.557636827057251
-0.38059336013153 -0.549328480353213
-0.425462625094242 -0.540472551073354
-0.469907779659446 -0.531077867000576
-0.513884519930494 -0.521153792960275
-0.557349008936294 -0.51071022148528
-0.600257920328953 -0.499757562954759
-0.642568481572428 -0.488306735216932
-0.684238516579116 -0.476369152705915
-0.725226487751896 -0.46395671506357
-0.765491537389711 -0.45108179527768
-0.804993528415412 -0.437757227348287
-0.84369308438527 -0.423996293494482
-0.881551628740268 -0.409812710914397
-0.918531423260059 -0.395220618111609
-0.954595605681243 -0.380234560801566
-0.989708226442469 -0.364869477412107
-1.02383428451975 -0.349140684192508
-1.05693976231623 -0.333063859945917
-1.08899165957167 -0.316655030400374
-1.11995802625784 -0.299930552234025
-1.14980799442697 -0.282907096770424
-1.17851180898164 -0.265601633360196
-1.20604085733525 -0.248031412465623
-1.23236769793371 -0.230213948465004
-1.25746608760976 -0.212167002193941
-1.28131100774275 -0.193908563240957
-1.30387868919772 -0.175456832015079
-1.32514663601899 -0.156830201603274
-1.34509364785465 -0.138047239435822
-1.3636998410895 -0.119126668777897
-1.38094666866549 -0.100087350065809
-1.39681693856979 -0.0809482621065102
-1.41129483097224 -0.0617284831591089
-1.42436591399483 -0.0424471719172452
-1.43601715809774 -0.0231235484112881
-1.44623694906745 -0.00377687484939154
-1.45501509959407 0.0155735635834929
-1.46234285942625 0.0349084779495089
-1.46821292409364 0.0542085947855274
-1.4726194421881 0.073454675315354
-1.47555802119657 0.0926275346273596
-1.47702573187959 0.111708060798417
-1.47702111119124 0.130677233945083
-1.4755441637375 0.149516145183029
-1.47259636177173 0.168206015475831
-1.46818064372703 0.186728214354321
-1.46230141128718 0.205064278487843
-1.45496452499894 0.223195930088901
-1.44617729843012 0.241105095132859
-1.43594849087926 0.258773921374519
-1.42428829864418 0.276184796143627
-1.41120834485808 0.293320363901568
-1.39672166790337 0.310163543541743
-1.38084270841467 0.326697545416387
-1.36358729488413 0.342905888072855
-1.34497262788324 0.358772414682694
-1.32501726291687 0.374281309147119
-1.30374109192681 0.389417111862846
-1.28116532346301 0.404164735132561
-1.25731246154244 0.418509478204661
-1.23220628321661 0.432437041927289
-1.20587181487008 0.445933543002032
-1.17833530727361 0.458985527823104
-1.14962420941681 0.471579985888184
-1.11976714114645 0.483704362767572
-1.08879386463748 0.49534657261872
-1.05673525472553 0.506495010233656
-1.02362326813016 0.517138562607314
-0.989490911599716 0.527266620015215
-0.954372209009483 0.536869086589474
-0.918302167445955 0.545936390382577
-0.881316742311 0.554459492908913
-0.843452801480739 0.562429898154525
-0.804748088554844 0.569839661046126
-0.765241185232896 0.576681395370922
-0.724971472855308 0.582948281139346
-0.68397909314715 0.588634071383373
-0.642304908204004 0.593733098383636
-0.599990459759736 0.598240279319129
-0.557077927776795 0.602151121333875
-0.51361008840031 0.605461726015502
-0.469630271317893 0.608168793281265
-0.42518231656767 0.610269624667643
-0.380310530837576 0.61176212602022
-0.335059643299482 0.612644809581189
-0.289474761022186 0.612916795472382
-0.243601324007706 0.612577812572347
-0.197485059895699 0.61162819878661
-0.151171938381151 0.610068900710845
-0.10470812539079 0.607901472687285
-0.058139937063884 0.605128075255321
-0.0115137935833054 0.601751472997826
0.0351238270971158 0.597775031785366
0.081726435582825 0.593202715421019
0.128247577380187 0.588039081689173
0.174640879203333 0.58228927781223
0.22086009520006 0.575959035319739
0.266859153050703 0.569054664335084
0.312592199894025 0.561583047285417
0.358013648034346 0.553551632041103
0.40307822038435 0.544968424491521
0.44774099559827 0.535841980564623
0.491957452850461 0.526181397698193
0.535683516214733 0.515996305771327
0.578875598600182 0.50529685750516
0.62149064519976 0.494093718342411
0.663486176408242 0.48239805581584
0.704820330166812 0.470221528416206
0.745451903692086 0.457576273970833
0.785340394547929 0.444474897544356
0.824446041019178 0.430930458873721
0.862729861746972 0.416956459349955
0.900153694586228 0.402566828559685
0.936680234646491 0.387775910399823
0.972273071478262 0.372598448779259
1.00689672536773 0.357049572921809
1.0405166827037 0.341144782285082
1.07309943038157 0.324899931110275
1.10461248920983 0.308331212618326
1.13502444628613 0.291455142868158
1.16430498631024 0.274288544293116
1.19242492180311 0.256848528931989
1.21935622220147 0.23915248137137
1.24507204179942 0.221218041416316
1.2695467465087 0.203063086506605
1.2927559394114 0.184705713896119
1.31467648507929 0.166164222613105
1.33528653263572 0.147457095219308
1.35456553753708 0.128602979386152
1.37249428205199 0.10962066930633
1.38905489441803 0.0905290869593517
1.40423086665663 0.071347263249697
1.41800707102863 0.0520943190363993
1.43036977511389 0.0327894460729577
1.44130665550014 0.0134518878765764
1.45080681006715 -0.00589907945418661
1.45886076885424 -0.0252441664542591
1.46546050350014 -0.0445640895202087
1.47059943524583 -0.0638395901325596
1.47427244149238 -0.0830514540531058
1.47647586090728 -0.102180530478077
1.47720749707415 -0.121207751128082
1.47646662068213 -0.140114149255777
1.47425397025294 -0.158880878552341
1.47057175140464 -0.177489231933878
1.46542363465306 -0.195920660189049
1.45881475175293 -0.214156790469329
1.45075169058245 -0.232179444603462
1.44124248857636 -0.249970657217856
1.43029662471402 -0.267512693644861
1.41792501007056 -0.284788067601067
1.40413997694053 -0.301779558618017
1.38895526654474 -0.318470229207938
1.37238601533273 -0.334843441747406
1.35444873989447 -0.350882875062082
1.3351613204962 -0.366572540696016
1.31454298325704 -0.381896798849293
1.29261428098394 -0.396840373968117
1.26939707268423 -0.411388369971824
1.24491450177607 -0.425526285101611
1.21919097301856 -0.439240026376211
1.19225212818461 -0.452515923640081
1.1641248205006 -0.465340743190106
1.13483708787853 -0.477701700967251
1.10441812496722 -0.489586475299982
1.07289825405045 -0.500983219186776
1.04030889482102 -0.511880572105477
1.00668253306095 -0.522267671337702
0.972052688258946 -0.532134162797041
0.936453880197423 -0.541470211350236
0.899921594542444 -0.550266510621051
0.862492247470844 -0.558514292267074
0.824203149369783 -0.56620533472019
0.785092467644947 -0.573331971382021
0.745199188674445 -0.579887098266153
0.704563078946338 -0.585864181079551
0.663224645418544 -0.591257261736078
0.621225095140619 -0.596060964295648
0.578606294177669 -0.600270500323083
0.535410725877356 -0.603881673661321
0.49168144852156 -0.606890884614245
0.447462052404961 -0.609295133534933
0.402796616383277 -0.611092023815775
0.3577296639345 -0.61227976427746
0.312306118776934 -0.612857170954466
0.266571260088244 -0.612823668275259
0.2205706773702 -0.61217928963604
0.17435022500406 -0.610924677367449
0.127955976541926 -0.609061082094281
0.0814341787796345 -0.606590361488833
0.034831205656932 -0.603514978419135
-0.0118064879690663 -0.599837998493916
-0.0584324126310923 -0.59556308700673
};
\addplot [line width=0.36pt, steelblue31119180, opacity=0.7, forget plot]
table {%
0.0529110204014957 -0.468509643111607
0.0107070430337663 -0.473039450281993
-0.0315076073473123 -0.47709772144084
-0.0736908502601705 -0.4806804112152
-0.115800636530886 -0.483783948301989
-0.157794990208716 -0.486405239027946
-0.199632050408636 -0.488541670433468
-0.241270113039179 -0.490191112877267
-0.282667672373987 -0.491351922159245
-0.323783462425613 -0.492022941159463
-0.364576498080364 -0.492203500991589
-0.405006115953145 -0.491893421669656
-0.445032014921612 -0.491093012287476
-0.484614296299203 -0.489803070710528
-0.523713503607021 -0.488024882780628
-0.562290661904903 -0.485760221034177
-0.600307316642496 -0.483011342935255
-0.637725571991586 -0.479780988625332
-0.67450812862148 -0.476072378191837
-0.710618320879796 -0.471889208458304
-0.746020153341579 -0.467235649299297
-0.780678336690323 -0.46211633948379
-0.814558322895134 -0.45653638205114
-0.847626339648951 -0.450501339224266
-0.879849424033524 -0.444017226865108
-0.91119545537756 -0.437090508477878
-0.941633187275314 -0.429728088766105
-0.97113227873368 -0.421937306749873
-0.999663324416757 -0.413725928450126
-1.02719788395773 -0.405102139147333
-1.05370851030882 -0.396074535222217
-1.07916877710116 -0.3866521155867
-1.10355330498709 -0.376844272713591
-1.12683778693892 -0.366660783273962
-1.14899901247865 -0.356111798391555
-1.17001489081471 -0.345207833523915
-1.18986447286253 -0.333959757980359
-1.20852797212704 -0.322378784087207
-1.22598678442628 -0.310476456011095
-1.24222350643648 -0.298264638251497
-1.25722195304002 -0.285755503813935
-1.27096717345914 -0.272961522075659
-1.2834454661592 -0.259895446355906
-1.29464439250667 -0.24657030120311
-1.30455278916821 -0.232999369411749
-1.31316077923855 -0.219196178781769
-1.32045978208598 -0.205174488633777
-1.32644252190573 -0.190948276093452
-1.33110303497263 -0.17653172215884
-1.3344366755859 -0.161939197564429
-1.33644012070009 -0.147185248456089
-1.33711137323752 -0.132284581891154
-1.33644976407908 -0.117252051178111
-1.33445595273116 -0.102102641070497
-1.33113192666827 -0.0868514528297705
-1.32648099935189 -0.0715136891720485
-1.3205078069275 -0.0561046391137049
-1.31321830360322 -0.0406396627309486
-1.3046197557145 -0.0251341758485647
-1.29472073448086 -0.00960363467308584
-1.28353110746196 0.00593647961429175
-1.27106202872134 0.0214706762896854
-1.25732592770789 0.0369834705280166
-1.24233649686581 0.0524593988386219
-1.22610867798575 0.067883034479596
-1.20865864731048 0.0832390028355041
-1.1900037994101 0.0985119967431329
-1.17016272984266 0.113686791750004
-1.14915521661783 0.128748261290437
-1.12700220048165 0.143681391764044
-1.10372576404238 0.158471297501607
-1.07934910975806 0.173103235603442
-1.05389653680783 0.187562620635439
-1.02739341686993 0.201835039168139
-0.999866168830699 0.215906264144358
-0.971342232449637 0.229762269061024
-0.941850041006856 0.2433892419511
-0.911418992960171 0.256773599151654
-0.880079422640081 0.269901998844351
-0.847862570011842 0.282761354354865
-0.814800549534794 0.295338847197962
-0.780926318149971 0.307621939855246
-0.746273642427902 0.319598388272827
-0.71087706490936 0.331256254066463
-0.67477186967261 0.342583916422
-0.637994047161466 0.353570083679251
-0.600580258309239 0.364203804587772
-0.562567797994318 0.374474479223302
-0.523994557863828 0.384371869554001
-0.484898988562407 0.393886109645943
-0.445320061403773 0.403007715497691
-0.405297229523265 0.41172759449416
-0.364870388550106 0.420037054470328
-0.324079836838565 0.42792781237578
-0.282966235297688 0.435392002531431
-0.241570566859614 0.442422184470205
-0.199934095626903 0.449011350353854
-0.158098325739573 0.455152931958529
-0.116104960002879 0.460840807222119
-0.0739958583170431 0.466069306346864
-0.031812995950401 0.47083321745113
0.0104015783024624 0.475127791764723
0.052605784035864 0.478948748362571
0.0947575511796863 0.48229227843204
0.136814861935739 0.485155049069643
0.178735792662012 0.487534206603351
0.220478555663075 0.489427379437191
0.262001540844962 0.490832680415312
0.303263357193016 0.491748708703137
0.344222874031355 0.492174551183746
0.38483926202282 0.492109783368094
0.425072033868549 0.491554469818148
0.4648810846666 0.490509164082532
0.504226731889383 0.488974908144735
0.54306975494007 0.486953231384444
0.581371434248542 0.484446149053018
0.61909358986789 0.481456160264651
0.656198619533026 0.47798624550519
0.692649536143431 0.474039863661131
0.728410004632708 0.46962094857172
0.763444378188158 0.46473390510762
0.797717733784309 0.459383604780037
0.831195906994935 0.453575380884694
0.863845526048908 0.447315023185486
0.895634045095903 0.440608772143114
0.92652977664881 0.433463312694464
0.956501923170505 0.425885767588911
0.985520607773512 0.417883690288209
1.01355690400193 0.409465057437039
1.04058286466596 0.400638260911712
1.06657154970025 0.391412099454966
1.09149705301837 0.38179576990519
1.11533452833655 0.37179885802881
1.13806021394098 0.361431328964994
1.159651456374 0.350703517292184
1.18008673301562 0.339626116726358
1.19934567353766 0.328210169461309
1.21740908020935 0.31646705516154
1.23425894703411 0.30440847961877
1.24987847769823 0.292046463083347
1.26425210231385 0.279393328282197
1.27736549293937 0.266461688135267
1.28920557786183 0.253264433182687
1.29976055462708 0.239814718735204
1.30901990180474 0.226125951760678
1.31697438947617 0.212211777519729
1.32361608843507 0.198086065963839
1.32893837809148 0.183762897909491
1.33293595307131 0.169256551002094
1.33560482850489 0.154581485483727
1.33694234399915 0.139752329778844
1.33694716628956 0.124783865912352
1.33561929056914 0.109691014774564
1.3329600404933 0.0944888212477252
1.32897206686035 0.0791924392089574
1.32365934496913 0.0638171164245361
1.31702717065635 0.0483781793505903
1.30908215501762 0.0328910178553565
1.29983221781731 0.0173710698782216
1.28928657959397 0.00183380604085251
1.27745575246914 -0.0137052857742595
1.26435152966858 -0.029230715862454
1.24998697376657 -0.0447270081373769
1.23437640366483 -0.0601787155578814
1.21753538031913 -0.0755704355259716
1.19948069122778 -0.0908868252404504
1.18023033369754 -0.106112616990954
1.15980349690342 -0.121232633377141
1.13822054276059 -0.136231802437853
1.11550298562711 -0.151095172675172
1.09167347085798 -0.165807927958404
1.06675575223177 -0.180355402293113
1.0407746682723 -0.19472309444051
1.01375611748909 -0.208896682372599
0.985727032561153 -0.222862037548686
0.956715353489909 -0.236605238999018
0.926749999747994 -0.250112587201508
0.895860841451692 -0.263370617737714
0.86407866958576 -0.276366114714466
0.831435165310304 -0.289086123937756
0.797962868380303 -0.301517965825761
0.763695144709282 -0.313649248048122
0.728666153109437 -0.325467877878895
0.692910811241394 -0.336962074250835
0.656464760807524 -0.348120379499028
0.619364332023512 -0.358931670782136
0.581646507403619 -0.369385171169886
0.5433488848957 -0.379470460385755
0.504509640402754 -0.389177485194119
0.465167489728355 -0.398496569421544
0.425361649983893 -0.407418423602206
0.385131800496111 -0.415934154237826
0.344518043253881 -0.424035272662911
0.303560862933676 -0.431713703506434
0.262301086543565 -0.438961792741541
0.220779842725952 -0.445772315315256
0.179038520759668 -0.452138482350563
0.137118729302223 -0.458053947913712
0.095062254913398 -0.463512815339983
0.052911020401496 -0.468509643111607
};
\addplot [line width=0.36pt, steelblue31119180, opacity=0.7, forget plot]
table {%
-0.0709702130985577 -0.911939103342883
-0.107954731003533 -0.908603568785947
-0.144831637284883 -0.904362318431649
-0.181564172240438 -0.899219580050582
-0.21811572008051 -0.893180480036206
-0.254449845427301 -0.88625103829475
-0.290530329634483 -0.878438162244443
-0.326321206890732 -0.869749639930043
-0.361786800071233 -0.860194132259547
-0.396891756301409 -0.849781164370813
-0.43160108219744 -0.838521116136688
-0.465880178748426 -0.826425211818132
-0.499694875805436 -0.813505508875627
-0.533011466143052 -0.799774885950041
-0.565796739059469 -0.785247030024912
-0.59801801348164 -0.769936422782976
-0.629643170542485 -0.7538583261705
-0.660640685597672 -0.737028767183848
-0.690979659650072 -0.719464521893422
-0.720629850150546 -0.701183098720908
-0.749561701144384 -0.682202720986502
-0.777746372733313 -0.66254230874351
-0.805155769823743 -0.642221459918422
-0.831762570132566 -0.621260430775278
-0.857540251422597 -0.599680115723782
-0.882463117940523 -0.577502026491298
-0.906506326030984 -0.554748270679493
-0.929645908901273 -0.531441529727003
-0.951858800511953 -0.507605036300075
-0.973122858569589 -0.483262551133745
-0.993416886598665 -0.458438339346619
-1.01272065507069 -0.433157146252875
-1.03101492156945 -0.407444172695595
-1.04828144997223 -0.381325049926017
-1.06450302862802 -0.354825814053749
-1.07966348751444 -0.327972880093403
-1.09374771435635 -0.30079301563354
-1.10674166969014 -0.273313314154152
-1.11863240085856 -0.245561168019292
-1.12940805492217 -0.217564241171771
-1.13905789047466 -0.189350441557138
-1.14757228835009 -0.160947893304434
-1.15494276121149 -0.132384908691446
-1.16116196201122 -0.103689959922412
-1.16622369131467 -0.0748916507463149
-1.17012290347999 -0.046018687944036
-1.17285571168769 -0.0170998527128161
-1.17441939181513 0.0118360280234697
-1.17481238515194 0.0407601103496156
-1.17403429995382 0.0696435621113284
-1.17208591183301 0.0984575916557448
-1.16896916298518 0.127173476531577
-1.16468716025334 0.155762592120278
-1.15924417203095 0.184196440169683
-1.15264562400705 0.212446677201687
-1.14489809375779 0.240485142765639
-1.13600930418983 0.268283887509296
-1.12598811584193 0.295815201039338
-1.11484451805255 0.323051639543698
-1.10258961900232 0.349966053148146
-1.08923563464114 0.376531612979876
-1.07479587651103 0.402721837911109
-1.05928473847692 0.428510620956058
-1.04271768237855 0.453872255294938
-1.02511122261772 0.478781459899083
-1.00648290969642 0.503213404731628
-0.986851312722095 0.527143735498627
-0.966236000897549 0.550548597925942
-0.944657524013944 0.573404661537704
-0.922137391966318 0.595689142912641
-0.898698053312057 0.617379828395082
-0.874362872893682 0.638455096238016
-0.849156108548265 0.658893938156121
-0.823102886926697 0.678675980267276
-0.796229178446884 0.697781503401691
-0.76856177140588 0.716191462758405
-0.740128245276724 0.733887506889559
-0.710956943216625 0.750851995993515
-0.681076943813882 0.767068019498596
-0.650518032101721 0.782519412919912
-0.619310669867921 0.79719077397247
-0.587485965289844 0.811067477924504
-0.555075641925125 0.824135692175729
-0.522112007088943 0.836382390045976
-0.488627919649384 0.847795363760473
-0.454656757273009 0.858363236618814
-0.420232383153267 0.868075474335508
-0.38538911225493 0.87692239554078
-0.350161677108188 0.884895181431174
-0.31458519318651 0.891985884560325
-0.278695123902775 0.898187436761147
-0.242527245258569 0.903493656191537
-0.206117610181891 0.907899253496569
-0.169502512588809 0.911399837081038
-0.132718451204897 0.913991917487099
-0.0958020931825096 0.915672910872629
-0.058790237550167 0.916441141586865
-0.0217197785304848 0.916295843840725
0.0153723312367907 0.91523716247016
0.0524491175299779 0.913266152791785
0.0894736214021655 0.910384779550912
0.126408936022698 0.906595914963051
0.163218243466706 0.901903335850823
0.199864851416012 0.896311719879138
0.236312229734836 0.889826640892401
0.272524046883824 0.882454563358382
0.308464206136119 0.874202835924293
0.344096881559366 0.865079684091495
0.379386553727775 0.855094202016142
0.414298045128661 0.844256343443928
0.448796555228154 0.832576911787975
0.482847695161123 0.820067549359759
0.516417522010752 0.806740725763797
0.549472572643571 0.792609725467676
0.581979897066241 0.777688634559809
0.613907091270819 0.76199232670811
0.64522232953578 0.745536448333605
0.675894396150589 0.728337403013733
0.705892716532198 0.710412335130908
0.735187387702459 0.691779112782623
0.763749208096058 0.672456309970146
0.791549706669274 0.652463188083544
0.818561171280535 0.631819676701518
0.844756676314479 0.610546353725161
0.870110109521999 0.588664424865461
0.894596198049494 0.566195702504985
0.918190533631407 0.543162583954825
0.940869596920907 0.519588029128468
0.962610780934492 0.495495537654856
0.983392413587116 0.470909125453444
1.0031937792954 0.445853300794614
1.02199513962737 0.420353039869287
1.03977775297817 0.394433761892128
1.05652389325207 0.368121303763106
1.07221686753227 0.341441894312718
1.08684103272071 0.314422128156521
1.10038181113149 0.287088939185033
1.11282570502222 0.259469573715456
1.12416031004883 0.231591563331945
1.13437432763054 0.203482697441531
1.14345757621248 0.175170995573037
1.15140100141494 0.146684679446598
1.15819668505892 0.118052144841639
1.16383785305918 0.089301933291345
1.16831888217676 0.0604627036318384
1.1716353056244 0.0315632034344304
1.17378381751908 0.00263224034941263
1.17476227617741 -0.0263013466100297
1.17456970625052 -0.0552087158151794
1.17320629969631 -0.0840610517717437
1.17067341558808 -0.112829593843772
1.16697357875979 -0.141485664922879
1.16211047728926 -0.170000700014208
1.15608895882179 -0.198346274710648
1.14891502573795 -0.226494133526889
1.14059582917024 -0.254416218065116
1.13113966187474 -0.28208469498422
1.12055594996465 -0.309471983744684
1.10885524351418 -0.336550784101462
1.09604920604197 -0.363294103317457
1.08215060288468 -0.389675283070473
1.0671732884722 -0.415668026026801
1.05113219251725 -0.441246422054982
1.03404330513315 -0.466384974053585
1.01592366089448 -0.491058623367272
0.996791321856669 -0.515242774765815
0.976665359551363 -0.53891332096115
0.955565835975478 -0.562046666638052
0.933513783592978 -0.584619751974453
0.91053118436926 -0.606610075627979
0.886640947859047 -0.627995717165774
0.861866888369648 -0.648755358915266
0.836233701222342 -0.668868307214084
0.809766938135536 -0.688314513037957
0.782492981754268 -0.707074591986007
0.754439019351401 -0.725129843603555
0.725633015726768 -0.742462270023125
0.696103685331243 -0.759054593905126
0.665880463643559 -0.774890275660267
0.634993477828378 -0.789953529936585
0.60347351670488 -0.804229341354627
0.5713520000558 -0.817703479475107
0.538660947307497 -0.830362512984115
0.505432945612308 -0.842193823081745
0.471701117364953 -0.853185616060788
0.437499087185416 -0.863326935062958
0.402860948401196 -0.872607671000923
0.367821229062328 -0.881018572635266
0.332414857523083 -0.88855125579632
0.296677127624616 -0.895198211741693
0.260643663513302 -0.900952814641142
0.224350384129811 -0.905809328181348
0.187833467404323 -0.909762911283996
0.151129314193591 -0.912809622931469
0.114274511995763 -0.914946426095337
0.0773057984791746 -0.916171190763737
0.0402600248614434 -0.91648269606461
0.00317411917536838 -0.915880631482699
-0.0339149505417303 -0.914365597169073
-0.0709702130985575 -0.911939103342883
};
\addplot [line width=0.36pt, steelblue31119180, opacity=0.7, forget plot]
table {%
-0.02135636907131 -0.76501492876425
-0.0668329553124619 -0.763363802854431
-0.112242921010298 -0.760951739304556
-0.157541000523452 -0.757781142512028
-0.202682039741137 -0.753855172996683
-0.247621041093749 -0.749177744250315
-0.292313208407422 -0.743753518835621
-0.336713991557829 -0.737587903738459
-0.380779130878713 -0.730687044978045
-0.424464701280884 -0.723057821480468
-0.467727156037708 -0.71470783822163
-0.510523370193425 -0.705645418646435
-0.552810683551049 -0.695879596371802
-0.594546943196975 -0.685420106181757
-0.635690545519921 -0.674277374323592
-0.676200477682317 -0.662462508114752
-0.716036358502785 -0.649987284870818
-0.755158478708988 -0.636864140165628
-0.79352784052069 -0.623106155435221
-0.831106196523593 -0.608727044937979
-0.867856087795193 -0.593741142083959
-0.903740881244646 -0.578163385147034
-0.938724806129431 -0.562009302374098
-0.972772989712406 -0.545294996506173
-1.00585149202371 -0.528037128726844
-1.03792733969285 -0.510252902054032
-1.0689685588173 -0.49196004419165
-1.09894420683477 -0.473176789858241
-1.12782440336741 -0.453921862610218
-1.15558036000724 -0.43421445617781
-1.18218440901305 -0.414074215332338
-1.20761003089023 -0.393521216303881
-1.23183188082592 -0.37257594676885
-1.25482581395336 -0.351259285427428
-1.27656890941994 -0.329592481191226
-1.29703949323523 -0.307597132001911
-1.3162171598761 -0.2852951633019
-1.33408279262737 -0.26270880617861
-1.35061858263778 -0.239860575204021
-1.36580804667223 -0.216773245991665
-1.37963604354264 -0.193469832493394
-1.392088789201 -0.16997356405858
-1.40315387047963 -0.146307862278586
-1.41282025746485 -0.12249631763962
-1.42107831449192 -0.0985626660072243
-1.42791980974998 -0.0745307649658453
-1.43333792348777 -0.0504245700370794
-1.4373272548117 -0.0262681108002843
-1.43988382706954 -0.00208546693937388
-1.44100509181447 0.0220992557603357
-1.44068993134539 0.0462619494412966
-1.43893865982113 0.0703785282049654
-1.43575302294722 0.094424952121152
-1.43113619623577 0.118377251191543
-1.42509278184007 0.142211549243519
-1.417628803967 0.165904087730442
-1.40875170287205 0.18943124941469
-1.39847032744262 0.212769581909837
-1.38679492637736 0.235895821058498
-1.37373713796995 0.258786914122546
-1.35930997850785 0.28142004276259
-1.34352782929738 0.303772645783783
-1.32640642232808 0.325822441625316
-1.30796282459078 0.347547450571164
-1.28821542106481 0.368926016659942
-1.26718389639146 0.389936829272038
-1.24488921525192 0.410558944372505
-1.22135360146919 0.430771805388529
-1.19660051585494 0.450555263700665
-1.17065463282324 0.469889598727426
-1.14354181579455 0.488755537583179
-1.11528909141452 0.507134274289784
-1.08592462261325 0.525007488522797
-1.05547768053186 0.542357363873569
-1.02397861534434 0.559166605609034
-0.99145882600392 0.575418457911466
-0.957950728943836 0.591096720581051
-0.923487725763961 0.606185765184595
-0.888104169935354 0.620670550634286
-0.851835332555976 0.634536638180977
-0.814717367191687 0.647770205807047
-0.776787273837591 0.660358062004487
-0.738082862035633 0.672287658924483
-0.698642713185229 0.683547104885386
-0.658506142084484 0.694125176226593
-0.617713157740346 0.704011328496545
-0.576304423486761 0.713195706963658
-0.534321216450576 0.721669156439737
-0.491805386405601 0.729423230406067
-0.448799314055844 0.736450199433085
-0.405345868789505 0.742743058885245
-0.361488365945836 0.748295535903391
-0.317270523637468 0.753102095657676
-0.27273641917125 0.757157946864808
-0.227930445111024 0.760459046564097
-0.18289726502616 0.763002104147572
-0.137681768969932 0.764784584640127
-0.092329028732139 0.765804711226444
-0.0468842529105664 0.766061467022156
-0.00139274184606958 0.765554596087502
0.0441001575337877 0.764284603682452
0.0895490969175403 0.762252755763056
0.134908771814003 0.759461077719507
0.180133966712772 0.755912352357192
0.225179600156032 0.751610117122728
0.270000769676728 0.74655866057776
0.314552796558327 0.74076301812403
0.358791270371532 0.734228966983978
0.402672093243563 0.726963020441879
0.446151523815876 0.718972421351265
0.489186220846502 0.710265134915085
0.531733286413535 0.700849840745823
0.57375030867672 0.690735924213467
0.615195404154495 0.67993346708997
0.656027259474352 0.668453237499521
0.696205172554913 0.656306679184643
0.735689093178642 0.643505900098825
0.77443966291478 0.630063660337046
0.81241825435268 0.615993359416237
0.849587009606452 0.601309022918352
0.885908878052525 0.586025288509363
0.921347653262517 0.570157391348113
0.955868009094588 0.553721148899583
0.989435534907308 0.536732945167695
1.02201676986093 0.519209714363382
1.05357923627189 0.501168924024199
1.08409147198727 0.482628557602301
1.11352306174691 0.463607096538147
1.1418446675021 0.444123501837798
1.16902805766028 0.424197195172173
1.19504613522697 0.403848039517107
1.21987296481662 0.383096319353504
1.24348379850556 0.361962720447324
1.2658551005013 0.340468309229563
1.28696457060349 0.318634511796774
1.30679116643323 0.296483092553075
1.32531512440864 0.274036132514907
1.34251797944554 0.251316007300206
1.35838258336392 0.228345364823894
1.37289312198157 0.205147102721951
1.38603513087802 0.181744345526547
1.39779550981301 0.158160421615008
1.40816253578504 0.134418839955586
1.41712587471717 0.110543266673202
1.42467659175823 0.0865575014585416
1.43080716018925 0.0624854538440053
1.43551146892629 0.0383511193701647
1.4387848286121 0.0141785556664899
1.44062397629056 -0.0100081415298244
1.44102707865927 -0.0341848623930074
1.43999373389701 -0.0583275070419173
1.43752497206432 -0.0824120095633214
1.43362325407668 -0.106414362001324
1.42829246925144 -0.130310638289023
1.42153793143084 -0.154077018098523
1.41336637368508 -0.17768981058557
1.40378594160062 -0.201125478005089
1.39280618516053 -0.224360659174126
1.38043804922481 -0.24737219275878
1.36669386262034 -0.270137140361927
1.35158732585122 -0.292632809388714
1.33513349744184 -0.314836775667025
1.31734877892619 -0.33672690580039
1.29825089849846 -0.358281379231033
1.2778588933412 -0.37947870999108
1.25619309064863 -0.400297768120238
1.23327508736405 -0.420717800728596
1.20912772865155 -0.440718452683563
1.18377508512343 -0.460279786900304
1.15724242884614 -0.479382304215457
1.12955620814846 -0.498006962824329
1.10074402125731 -0.516135197262175
1.07083458878717 -0.533748936910661
1.03985772511079 -0.550830624011049
1.00784430863956 -0.567363231166149
0.97482625104322 -0.583330278313604
0.940836465439663 -0.598715849153572
0.90590883358637 -0.61350460701444
0.870078172106367 -0.627681810140751
0.833380197782252 -0.641233326388107
0.795851491952924 -0.654145647310401
0.757529464048516 -0.666405901625323
0.718452314299857 -0.678001868044741
0.678658995659652 -0.688921987457143
0.638189174973334 -0.699155374450012
0.597083193438286 -0.708691828160641
0.555382026390862 -0.717521842444576
0.513127242461275 -0.725636615351546
0.470360962137089 -0.733028057899441
0.4271258157766 -0.739688802137583
0.38346490111396 -0.745612208491259
0.339421740298429 -0.750792372380195
0.295040236510541 -0.755224130104371
0.250364630198467 -0.758903063991308
0.205439454978169 -0.761825506799697
0.160309493241323 -0.763988545374993
0.115019731515259 -0.765390023553301
0.0696153156194103 -0.766028544310695
0.0241415056629827 -0.765903471155798
-0.0213563690713096 -0.76501492876425
};
\addplot [line width=0.36pt, steelblue31119180, opacity=0.7, forget plot]
table {%
-0.0374651638482276 -0.729004766890348
-0.0811665450321947 -0.726394552912616
-0.124787017638892 -0.723060253047164
-0.168283099832582 -0.719005190996612
-0.211611433772596 -0.714233408935
-0.254728828833375 -0.708749663478461
-0.297592304657818 -0.702559420943725
-0.340159134001044 -0.695668851899175
-0.382386885321835 -0.688084825013887
-0.42423346507932 -0.679814900210796
-0.465657159692732 -0.670867321130796
-0.506616677122413 -0.661251006915303
-0.547071188030614 -0.650975543315461
-0.586980366481073 -0.640051173136852
-0.626304430136778 -0.628488786029249
-0.665004179915871 -0.616299907631577
-0.703041039066145 -0.603496688082902
-0.740377091619185 -0.590091889910909
-0.776975120185836 -0.576098875309939
-0.8127986430553 -0.561531592821263
-0.847811950560905 -0.546404563428868
-0.881980140676271 -0.530732866084631
-0.915269153806413 -0.514532122677287
-0.947645806739086 -0.497818482460192
-0.979077825722529 -0.480608605953398
-1.00953387863665 -0.46291964833608
-1.03898360622555 -0.444769242345886
-1.06739765236031 -0.426175480702234
-1.09474769330181 -0.407156898071102
-1.12100646593443 -0.387732452589266
-1.14614779494255 -0.367921506966418
-1.17014661890264 -0.347743809183997
-1.19297901526507 -0.32721947280997
-1.21462222420059 -0.30636895694919
-1.23505467128785 -0.285213045849318
-1.25425598901923 -0.26377282818263
-1.27220703710363 -0.242069676024369
-1.28888992154592 -0.220125223548599
-1.30428801248409 -0.19796134546279
-1.31838596076622 -0.175600135202629
-1.33116971325085 -0.153063882908807
-1.3426265268155 -0.130375053207721
-1.35274498105924 -0.107556262818247
-1.36151498968684 -0.084630258006899
-1.36892781056296 -0.0616198919138605
-1.37497605442655 -0.0385481017724741
-1.37965369225658 -0.0154378860449103
-1.38295606128194 0.00768771850320092
-1.38487986962937 0.0308056597663182
-1.38542319960486 0.0538928932778296
-1.38458550960529 0.0769264051812779
-1.38236763465826 0.0998832351710727
-1.37877178558975 0.122740499379823
-1.37380154682031 0.145475413189474
-1.36746187279205 0.168065313943509
-1.3597590830299 0.190487683537579
-1.35070085584221 0.212720170866039
-1.34029622066686 0.234740614102015
-1.32855554907044 0.256527062788797
-1.31549054440975 0.278057799720529
-1.30111423016556 0.299311362590396
-1.28544093696055 0.320266565384713
-1.26848628827427 0.340902519501609
-1.25026718486923 0.361198654573236
-1.23080178794394 0.381134738970762
-1.21010950102938 0.400690899971695
-1.18821095064715 0.419847643569451
-1.16512796574851 0.438585873905396
-1.14088355595482 0.456886912304019
-1.11550188862102 0.474732515892239
-1.08900826474509 0.4921048957843
-1.06142909374749 0.508986734814117
-1.03279186714565 0.525361204797405
-1.00312513114984 0.541211983306378
-0.972458458207686 0.556523269940301
-0.940822417525701 0.571279802075673
-0.908248544597246 0.585466870080341
-0.874769309767264 0.599070331976383
-0.840418085865127 0.612076627537138
-0.805229114937878 0.624472791804334
-0.769237474117017 0.636246468011842
-0.732479040652847 0.647385919903169
-0.694990456151255 0.657880043430412
-0.656809090048559 0.667718377823017
-0.617973002360848 0.676891116015301
-0.578520905744921 0.685389114422349
-0.538492126908677 0.693203902054535
-0.49792656740939 0.700327688961596
-0.456864663878975 0.706753373997818
-0.415347347715871 0.712474551900611
-0.373416004283735 0.717485519675421
-0.331112431657608 0.721781282280593
-0.288478798958686 0.725357557606539
-0.245557604319225 0.728210780744239
-0.202391632519477 0.730338107538817
-0.159023912338892 0.731737417424662
-0.11549767366409 0.732407315539247
-0.0718563043963761 0.732347134113563
-0.0281433072017343 0.731556933137763
0.0155977438535776 0.730037500301365
0.0593232467386014 0.727790350208066
0.102989614921135 0.724817722865947
0.146553320815732 0.721122581454592
0.18997093917295 0.716708609371319
0.233199190366579 0.711580206559495
0.276194983535715 0.705742485122574
0.31891545953867 0.69920126422824
0.361318033675895 0.691963064307738
0.403360438139332 0.684035100556161
0.445000764145889 0.675425275740192
0.486197503713024 0.666142172320451
0.526909591034803 0.656195043896313
0.567096443417192 0.645593805981717
0.606718001731767 0.634349026121169
0.645734770347523 0.622471913355781
0.684107856500981 0.609974307049857
0.721799009065339 0.596868665089157
0.758770656680027 0.583168051462604
0.794985945202655 0.568886123239816
0.830408774446028 0.554037116957445
0.865003834163591 0.538635834427885
0.898736639247458 0.522697627984502
0.93157356410391 0.506238385178097
0.963481876172126 0.489274512939851
0.994429768552703 0.471822921226534
1.02438639171348 0.453901006164297
1.05332188424099 0.435526632707835
1.08120740260701 0.416718116832207
1.10801514992039 0.397494207275081
1.1337184036356 0.377874066847576
1.15829154219036 0.357877253332365
1.18171007054574 0.337523699988041
1.20395064460337 0.316833695679221
1.22499109447537 0.295827864652159
1.24481044658371 0.27452714597605
1.26338894456719 0.252952772670512
1.28070806897499 0.231126250540046
1.29675055572724 0.209069336736589
1.31150041332425 0.186804018071505
1.32494293878721 0.164352489098658
1.33706473231439 0.141737129990398
1.34785371063838 0.118980484228513
1.35729911907095 0.0961052361324032
1.36539154222356 0.0731341882468528
1.37212291339282 0.050090238611961
1.37748652260156 0.0269963579378764
1.38147702328748 0.00387556670709127
1.38409043763269 -0.0192490877728713
1.38532416052896 -0.0423545543435201
1.38517696217446 -0.0654178009732865
1.38364898929972 -0.0884158377163598
1.38074176502133 -0.11132573962956
1.37645818732368 -0.134124669624424
1.37080252617016 -0.15678990123171
1.36378041924682 -0.179298841255636
1.35539886634256 -0.201629052295269
1.34566622237158 -0.223758275110616
1.33459218904506 -0.245664450811115
1.32218780520025 -0.267325742844421
1.30846543579673 -0.288720558763552
1.29343875959077 -0.309827571750714
1.27712275549998 -0.330625741876332
1.2595336876721 -0.351094337072113
1.24068908927244 -0.371212953797224
1.22060774500652 -0.390961537376978
1.19930967239502 -0.410320401993781
1.17681610181991 -0.42927025031037
1.15314945536152 -0.44779219270583
1.12833332444772 -0.465867766105171
1.10239244633752 -0.483478952383729
1.07535267946237 -0.500608196328018
1.04724097764992 -0.517238423135155
1.01808536325593 -0.533353055433388
0.98791489923087 -0.548936029806789
0.956759660149431 -0.563971812807609
0.924650702231477 -0.578445416440359
0.891620032384558 -0.592342413102171
0.857700576298714 -0.605648949964546
0.822926145625441 -0.618351762782154
0.787331404273498 -0.630438189114921
0.750951833855165 -0.641896180950227
0.713823698317406 -0.652714316712628
0.675984007793174 -0.662881812649125
0.637470481708908 -0.67238853357865
0.59832151118499 -0.681225002995032
0.558576120766633 -0.689382412513384
0.518273929523377 -0.696852630650493
0.47745511155593 -0.703628210930461
0.436160355949757 -0.709702399307512
0.394430826215309 -0.715069140898572
0.352308119255345 -0.719723086018909
0.309834223900234 -0.723659595514815
0.267051479052571 -0.726874745388014
0.224002531482845 -0.729365330707191
0.180730293318209 -0.731128868802736
0.137277899266733 -0.732163601741521
0.0936886636197971 -0.732468498079251
0.0500060370754551 -0.732043253888625
0.00627356342584241 -0.730888293062301
-0.0374651638482272 -0.729004766890348
};
\addplot [line width=0.64pt, darkorange25512714]
table {%
-0.0162129299680999 -0.815376068647767
-0.0607193514478565 -0.81408455041233
-0.105165246557586 -0.811981534992614
-0.14950631066245 -0.809069118720125
-0.193698343625435 -0.805350204754742
-0.237697293867009 -0.800828500190783
-0.281459302276699 -0.795508512361694
-0.3249407459328 -0.78939554434704
-0.368098281586657 -0.782495689686281
-0.410888888868155 -0.774815826304599
-0.45326991316936 -0.76636360965683
-0.495199108163561 -0.757147465096347
-0.536634677917327 -0.747176579476475
-0.577535318553601 -0.736460891992845
-0.617860259424303 -0.72501108427578
-0.6575693037514 -0.712838569742621
-0.69662286869593 -0.699955482220576
-0.734982024815031 -0.686374663851466
-0.772608534867663 -0.67210965229039
-0.809464891930317 -0.657174667211105
-0.845514356784734 -0.641584596131537
-0.880720994540355 -0.625354979573587
-0.915049710455004 -0.608501995571994
-0.948466284918089 -0.591042443547726
-0.980937407561452 -0.572993727561943
-1.01243071046387 -0.554373838967259
-1.04291480041611 -0.535201338473562
-1.07235929021435 -0.515495337646302
-1.1007348289508 -0.495275479855663
-1.12801313127135 -0.474561920695625
-1.15416700557102 -0.453375307892429
-1.17917038109911 -0.431736760722478
-1.20299833394714 -0.409667848960177
-1.22562711189349 -0.38719057137672
-1.2470341580802 -0.364327333811227
-1.26719813349807 -0.341100926836128
-1.28609893825788 -0.317534503039013
-1.3037177316264 -0.29365155394364
-1.32003695080724 -0.269475886593064
-1.33504032844781 -0.245031599818263
-1.34871290885499 -0.220343060215902
-1.36104106290328 -0.195434877859173
-1.37201250162062 -0.170331881765945
-1.3816162884383 -0.145059095148653
-1.38984285009275 -0.119641710470623
-1.39668398616841 -0.0941050643336687
-1.40213287727206 -0.0684746122220228
-1.40618409183054 -0.0427759031277518
-1.40883359150508 -0.0170345540829677
-1.41007873521679 0.0087237753757878
-1.40991828177934 0.0344734087855273
-1.40835239113619 0.0601886783516735
-1.4053826242012 0.0858439505342764
-1.40101194130263 0.111413651600082
-1.39524469923225 0.136872293114988
-1.38808664690238 0.162194497351465
-1.37954491961526 0.187355022585633
-1.36962803195039 0.212328788258756
-1.35834586927708 0.23709089997809
-1.34570967790045 0.261616674332154
-1.33173205385085 0.285881663495697
-1.31642693032794 0.309861679599811
-1.29980956381167 0.333532818842937
-1.28189651885439 0.356871485318684
-1.26270565156888 0.379854414536748
-1.242256091829 0.40245869661346
-1.22056822420063 0.424661799108856
-1.19766366762189 0.446441589487503
-1.17356525385288 0.467776357180693
-1.14829700471655 0.488644835228007
-1.12188410815315 0.509026221476688
-1.0943528931124 0.528900199317675
-1.06573080330818 0.548246957937645
-1.03604636986209 0.567047212066858
-1.00532918286294 0.585282221203137
-0.973609861870762 0.602933808292807
-0.940920025394543 0.619984377849976
-0.90729225937421 0.636416933496098
-0.872760084698256 0.652215094902338
-0.837357923789391 0.667363114117832
-0.801121066291525 0.681845891267592
-0.764085633892276 0.695648989604389
-0.726288544316095 0.708758649899616
-0.687767474523868 0.721161804158786
-0.648560823155708 0.732846088647995
-0.608707672254351 0.743799856218357
-0.568247748307328 0.75401218791614
-0.527221382646738 0.763472903867012
-0.485669471246098 0.772172573423567
-0.443633433954349 0.780102524565993
-0.401155173207648 0.787254852546533
-0.358277032260111 0.793622427769109
-0.315041752975133 0.79919890289626
-0.271492433219371 0.803978719176305
-0.227672483901852 0.807957111984424
-0.183625585701038 0.811130115572147
-0.139395645522972 0.813494567020492
-0.0950267527339201 0.815048109392834
-0.0505631352111346 0.81578919408435
-0.00604911525554054 0.815717082365703
0.0384709345896937 0.814831846119422
0.0829526357706764 0.813134367768249
0.127351647960298 0.810626339395521
0.171623713257633 0.807310261058464
0.215724700305181 0.803189438296081
0.259610648279958 0.79826797883412
0.303237810714607 0.792550788490402
0.346562699104826 0.786043566284598
0.389542126259652 0.778752798757319
0.432133249351395 0.770685753504193
0.474293612622302 0.761850471931369
0.515981189705374 0.752255761239664
0.557154425517167 0.741911185645359
0.597772277680805 0.730827056846375
0.637794257437906 0.719014423743355
0.677180470008666 0.706485061425871
0.715891654359835 0.693251459434766
0.753889222340976 0.679326809312305
0.791135297149966 0.664724991452558
0.827592751089422 0.649460561265123
0.863225242576392 0.633548734665976
0.89799725236844 0.617005372909917
0.931874118969994 0.599846966779722
0.964822073183685 0.582090620147778
0.996808271772211 0.563754032926567
1.02780083019719 0.544855483425007
1.05776885440238 0.525413810128242
1.08668247160951 0.505448392919024
1.11451286009615 0.484979133759428
1.14123227792581 0.464026436852143
1.1668140906017 0.442611188301118
1.19123279761662 0.42075473529184
1.21446405787235 0.398478864811996
1.23648471394346 0.375805781933727
1.25727281516107 0.352758087679133
1.27680763949379 0.329358756491081
1.29506971420389 0.305631113331778
1.31204083525815 0.281598810431947
1.32770408547402 0.257285803713764
1.34204385138304 0.232716328911077
1.3550458387947 0.207914877410698
1.36669708704518 0.182906171838855
1.37698598191685 0.157715141417144
1.38590226721552 0.132366897112533
1.39343705499412 0.106886706606212
1.39958283441231 0.0812999691062106
1.40433347922351 0.0556321900289204
1.40768425388167 0.0299089555747406
1.40963181826174 0.00415590722319353
1.41017423098918 -0.0216012838270541
1.40931095137518 -0.0473369422478064
1.40704283995566 -0.0730254141750588
1.40337215763338 -0.0986410927813294
1.39830256342433 -0.124158443801104
1.39183911081026 -0.149552030983942
1.3839882427013 -0.174796541449891
1.37475778501352 -0.199866810921908
1.36415693886788 -0.224737848810163
1.35219627141837 -0.249384863123205
1.33888770531841 -0.273783285181149
1.32424450683613 -0.297908794106282
1.30828127263023 -0.321737341066632
1.29101391519972 -0.345245173248375
1.27245964702201 -0.368408857533152
1.25263696339515 -0.391205303856712
1.23156562400126 -0.413611788225592
1.2092666332097 -0.435605975368882
1.18576221913946 -0.457165941002517
1.16107581150162 -0.478270193683875
1.13523201824421 -0.498897696234914
1.10825660102238 -0.519027886712494
1.08017644951868 -0.53864069890496
1.05101955463879 -0.557716582334587
1.02081498060963 -0.576236521745916
0.989592836007468 -0.594182056060577
0.957384243745053 -0.611535296779693
0.924221310047614 -0.628278945815521
0.890137092448674 -0.644396312734568
0.85516556683759 -0.659871331394974
0.819341593591639 -0.674688575961594
0.782700882826439 -0.688833276282802
0.745279958799332 -0.702291332613704
0.707116123501203 -0.71504932967107
0.668247419473052 -0.727094550005976
0.628712591884359 -0.738414986680835
0.588551049911052 -0.748999355238179
0.547802827451591 -0.758837104949243
0.506508543220299 -0.767918429331169
0.464709360257754 -0.776234275922314
0.422446944898567 -0.783776355305946
0.379763425237475 -0.790537149373314
0.336701349135144 -0.796509918817859
0.29330364180553 -0.801688709853106
0.249613563027099 -0.806068360147524
0.205674664020534 -0.809644503970452
0.161530744035932 -0.812413576543955
0.117225806692761 -0.81437281759627
0.0728040161160919 -0.815520274113309
0.0283096529128415 -0.815854802285461
-0.0162129299680996 -0.815376068647767
};
\addlegendentry{arithmétique}
\addplot [line width=0.64pt, mediumpurple148103189]
table {%
-0.000375422082854197 -0.68421180249173
-0.0432278405944037 -0.683847270973628
-0.0860371686543447 -0.682801065623748
-0.128760732993465 -0.681074229322333
-0.171355945833693 -0.678668483417364
-0.213780347340525 -0.675586226008678
-0.255991647947907 -0.671830529557493
-0.297947770513397 -0.667405137823707
-0.339606892261575 -0.662314462134044
-0.380927486473893 -0.656563576984741
-0.421868363883416 -0.650158214983183
-0.462388713733174 -0.643104761133515
-0.502448144457213 -0.635410246471931
-0.54200672394379 -0.627082341057981
-0.581025019340565 -0.618129346328897
-0.619464136362132 -0.608560186824529
-0.657285758060682 -0.59838440129118
-0.694452183021175 -0.587612133173167
-0.730926362942928 -0.576254120501623
-0.766671939570167 -0.564321685190589
-0.801653280934727 -0.551826721751084
-0.835835516874775 -0.538781685434403
-0.869184573794141 -0.525199579816443
-0.901667208627626 -0.511093943835465
-0.933251041978405 -0.496478838296178
-0.963904590394521 -0.481368831853628
-0.99359729775227 -0.465778986490844
-1.02229956571521 -0.449724842504727
-1.04998278323843 -0.433222403015141
-1.07661935508866 -0.41628811801266
-1.10218272935179 -0.398938867960854
-1.12664742390042 -0.381191946969481
-1.14998905179496 -0.363065045555339
-1.17218434559312 -0.344576233007975
-1.19321118054332 -0.325743939377824
-1.21304859663923 -0.306586937104736
-1.23167681951304 -0.287124322305195
-1.24907728014708 -0.267375495736901
-1.26523263338377 -0.247360143459667
-1.28012677521567 -0.22709821721193
-1.29374485883828 -0.206609914522422
-1.30607330944969 -0.185915658576839
-1.3170998377822 -0.165036077859556
-1.32681345235259 -0.143991985590716
-1.3352044704186 -0.122804358979152
-1.34226452763092 -0.10149431831185
-1.34798658637101 -0.0800831059007925
-1.35236494276625 -0.0585920649081531
-1.35539523237579 -0.0370426180709735
-1.35707443454102 -0.0154562463465108
-1.35740087539672 0.00614553250044727
-1.35637422953955 0.0277411853469579
-1.35399552035241 0.0493091851766092
-1.35026711898436 0.0708280325381172
-1.34519274198697 0.0922762769764478
-1.33877744760959 0.113632538415096
-1.33102763075715 0.134875528468211
-1.32195101661566 0.155984071661323
-1.31155665295147 0.176937126539515
-1.29985490109234 0.197713806642006
-1.286857425599 0.218293401322225
-1.27257718263771 0.238655396392636
-1.25702840706519 0.258779494573729
-1.24022659823909 0.278645635726785
-1.22218850456785 0.298234016850259
-1.20293210681553 0.317525111819847
-1.18247660017819 0.336499690852548
-1.16084237514969 0.355138839675328
-1.13805099719605 0.373423978379284
-1.11412518525847 0.391336879940497
-1.08908878910662 0.408859688389126
-1.06296676556466 0.425974936608631
-1.03578515363375 0.44266556374737
-1.0075710485358 0.458914932225228
-0.978352574704374 0.474706844318321
-0.948158857749606 0.490025558305239
-0.917019995425205 0.504855804158732
-0.884967027626331 0.519182798767212
-0.852031905448364 0.532992260670873
-0.818247459337352 0.546270424297767
-0.783647366363912 0.559004053685622
-0.748266116653189 0.571180455675736
-0.712138979004355 0.5827874925658
-0.675301965733898 0.593813594209014
-0.637791796777766 0.604247769547468
-0.599645863088131 0.614079617568268
-0.56090218936128 0.623299337671485
-0.521599396133768 0.631897739439613
-0.481776661284627 0.639866251798776
-0.441473680982006 0.647196931562557
-0.400730630113166 0.653882471349943
-0.359588122237279 0.659916206869478
-0.318087169100948 0.665292123562377
-0.276269139756809 0.670004862597968
-0.234175719325954 0.674049726215488
-0.191848867445303 0.677422682406917
-0.14933077644132 0.68012036893617
-0.106663829271787 0.682140096690646
-0.0638905572775477 0.683479852361795
-0.0210535977863439 0.684138300452027
0.021804348388999 0.684114784605967
0.0646405595156743 0.683409328264721
0.107412335526844 0.682022634642512
0.150077040585994 0.6799560860257
0.192592145587266 0.677211742394889
0.234915270549389 0.673792339371493
0.277004226860972 0.669701285490811
0.318817059335028 0.66494265880432
0.360312088030817 0.659521202814591
0.401447949801315 0.653442321746858
0.442183639524901 0.646712075161972
0.482478550980151 0.639337171916095
0.522292517323009 0.631324963473169
0.561585851125963 0.622683436576818
0.600319383939341 0.613421205288984
0.638454505335271 0.603547502403246
0.675953201395396 0.593072170241371
0.712778092603974 0.582005650842273
0.748892471108595 0.570358975553158
0.784260337311362 0.558143754033243
0.818846435754075 0.545372162680992
0.852616290261641 0.532056932496415
0.885536238308671 0.518211336390537
0.917573464575021 0.503849175954665
0.948696033656814 0.488984767702664
0.978872921900346 0.473632928799947
1.00807404832714 0.457808962293398
1.03627030461932 0.441528641856964
1.0634335841354 0.424808196068116
1.08953680992764 0.407664292230846
1.11455396173287 0.390114019761339
1.13846010191012 0.37217487315287
1.16123140029894 0.353864734536912
1.18284515797389 0.33520185585784
1.20327982987127 0.316204840678996
1.22251504626572 0.296892625638251
1.24053163307522 0.277284461571551
1.25731163097422 0.257399894323271
1.27283831329589 0.237258745262485
1.28709620270562 0.216881091524602
1.30007108662917 0.196287245998032
1.31175003142007 0.175497737075864
1.32212139525218 0.154533288192716
1.33117483972445 0.133414797167161
1.33890134016657 0.11216331537033
1.34529319463484 0.0908000267414505
1.35034403158974 0.0693462266712268
1.35404881624713 0.0478233007741402
1.3564038555971 0.0262527035707969
1.35740680208522 0.00465593710160009
1.3570566559526 -0.0169454705069559
1.35535376623254 -0.0385299865019863
1.35229983040253 -0.0600760949685342
1.34789789269223 -0.0815623182771041
1.34215234104888 -0.102967238493026
1.3350689027633 -0.124269518726322
1.32665463876084 -0.145447924400782
1.31691793656284 -0.166481344421048
1.3058685019258 -0.187348812216617
1.29351734916647 -0.208029526641771
1.27987679018252 -0.228502872710616
1.2649604221798 -0.248748442146541
1.24878311411833 -0.268746053725639
1.23136099189061 -0.288475773393789
1.21271142224696 -0.307917934137348
1.19285299548399 -0.327053155587667
1.17180550691337 -0.345862363339858
1.14958993712944 -0.364326807966577
1.12622843109529 -0.382428083707863
1.1017442760682 -0.400148146818395
1.07616187838638 -0.417469333553896
1.04950673914019 -0.434374377778733
1.02180542875211 -0.450846428177176
0.993085560490789 -0.466869065051162
0.963375762945488 -0.482426316687797
0.932705651488521 -0.497502675280319
0.901105798753993 -0.512083112386606
0.86860770416233 -0.52615309390986
0.835243762520979 -0.539698594586511
0.801047231732558 -0.5527061119669
0.766052199642647 -0.56516267987482
0.730293550060298 -0.57705588133248
0.693806927985076 -0.588373860938016
0.656628704075367 -0.599105336683219
0.618795938393304 -0.609239611199682
0.58034634346249 -0.618766582422169
0.541318246675343 -0.627676753658578
0.501750552087509 -0.635961243056443
0.461682701637456 -0.643611792456566
0.421154635829898 -0.650620775624923
0.380206753922219 -0.656981205854662
0.338879873653626 -0.662686742930606
0.297215190557133 -0.667731699449317
0.25525423689496 -0.672111046488419
0.213038840258271 -0.675820418619544
0.170611081872517 -0.678856118259881
0.128013254649962 -0.68121511935801
0.0852878210311838 -0.682895070410333
0.0424773706575917 -0.683894296805107
-0.000375422082853864 -0.68421180249173
};
\addlegendentry{harmonique}
\addplot [line width=0.64pt, forestgreen4416044]
table {%
-0.00486618432593241 -0.732246557935913
-0.0486398214252344 -0.731590682085171
-0.0923649732927378 -0.730205540732146
-0.135998053746142 -0.728092514615864
-0.179495568381933 -0.725253710046709
-0.222814157931554 -0.721691956806814
-0.265910641482864 -0.717410805329268
-0.308742059523811 -0.71241452315897
-0.35126571676541 -0.706708090698642
-0.393439224701333 -0.700297196244265
-0.435220543861694 -0.693188230314851
-0.476568025718908 -0.685388279282243
-0.51744045420385 -0.67690511830726
-0.557797086790927 -0.66774720358925
-0.597597695111113 -0.657923663936769
-0.636802605052465 -0.647444291667781
-0.675372736308139 -0.636319532848465
-0.71326964133249 -0.624560476880349
-0.75045554366642 -0.612178845446155
-0.786893375593778 -0.599186980825368
-0.822546815091259 -0.585597833591183
-0.857380322034996 -0.571424949701097
-0.891359173627718 -0.556682456993997
-0.924449499011199 -0.541385051107229
-0.95661831302946 -0.525547980827664
-0.987833549109089 -0.509187032891374
-1.01806409122391 -0.492318516247076
-1.04727980491209 -0.474959245799011
-1.07545156731486 -0.457126525645484
-1.10255129620678 -0.438838131829766
-1.12855197798874 -0.42011229462054
-1.15342769461568 -0.400967680339581
-1.17715364943229 -0.381423372754753
-1.19970619189081 -0.3614988540569
-1.22106284112647 -0.341213985439572
-1.24120230836684 -0.320588987300955
-1.26010451815299 -0.299644419087738
-1.27775062835113 -0.278401158801011
-1.29412304893486 -0.256880382184621
-1.30920545951927 -0.235103541616723
-1.32298282562948 -0.213092344725592
-1.33544141368734 -0.19086873275098
-1.34656880470131 -0.168454858672615
-1.35635390664605 -0.145873065127627
-1.36478696551914 -0.123145862138923
-1.37185957506408 -0.100295904676706
-1.37756468514987 -0.0773459700755069
-1.3818966087987 -0.054318935329243
-1.38485102785488 -0.0312377542869337
-1.38642499728924 -0.0081254347718033
-1.3866169481348 0.0149949843534169
-1.38542668905081 0.0381004561521277
-1.38285540651338 0.0611679485875508
-1.37890566363289 0.0841744674815623
-1.37358139759893 0.107097079435791
-1.36688791575567 0.129912934692129
-1.35883189031138 0.152599289909863
-1.34942135168742 0.17513353083674
-1.33866568051336 0.197493194851335
-1.32657559827612 0.219655993354289
-1.31316315663262 0.241599833986065
-1.29844172539637 0.263302842649094
-1.28242597921018 0.284743385312353
-1.26513188291817 0.305900089576631
-1.24657667565161 0.326751865979011
-1.22677885364465 0.347277929015297
-1.20575815179681 0.367457817859461
-1.18353552400086 0.387271416759438
-1.16013312225546 0.406698975088945
-1.13557427458362 0.425721127035334
-1.1098834617788 0.444318910903861
-1.08308629300193 0.462473788019104
-1.0552094802536 0.480167661204732
-1.02628081174703 0.497382892823151
-0.996329124208113 0.514102322357093
-0.965384274130364 0.530309283515591
-0.933477108013321 0.545987620847299
-0.900639431614074 0.561121705844597
-0.8669039782426 0.57569645252243
-0.832304376132488 0.589697332456342
-0.796875114919581 0.603110389264729
-0.760651511261959 0.615922252520861
-0.723669673635526 0.628120151080818
-0.685966466340296 0.639691925814038
-0.647579472753259 0.65062604172381
-0.608546957864447 0.660911599445602
-0.568907830133569 0.670538346111791
-0.528701602705204 0.67949668557194
-0.487968354021243 0.687777687958454
-0.446748687869828 0.695373098588061
-0.405083692910609 0.702275346190263
-0.363014901716677 0.708477550454542
-0.320584249373993 0.713973528888804
-0.277834031679577 0.718757802982222
-0.234806862980145 0.722825603666332
-0.191545633693197 0.726172876068948
-0.148093467552915 0.728796283556143
-0.104493678623492 0.730693211058281
-0.0607897281227241 0.731861767676773
-0.0170251810989292 0.732300788568971
0.0267563369956422 0.732009836109299
0.0705111837916089 0.730989200325499
0.114195743506119 0.729239898609514
0.157766470420013 0.726763674703341
0.201179932285143 0.723562996960818
0.244392853618579 0.719641055887129
0.287362158840558 0.715001760958428
0.330045015213151 0.709649736724794
0.372398875536876 0.703590318200376
0.414381520562664 0.696829545545339
0.455951101076931 0.689374158044895
0.497066179617778 0.681231587391443
0.537685771780759 0.672409950276491
0.577769387073026 0.662918040299765
0.61727706927513 0.652765319203558
0.656169436270255 0.641961907441055
0.69440771930116 0.63051857408805
0.731953801615727 0.618446726108096
0.768770256462567 0.605758396981793
0.804820384398818 0.592466234711555
0.840068249872945 0.578583489213806
0.874478717046087 0.564123999111177
0.908017484816217 0.549102177937861
0.940651121010235 0.533532999771894
0.972347095709875 0.517431984308663
1.00307381367825 0.500815181390543
1.03280064585464 0.483699155008057
1.06149795988626 0.46610096678854
1.08913714966637 0.448038158988725
1.11569066384952 0.429528737008251
1.1411320333153 0.410591151441487
1.16543589755332 0.391244279685572
1.18857802994316 0.371507407123022
1.21053536190398 0.351400207897637
1.23128600588975 0.330942725302891
1.25080927720726 0.310155351802335
1.26908571463509 0.289058808701952
1.28609709982294 0.267674125494705
1.30182647545209 0.24602261889788
1.31625816213884 0.224125871604123
1.32937777406405 0.202005710767341
1.34117223331322 0.179684186244914
1.35162978291288 0.157183548617931
1.36073999855016 0.134526227011311
1.36849379896398 0.111734806735976
1.37488345499747 0.0888320067753184
1.37990259730252 0.0658406571384186
1.38354622268895 0.0427836761026041
1.38581069911171 0.0196840473680054
1.38669376929147 -0.00343520285309093
1.38619455296467 -0.0265510287892763
1.38431354776099 -0.049640388082542
1.38105262870734 -0.0726802647574045
1.37641504635874 -0.0956476921637323
1.37040542355813 -0.118519775870398
1.36302975082822 -0.14127371648695
1.35429538039998 -0.163886832390535
1.34421101888382 -0.18633658233544
1.3327867185906 -0.20860058792269
1.32003386751131 -0.230656655907325
1.30596517796524 -0.252482800321111
1.2905946739281 -0.274057264388629
1.27393767705259 -0.295358542214901
1.25601079139543 -0.31636540022294
1.23683188686613 -0.337056898319831
1.21642008141381 -0.357412410770281
1.19479572197005 -0.377411646756794
1.17198036416663 -0.397034670606002
1.14799675084839 -0.416261921660979
1.12286878940267 -0.435074233779724
1.09662152792791 -0.453452854440394
1.06928113026517 -0.471379463434223
1.04087484991742 -0.488836191127504
1.0114310028827 -0.50580563627443
0.980978939428049 -0.522270883363031
0.949549014832589 -0.538215519476932
0.917172559128682 -0.553623650656093
0.883881845871485 -0.568479917740266
0.849710059967968 -0.582769511679331
0.814691264597477 -0.596478188295286
0.778860367256811 -0.609592282481147
0.74225308496368 -0.62209872182263
0.704905908653195 -0.633985039629014
0.666856066802915 -0.64523938736021
0.628141488322693 -0.655850546437642
0.588800764746307 -0.665807939427176
0.548873111762589 -0.675101640582929
0.508398330124363 -0.683722385741473
0.467416765974198 -0.691661581556553
0.425969270626495 -0.698911314065122
0.384097159845998 -0.705464356576151
0.341842172663356 -0.71131417687435
0.299246429768738 -0.716454943731628
0.256352391525018 -0.720881532719781
0.213202815642366 -0.724589531318638
0.169840714556418 -0.727575243314555
0.126309312552566 -0.729835692484883
0.0826520026790282 -0.731368625564729
0.0389123034917334 -0.732172514493061
-0.00486618432593207 -0.732246557935913
};
\addlegendentry{de Fréchet}
\addplot [line width=0.64pt, sienna1408675, dashed]
table {%
-0.00672853428885439 -0.747017791698577
-0.050400225894203 -0.746252080296484
-0.0940216774582937 -0.74474248859152
-0.137549406169543 -0.742490521377527
-0.180940022641346 -0.7394984234644
-0.22415027416356 -0.735769177440408
-0.267137087817744 -0.731306500699092
-0.309857613413188 -0.726114841733679
-0.352269266200906 -0.720199375702732
-0.394329769323045 -0.71356599927144
-0.435997195955376 -0.706221324733688
-0.477230011100863 -0.698172673420786
-0.517987112992653 -0.689428068403401
-0.558227874065207 -0.67999622649399
-0.597912181452748 -0.669886549557688
-0.63700047697463 -0.659109115140327
-0.675453796567799 -0.647674666422918
-0.713233809127015 -0.635594601512615
-0.750302854714132 -0.622880962080834
-0.786623982098344 -0.609546421359852
-0.822160985589975 -0.595604271509855
-0.856878441131099 -0.581068410369022
-0.890741741607009 -0.565953327599859
-0.923717131343344 -0.550274090245589
-0.955771739754476 -0.534046327710992
-0.986873614109629 -0.517286216182679
-1.01699175138405 -0.500010462504314
-1.04609612916351 -0.482236287522871
-1.07415773557129 -0.463981408922518
-1.10114859818784 -0.445264023563244
-1.12704181193432 -0.426102789341836
-1.15181156589215 -0.406516806593276
-1.1754331690319 -0.386525599051118
-1.19788307482585 -0.366149094385801
-1.21913890471966 -0.345407604340317
-1.23917947043985 -0.324321804483023
-1.25798479511465 -0.302912713597781
-1.27553613318747 -0.281201672731985
-1.2918159891028 -0.259210323923334
-1.30680813474624 -0.23696058862658
-1.32049762562095 -0.214474645861746
-1.33287081574472 -0.191774910105597
-1.34391537125255 -0.168884008948401
-1.3536202826913 -0.145824760538254
-1.36197587599421 -0.122620150835457
-1.36897382212414 -0.099293310699606
-1.37460714537622 -0.0758674928322536
-1.3788702303313 -0.052366048598106
-1.38175882745361 -0.0288124047478774
-1.38327005732669 -0.00523004006599388
-1.38340241352376 0.0183575380335687
-1.38215576410927 0.0419268169399858
-1.37953135177047 0.0654543022834675
-1.37553179257866 0.0889165413549555
-1.37016107338142 0.112290146484295
-1.36342454782844 0.135551818353572
-1.35532893103491 0.15867836922238
-1.3458822928877 0.181646746041865
-1.33509405000116 0.204434053434505
-1.32297495633041 0.227017576516721
-1.3095370924516 0.249374803541566
-1.29479385351967 0.271483448338924
-1.27875993591583 0.29332147253085
-1.26145132259784 0.314867107499904
-1.24288526716788 0.336098876088582
-1.22308027667377 0.356995614008213
-1.20205609316074 0.377536490935984
-1.17983367399216 0.397701031279057
-1.15643517095873 0.417469134585088
-1.13188390819713 0.436821095578791
-1.10620435894003 0.455737623804587
-1.07942212112059 0.474199862855752
-1.05156389185592 0.492189409170888
-1.02265744083479 0.509688330378999
-0.992731582636251 0.526679183174863
-0.961816148006555 0.543145030706902
-0.929941954123282 0.559069459460203
-0.897140773876085 0.574436595617866
-0.863445304194781 0.58923112088437
-0.828889133456349 0.603438287755183
-0.793506708003294 0.617043934217395
-0.757333297806794 0.630034497866722
-0.720404961308818 0.642397029426805
-0.682758509478293 0.654119205657327
-0.644431469117125 0.665189341638093
-0.60546204545267 0.675596402416807
-0.565889084053929 0.685330014008957
-0.525752032109444 0.694380473738823
-0.48509089910548 0.702738759911322
-0.443946216943699 0.710396540805022
-0.402358999538081 0.717346182977391
-0.36037070193136 0.723580758873973
-0.318023178971737 0.729094053733926
-0.275358643591057 0.733880571785033
-0.232419624726041 0.737935541722006
-0.189248924924517 0.741254921462624
-0.145889577678909 0.743835402176973
-0.102384804529514 0.745674411585759
-0.0587779719803357 0.746770116524408
-0.0151125482704036 0.747121424770409
0.0285679399563063 0.746727986132061
0.0722199510391821 0.745590192797556
0.115799971704248 0.743709178944035
0.159264560439141 0.741086819607013
0.202570390796559 0.737725728811308
0.245674294582999 0.733629256965321
0.288533304889761 0.728801487521277
0.331104698923286 0.723247232904751
0.373346040592174 0.716972029717535
0.4152152228084 0.709982133218639
0.456670509460575 0.702284511088905
0.497670577017408 0.693886836485479
0.538174555719903 0.684797480393034
0.578142070321223 0.675025503279396
0.617533280333614 0.664580646063869
0.65630891974227 0.653473320407276
0.694430336146548 0.641714598333391
0.731859529289525 0.629316201192101
0.76855918893748 0.616290487975313
0.804492732071544 0.602650442997239
0.839624339354451 0.588409662951354
0.873918990836031 0.573582343356909
0.907342500861859 0.558183264408538
0.939861552150259 0.542227776243028
0.971443729003692 0.525731783637977
1.00205754962143 0.508711730157565
1.03167249748132 0.491184581761245
1.06025905175925 0.473167809891714
1.08778871675625 0.454679374058994
1.11423405030353 0.435737703938003
1.1395686911175 0.416361680997454
1.16376738507721 0.396570619678398
1.18680601039828 0.376384248141167
1.20866160167793 0.355822688599914
1.22931237278748 0.334906437264357
1.24873773858917 0.313656343908707
1.26691833545593 0.292093591088158
1.2838360405734 0.270239673023664
1.29947399000521 0.24811637417602
1.31381659550327 0.225745747530649
1.32684956004651 0.2031500926147
1.33855989209243 0.180351933268395
1.34893591852738 0.157373995192777
1.35796729630255 0.134239183296235
1.36564502274417 0.110970558862397
1.37196144452754 0.0875913165621358
1.37691026530606 0.0641247613326195
1.38048655198757 0.0405942851464444
1.38268673965171 0.0170233436940044
1.38350863510357 -0.00656456699764701
1.3829514190599 -0.0301459339861518
1.38101564696576 -0.053697250852056
1.37770324844088 -0.077195041130521
1.37301752535614 -0.100615881713177
1.36696314854221 -0.123936426196781
1.35954615313359 -0.147133428155427
1.35077393255258 -0.17018376431308
1.34065523113944 -0.193064457593369
1.32920013543579 -0.21575270002363
1.31642006413016 -0.238225875470392
1.30232775667556 -0.260461582183633
1.28693726059054 -0.282437655127325
1.27026391745629 -0.304132188074029
1.25232434762387 -0.325523555441501
1.23313643364661 -0.346590433849534
1.21271930245448 -0.367311823375579
1.19109330628786 -0.387667068487913
1.16828000241006 -0.40763587863553
1.1443021316185 -0.427198348474199
1.11918359557626 -0.446334977708548
1.09294943298633 -0.465026690530379
1.06562579463252 -0.483254854633851
1.03723991731176 -0.501001299788564
1.00782009668389 -0.518248335952047
0.977395659065919 -0.534978770903567
0.945996932198934 -0.551175927381719
0.913655215016722 -0.566823659708674
0.880402746446324 -0.581906369884548
0.846272673271575 -0.596409023135823
0.811299017091666 -0.61031716290234
0.77551664040767 -0.623616925247911
0.73896121187085 -0.63629505268019
0.701669170727365 -0.648338907366034
0.663677690494841 -0.659736483729166
0.625024641907002 -0.670476420417598
0.58574855516329 -0.680548011628869
0.545888581521136 -0.689941217781826
0.505484454269125 -0.698646675524297
0.464576449119999 -0.70665570706668
0.423205344062944 -0.713960328832156
0.381412378715204 -0.720553259414892
0.339239213213537 -0.7264279268383
0.296727886686481 -0.731578475106134
0.253920775348841 -0.735999770039868
0.210860550260154 -0.739687404396559
0.167590134789245 -0.742637702262083
0.124152661827286 -0.744847722715365
0.0805914307919886 -0.746315262759956
0.0369498644658028 -0.747038859520022
-0.00672853428885405 -0.747017791698577
};
\addlegendentry{\textsc{gah}}
\addplot [line width=0.64pt, crimson2143940, dashed]
table {%
-0.0057400872837216 -0.730980162994304
-0.0495883618092768 -0.730271489558875
-0.09338720557735 -0.728834865620963
-0.137092958947951 -0.726671723238553
-0.180662055075623 -0.723784218679014
-0.22405106333782 -0.720175230269688
-0.267216732627497 -0.715848355528709
-0.310116034466754 -0.710807907578918
-0.35270620589856 -0.705058910848449
-0.394944792113793 -0.698607096062269
-0.436789688771125 -0.691458894529671
-0.47819918396754 -0.683621431733398
-0.519131999817667 -0.675102520226811
-0.559547333600473 -0.665910651846168
-0.599404898432305 -0.65605498924577
-0.63866496342573 -0.645545356764437
-0.677288393294149 -0.634392230632385
-0.7152366873627 -0.622606728528296
-0.752472017946568 -0.610200598496976
-0.788957268058448 -0.597186207238644
-0.824656068407558 -0.583576527781538
-0.859532833653346 -0.56938512655012
-0.893552797877727 -0.554626149841767
-0.926682049240508 -0.539314309725434
-0.958887563783447 -0.523464869376348
-0.990137238349253 -0.507093627861335
-1.02039992258271 -0.490216904389973
-1.04964544998204 -0.472851522047249
-1.0778446679695 -0.455014791023938
-1.10496946695133 -0.43672449136143
-1.13099280833802 -0.417998855228202
-1.15588875149692 -0.398856548745587
-1.17963247961045 -0.37931665338099
-1.20220032441408 -0.359398646927055
-1.22356978978929 -0.339122384085787
-1.24371957418827 -0.318508076676946
-1.26262959186772 -0.297576273490471
-1.28028099291077 -0.276347839802995
-1.29665618201698 -0.254843936578887
-1.31173883604172 -0.233085999376539
-1.32551392026746 -0.21109571698094
-1.33796770339066 -0.188895009783817
-1.34908777120946 -0.166506007932919
-1.35886303899843 -0.143951029272199
-1.36728376255804 -0.121252557094904
-1.37434154792793 -0.0984332177317374
-1.38002935975413 -0.0755157579964448
-1.38434152830214 -0.0525230225112907
-1.38727375510855 -0.029477930935045
-1.38882311726593 -0.00640345511616751
-1.38898807033637 0.0166774038060301
-1.38776844989107 0.0397416383294193
-1.38516547167422 0.0627662575234345
-1.3811817303911 0.0857283099469315
-1.37582119712167 0.108604906526684
-1.36908921536206 0.131373243373709
-1.36099249569811 0.154010624514679
-1.35153910911605 0.176494484515764
-1.34073847895723 0.198802410976342
-1.32860137152466 0.220912166870174
-1.315139885351 0.242801712711748
-1.30036743913841 0.264449228525724
-1.28429875838255 0.285833135597556
-1.26694986069385 0.306932117983624
-1.24833803983083 0.327725143759433
-1.22848184846129 0.34819148598469
-1.2074010796686 0.368310743364371
-1.18511674722155 0.388062860585172
-1.16165106462737 0.407428148307085
-1.13702742298879 0.426387302790159
-1.11127036768729 0.44492142513689
-1.08440557391576 0.463012040131048
-1.05645982108488 0.480641114654178
-1.02746096612882 0.497791075661405
-0.997437915736854 0.514444827698625
-0.966420597538551 0.530585769943628
-0.934439930271242 0.546197812754157
-0.901527792959564 0.561265393706411
-0.86771699313775 0.575773493108015
-0.833041234146369 0.589707648969967
-0.7975350815361 0.603053971422668
-0.76123392861204 0.615799156561647
-0.724173961152889 0.627930499709179
-0.686392121340177 0.639435908078587
-0.647926070933495 0.650303912828593
-0.608814153728436 0.660523680495707
-0.569095357334662 0.670085023793264
-0.528809274312209 0.678978411766327
-0.487996062704761 0.687194979292355
-0.44669640600924 0.694726535918147
-0.40495147262161 0.701565574024264
-0.362802874799328 0.70770527630879
-0.320292627181335 0.713139522582963
-0.277463104906955 0.717862895871914
-0.234357001375432 0.721870687814428
-0.191017285688224 0.725158903356336
-0.147487159816464 0.727724264732874
-0.1038100155363 0.729564214736032
-0.0600293911750249 0.730676919263632
-0.0161889282111277 0.731061269147602
0.0276676722284809 0.730716881259619
0.0714966929307053 0.729644098893021
0.115254444174531 0.727843991420599
0.15889730728189 0.725318353228631
0.202381778097708 0.722069701928193
0.245664510355789 0.718101275845553
0.28870235888732 0.713417030794132
0.331452422628914 0.708021636131271
0.37387208738732 0.701920470103705
0.415919068318178 0.695119614486417
0.457551452076473 0.68762584852019
0.498727738596668 0.679446642153909
0.539406882460872 0.670590148598361
0.579548333813806 0.661065196198927
0.619112078783778 0.650881279635307
0.658058679369386 0.640048550457006
0.696349312752165 0.628577806964056
0.733945809996031 0.616480483443033
0.770810694094894 0.603768638769111
0.80690721733056 0.590454944385513
0.842199397903653 0.576552671672342
0.876652055801057 0.562075678717381
0.910230847864113 0.547038396502049
0.942902302022634 0.531455814516282
0.974633850660588 0.515343465816689
1.0053938630802 0.498717411542859
1.03515167703215 0.48159422490727
1.0638776292803 0.463990974674751
1.09154308517076 0.44592520814796
1.11812046717541 0.427414933675862
1.14358328238188 0.408478602702604
1.16790614890223 0.389135091374721
1.19106482117413 0.369403681724982
1.21303621412943 0.349304042451635
1.23379842620579 0.32885620931223
1.25333076117868 0.308080565151533
1.27161374879184 0.286997819583462
1.28862916416558 0.265628988347303
1.30436004596382 0.243995372358752
1.31879071330148 0.222118536476707
1.33190678137552 0.200020288006945
1.34369517580404 0.177722654964127
1.3541441456591 0.15524786411379
1.36324327518031 0.132618318816227
1.37098349415748 0.109856576694331
1.377357086972 0.0869853271476638
1.38235770028794 0.0640273687351688
1.3859803493852 0.0410055864490743
1.38822142312838 0.0179429289026314
1.38907868756642 -0.00513761454555756
1.3885512881595 -0.0282130367078371
1.38663975063081 -0.0512603355015589
1.38334598044253 -0.0742565368780602
1.37867326089642 -0.0971787177236967
1.37262624986094 -0.120004028710097
1.3652109751282 -0.142709717070873
1.35643482840532 -0.165273149282072
1.34630655794621 -0.187671833623763
1.33483625983113 -0.209883442600276
1.3220353679027 -0.231885835196733
1.30791664236837 -0.253657078949694
1.29249415708082 -0.275175471809911
1.27578328550878 -0.296419563775408
1.25780068541244 -0.317368178273304
1.23856428223868 -0.338000433269083
1.21809325125252 -0.358295762082265
1.19640799842284 -0.378233933887711
1.17353014008126 -0.397795073882158
1.14948248137453 -0.41695968309585
1.12428899353182 -0.435708657829538
1.09797478996971 -0.454023308697461
1.07056610125858 -0.471885379257333
1.04209024897542 -0.489277064208762
1.01257561846906 -0.50618102714196
0.982051630565073 -0.522580417819054
0.950548712238418 -0.538458888970769
0.918098266283203 -0.553800612591741
0.884732640009673 -0.568590295718218
0.850485092999724 -0.582813195672419
0.815389763953019 -0.596455134758355
0.779481636656792 -0.609502514394465
0.74279650511325 -0.621942328668978
0.70537093785933 -0.633762177304488
0.66724224151439 -0.644950278018821
0.628448423592157 -0.65549547826987
0.589028154614009 -0.665387266372693
0.549020729561363 -0.674615781977783
0.508466028705577 -0.683171825900089
0.467404477854435 -0.691046869288954
0.425877008054818 -0.698233062129865
0.383925014791746 -0.704723241069515
0.341590316724463 -0.710510936556391
0.29891511400068 -0.715590379289762
0.25594194619056 -0.719956505970646
0.212713649882341 -0.723604964349015
0.169273315981896 -0.726532117562212
0.125664246758784 -0.728735047760254
0.0819299126815998 -0.730211559014413
0.038113909085671 -0.730960179506158
-0.00574008728372126 -0.730980162994304
};
\addlegendentry{log-euclidienne}

\nextgroupplot[
height=4.6cm,
label style={font=\footnotesize},
legend style={draw=none, fill=white, fill opacity=0.9, text opacity=1, font=\footnotesize},
log basis y={10},
tick align=outside,
tick label style={font=\scriptsize},
tick pos=left,
title style={font=\footnotesize, yshift=-2pt},
title={déterminants},
width=0.45\textwidth,
x grid style={darkgray176},
xlabel={matrices triées},
xmin=0.3, xmax=15.7,
xtick style={color=black},
y grid style={darkgray176},
ylabel={\(\displaystyle \det\)},
ymin=0.349511235857449, ymax=7.54244028810198,
ymode=log,
ytick style={color=black}
]
\addplot [semithick, steelblue31119180, mark=*, mark size=1.5, mark options={solid}, only marks]
table {%
1 0.401883162979835
2 0.453860550127679
3 0.711936822021271
4 0.734112973361792
5 0.787734131840946
6 0.860462142381226
7 1.02471693850366
8 1.03265102678936
9 1.05860278331886
10 1.15728235079122
11 1.15754370135465
12 1.21698771074158
13 1.26924071070826
14 1.6123109079403
15 6.55953737133251
};
\addplot [line width=0.56pt, darkorange25512714]
table {%
0.300000000000001 1.32297930186553
15.7 1.32297930186553
};
\addplot [line width=0.56pt, mediumpurple148103189]
table {%
0.300000000000001 0.862634036808105
15.7 0.862634036808105
};
\addplot [line width=0.56pt, forestgreen4416044]
table {%
0.300000000000001 1.03113888574665
15.7 1.03113888574665
};
\addplot [line width=0.56pt, sienna1408675, dashed]
table {%
0.300000000000001 1.06829161551602
15.7 1.06829161551602
};
\addplot [line width=0.56pt, crimson2143940, dashed]
table {%
0.300000000000001 1.03113888574665
15.7 1.03113888574665
};
\end{groupplot}

\end{tikzpicture}

  \end{center}
  \provMoyennesQR[.5cm]
  \caption{Moyennes d'une collection de matrices de covariance $2\times 2$
    tirées le long de géodésiques issues d'un même centre. À gauche, les
    ellipses de concentration: la moyenne arithmétique est sensiblement plus
    large que les deux autres, ce qui est l'effet de gonflement. À droite, les
    déterminants: celui de la moyenne arithmétique dépasse la plupart de ceux
    de l'échantillon, tandis que les moyennes log-euclidienne et de Fréchet
    ont toutes deux pour déterminant la moyenne géométrique des déterminants,
    conséquence de la $g$-linéarité de $\log\Det$ le long des géodésiques.
    \provMoyennes}
  \label{fig:context-riemann-moyennes}
\end{figure}


\section{Les matrices orthogonales}
\label{sec:context-riemann-orthogonal}

Le cône n'est pas la seule variété que nous rencontrerons. Deux besoins font apparaître les matrices à colonnes orthonormées. Le premier vient de l'apprentissage profond sur variété \textsc{spd}: dans une architecture de type \textsc{spdnet} \sidecite[0cm]{huang2017riemannian}, la transformation élémentaire d'une couche est $\MCov\mapsto\mat{W}^\T\MCov\mat{W}$ avec $\mat{W}$ à colonnes orthonormées, de sorte que l'apprentissage des poids \emph{est} un problème d'optimisation sur la variété correspondante. Le second est la réduction de dimension: dès que l'on cherche à représenter les données dans un sous-espace, c'est ce sous-espace lui-même qui devient le paramètre.

\begin{definition}[{Variétés de Stiefel et de Grassmann}]
\label{def:context-riemann-stiefel}
  Pour $k\leq\Nvar$, la \textbf{variété de Stiefel} est l'ensemble des matrices à colonnes orthonormées
  \begin{equation}
    \Stiefel{\Nvar}{k}
    \eqdef
    \left\{ \mat{W}\in\R^{\Nvar\times k};\; \mat{W}^\T\mat{W}=\Id{k} \right\},
    \label{eq:context-riemann-stiefel}
  \end{equation}
  de dimension $\Nvar k - k(k+1)/2$. Le cas $k=\Nvar$ redonne le \textbf{groupe orthogonal} $\Orthgroup{\Nvar}$, de dimension $\Nvar(\Nvar-1)/2$. La \textbf{variété de Grassmann} $\Grassmann{\Nvar}{k}$ est l'ensemble des sous-espaces de dimension $k$ de $\R^\Nvar$; elle s'obtient à partir de $\Stiefel{\Nvar}{k}$ en identifiant les bases qui engendrent le même sous-espace.
\end{definition}

Contrairement au cône, il s'agit ici du cas \emph{contraint} annoncé en \cref{def:context-riemann-manifold}, et la différence se voit immédiatement sur les objets de base. L'espace tangent en $\mat{W}$ s'obtient en différentiant la contrainte $\mat{W}^\T\mat{W}=\Id{k}$, ce qui donne $\Tangent{\mat{W}}{\Stiefel{\Nvar}{k}}=\{\tvec;\,\mat{W}^\T\tvec+\tvec^\T\mat{W}=\mat{0}\}$, et la projection orthogonale associée à la métrique héritée de $\R^{\Nvar\times k}$ n'est plus l'identité mais
\begin{equation}
  \projman_{\mat{W}}(\mat{Z})
  =
  \mat{Z} - \mat{W}\,\symop\!\left(\mat{W}^\T\mat{Z}\right).
  \label{eq:context-riemann-projstiefel}
\end{equation}
La géométrie complète de ces variétés est développée dans l'article de référence \sidecite[-0.5cm]{edelman1998geometry}.

\begin{remark}[{Deux différences pratiques avec le cône}]
\label{rem:context-riemann-stiefelvscone}
  Deux points distinguent ces variétés du cône et auront des conséquences algorithmiques.

  D'une part, $\Stiefel{\Nvar}{k}$ est \textbf{compacte}. Aucune divergence n'est possible, et la question de la complétude qui nous occupait tant sur le cône ne se pose pas. En revanche l'exponentielle riemannienne y fait intervenir une exponentielle de matrice coûteuse, ce qui motive l'emploi de \emph{rétractions} approchées --- typiquement le facteur $\mat{Q}$ d'une décomposition \textsc{qr} --- que nous introduirons en \cref{sec:context-riemann-optimization}.

  D'autre part, $\Orthgroup{\Nvar}$ n'est \textbf{pas connexe}: elle possède deux composantes, selon que le déterminant vaut $+1$ ou $-1$. Aucun algorithme itératif procédant par petits déplacements le long de la variété ne peut passer de l'une à l'autre, et le signe du déterminant est donc entièrement fixé par l'initialisation. C'est une différence de nature avec le cône, qui est connexe et même convexe, et c'est un piège classique en pratique.
\end{remark}


\section{Optimisation en contexte riemannien}
\label{sec:context-riemann-optimization}

Nous en venons à ce qui motive l'essentiel de ce chapitre. Presque tous les estimateurs des \cref{ch:context-definition,ch:context-distributions} sont définis comme minimiseurs d'un critère sous contrainte de positivité. Nous disposons maintenant du cadre pour traiter cette contrainte non pas comme une gêne mais comme la structure même du problème.

Le principe général est simple, et se décline en deux étapes qui rejouent la descente de gradient euclidienne \sidecite[0cm]{absil2008optimization,boumal2023introduction}: en un itéré $\MCov_{(\iiter)}$, on détermine d'abord une direction de descente \emph{dans l'espace tangent}, puis on l'utilise pour produire un nouvel itéré \emph{sur la variété}.

\subsection{Gradient riemannien}

\begin{definition}[{Gradient riemannien}]
\label{def:context-riemann-gradient}
  Soit $f:\Manifold\to\R$ et $\param\in\Manifold$. Le \textbf{gradient riemannien} $\Rgrad f(\param)$ est l'unique vecteur tangent vérifiant
  \begin{equation}
    \metric{\Rgrad f(\param)}{\tvec}{\param}
    =
    \dirderiv f(\param)[\tvec],
    \qquad
    \forall\tvec\in\Tangent{\param}{\Manifold}.
    \label{eq:context-riemann-defgradient}
  \end{equation}
\end{definition}

La définition mérite qu'on s'y arrête, car elle contient tout. Le membre de droite, la dérivée directionnelle, ne dépend que de la fonction: c'est une donnée du problème. Le membre de gauche fait intervenir la métrique. Le gradient riemannien est donc le \emph{même} objet analytique lu à travers une géométrie différente, et changer de métrique change la direction de plus forte pente sans changer la fonction. On en déduit une recette systématique en deux temps: calculer le gradient euclidien, puis le convertir.

\begin{property}[{Gradient sur le cône}]
\label{prop:context-riemann-gradientspd}
  Soit $f:\SPD{\Nvar}\to\R$ et $\nabla^{\mathrm{e}}f(\MCov)\in\Sym{\Nvar}$ son gradient euclidien, défini par $\dirderiv f(\MCov)[\tvec]=\trace{\nabla^{\mathrm{e}}f(\MCov)\,\tvec}$. Alors le gradient riemannien associé à la métrique affine invariante naturelle est
  \begin{equation}
    \Rgrad f(\MCov)
    =
    \MCov\,\nabla^{\mathrm{e}}f(\MCov)\,\MCov.
    \label{eq:context-riemann-gradientspd}
  \end{equation}
\end{property}

\begin{remark}[{Un préconditionnement gratuit}]
  Le facteur $\MCov\cdot\MCov$ de \cref{eq:context-riemann-gradientspd} n'est pas un simple changement d'échelle. Il agit comme un préconditionneur qui adapte le pas à l'anisotropie locale du problème: dans les directions où $\MCov$ est de faible énergie, le pas effectif est réduit, ce qui empêche l'itéré de sortir du cône. C'est la traduction algorithmique de la propriété de complétude, et c'est la raison pour laquelle les méthodes riemanniennes se comportent bien sur les problèmes mal conditionnés évoqués en \cref{sec:context-covariance-regularized-estimation}.
\end{remark}

\subsection{Rétractions}

Une fois la direction $\tvec$ obtenue, il faut se déplacer. L'exponentielle riemannienne \cref{eq:context-riemann-explogspd} est le déplacement idéal, mais elle demande une exponentielle de matrice à chaque itération. On se contente en pratique d'une approximation, à condition qu'elle reste sur la variété.

\begin{definition}[{Rétraction}]
\label{def:context-riemann-retraction}
  Une \textbf{rétraction} en $\param\in\Manifold$ est une application $\Retr_\param:\Tangent{\param}{\Manifold}\to\Manifold$ vérifiant $\Retr_\param(\tvec)=\param+\tvec+o(\norm{\tvec})$.
\end{definition}

Sur le cône, qui est ouvert dans $\Sym{\Nvar}$, l'approximation au premier ordre est simplement $\Retr^{(1)}_{\MCov}(\tvec)=\MCov+\tvec$ --- mais elle ne garantit pas de rester dans $\SPD{\Nvar}$. On lui préfère souvent l'approximation au second ordre suivante, qui n'a pas ce défaut \sidecite[0cm]{jeuris2012survey,boumal2023introduction}:
\begin{equation}
  \Retr^{(2)}_{\MCov}(\tvec)
  =
  \MCov + \tvec + \tfrac{1}{2}\tvec\inv{\MCov}\tvec
  \;\in\;\SPD{\Nvar},
  \qquad \forall\tvec\in\Sym{\Nvar}.
  \label{eq:context-riemann-retraction2}
\end{equation}
La positivité se vérifie en écrivant $\Retr^{(2)}_{\MCov}(\tvec)=\MCov^{1/2}(\Id{\Nvar}+\mat{L}+\tfrac{1}{2}\mat{L}^2)\MCov^{1/2}$ avec $\mat{L}=\MCov^{-1/2}\tvec\MCov^{-1/2}$, et en constatant que le polynôme $\eig\mapsto 1+\eig+\eig^2/2$ est strictement positif sur $\R$.

L'algorithme de descente de gradient riemannien s'écrit alors
\begin{equation}
  \MCov_{(\iiter+1)}
  =
  \Retr_{\MCov_{(\iiter)}}\!\left(-\gamma_\iiter\,\Rgrad f(\MCov_{(\iiter)})\right),
  \label{eq:context-riemann-descente}
\end{equation}
où le pas $\gamma_\iiter$ est fixé ou obtenu par recherche linéaire. Un pas de cet algorithme est illustré en \cref{fig:context-riemann-descente}.

\begin{figure}
  \centering
  % Un pas de descente de gradient riemannienne : la direction est determinee
% dans l'espace tangent, la retraction la ramene sur la variete.
% Utilise dans la \cref{sec:context-riemann-optimization}.
\begin{tikzpicture}[
    >={Stealth[length=5pt,width=4pt]},
    line join=round,
    font=\footnotesize,
  ]
  \definecolor{rmanifold}{RGB}{90,90,90}
  \definecolor{rtangent}{RGB}{50,110,120}
  \definecolor{rgeodesic}{RGB}{40,70,140}
  \definecolor{rvector}{RGB}{175,95,45}

  \coordinate (Sn)  at (-2.95,1.35);
  \coordinate (Sn1) at ( 0.75,2.10);
  \coordinate (Sst) at ( 3.55,1.70);

  % ---------------------------------------------------------------- variete
  \fill[rmanifold!7]
    (-5.0,0.2) .. controls (-4.6,2.45) and (-2.2,3.3) .. (-0.2,3.3)
    .. controls ( 2.1,3.3) and ( 4.4,2.35) .. ( 5.0,0.2)
    .. controls ( 3.0,-0.9) and (-3.0,-0.95) .. (-5.0,0.2) -- cycle;
  \draw[rmanifold!70, thick]
    (-5.0,0.2) .. controls (-4.6,2.45) and (-2.2,3.3) .. (-0.2,3.3)
    .. controls ( 2.1,3.3) and ( 4.4,2.35) .. ( 5.0,0.2);
  \draw[rmanifold!70, thick]
    (-5.0,0.2) .. controls (-3.0,-0.9) and (3.0,-0.95) .. (5.0,0.2);
  \node[rmanifold!85, font=\normalsize] at (4.30,2.65) {$\SPD{\Nvar}$};

  % ---------------------------------------------------- plan tangent en Sn
  \fill[rtangent!10, draw=rtangent!65]
    ($(Sn)+(-1.55,-0.45)$) -- ($(Sn)+(2.10,-1.05)$)
    -- ($(Sn)+(2.75,0.25)$) -- ($(Sn)+(-0.90,0.85)$) -- cycle;
  \node[rtangent!90, anchor=east] at ($(Sn)+(-1.65,0.55)$)
    {$\Tangent{\MCov_{(\iiter)}}{\SPD{\Nvar}}$};

  % -------------------------------------------------- direction de descente
  \draw[rvector, very thick, ->] (Sn) -- ($(Sn)+(2.30,-0.28)$);
  \node[rvector, anchor=north] at ($(Sn)+(1.15,-0.42)$)
    {$-\gamma_\iiter\,\Rgrad f(\MCov_{(\iiter)})$};

  % ---------------------------------------------------------- retraction
  \draw[rvector!60, dotted, very thick]
    ($(Sn)+(2.33,-0.26)$) .. controls (-0.30,1.05) and (0.20,1.55) .. (Sn1);
  \node[rvector!85, anchor=north west] at (-0.35,1.20)
    {$\Retr_{\MCov_{(\iiter)}}$};

  % ------------------------------------------------ suite des iteres et opt.
  \draw[rgeodesic!50, thick, dashed]
    (Sn1) .. controls (2.05,2.45) .. (Sst);
  \node[rgeodesic!70, anchor=south] at (2.15,2.44) {itérations suivantes};

  \fill (Sn) circle (1.8pt);
  \node[anchor=south east, inner sep=3pt] at (Sn) {$\MCov_{(\iiter)}$};
  \fill (Sn1) circle (1.8pt);
  \node[anchor=east, inner sep=4pt] at (Sn1) {$\MCov_{(\iiter+1)}$};
  \fill[rgeodesic] (Sst) circle (2.1pt);
  \node[anchor=west, inner sep=4pt] at (Sst) {$\MCov^{\opt}$};
\end{tikzpicture}

  \caption{Un pas de descente de gradient riemannienne sur le cône. La direction
    de descente $-\gamma_\iiter\Rgrad f(\MCov_{(\iiter)})$ est déterminée dans
    l'espace tangent, puis ramenée sur la variété par la rétraction
    $\Retr_{\MCov_{(\iiter)}}$. L'itéré ne peut donc jamais sortir de
    $\SPD{\Nvar}$, contrairement à ce qui se produirait avec une descente
    euclidienne.}
  \label{fig:context-riemann-descente}
\end{figure}

\subsection{Ce que cela révèle sur les estimateurs du chapitre}

Nous pouvons maintenant relire la \cref{subsec:context-robust-m-estimation} sous cet éclairage. Rappelons que l'estimateur du maximum de vraisemblance d'un modèle elliptique était obtenu comme solution de l'équation de point fixe \cref{eq:context-robust-fixedpoint}, et calculé par l'itération correspondante --- une recette dont nous n'avions justifié que la convergence, sans dire d'où elle venait.

\begin{property}[{L'algorithme du point fixe est une descente de gradient riemannienne}]
\label{prop:context-riemann-pointfixe}
  Soit $L(\MScatter)=-\likelihood(\X;\MScatter)$ l'opposé de la log-vraisemblance d'un échantillon elliptique. Son gradient riemannien pour la métrique affine invariante naturelle vaut
  \begin{equation}
    \Rgrad L(\MScatter)
    =
    \Nobs\,\MScatter
    -
    \sum_{\iobs=1}^{\Nobs}
      \scorefun\!\left(\quadratic{\x_\iobs}{\MScatter}\right)\x_\iobs\x_\iobs^\T .
    \label{eq:context-riemann-gradientvraisemblance}
  \end{equation}
  L'itération de point fixe \cref{eq:context-robust-iterations} est exactement l'itération \cref{eq:context-riemann-descente} obtenue en choisissant la rétraction du premier ordre $\Retr^{(1)}$ et le pas constant $\gamma_\iiter=1/\Nobs$.
\end{property}

\begin{proof}
  Il suffit de reporter \cref{eq:context-riemann-gradientvraisemblance} dans $\Retr^{(1)}_{\MScatter}(-\gamma\Rgrad L(\MScatter))=\MScatter-\gamma\Rgrad L(\MScatter)$ avec $\gamma=1/\Nobs$, ce qui donne directement $\frac{1}{\Nobs}\sum_\iobs\scorefun(\quadratic{\x_\iobs}{\MScatter})\x_\iobs\x_\iobs^\T$.
\end{proof}

\begin{remark}
  Cette lecture, due à \sidecite[0cm]{bouchard2024fisherrao}, éclaire plusieurs choses d'un coup. Elle explique d'abord pourquoi la rétraction du premier ordre suffit ici alors qu'elle n'est pas sûre en général: la structure particulière de \cref{eq:context-riemann-gradientvraisemblance} fait que l'itéré est une somme de matrices de rang un pondérées positivement, donc automatiquement dans le cône. Elle indique ensuite que le point fixe n'est qu'un choix parmi d'autres --- métrique $(\alpha,\beta)$, pas, rétraction sont autant de degrés de liberté --- et ouvre donc la porte à des variantes plus rapides, ou adaptées à des contraintes structurelles supplémentaires. Elle permet enfin d'aborder des problèmes où le point fixe n'existe pas, typiquement les critères pénalisés de la \cref{sec:context-covariance-regularized-estimation}, pour lesquels la descente riemannienne reste parfaitement définie.
\end{remark}

\subsection{Méthodes d'ordre supérieur et transport}

Le schéma \cref{eq:context-riemann-descente} se prête aux mêmes raffinements que son homologue euclidien, à une nuance près. Une méthode de gradient conjugué construit la direction de descente comme combinaison de la direction courante et de la précédente; or celles-ci vivent dans des espaces tangents attachés à deux points distincts, et ne peuvent pas être additionnées directement. C'est ici que sert le transport parallèle introduit en \cref{sec:context-riemann-notion}: on transporte la direction précédente en l'itéré courant avant de combiner. La même remarque vaut pour les méthodes de quasi-Newton. Les méthodes du second ordre demandent quant à elles le hessien riemannien, que nous n'introduisons pas.

Signalons enfin, car nous en aurons directement besoin au \cref{ch:robust-detection}, que la descente de gradient \emph{stochastique} s'étend également à ce cadre \sidecite[0cm]{bonnabel2013stochastic,Zhou_2019}. C'est ce qui rend possible l'estimation récursive, où les observations arrivent une à une et où l'on ne peut pas se permettre de recalculer un point fixe à chaque nouvel échantillon.

\subsection{Convexité géodésique}

Il reste une question de fond. Les algorithmes ci-dessus convergent vers un point critique; qu'est-ce qui garantit que ce point critique est le minimum global cherché ? Les critères que nous manipulons ne sont pas convexes au sens euclidien --- celui de Tyler ne l'est pas --- et pourtant leurs minimiseurs sont uniques. La géométrie fournit l'explication.

\begin{definition}[{Convexité géodésique}]
\label{def:context-riemann-gconvexite}
  Une fonction $f:\Manifold\to\R$ est dite \textbf{géodésiquement convexe}, ou \textbf{$g$-convexe}, si pour tout couple de points et toute géodésique $\geod$ les joignant,
  \begin{equation}
    f(\geod(t)) \leq (1-t)f(\geod(0)) + t\,f(\geod(1)),
    \qquad \forall t\in[0,1].
    \label{eq:context-riemann-gconvexite}
  \end{equation}
  Si l'inégalité est stricte, $f$ est dite strictement $g$-convexe.
\end{definition}

Il s'agit de la convexité usuelle, où les segments ont été remplacés par les géodésiques, et elle en conserve la conséquence essentielle: \textbf{tout minimum local d'une fonction $g$-convexe est global, et il est unique si la fonction est strictement $g$-convexe}. La \cref{fig:context-riemann-gconvexite} en donne une illustration frappante sur le critère de Tyler, non convexe au sens euclidien mais $g$-convexe.

Sur le cône muni de la métrique affine invariante, deux briques élémentaires suffisent à traiter la plupart de nos critères:
\begin{itemize}
  \item $\MCov\mapsto\log\Det(\MCov)$ est $g$-\emph{linéaire}: nous l'avons déjà constaté en \cref{rem:context-riemann-normalisation}, puisque $\log\Det(\MCov_0\geodmean_t\MCov_1)=(1-t)\log\Det(\MCov_0)+t\log\Det(\MCov_1)$;
  \item $\MCov\mapsto\trace{\inv{\MCov}\mat{S}}$ est $g$-convexe pour toute $\mat{S}\in\SPD{\Nvar}$.
\end{itemize}

\begin{remark}[{Retour sur l'unicité des M-estimateurs}]
\label{rem:context-riemann-gconvexiteM}
  Ces deux briques donnent immédiatement la $g$-convexité de la log-vraisemblance gaussienne négative $\trace{\inv{\MCov}\Smat}+\Nobs\log\Det(\MCov)$, et par extension celle des critères associés aux M-estimateurs sous les hypothèses de \sidecite[0cm]{6298979,wiesel2012unified}. Nous tenons donc la promesse faite en \cref{subsec:context-robust-m-estimation}: les conditions d'existence et d'unicité du point fixe de Tyler, que nous avions énoncées de manière technique, s'expliquent par la $g$-convexité stricte du critère associé. Le cas de Tyler est d'ailleurs le plus instructif: son critère est $g$-convexe mais pas strictement, la direction $\MCov\mapsto c\MCov$ le long de laquelle il est constant étant précisément celle qui rend la solution non identifiable en échelle. L'unicité «~à un facteur multiplicatif près~» n'est donc pas un accident de calcul, c'est la lecture géométrique exacte de la $g$-convexité non stricte. Un panorama des critères pénalisés préservant la $g$-convexité est donné dans \sidecite[1cm]{wiesel2015structured,sun2016robust}.
\end{remark}

\begin{figure}[h]
  \begin{center}
    % This file was created with matplot2tikz v0.5.3.
\begin{tikzpicture}

\definecolor{darkgray176}{RGB}{176,176,176}
\definecolor{darkorange25512714}{RGB}{255,127,14}
\definecolor{forestgreen4416044}{RGB}{44,160,44}
\definecolor{gray127}{RGB}{127,127,127}

\begin{groupplot}[group style={group size=2 by 1, vertical sep=1.5cm, horizontal sep=1.4cm}]
\nextgroupplot[
height=4.6cm,
label style={font=\footnotesize},
legend style={draw=none, fill=white, fill opacity=0.9, text opacity=1, font=\footnotesize},
tick align=outside,
tick label style={font=\scriptsize},
tick pos=left,
title style={font=\footnotesize, yshift=-2pt},
title={segment euclidien},
width=0.45\textwidth,
x grid style={darkgray176},
xlabel={\(\displaystyle t\)},
xmin=-0.05, xmax=1.05,
xtick style={color=black},
y grid style={darkgray176},
ylabel={\(\displaystyle L\)},
ymin=0.38937564836453, ymax=9.69140677310581,
ytick style={color=black}
]
\addplot [semithick, darkorange25512714]
table {%
0 7.22321310791446
0.005 3.66392482041289
0.01 3.51977126800225
0.015 3.50057139068436
0.02 3.51491899912505
0.025 3.54135720837481
0.03 3.57246801517583
0.035 3.60516211580002
0.04 3.63801529420469
0.045 3.67033711640028
0.05 3.70179095596579
0.055 3.73222144564196
0.06 3.76156977670936
0.065 3.78982963623954
0.07 3.81702326669366
0.075 3.84318803759633
0.08 3.8683687480987
0.085 3.89261316198124
0.09 3.9159694154759
0.095 3.93848453266715
0.1 3.96020360552759
0.105 3.98116937618751
0.11 4.00142206303059
0.115 4.02099933355579
0.12 4.03993636387945
0.125 4.05826594737571
0.13 4.07601862901224
0.135 4.09322285077584
0.14 4.10990509918842
0.145 4.12609004948971
0.15 4.14180070335066
0.155 4.15705851844339
0.16 4.17188352912091
0.165 4.18629445804178
0.17 4.20030881893048
0.175 4.21394301087318
0.18 4.2272124046621
0.185 4.24013142175401
0.19 4.25271360642096
0.195 4.26497169166111
0.2 4.2769176594114
0.205 4.28856279557098
0.21 4.29991774030714
0.215 4.31099253407735
0.22 4.32179665976365
0.225 4.33233908127983
0.23 4.34262827897809
0.235 4.35267228215105
0.24 4.36247869889601
0.245 4.37205474358279
0.25 4.38140726214249
0.255 4.39054275537342
0.26 4.39946740044104
0.265 4.40818707073166
0.27 4.41670735420385
0.275 4.42503357036776
0.28 4.43317078600983
0.285 4.44112382976938
0.29 4.44889730566305
0.295 4.45649560564451
0.3 4.46392292127829
0.305 4.47118325459953
0.31 4.47828042822463
0.315 4.48521809477216
0.32 4.49199974564784
0.325 4.49862871924252
0.33 4.50510820858815
0.335 4.51144126851219
0.34 4.51763082232812
0.345 4.52367966809574
0.35 4.52959048448272
0.355 4.53536583625574
0.36 4.54100817942753
0.365 4.54651986608363
0.37 4.55190314891108
0.375 4.55716018544909
0.38 4.56229304208046
0.385 4.56730369778067
0.39 4.57219404764064
0.395 4.57696590617736
0.4 4.58162101044613
0.405 4.58616102296637
0.41 4.59058753447273
0.415 4.59490206650173
0.42 4.59910607382388
0.425 4.60320094673005
0.43 4.60718801318049
0.435 4.61106854082417
0.44 4.61484373889541
0.445 4.6185147599946
0.45 4.62208270175874
0.455 4.62554860842766
0.46 4.62891347231099
0.465 4.63217823516073
0.47 4.63534378945368
0.475 4.63841097958803
0.48 4.64138060299767
0.485 4.64425341118787
0.49 4.64703011069532
0.495 4.6497113639756
0.5 4.65229779022069
0.505 4.65478996610902
0.51 4.65718842649024
0.515 4.65949366500681
0.52 4.66170613465417
0.525 4.66382624828122
0.53 4.66585437903259
0.535 4.66779086073393
0.54 4.66963598822155
0.545 4.67139001761726
0.55 4.67305316654944
0.555 4.67462561432094
0.56 4.67610750202459
0.565 4.67749893260664
0.57 4.67879997087867
0.575 4.68001064347814
0.58 4.68113093877784
0.585 4.68216080674415
0.59 4.68310015874435
0.595 4.68394886730265
0.6 4.68470676580487
0.605 4.68537364815162
0.61 4.68594926835948
0.615 4.68643334011012
0.62 4.68682553624679
0.625 4.68712548821791
0.63 4.68733278546749
0.635 4.68744697477188
0.64 4.68746755952279
0.645 4.68739399895627
0.65 4.68722570732774
0.655 4.68696205303313
0.66 4.68660235767659
0.665 4.68614589508542
0.67 4.68559189027337
0.675 4.68493951835372
0.68 4.68418790340445
0.685 4.68333611728826
0.69 4.6823831784312
0.695 4.68132805056491
0.7 4.6801696414384
0.705 4.67890680150741
0.71 4.67753832261095
0.715 4.67606293664711
0.72 4.67447931426327
0.725 4.67278606357911
0.73 4.67098172896496
0.735 4.6690647899036
0.74 4.66703365996928
0.745 4.66488668596563
0.75 4.66262214727318
0.755 4.66023825546844
0.76 4.65773315429019
0.765 4.65510492004516
0.77 4.65235156256599
0.775 4.64947102685935
0.78 4.6464611956131
0.785 4.64331989276973
0.79 4.64004488842054
0.795 4.63663390533418
0.8 4.6330846275063
0.805 4.62939471120936
0.81 4.62556179913673
0.815 4.62158353838144
0.82 4.61745760317416
0.825 4.61318172354042
0.83 4.60875372133768
0.835 4.60417155552035
0.84 4.59943337898095
0.845 4.59453760996777
0.85 4.58948302193282
0.855 4.58426885678983
0.86 4.57889496805786
0.865 4.57336200236933
0.87 4.56767163052643
0.875 4.56182684297535
0.88 4.55583232963992
0.885 4.54969497110644
0.89 4.54342447806702
0.895 4.5370342300407
0.9 4.53054238474922
0.905 4.52397335931567
0.91 4.51735982874991
0.915 4.51074545419789
0.92 4.50418865677756
0.925 4.49776791562069
0.93 4.49158933132982
0.935 4.48579763101359
0.94 4.48059253345419
0.945 4.4762537041259
0.95 4.47317993827007
0.955 4.47195284125547
0.96 4.47344467292383
0.965 4.47901036505304
0.97 4.49085138041194
0.975 4.51276242556832
0.98 4.55183555055335
0.985 4.62297779362201
0.99 4.76401720739564
0.995 5.11363394585526
1 9.26858717574037
};
\addplot [semithick, gray127, mark=*, mark size=2, mark options={solid}, only marks]
table {%
0.015 3.50057139068436
0.955 4.47195284125547
};

\nextgroupplot[
height=4.6cm,
label style={font=\footnotesize},
legend style={draw=none, fill=white, fill opacity=0.9, text opacity=1, font=\footnotesize},
scaled y ticks=manual:{}{\pgfmathparse{#1}},
tick align=outside,
tick label style={font=\scriptsize},
tick pos=left,
title style={font=\footnotesize, yshift=-2pt},
title={géodésique},
width=0.45\textwidth,
x grid style={darkgray176},
xlabel={\(\displaystyle t\)},
xmin=-0.05, xmax=1.05,
xtick style={color=black},
y grid style={darkgray176},
ymin=0.38937564836453, ymax=9.69140677310581,
ytick style={color=black},
yticklabels={}
]
\addplot [semithick, forestgreen4416044]
table {%
0 7.22321310791556
0.005 7.11245407161653
0.01 7.00216148493118
0.015 6.89234189954155
0.02 6.78300186270914
0.025 6.67414792121277
0.03 6.56578662713859
0.035 6.45792454561291
0.04 6.35056826455989
0.045 6.24372440655611
0.05 6.13739964284505
0.055 6.03160070956459
0.06 5.92633442623191
0.065 5.82160771651879
0.07 5.71742763134354
0.075 5.61380137429515
0.08 5.51073632939679
0.085 5.40824009120882
0.09 5.3063204972618
0.095 5.20498566280297
0.1 5.10424401783189
0.105 5.00410434639131
0.11 4.90457582807113
0.115 4.80566808167191
0.12 4.70739121096235
0.125 4.60975585245
0.13 4.51277322506577
0.135 4.41645518163975
0.14 4.32081426201969
0.145 4.22586374764613
0.15 4.13161771736037
0.155 4.03809110416727
0.16 3.94529975262029
0.165 3.85326047642013
0.17 3.76199111574268
0.175 3.67151059371487
0.18 3.58183897135383
0.185 3.49299750016866
0.19 3.40500867149774
0.195 3.31789626151982
0.2 3.23168537074173
0.205 3.14640245662501
0.21 3.06207535788543
0.215 2.97873330888378
0.22 2.89640694243362
0.225 2.81512827929982
0.23 2.73493070265417
0.235 2.65584891581257
0.24 2.57791888171117
0.245 2.50117774280382
0.25 2.42566372038676
0.255 2.35141599278988
0.26 2.27847455241566
0.265 2.20688004225677
0.27 2.13667357325829
0.275 2.06789652469891
0.28 2.00059033060209
0.285 1.93479625601543
0.29 1.87055516775951
0.295 1.80790730489352
0.3 1.74689205461029
0.305 1.68754773951385
0.31 1.62991142219899
0.315 1.57401873272386
0.32 1.5199037239375
0.325 1.46759875871119
0.33 1.4171344319724
0.335 1.36853952910789
0.34 1.32184102088247
0.345 1.277064093585
0.35 1.23423221176893
0.355 1.19336720976989
0.36 1.15448940723377
0.365 1.11761774321263
0.37 1.08276992300598
0.375 1.04996257183169
0.38 1.01921138957875
0.385 0.990531301273453
0.39 0.963936598423709
0.395 0.939441067034032
0.4 0.917058098742982
0.405 0.896800782179987
0.41 0.878681972231518
0.415 0.862714335428332
0.42 0.848910370114081
0.425 0.837282400440903
0.43 0.827842543588222
0.435 0.820602649943341
0.44 0.815574216355273
0.445 0.812768273004428
0.45 0.812195244943679
0.455 0.813864789968332
0.46 0.81778561516377
0.465 0.823965275225779
0.47 0.832409956428971
0.475 0.843124250864511
0.48 0.856110926238971
0.485 0.871370697052128
0.49 0.888902003303759
0.495 0.908700802977655
0.5 0.930760384391163
0.505 0.955071204069821
0.51 0.981620755143483
0.515 1.01039347037282
0.52 1.04137066288774
0.525 1.07453050658387
0.53 1.10984805695875
0.535 1.1472953120434
0.54 1.18684131202218
0.545 1.22845227521176
0.55 1.27209176727004
0.555 1.31772089989079
0.56 1.36529855475004
0.565 1.41478162816427
0.57 1.46612529172809
0.575 1.5192832641448
0.58 1.57420808951909
0.585 1.63085141752397
0.59 1.68916428109973
0.595 1.74909736764519
0.6 1.81060128005092
0.605 1.8736267843405
0.61 1.93812504116995
0.615 2.00404781893789
0.62 2.07134768676785
0.625 2.13997818615492
0.63 2.20989398055394
0.635 2.28105098268496
0.64 2.35340645973738
0.645 2.42691911707457
0.65 2.50154916133557
0.655 2.57725834413119
0.66 2.65400998772631
0.665 2.73176899425855
0.67 2.81050184014164
0.675 2.89017655734616
0.68 2.97076270326406
0.685 3.05223132082542
0.69 3.13455489046728
0.695 3.21770727549385
0.7 3.30166366223958
0.705 3.38640049635702
0.71 3.47189541639789
0.715 3.55812718580222
0.72 3.64507562420231
0.725 3.73272153889405
0.73 3.82104665721171
0.735 3.91003356040437
0.74 3.99966561950917
0.745 4.08992693372199
0.75 4.18080227153768
0.755 4.27227701495974
0.76 4.36433710698961
0.765 4.45696900250649
0.77 4.55015962267916
0.775 4.64389631288918
0.78 4.73816680422469
0.785 4.8329591784597
0.79 4.92826183647963
0.795 5.02406347005505
0.8 5.12035303681532
0.805 5.21711973839759
0.81 5.31435300141061
0.815 5.41204246129137
0.82 5.51017794873549
0.825 5.60874947846818
0.83 5.70774724051283
0.835 5.80716159305824
0.84 5.90698305788952
0.845 6.00720231679151
0.85 6.10781020996654
0.855 6.20879773556425
0.86 6.31015605062104
0.865 6.41187647254648
0.87 6.51395048221497
0.875 6.61636972692509
0.88 6.71912602473259
0.885 6.82221136867431
0.89 6.92561793159581
0.895 7.02933807126484
0.9 7.13336433566525
0.905 7.23768946797733
0.91 7.34230641196603
0.915 7.44720831709702
0.92 7.55238854377697
0.925 7.65784066812314
0.93 7.76355848615644
0.935 7.86953601883019
0.94 7.97576751555928
0.945 8.08224745853654
0.95 8.18897056509217
0.955 8.29593179212433
0.96 8.40312633772356
0.965 8.51054964399026
0.97 8.61819739868259
0.975 8.72606553721333
0.98 8.83415024434738
0.985 8.94244795258038
0.99 9.05095534607926
0.995 9.15966936027299
1 9.26858717652666
};
\addplot [semithick, gray127, mark=*, mark size=2, mark options={solid}, only marks]
table {%
0.45 0.812195244943679
};
\end{groupplot}

\end{tikzpicture}

  \end{center}
  \provGconvexiteQR[.5cm]
  \caption{Fonction de coût de Tyler évaluée le long du segment euclidien
    $(1-t)\MCov_0+t\MCov_1$ et le long de la géodésique
    $\MCov_0\geodmean_t\MCov_1$, pour les deux mêmes extrémités: une matrice
    mal conditionnée et son inverse, qui encadrent la vraie matrice de forme.
    La lecture euclidienne présente un maximum intérieur et donc deux minima
    locaux, aux extrémités; la lecture géodésique n'en a qu'un, au milieu du
    chemin. Les minima locaux de chaque courbe sont marqués d'un point. C'est le même critère dans les deux panneaux: seule change la
    notion de droite le long de laquelle on le lit. \provGconvexite}
  \label{fig:context-riemann-gconvexite}
\end{figure}

\subsection{Considérations numériques}

Ces algorithmes sont disponibles dans des bibliothèques spécialisées qui prennent en charge les variétés usuelles et leurs opérations: \texttt{Manopt} \sidecite[0cm]{boumal2014manopt}, \texttt{Pymanopt} \sidecite[1cm]{townsend2016pymanopt} et \texttt{Geomstats} \sidecite[2cm]{miolane2020geomstats}. En pratique, l'utilisateur ne fournit que la fonction de coût et, éventuellement, son gradient euclidien; la conversion \cref{eq:context-riemann-gradientspd} et la rétraction sont prises en charge par la bibliothèque.

Le coût par itération est dominé par les décompositions spectrales, en $\bigO(\Nvar^3)$, ce qui est le prix de l'exponentielle et du logarithme de matrice. C'est ici que se joue l'arbitrage évoqué en \cref{rem:context-riemann-LEbilan}: la métrique affine invariante paye ce coût à chaque évaluation de distance, là où la métrique log-euclidienne ne le paye qu'une fois par matrice, ce qui la rend préférable dès que le nombre de matrices manipulées devient grand devant leur dimension. La \cref{fig:context-riemann-convergence} compare enfin les vitesses de convergence des différents algorithmes introduits dans cette section.

\begin{figure}[h]
  \begin{center}
    % This file was created with matplot2tikz v0.5.3.
\begin{tikzpicture}

\definecolor{crimson2143940}{RGB}{214,39,40}
\definecolor{darkgray176}{RGB}{176,176,176}
\definecolor{darkorange25512714}{RGB}{255,127,14}
\definecolor{forestgreen4416044}{RGB}{44,160,44}

\begin{groupplot}[group style={group size=2 by 1, vertical sep=1.5cm, horizontal sep=1.4cm}]
\nextgroupplot[
height=4.6cm,
label style={font=\footnotesize},
legend cell align={left},
legend columns=2,
legend style={
  fill opacity=0.8,
  draw opacity=1,
  text opacity=1,
  at={(0,-0.45)},
  anchor=north west,
  draw=none
},
legend style={draw=none, fill=white, fill opacity=0.9, text opacity=1, font=\footnotesize},
log basis y={10},
tick align=outside,
tick label style={font=\scriptsize},
tick pos=left,
title style={font=\footnotesize, yshift=-2pt},
width=0.45\textwidth,
x grid style={darkgray176},
xlabel={itération},
xmin=-7.5, xmax=157.5,
xtick style={color=black},
y grid style={darkgray176},
ylabel={\(\displaystyle L - L^{\star}\)},
ymin=5e-15, ymax=43.685699240797,
ymode=log,
ytick style={color=black}
]
\addplot [line width=0.56pt, darkorange25512714, mark=*, mark size=1.5, mark repeat=10, mark options={solid}]
table {%
0 7.86251034725154
1 0.363424692178823
2 0.0371536351874333
3 0.00418682106529999
4 0.000502501383412124
5 6.28803955251556e-05
6 8.09665446332986e-06
7 1.06369673424922e-06
8 1.41783775831072e-07
9 1.91033464602697e-08
10 2.59511168110294e-09
11 3.5479708060393e-10
12 4.87556661710187e-11
13 6.72883970764815e-12
14 9.43245481721533e-13
15 1.31450406115619e-13
16 2.48689957516035e-14
17 1.06581410364015e-14
18 1e-14
19 1.06581410364015e-14
20 1e-14
21 1e-14
22 1.06581410364015e-14
23 1e-14
24 1e-14
25 1e-14
26 1e-14
27 1.77635683940025e-14
28 1e-14
29 1e-14
};
\addlegendentry{point fixe}
\addplot [line width=0.56pt, forestgreen4416044, mark=square*, mark size=1.5, mark repeat=10, mark options={solid}]
table {%
0 7.86251034725154
1 2.40206897416283
2 0.69500149712794
3 0.131134962925554
4 0.0160210430620431
5 0.00180659120125171
6 0.000224584913846471
7 2.95953945688154e-05
8 4.01079660505843e-06
9 5.51553277716721e-07
10 7.64736647340669e-08
11 1.06555244627771e-08
12 1.48935264121519e-09
13 2.08615347219165e-10
14 2.92654789291191e-11
15 4.11937151056918e-12
16 5.80868686483882e-13
17 8.34887714518118e-14
18 1.77635683940025e-14
19 1.06581410364015e-14
20 1e-14
21 1e-14
22 1e-14
23 1e-14
24 1e-14
25 1e-14
26 1e-14
27 1e-14
28 1e-14
};
\addlegendentry{gradient riemannien}
\addplot [line width=0.56pt, crimson2143940, mark=triangle*, mark size=1.5, mark repeat=10, mark options={solid}]
table {%
0 7.86251034725154
1 0.363424692178823
2 0.315212558531327
3 0.293013398365204
4 0.276245280407656
5 0.259499903380025
6 0.25204558688854
7 0.237746105232661
8 0.237015614501507
9 0.224412981717885
10 0.205180641085235
11 0.179758235874901
12 0.175274750352578
13 0.166904914017749
14 0.156525184015432
15 0.154396671367657
16 0.148877340275869
17 0.147110487849453
18 0.146870882781526
19 0.143836059753438
20 0.141736835098374
21 0.140147566239449
22 0.136574117920819
23 0.135215746772616
24 0.135157702605824
25 0.132572576405309
26 0.13086490610638
27 0.12974601733287
28 0.126688264046352
29 0.125640599407124
30 0.124209413859475
31 0.123187632859699
32 0.120565337681363
33 0.11966230507433
34 0.118375001629614
35 0.117567585152159
36 0.115151125714014
37 0.114378242441122
38 0.11318834117457
39 0.112590053183945
40 0.110303204255342
41 0.109640214281296
42 0.108524860077884
43 0.108107353943943
44 0.105917821558229
45 0.105346560855129
46 0.104292238590146
47 0.104026201386525
48 0.101917478744951
49 0.101423122122572
50 0.100420761639768
51 0.1002811456061
52 0.0982431352202831
53 0.0978139200962929
54 0.096856869823668
55 0.0968235001856605
56 0.0948489908846817
57 0.0944756301311376
58 0.0935587120166055
59 0.092977347692818
60 0.0928666367476652
61 0.0918104810126366
62 0.0916823505397044
63 0.0906717321494899
64 0.0905219467474137
65 0.0895601852096455
66 0.0893846931023088
67 0.0884748928811288
68 0.0882702755587896
69 0.0874150788860213
70 0.0871787127478889
71 0.0863801088193803
72 0.0861103318098007
73 0.0853694668075757
74 0.0850657061406501
75 0.0843827217747712
76 0.0840455604741521
77 0.0834194866427698
78 0.0830506494469301
79 0.0824793748678658
80 0.0820816235029902
81 0.0815619603272797
82 0.0811389019422286
83 0.0806667468566786
84 0.0802225748109269
85 0.0801485055167372
86 0.0792709150379736
87 0.0789351317794189
88 0.0783997983701301
89 0.0782324573164956
90 0.0771168970648759
91 0.0768757542276468
92 0.0762784771771727
93 0.0759045295910781
94 0.07585270263003
95 0.0751241966941922
96 0.0750266589790822
97 0.0743598156151819
98 0.0742172653320328
99 0.0736115910939539
100 0.0734259064266318
101 0.0728796022734972
102 0.0726539003914173
103 0.0721638612773727
104 0.0719023020927523
105 0.071464269100856
106 0.0711717754414352
107 0.0707805942233151
108 0.0704625127069782
109 0.0701124704350224
110 0.070099053673701
111 0.069246388213875
112 0.0688323039586329
113 0.0687444617752124
114 0.0678617343622641
115 0.0676975139241947
116 0.067216825682582
117 0.0669412919041079
118 0.0669372248655957
119 0.0663357614727769
120 0.0662804190171791
121 0.0657424668565927
122 0.0656391346095226
123 0.0651617156275464
124 0.0650148134297641
125 0.0645936175883381
126 0.0644085094564542
127 0.0640381293586429
128 0.0638207470077745
129 0.0634950420472293
130 0.0632514848892409
131 0.0629639972494083
132 0.0627001582207143
133 0.0626461838440306
134 0.0621386039152263
135 0.061942952199491
136 0.0616176960647401
137 0.0615599429585068
138 0.0608571740548136
139 0.0607372831710222
140 0.0603460018531141
141 0.0601340157277246
142 0.0598629819740992
143 0.0596991832877141
144 0.0591432431165568
145 0.059001823664353
146 0.0586692052158178
147 0.0584574135005145
148 0.0584328427575205
149 0.0580027117771742
150 0.0579308079580976
};
\addlegendentry{gradient euclidien projeté}

\nextgroupplot[
height=4.6cm,
label style={font=\footnotesize},
legend style={draw=none, fill=white, fill opacity=0.9, text opacity=1, font=\footnotesize},
log basis x={10},
log basis y={10},
scaled y ticks=manual:{}{\pgfmathparse{#1}},
tick align=outside,
tick label style={font=\scriptsize},
tick pos=left,
title style={font=\footnotesize, yshift=-2pt},
width=0.45\textwidth,
x grid style={darkgray176},
xlabel={temps (s)},
xmin=4.58657733252005e-05, xmax=0.095636132912106,
xmode=log,
xtick style={color=black},
y grid style={darkgray176},
ymin=5e-15, ymax=43.685699240797,
ymode=log,
ytick style={color=black},
yticklabels={}
]
\addplot [line width=0.56pt, darkorange25512714, mark=*, mark size=1.5, mark repeat=10, mark options={solid}]
table {%
0.000134457994136028 7.86251034725154
0.000262791989371181 0.363424692178823
0.000370291993021965 0.0371536351874333
0.000484249991131946 0.00418682106529999
0.000583499990170822 0.000502501383412124
0.000681624995195307 6.28803955251556e-05
0.000790416990639642 8.09665446332986e-06
0.000887499991222285 1.06369673424922e-06
0.000995458001852967 1.41783775831072e-07
0.00109787499241065 1.91033464602697e-08
0.00120604199764784 2.59511168110294e-09
0.00130241698934697 3.5479708060393e-10
0.00140970799839124 4.87556661710187e-11
0.00151824999193195 6.72883970764815e-12
0.00162249999993946 9.43245481721533e-13
0.00171820799005218 1.31450406115619e-13
0.00182520798989572 2.48689957516035e-14
0.00192112498916686 1.06581410364015e-14
0.00203020799381193 1e-14
0.00212629199086223 1.06581410364015e-14
0.00223574999836273 1e-14
0.0023313749989029 1e-14
0.0024367500009248 1.06581410364015e-14
0.0025357919948874 1e-14
0.00263254198944196 1e-14
0.00273804199241567 1e-14
0.00283649998891633 1e-14
0.00294233299791813 1.77635683940025e-14
0.0030391670006793 1e-14
0.00314441700174939 1e-14
};
\addplot [line width=0.56pt, forestgreen4416044, mark=square*, mark size=1.5, mark repeat=10, mark options={solid}]
table {%
7.7666001743637e-05 7.86251034725154
0.00024933299573604 2.40206897416283
0.000396040995838121 0.69500149712794
0.000540583001566119 0.131134962925554
0.000683374993968755 0.0160210430620431
0.000815499995951541 0.00180659120125171
0.000954290997469798 0.000224584913846471
0.00109595799585804 2.95953945688154e-05
0.00123816600535065 4.01079660505843e-06
0.00138091600092594 5.51553277716721e-07
0.00152291600534227 7.64736647340669e-08
0.00166424999770243 1.06555244627771e-08
0.00181429099757224 1.48935264121519e-09
0.00195575000543613 2.08615347219165e-10
0.00209816600545309 2.92654789291191e-11
0.0022575000039069 4.11937151056918e-12
0.0023952079936862 5.80868686483882e-13
0.00253616599366069 8.34887714518118e-14
0.00267712499771733 1.77635683940025e-14
0.00281845800054725 1.06581410364015e-14
0.00295979100337718 1e-14
0.00310087499383371 1e-14
0.00327079099952243 1e-14
0.00346995799918659 1e-14
0.00366987500456162 1e-14
0.00403991600614972 1e-14
0.0046931249962654 1e-14
0.00535612499515992 1e-14
0.00570204100222327 1e-14
};
\addplot [line width=0.56pt, crimson2143940, mark=triangle*, mark size=1.5, mark repeat=10, mark options={solid}]
table {%
6.49169960524887e-05 7.86251034725154
0.000266250004642643 0.363424692178823
0.000719250005204231 0.315212558531327
0.00118812499567866 0.293013398365204
0.00160812500689644 0.276245280407656
0.00209729200287256 0.259499903380025
0.00252066699613351 0.25204558688854
0.00300966699433047 0.237746105232661
0.00345212500542402 0.237015614501507
0.00396375000127591 0.224412981717885
0.00445770900114439 0.205180641085235
0.00472087500384077 0.179758235874901
0.00514262499928009 0.175274750352578
0.00561179200303741 0.166904914017749
0.00609412499761675 0.156525184015432
0.00647575000766665 0.154396671367657
0.00697062500694301 0.148877340275869
0.00741529199876823 0.147110487849453
0.00784583399945404 0.146870882781526
0.00834787500207312 0.143836059753438
0.0088305420067627 0.141736835098374
0.00919379200786352 0.140147566239449
0.00966341700404882 0.136574117920819
0.0100786669936497 0.135215746772616
0.0104992499982473 0.135157702605824
0.0109821670048404 0.132572576405309
0.0114731670037145 0.13086490610638
0.0118489589949604 0.12974601733287
0.0123411670065252 0.126688264046352
0.01276329200482 0.125640599407124
0.0132311670022318 0.124209413859475
0.013588333997177 0.123187632859699
0.0140965000027791 0.120565337681363
0.0145189589966321 0.11966230507433
0.0150228749989765 0.118375001629614
0.0154157090000808 0.117567585152159
0.0159077499993145 0.115151125714014
0.0163347500056261 0.114378242441122
0.016818208998302 0.11318834117457
0.0171724999963772 0.112590053183945
0.0176621249993332 0.110303204255342
0.0180753750028089 0.109640214281296
0.0185478750063339 0.108524860077884
0.0189249999966705 0.108107353943943
0.0194086249975953 0.105917821558229
0.0198375839972869 0.105346560855129
0.0203292499936651 0.104292238590146
0.0207349590054946 0.104026201386525
0.0212041250051698 0.101917478744951
0.0216205000033369 0.101423122122572
0.022093000006862 0.100420761639768
0.0224668750015553 0.1002811456061
0.0230224999977509 0.0982431352202831
0.0234589169995161 0.0978139200962929
0.0239622499939287 0.096856869823668
0.0243366670038085 0.0968235001856605
0.0248165420052828 0.0948489908846817
0.0252319590072148 0.0944756301311376
0.025729125001817 0.0935587120166055
0.0261490419943584 0.092977347692818
0.0265664589969674 0.0928666367476652
0.0270409999939147 0.0918104810126366
0.0274695419939235 0.0916823505397044
0.0279658340004971 0.0906717321494899
0.0284127920022001 0.0905219467474137
0.0289122090034653 0.0895601852096455
0.029341834000661 0.0893846931023088
0.0298279590060702 0.0884748928811288
0.0302563339937478 0.0882702755587896
0.0307225419965107 0.0874150788860213
0.0311339170002611 0.0871787127478889
0.031600874994183 0.0863801088193803
0.0320291249954607 0.0861103318098007
0.0326924170076381 0.0853694668075757
0.0331438749999506 0.0850657061406501
0.0336594170075841 0.0843827217747712
0.0340843750018394 0.0840455604741521
0.0345623339962913 0.0834194866427698
0.0349853340012487 0.0830506494469301
0.0354606670007342 0.0824793748678658
0.03588204200787 0.0820816235029902
0.036365542007843 0.0815619603272797
0.0368036669970024 0.0811389019422286
0.0372961249959189 0.0806667468566786
0.037716916995123 0.0802225748109269
0.0381350840034429 0.0801485055167372
0.0386037090065656 0.0792709150379736
0.0390269589988748 0.0789351317794189
0.039506834000349 0.0783997983701301
0.0398914590041386 0.0782324573164956
0.0403776669991203 0.0771168970648759
0.0408178750076331 0.0768757542276468
0.0412741670006653 0.0762784771771727
0.0416995419946034 0.0759045295910781
0.0421227500046371 0.07585270263003
0.0425934590020915 0.0751241966941922
0.0430192920030095 0.0750266589790822
0.0434992499940563 0.0743598156151819
0.0439396669971757 0.0742172653320328
0.0444252500019502 0.0736115910939539
0.0448744589957641 0.0734259064266318
0.0453581670008134 0.0728796022734972
0.0457867500081193 0.0726539003914173
0.0462557919963729 0.0721638612773727
0.0466624999971827 0.0719023020927523
0.0471333750028862 0.071464269100856
0.0475515419966541 0.0711717754414352
0.0480314169981284 0.0707805942233151
0.0484688340075081 0.0704625127069782
0.0489790419960627 0.0701124704350224
0.0493728340079542 0.070099053673701
0.0498717500013299 0.069246388213875
0.0503542499936884 0.0688323039586329
0.0507314169954043 0.0687444617752124
0.051194208994275 0.0678617343622641
0.0516052919992944 0.0676975139241947
0.052076499996474 0.067216825682582
0.0525094589975197 0.0669412919041079
0.0529469169996446 0.0669372248655957
0.0534438750037225 0.0663357614727769
0.0538812919985503 0.0662804190171791
0.0543647919985233 0.0657424668565927
0.0547908339940477 0.0656391346095226
0.0552597500063712 0.0651617156275464
0.055672250004136 0.0650148134297641
0.0561486250080634 0.0645936175883381
0.0565921670058742 0.0644085094564542
0.0570723749988247 0.0640381293586429
0.0575015840004198 0.0638207470077745
0.057989249995444 0.0634950420472293
0.0584279589966172 0.0632514848892409
0.0589200420072302 0.0629639972494083
0.0593696669966448 0.0627001582207143
0.0598005000065314 0.0626461838440306
0.0602748340024846 0.0621386039152263
0.0606953749957029 0.061942952199491
0.0611640840070322 0.0616176960647401
0.061557624998386 0.0615599429585068
0.062036208997597 0.0608571740548136
0.0624952500074869 0.0607372831710222
0.0630382090021158 0.0603460018531141
0.0635351249948144 0.0601340157277246
0.0640229169948725 0.0598629819740992
0.0644122090016026 0.0596991832877141
0.0649117919965647 0.0591432431165568
0.0653397090063663 0.059001823664353
0.0658235839946428 0.0586692052158178
0.0662376670079539 0.0584574135005145
0.0666732920071809 0.0584328427575205
0.0671417090052273 0.0580027117771742
0.0675697499973467 0.0579308079580976
};
\end{groupplot}

\end{tikzpicture}

  \end{center}
  \provConvergenceQR[.5cm]
  \caption{Convergence comparée de trois algorithmes minimisant le critère de
    Tyler depuis le même point de départ: l'itération de point fixe, la
    descente de gradient riemannienne à pas déterminé par recherche linéaire,
    et une descente de gradient euclidienne avec projection sur le cône. Les
    deux premières courbes se superposent, comme l'annonce la
    \cref{prop:context-riemann-pointfixe}: le point fixe \emph{est} un pas de
    gradient riemannien. La troisième minimise le même critère et converge vers
    le même point, mais n'y est pas parvenue au bout de son budget.
    \provConvergence}
  \label{fig:context-riemann-convergence}
\end{figure}


\section{Le cas complexe}
\label{sec:context-riemann-complexe}

Il nous reste à transposer ce qui précède au cas complexe circulaire, que nous emploierons à partir de la \cref{ch:context-geometry} et dans tout le \cref{ch:robust-detection}.

Les \cref{sec:context-definition-complex,sec:context-distributions-complex} s'ouvraient sur une mise en garde: la transposition naïve du cas réel au cas complexe y était insuffisante, une matrice de pseudo-covariance venant s'ajouter et les constantes de normalisation changeant. Il faut ici dire exactement l'inverse, et il vaut la peine de dire pourquoi. La géométrie riemannienne ne voit jamais l'espace des matrices hermitiennes comme un espace vectoriel \emph{complexe}: elle le voit comme un espace vectoriel \emph{réel}, muni du produit scalaire $\Realp\trace{\cdot}$. Il n'existe donc pas d'analogue géométrique de la pseudo-covariance, et la transposition est essentiellement une réécriture.

\begin{property}[{Géométrie du cône hermitien}]
\label{prop:context-riemann-conehermitien}
  L'ensemble $\HPD{\Nvar}$ est un ouvert de l'espace $\Herm{\Nvar}$ des matrices hermitiennes, vu comme espace vectoriel \emph{réel}. C'est donc une variété de dimension réelle $\Nvar^2$, d'espace tangent $\Tangent{\MCov}{\HPD{\Nvar}}=\Herm{\Nvar}$ en tout point. La métrique affine invariante s'y écrit, dans sa version à deux paramètres de la \cref{def:context-riemann-affineinvariante},
  \begin{equation}
    \metric{\tvec}{\tvecbis}{\MCov}
    =
    \Realp\left(
      \alpha\,\trace{\inv{\MCov}\tvec\inv{\MCov}\tvecbis}
      +
      \beta\,\trace{\inv{\MCov}\tvec}\trace{\inv{\MCov}\tvecbis}
    \right),
    \label{eq:context-riemann-metriquecomplexe}
  \end{equation}
  et l'ensemble des résultats des \cref{prop:context-riemann-projspd,prop:context-riemann-geodspd,prop:context-riemann-explogspd,prop:context-riemann-transportspd,prop:context-riemann-distspd} reste valable en remplaçant partout la transposée $(\cdot)^\T$ par la transposée conjuguée $(\cdot)^\herm$ et l'opérateur $\symop$ par $\hermop(\mat{Z})=\tfrac{1}{2}(\mat{Z}+\mat{Z}^\herm)$.
\end{property}

Trois points seulement demandent une attention particulière.

\paragraph{La dimension.} Elle passe de $\Nvar(\Nvar+1)/2$ à $\Nvar^2$. Ce n'est pas un détail cosmétique: cette dimension est celle qui apparaîtra dans les bornes de Cramér-Rao de la \cref{sec:context-geometry-icrb}, et c'est elle qui gouverne le nombre d'observations nécessaires à une estimation bien posée.

\paragraph{Le gradient.} La \cref{def:context-riemann-gradient} demande une dérivée directionnelle réelle, or la log-vraisemblance complexe n'est pas holomorphe. C'est exactement la difficulté que nous avons traitée en \cref{def:context-complex-wirtinger}: le gradient euclidien à employer dans \cref{eq:context-riemann-gradientspd} est $\nabla^{\mathrm{e}}f=2\,\partial f/\partial\conj{\MCov}$ au sens de Wirtinger, et non $\partial f/\partial\MCov$. Le facteur $2$ annoncé alors trouve ici son usage.

\paragraph{Le groupe unitaire.} La transposition de la \cref{sec:context-riemann-orthogonal} est la seule à avoir une conséquence algorithmique. Le groupe unitaire $\Unitgroup{\Nvar}=\{\Unit;\,\Unit^\herm\Unit=\Id{\Nvar}\}$ est de dimension réelle $\Nvar^2$, contre $\Nvar(\Nvar-1)/2$ pour $\Orthgroup{\Nvar}$, et il est surtout \textbf{connexe}, là où $\Orthgroup{\Nvar}$ possède deux composantes. Le piège d'initialisation signalé en \cref{rem:context-riemann-stiefelvscone} disparaît donc dans le cas complexe.

Le \cref{tab:context-riemann-complexe} rassemble ces correspondances, dans le même esprit que le \cref{tab:context-complex-summary}.

\begin{table}
  \centering
  \footnotesize
  \begin{tabular}{@{}lll@{}}
    \toprule
     & Cas réel & Cas complexe circulaire \\
    \midrule
    Variété                 & $\SPD{\Nvar}$ & $\HPD{\Nvar}$ \\
    Espace ambiant          & $\Sym{\Nvar}$ & $\Herm{\Nvar}$, e.v. \emph{réel} \\
    Dimension (réelle)      & $\Nvar(\Nvar+1)/2$ & $\Nvar^2$ \\
    Métrique                & $\alpha\trace{\cdot}+\beta\trace{\cdot}\trace{\cdot}$ & idem, précédé de $\Realp$ \\
    Projection              & $\symop$ & $\hermop$ \\
    Géodésique, $\exp$, $\log$, $\geodist$ & \cref{eq:context-riemann-geodspd}--\cref{eq:context-riemann-distspd} & identiques, $(\cdot)^\T\to(\cdot)^\herm$ \\
    Gradient euclidien      & $\nabla^{\mathrm{e}}f$ & $2\,\partial f/\partial\conj{\MCov}$ \\
    Groupe associé          & $\Orthgroup{\Nvar}$, $\Nvar(\Nvar-1)/2$, 2 composantes & $\Unitgroup{\Nvar}$, $\Nvar^2$, connexe \\
    \bottomrule
  \end{tabular}
  \vspace{0.5em}
  \caption{Correspondances entre les cas réel et complexe circulaire pour la géométrie riemannienne. À la différence de ce que nous avions observé en \cref{sec:context-distributions-complex} pour les distributions, la transposition est ici essentiellement mécanique.}
  \label{tab:context-riemann-complexe}
\end{table}




\section{Bilan}
\label{sec:context-riemann-bilan}

% ============================ TODO texte (~1,5 page) ===============================
% Bilan absent lui aussi, alors que \cref{ch:commentlire} le promet. Meme cahier
% des charges qu'en \cref{sec:context-definition-bilan}: une declaration de
% dette, pas un resume.
%
%  1. CE QUE LE CHAPITRE ACHETE AUX CONTRIBUTIONS:
%       - le cone SPD et sa distance (\cref{sec:context-riemann-spd}) portent la
%         recursion en ligne du \cref{sec:detection-online-change-detection};
%       - les moyennes du cone (\cref{sec:context-riemann-moyennes}) sont
%         exactement l'objet du \cref{sec:spdnet-batchnorm}, et la moyenne de
%         Frechet ponderee est ce que devient l'attention en
%         \cref{sec:perspectives-transformers};
%       - la variete de Stiefel contraint les poids \BiMap{} et fonde
%         \cref{sec:spdnet-federe};
%       - l'optimisation riemannienne (\cref{sec:context-riemann-optim}) est
%         l'outil commun aux deux chapitres de contribution.
%
%  2. CE QUI RESTE INACCESSIBLE, et qu'il vaut mieux annoncer ici que laisser
%     decouvrir: ni logarithme ni distance en forme close sur Stiefel, d'ou les
%     approximations du \cref{ch:spdnet}; c'est le fil « l'objet juste est
%     toujours iteratif » que le memoire suit ensuite.
%
%  3. RENVOYER a \cref{fig:commentlire-carte}.
% ====================================================================================

\wip{Bilan du chapitre 3}{Tableau consolidant les géométries du cône et leurs outils (métriques, géodésiques, moyennes, rétractions), et ce qui reste inaccessible en forme close --- le fil « l'objet juste est toujours itératif ». Une phrase de liaison. Environ 1,5 page.}

